% Generated by scripts/build.py - do not edit.

\documentclass[11pt]{article}

% Preamble for the English bank (tex/bank-en.tex); compiles with pdfLaTeX or XeLaTeX.
\usepackage{amsmath,amssymb,amsthm}
\newcommand{\R}{\mathbb{R}}
\newcommand{\N}{\mathbb{N}}
\newcommand{\Q}{\mathbb{Q}}
\newcommand{\Z}{\mathbb{Z}}
\newcommand{\eps}{\varepsilon}
\newenvironment{question}{\par\noindent\textbf{Question.}\ }{\par}
\newenvironment{hint}{\par\noindent\textbf{Hint.}\ }{\par}
\newenvironment{solution}{\par\noindent\textbf{Solution.}\ }{\par}


\begin{document}

\title{Exam Question Bank}\date{}\maketitle

\section*{2015/16 Semester B Moed B}

\subsection*{Question 1}

\begin{question}
\textbf{Prove or disprove:} Let $f$ be a strictly increasing function on the interval $[a,b]$; then $f$ is convex.
\end{question}

\begin{hint}
Look for a strictly increasing function whose graph ``bends'' in both directions, for example an odd polynomial.
\end{hint}

\begin{solution}
\textbf{The statement is false.} Monotonicity is unrelated to convexity.

Counterexample: $f(x)=x^3$ on the interval $[a,b]=[-1,1]$. Since $f'(x)=3x^2\ge 0$ and it vanishes only at the single point $x=0$, the function is strictly increasing on the interval (one can also see this directly: if $x<y$ then $y^3-x^3=(y-x)(x^2+xy+y^2)>0$, since $x^2+xy+y^2=(x+\tfrac y2)^2+\tfrac34y^2>0$ when $x\ne y$).

However, $f''(x)=6x$, so $f''<0$ on $(-1,0)$ and $f''>0$ on $(0,1)$. That is, $f$ is concave on $[-1,0]$ and convex on $[0,1]$, and hence it is not convex (and not concave either) on the whole interval.

A direct check by the definition: for the points $x_1=-1,\ x_2=0$ and the midpoint $-\tfrac12$:
\[ f\!\left(-\tfrac12\right)=-\tfrac18 \;>\; -\tfrac12=\frac{f(-1)+f(0)}{2}, \]
i.e., the chord between $(-1,-1)$ and $(0,0)$ lies \emph{below} the graph, contradicting convexity ($f(\tfrac{x_1+x_2}{2})\le \tfrac{f(x_1)+f(x_2)}{2}$). On the other hand, for $x_1=0,\ x_2=1$: $f(\tfrac12)=\tfrac18<\tfrac12=\tfrac{f(0)+f(1)}2$, so the function is not concave either. Hence the example disproves the statement under either convention for the word ``convex''.

(Another example: $f(x)=\sqrt{x}$ on $[0,1]$ is strictly increasing and strictly concave.)
\end{solution}

\subsection*{Question 2}

\begin{question}
\textbf{Prove or disprove:} The following function is differentiable at zero:
\[ f(x)=\begin{cases} x\sqrt{|x|}\cos(\frac{1}{x}) & x\neq 0\\ 0 & x=0\end{cases}. \]
\end{question}

\begin{hint}
Compute the derivative by the definition, as the limit of the difference quotient, and use the fact that ``bounded times tending to zero tends to zero''.
\end{hint}

\begin{solution}
\textbf{The statement is true.} We compute using the definition of the derivative. For $h\neq 0$:
\[ \frac{f(h)-f(0)}{h}=\frac{h\sqrt{|h|}\cos(\frac1h)-0}{h}=\sqrt{|h|}\cos\!\left(\frac1h\right). \]
We have $\left|\sqrt{|h|}\cos(\frac1h)\right|\le \sqrt{|h|}$, and $\sqrt{|h|}\to 0$ as $h\to 0$. By the Squeeze theorem (or: a bounded function, $\cos(\frac1h)$, times a function tending to zero, $\sqrt{|h|}$, tends to zero):
\[ f'(0)=\lim_{h\to 0}\frac{f(h)-f(0)}{h}=\lim_{h\to0}\sqrt{|h|}\cos\!\left(\frac1h\right)=0. \]
The limit exists and is finite, hence $f$ is differentiable at zero and $\boxed{f'(0)=0}$.
\end{solution}

\subsection*{Question 3}

\begin{question}
\textbf{Prove or disprove:} The integral $\int_2^\infty \frac{dx}{x\ln(x)}dx$ converges.
\end{question}

\begin{hint}
Find an antiderivative using the substitution $t=\ln x$.
\end{hint}

\begin{solution}
\textbf{The statement is false} -- the integral diverges.

The function $\frac{1}{x\ln x}$ is continuous and positive on $[2,\infty)$, so the improper integral is defined as $\lim_{R\to\infty}\int_2^R\frac{dx}{x\ln x}$.
With the substitution $t=\ln x$, $dt=\frac{dx}{x}$:
\[ \int\frac{dx}{x\ln x}=\int\frac{dt}{t}=\ln|t|+C=\ln(\ln x)+C \qquad (x>1). \]
By the Fundamental Theorem (the Newton--Leibniz formula):
\[ \int_2^R\frac{dx}{x\ln x}=\ln(\ln R)-\ln(\ln 2)\xrightarrow[R\to\infty]{}\infty, \]
since $\ln R\to\infty$ and hence also $\ln(\ln R)\to\infty$. The limit is not finite, so the integral \textbf{diverges}.
\end{solution}

\subsection*{Question 4a}

\begin{question}
Compute the following limit: $\lim_{x\to 0}\frac{1-\cos(x^2)}{x^2\sin^2(x)}$.
\end{question}

\begin{hint}
Use the Taylor expansion of $\cos t$ around $0$ with $t=x^2$, and the fundamental limit $\frac{\sin x}{x}\to1$.
\end{hint}

\begin{solution}
Write
\[ \frac{1-\cos(x^2)}{x^2\sin^2 x}=\frac{1-\cos(x^2)}{x^4}\cdot\left(\frac{x}{\sin x}\right)^2. \]
By the Maclaurin expansion $\cos t=1-\frac{t^2}{2}+o(t^2)$, and with $t=x^2$: $1-\cos(x^2)=\frac{x^4}{2}+o(x^4)$, hence $\frac{1-\cos(x^2)}{x^4}\to\frac12$.
(Alternatively: $1-\cos(x^2)=2\sin^2(\frac{x^2}{2})$, hence $\frac{1-\cos(x^2)}{x^4}=\frac12\left(\frac{\sin(x^2/2)}{x^2/2}\right)^2\to\frac12$.)
Also $\frac{x}{\sin x}\to 1$ by the fundamental limit. By the arithmetic of limits:
\[ \lim_{x\to0}\frac{1-\cos(x^2)}{x^2\sin^2 x}=\frac12\cdot 1=\boxed{\frac12}. \]
\end{solution}

\subsection*{Question 4b}

\begin{question}
Show that:
\[ \frac15\le \arctan(3)-\arctan(1)\le 1. \]
\end{question}

\begin{hint}
Apply Lagrange's theorem to $\arctan$ on the interval $[1,3]$ and bound the derivative.
\end{hint}

\begin{solution}
Define $g(t)=\arctan t$. The function is continuous on $[1,3]$ and differentiable on $(1,3)$ with $g'(t)=\frac{1}{1+t^2}$. By Lagrange's theorem there exists $c\in(1,3)$ such that
\[ \arctan(3)-\arctan(1)=g'(c)(3-1)=\frac{2}{1+c^2}. \]
Since $1<c<3$ we have $2<1+c^2<10$, and therefore
\[ \frac{2}{10}<\frac{2}{1+c^2}<\frac{2}{2}, \qquad\text{i.e.}\qquad \frac15<\arctan(3)-\arctan(1)<1, \]
and in particular the required weak inequality holds. (Numerical check: $\arctan 3-\arctan 1\approx 0.4636$.)
\end{solution}

\subsection*{Question 5}

\begin{question}
\begin{enumerate}
  \item (15 pts) Compute the following indefinite integral:
  \[ \int\frac{xdx}{(x-1)(x^2+2x+2)}. \]
  \item (5 pts) Compute the area enclosed under the graph of the function $f(x)=\frac{x}{(x-1)(x^2+2x+2)}$ on the interval $[-1,0]$.
\end{enumerate}
\end{question}

\begin{hint}
Decompose into partial fractions: $\frac{A}{x-1}+\frac{Bx+C}{x^2+2x+2}$. For the quadratic factor, complete the square: $x^2+2x+2=(x+1)^2+1$.
\end{hint}

\begin{hint}
In part (b), check the sign of the function on the interval before computing the area.
\end{hint}

\begin{solution}
\begin{enumerate}
  \item The factor $x^2+2x+2=(x+1)^2+1$ has no real roots (the discriminant is $4-8<0$), so we decompose:
  \[ \frac{x}{(x-1)(x^2+2x+2)}=\frac{A}{x-1}+\frac{Bx+C}{x^2+2x+2}
     \iff x=A(x^2+2x+2)+(Bx+C)(x-1). \]
  Substituting $x=1$: $1=5A$, i.e., $A=\frac15$. Comparing the coefficients of $x^2$: $0=A+B$, hence $B=-\frac15$. Comparing the constant terms: $0=2A-C$, hence $C=\frac25$. That is,
  \[ \frac{x}{(x-1)(x^2+2x+2)}=\frac{1}{5}\cdot\frac{1}{x-1}-\frac15\cdot\frac{x-2}{x^2+2x+2}. \]
  Write $x-2=\frac12(2x+2)-3$:
  \[ \frac{x-2}{x^2+2x+2}=\frac12\cdot\frac{2x+2}{x^2+2x+2}-\frac{3}{(x+1)^2+1}. \]
  Hence (using $\int\frac{u'}{u}=\ln|u|$ and $\int\frac{dt}{t^2+1}=\arctan t$):
  \[ \boxed{\int\frac{x\,dx}{(x-1)(x^2+2x+2)}=\frac15\ln|x-1|-\frac1{10}\ln(x^2+2x+2)+\frac35\arctan(x+1)+C.} \]
  (Check: differentiating the expression gives back the integrand.)

  \item On the interval $[-1,0]$: $x\le 0$, $x-1<0$, $x^2+2x+2>0$, hence $f(x)\ge0$ (equal to zero only at $x=0$). The function is continuous on the interval, so the area is $\int_{-1}^0 f(x)\,dx$. Denote by $F$ the antiderivative from part (a) and use the Newton--Leibniz formula:
  \[ F(0)=\frac15\ln1-\frac1{10}\ln2+\frac35\cdot\frac\pi4=-\frac{\ln2}{10}+\frac{3\pi}{20},\qquad
     F(-1)=\frac15\ln2-\frac1{10}\ln1+\frac35\arctan0=\frac{\ln 2}{5}. \]
  Therefore
  \[ S=F(0)-F(-1)=\boxed{\frac{3\pi}{20}-\frac{3\ln2}{10}}\approx 0.2633. \]
\end{enumerate}
\end{solution}

\subsection*{Question 6a}

\begin{question}
Find the domain, the intersection points with the axes, the intervals of increase and decrease, the extremum points, the asymptotes, and the intervals of convexity and concavity of the following function:
\[ f(x)=(x+2)e^{\frac1x}. \]
\end{question}

\begin{hint}
For the oblique asymptote compute $\lim\frac{f(x)}{x}$ and then $\lim (f(x)-x)$; use the substitution $t=\frac1x$ and the limit $\frac{e^t-1}{t}\to1$.
\end{hint}

\begin{solution}
\textbf{Domain:} $x\neq0$.

\textbf{Intersection with the axes:} There is no intersection with the $y$-axis (since $0$ is not in the domain). $f(x)=0\iff x+2=0$ (since $e^{1/x}>0$), so the intersection point with the $x$-axis is $(-2,0)$. In addition, $f<0$ for $x<-2$ and $f>0$ for $x>-2,\ x\ne0$.

\textbf{First derivative:}
\[ f'(x)=e^{1/x}+(x+2)e^{1/x}\cdot\left(-\frac1{x^2}\right)=e^{1/x}\,\frac{x^2-x-2}{x^2}=\frac{(x-2)(x+1)}{x^2}e^{1/x}. \]
The sign is determined by $(x-2)(x+1)$:
\begin{itemize}
  \item $f'>0$ on $(-\infty,-1)$ and on $(2,\infty)$ -- \textbf{increasing};
  \item $f'<0$ on $(-1,0)$ and on $(0,2)$ -- \textbf{decreasing}.
\end{itemize}
\textbf{Extrema:} At $x=-1$ the derivative changes sign from $+$ to $-$: a local maximum $f(-1)=e^{-1}=\frac1e$. At $x=2$ from $-$ to $+$: a local minimum $f(2)=4e^{1/2}=4\sqrt e$.

\textbf{Vertical asymptotes:} As $x\to0^+$, $\frac1x\to+\infty$, hence $e^{1/x}\to\infty$ and $f(x)\to 2\cdot\infty=+\infty$: $x=0$ is a vertical asymptote (from the right). As $x\to0^-$, $\frac1x\to-\infty$, $e^{1/x}\to0$ and $f(x)\to 0$ (no asymptote from the left; the graph has a ``hole'' at the origin).

\textbf{Oblique asymptotes:} $\frac{f(x)}{x}=\left(1+\frac2x\right)e^{1/x}\to1$ as $x\to\pm\infty$. Also
\[ f(x)-x=x\left(e^{1/x}-1\right)+2e^{1/x}=\frac{e^{t}-1}{t}+2e^{t}\Big|_{t=1/x}\;\xrightarrow[t\to0]{}\;1+2=3. \]
Hence $y=x+3$ is an oblique asymptote both at $+\infty$ and at $-\infty$. (There is no horizontal asymptote.)

\textbf{Second derivative:} Differentiate $f'(x)=\left(1-\frac1x-\frac2{x^2}\right)e^{1/x}$:
\[ f''(x)=\left(\frac1{x^2}+\frac4{x^3}\right)e^{1/x}-\frac1{x^2}\left(1-\frac1x-\frac2{x^2}\right)e^{1/x}=\frac{5x+2}{x^4}e^{1/x}. \]
The sign is determined by $5x+2$:
\begin{itemize}
  \item $f''<0$ on $(-\infty,-\frac25)$ -- \textbf{concave} ($\cap$);
  \item $f''>0$ on $(-\frac25,0)$ and on $(0,\infty)$ -- \textbf{convex} ($\cup$).
\end{itemize}
Inflection point: $x=-\frac25$, $f(-\frac25)=\frac85e^{-5/2}$.
\end{solution}

\subsection*{Question 6b}

\begin{question}
Compute the volume of the solid of revolution obtained by rotating the graph of the function:
\[ f(x)=\sin^2(x) \]
around the $x$-axis on the interval $[0,\frac\pi2]$.
\end{question}

\begin{hint}
The volume of the solid of revolution is $\pi\int_a^b f^2(x)\,dx$; to reduce the power, use $\sin^2x=\frac{1-\cos2x}{2}$.
\end{hint}

\begin{solution}
The volume of a solid of revolution around the $x$-axis: $V=\pi\int_0^{\pi/2}f^2(x)\,dx=\pi\int_0^{\pi/2}\sin^4x\,dx$. We use the identities
\[ \sin^4x=\left(\frac{1-\cos2x}{2}\right)^2=\frac14\left(1-2\cos2x+\cos^22x\right)=\frac14\left(1-2\cos2x+\frac{1+\cos4x}{2}\right)=\frac38-\frac12\cos2x+\frac18\cos4x. \]
Hence
\[ \int_0^{\pi/2}\sin^4x\,dx=\left[\frac38x-\frac14\sin2x+\frac1{32}\sin4x\right]_0^{\pi/2}=\frac{3\pi}{16}, \]
(since $\sin\pi=\sin2\pi=0$), and therefore
\[ \boxed{V=\frac{3\pi^2}{16}}. \]
\end{solution}

\section*{2015/16 Semester B Moed C}

\subsection*{Question 1}

\begin{question}
\textbf{Prove or disprove:} The following equation has a solution in $[0,\pi]$:
\[ \ln(x+1)-\sin(x)=1 \]
\end{question}

\begin{hint}
Define $g(x)=\ln(x+1)-\sin(x)-1$ and check its sign at the endpoints of the interval.
\end{hint}

\begin{solution}
\textbf{The statement is true.} Define $g(x)=\ln(x+1)-\sin x-1$. The function $g$ is continuous on $[0,\pi]$ as a sum of continuous functions ($\ln(x+1)$ is continuous for $x>-1$).

At the endpoints:
\[ g(0)=\ln1-\sin0-1=-1<0,\qquad g(\pi)=\ln(\pi+1)-\sin\pi-1=\ln(\pi+1)-1>0, \]
where the last inequality follows from $\pi+1>e$ (since $\pi+1>4>e$), and hence $\ln(\pi+1)>\ln e=1$. (Numerically $\ln(\pi+1)\approx1.421$.)

By the Intermediate Value Theorem (Bolzano's theorem) there exists $c\in(0,\pi)$ for which $g(c)=0$, i.e., $\ln(c+1)-\sin c=1$. Hence the equation has a solution in the interval.
\end{solution}

\subsection*{Question 2}

\begin{question}
\textbf{Prove or disprove:} The following function is continuous at zero:
\[ f(x)=\begin{cases} \frac{\sin(2x)}{\sqrt{x+4}-2} & x\neq 0\\ 5 & x=0\end{cases}. \]
\end{question}

\begin{hint}
Multiply the numerator and the denominator by the conjugate $\sqrt{x+4}+2$.
\end{hint}

\begin{solution}
\textbf{The statement is false.} $f$ is continuous at $0$ if and only if $\lim_{x\to0}f(x)=f(0)=5$. We compute the limit (for $x\ne0$, $x\ge-4$) by multiplying by the conjugate:
\[ \frac{\sin 2x}{\sqrt{x+4}-2}=\frac{\sin 2x\,(\sqrt{x+4}+2)}{(x+4)-4}=\frac{\sin2x}{x}\left(\sqrt{x+4}+2\right)=2\cdot\frac{\sin 2x}{2x}\cdot\left(\sqrt{x+4}+2\right). \]
By the fundamental limit $\frac{\sin2x}{2x}\to1$, and by continuity of the square root $\sqrt{x+4}+2\to4$. By the arithmetic of limits
\[ \lim_{x\to0}f(x)=2\cdot1\cdot4=8\neq5=f(0). \]
Hence $f$ is \textbf{not continuous} at zero (it has a removable discontinuity; it would be continuous had we defined $f(0)=8$).
\end{solution}

\subsection*{Question 3}

\begin{question}
\textbf{Prove or disprove:} The integral $\int_0^1\ln^2(x)dx$ converges.
\end{question}

\begin{hint}
Find an antiderivative using integration by parts (twice), and compute the limit $\lim_{x\to0^+}x\ln^k x$ using L'Hôpital's rule.
\end{hint}

\begin{solution}
\textbf{The statement is true}, and the value of the integral is $2$.

The function $\ln^2x$ is continuous on $(0,1]$ and unbounded near $0$, so this is an improper integral: $\int_0^1\ln^2x\,dx=\lim_{\eps\to0^+}\int_\eps^1\ln^2x\,dx$.

\textbf{Antiderivative} -- integration by parts with $u=\ln^2x,\ v'=1$:
\[ \int\ln^2x\,dx=x\ln^2x-\int x\cdot\frac{2\ln x}{x}dx=x\ln^2x-2\int\ln x\,dx=x\ln^2x-2(x\ln x-x)+C. \]
Hence
\[ \int_\eps^1\ln^2x\,dx=\big[x\ln^2x-2x\ln x+2x\big]_\eps^1=2-\left(\eps\ln^2\eps-2\eps\ln\eps+2\eps\right). \]
\textbf{The limits:} By L'Hôpital's rule ($\frac{\infty}{\infty}$):
\[ \lim_{\eps\to0^+}\eps\ln\eps=\lim\frac{\ln\eps}{1/\eps}=\lim\frac{1/\eps}{-1/\eps^2}=\lim(-\eps)=0, \]
\[ \lim_{\eps\to0^+}\eps\ln^2\eps=\lim\frac{\ln^2\eps}{1/\eps}=\lim\frac{2\ln\eps/\eps}{-1/\eps^2}=\lim(-2\eps\ln\eps)=0. \]
Hence $\int_0^1\ln^2x\,dx=2-0=2$: the integral \textbf{converges}.

(Another way: $\ln^2x\le\frac{C}{\sqrt x}$ near $0$ since $\sqrt x\ln^2x\to0$, and $\int_0^1\frac{dx}{\sqrt x}$ converges, so by the comparison test the integral converges.)
\end{solution}

\subsection*{Question 4a}

\begin{question}
Compute the following limit: $\lim_{x\to0}\frac{\cos(\sin(x))-\cos(x)}{x^4}$.
\end{question}

\begin{hint}
Expand $\cos t$ and $\sin x$ by Maclaurin up to order $4$ (note that $\sin^2x=x^2-\frac{x^4}{3}+o(x^4)$).
\end{hint}

\begin{solution}
We use the Maclaurin expansions
$\sin x=x-\frac{x^3}{6}+o(x^4)$ and $\cos t=1-\frac{t^2}{2}+\frac{t^4}{24}+o(t^4)$.
We compute:
\[ \sin^2x=\left(x-\frac{x^3}{6}+o(x^4)\right)^2=x^2-\frac{x^4}{3}+o(x^4),\qquad \sin^4x=x^4+o(x^4). \]
Since $\sin x\to0$ and $|\sin x|\le|x|$, we may substitute $t=\sin x$ (the remainder $o(\sin^4x)$ is $o(x^4)$):
\[ \cos(\sin x)=1-\frac{\sin^2x}{2}+\frac{\sin^4x}{24}+o(x^4)=1-\frac{x^2}{2}+\frac{x^4}{6}+\frac{x^4}{24}+o(x^4)=1-\frac{x^2}{2}+\frac{5x^4}{24}+o(x^4). \]
Also $\cos x=1-\frac{x^2}{2}+\frac{x^4}{24}+o(x^4)$. Hence
\[ \cos(\sin x)-\cos x=\frac{4x^4}{24}+o(x^4)=\frac{x^4}{6}+o(x^4), \]
and therefore
\[ \lim_{x\to0}\frac{\cos(\sin x)-\cos x}{x^4}=\boxed{\frac16}. \]
(Another way: $\cos A-\cos B=-2\sin\frac{A+B}{2}\sin\frac{A-B}{2}$ with $A=\sin x$, $B=x$; then $\frac{A+B}{2}\sim x$ and $\frac{A-B}{2}=\frac{\sin x-x}{2}\sim-\frac{x^3}{12}$, and we get $-2\cdot x\cdot(-\frac{x^3}{12})=\frac{x^4}{6}$.)
\end{solution}

\subsection*{Question 4b}

\begin{question}
Show that for every $0<a<b$ and every $p>1$:
\[ pa^{p-1}(b-a)\le b^p-a^p\le pb^{p-1}(b-a) \]
\end{question}

\begin{hint}
Apply Lagrange's theorem to $f(x)=x^p$ on the interval $[a,b]$, and use the monotonicity of $x^{p-1}$.
\end{hint}

\begin{solution}
Let $f(x)=x^p$. It is continuous on $[a,b]$ and differentiable on $(a,b)$ (since $a>0$), with $f'(x)=px^{p-1}$. By Lagrange's theorem there exists $c\in(a,b)$ such that
\[ b^p-a^p=f'(c)(b-a)=pc^{p-1}(b-a). \]
Since $p>1$, the exponent $p-1>0$, so the function $x\mapsto x^{p-1}$ is strictly increasing on $(0,\infty)$. From $0<a<c<b$ it follows that $a^{p-1}<c^{p-1}<b^{p-1}$. Multiply by $p(b-a)>0$:
\[ pa^{p-1}(b-a)<pc^{p-1}(b-a)=b^p-a^p<pb^{p-1}(b-a), \]
and in particular the required inequality holds (even strictly).
\end{solution}

\subsection*{Question 5}

\begin{question}
\begin{enumerate}
  \item (15 pts) Compute the following indefinite integral:
  \[ \int\frac{dx}{\cos^3(x)}. \]
  \item (5 pts) Compute the area enclosed under the graph of the function $f(x)=\frac{1}{\cos^3(x)}$ on the interval $[0,1]$.
\end{enumerate}
\end{question}

\begin{hint}
Write $\frac{1}{\cos^3x}=\frac{1}{\cos x}\cdot\frac1{\cos^2x}$ and integrate by parts with $v'=\frac1{\cos^2x}$, $v=\tan x$.
\end{hint}

\begin{hint}
You will also need $\int\frac{dx}{\cos x}=\ln\left|\frac1{\cos x}+\tan x\right|+C$ (or: multiply by $\cos x$ and substitute $t=\sin x$).
\end{hint}

\begin{solution}
\begin{enumerate}
  \item Denote $I=\int\frac{dx}{\cos^3x}$ and integrate by parts with $u=\frac1{\cos x}$, $v'=\frac1{\cos^2x}$; then $u'=\frac{\sin x}{\cos^2x}$, $v=\tan x$:
  \[ I=\frac{\tan x}{\cos x}-\int\frac{\sin x}{\cos^2x}\cdot\frac{\sin x}{\cos x}dx=\frac{\tan x}{\cos x}-\int\frac{1-\cos^2x}{\cos^3x}dx=\frac{\tan x}{\cos x}-I+\int\frac{dx}{\cos x}. \]
  We compute $\int\frac{dx}{\cos x}=\int\frac{\cos x\,dx}{1-\sin^2x}$; with the substitution $t=\sin x$:
  \[ \int\frac{dt}{1-t^2}=\frac12\ln\left|\frac{1+t}{1-t}\right|+C=\frac12\ln\left|\frac{1+\sin x}{1-\sin x}\right|+C=\ln\left|\frac{1+\sin x}{\cos x}\right|+C, \]
  (since $\frac{1+\sin x}{1-\sin x}=\frac{(1+\sin x)^2}{\cos^2x}$). Hence $2I=\frac{\sin x}{\cos^2x}+\ln\left|\frac{1+\sin x}{\cos x}\right|+C$, i.e.,
  \[ \boxed{\int\frac{dx}{\cos^3x}=\frac{\sin x}{2\cos^2x}+\frac12\ln\left|\frac{1}{\cos x}+\tan x\right|+C.} \]
  (Check: differentiating the result gives back $\frac1{\cos^3x}$.)

  \item On $[0,1]$ we have $0\le x\le1<\frac\pi2$, hence $\cos x>0$ and $f(x)>0$ is continuous. The area is
  \[ S=\int_0^1\frac{dx}{\cos^3x}=\left[\frac{\sin x}{2\cos^2x}+\frac12\ln\left(\frac{1+\sin x}{\cos x}\right)\right]_0^1
      =\boxed{\frac{\sin1}{2\cos^21}+\frac12\ln\left(\frac{1+\sin1}{\cos1}\right)}\approx2.054, \]
  since at the point $0$ the expression equals $0+\frac12\ln1=0$.
\end{enumerate}
\end{solution}

\subsection*{Question 6a}

\begin{question}
Sketch the region enclosed between the curves $x+y=0$ and $y=2x-x^2$ and find its area.
\end{question}

\begin{hint}
Find the intersection points of the line $y=-x$ with the parabola, and determine which curve lies above the other.
\end{hint}

\begin{solution}
The curve $x+y=0$ is the line $y=-x$, and the curve $y=2x-x^2$ is a downward-opening parabola with vertex $(1,1)$ that crosses the $x$-axis at $0$ and at $2$.
\textbf{Intersection points:} $-x=2x-x^2\iff x^2-3x=0\iff x=0$ or $x=3$, i.e., the points $(0,0)$ and $(3,-3)$.
On the interval $(0,3)$: $(2x-x^2)-(-x)=3x-x^2=x(3-x)>0$, i.e., the parabola is above the line. The region lies between the line (below) and the arc of the parabola (above) for $0\le x\le3$. The area:
\[ S=\int_0^3\big[(2x-x^2)-(-x)\big]dx=\int_0^3(3x-x^2)\,dx=\left[\frac{3x^2}{2}-\frac{x^3}{3}\right]_0^3=\frac{27}{2}-9=\boxed{\frac92}. \]
\end{solution}

\subsection*{Question 6b}

\begin{question}
Find the domain, the intersection points with the axes, the intervals of increase and decrease, the extremum points, the asymptotes, and the intervals of convexity/concavity of the following function:
\[ f(x)=\frac{\sin(x)}{2+\cos(x)}. \]
Sketch the graph of the function
\end{question}

\begin{hint}
The function is periodic with period $2\pi$ and odd, so it suffices to analyze it on $[0,2\pi]$ (or $[-\pi,\pi]$). Note that $2+\cos x\ge1>0$.
\end{hint}

\begin{solution}
\textbf{Domain:} $2+\cos x\ge1>0$ for every $x$, so $f$ is defined (and continuous and differentiable) on all of $\R$.
\textbf{Symmetry and periodicity:} $f(-x)=-f(x)$ (odd), and $f(x+2\pi)=f(x)$ (periodic with period $2\pi$).

\textbf{Intersection with the axes:} $f(x)=0\iff\sin x=0\iff x=k\pi$, $k\in\Z$. The intersection with the $y$-axis: $(0,0)$. $f>0$ on $(2k\pi,(2k+1)\pi)$ and $f<0$ on $((2k-1)\pi,2k\pi)$.

\textbf{First derivative:}
\[ f'(x)=\frac{\cos x(2+\cos x)-\sin x\cdot(-\sin x)}{(2+\cos x)^2}=\frac{2\cos x+\cos^2x+\sin^2x}{(2+\cos x)^2}=\frac{2\cos x+1}{(2+\cos x)^2}. \]
$f'(x)=0\iff\cos x=-\frac12\iff x=\pm\frac{2\pi}{3}+2k\pi$. Over the period $[-\pi,\pi]$:
\begin{itemize}
  \item $f'>0$ (\textbf{increasing}) for $\cos x>-\frac12$, i.e., on $\left(-\frac{2\pi}3+2k\pi,\ \frac{2\pi}3+2k\pi\right)$;
  \item $f'<0$ (\textbf{decreasing}) on $\left(\frac{2\pi}3+2k\pi,\ \frac{4\pi}3+2k\pi\right)$.
\end{itemize}
\textbf{Extrema:} A maximum at $x=\frac{2\pi}3+2k\pi$ with $f=\frac{\sqrt3/2}{2-1/2}=\frac{1}{\sqrt3}$; a minimum at $x=-\frac{2\pi}3+2k\pi$ with $f=-\frac1{\sqrt3}$. These are also the global extrema: $-\frac1{\sqrt3}\le f(x)\le\frac1{\sqrt3}$ for every $x$.

\textbf{Asymptotes:} There are no vertical asymptotes ($f$ is continuous on all of $\R$). There is no horizontal or oblique asymptote: $f$ is periodic and not constant, so $\lim_{x\to\pm\infty}f(x)$ does not exist (for example $f(2k\pi)=0$ and $f(\frac{2\pi}3+2k\pi)=\frac1{\sqrt3}$), and also $\frac{f(x)}{x}\to0$, so an oblique asymptote would have to be horizontal.

\textbf{Second derivative:} Differentiate $f'=\frac{2\cos x+1}{(2+\cos x)^2}$:
\[ f''(x)=\frac{-2\sin x(2+\cos x)^2+(2\cos x+1)\cdot2(2+\cos x)\sin x}{(2+\cos x)^4}=\frac{2\sin x\,\big[-(2+\cos x)+2\cos x+1\big]}{(2+\cos x)^3}=\frac{2\sin x(\cos x-1)}{(2+\cos x)^3}. \]
Since $\cos x-1\le0$ (and equals $0$ only at $x=2k\pi$) and the denominator is positive, the sign of $f''$ is opposite to the sign of $\sin x$:
\begin{itemize}
  \item On $(2k\pi,(2k+1)\pi)$: $\sin x>0$, hence $f''<0$ -- \textbf{concave} ($\cap$);
  \item On $((2k+1)\pi,(2k+2)\pi)$: $\sin x<0$, hence $f''>0$ -- \textbf{convex} ($\cup$).
\end{itemize}
\textbf{Inflection points:} $x=k\pi$ (the points $(k\pi,0)$), since $f''$ changes sign there.

\textbf{Sketch:} A periodic sine-like wave that passes through the origin, rises to the maximum $\frac1{\sqrt3}\approx0.577$ at $x=\frac{2\pi}{3}$, descends through $(\pi,0)$ (an inflection point) to the minimum $-\frac1{\sqrt3}$ at $x=\frac{4\pi}3$, and returns to $(2\pi,0)$; compared with $\sin x$ the graph is ``tilted to the right'' -- the maximum is shifted from $\frac\pi2$ to $\frac{2\pi}3$. The pattern repeats in every period of length $2\pi$.
\end{solution}

\section*{2016/17 Semester A Moed A}

\subsection*{Question 1}

\begin{question}
Use Taylor's formula to compute $\sqrt[3]{129}$ with an error not greater than $10^{-3}$.
Explain how you estimated the error.
\end{question}

\begin{hint}
$129$ is close to $125=5^3$. Expand $f(x)=\sqrt[3]{x}$ around $x_0=125$.
\end{hint}

\begin{hint}
Estimate the remainder in Lagrange form $R_n(x)=\frac{f^{(n+1)}(c)}{(n+1)!}(x-x_0)^{n+1}$ with $125<c<129$, and check which $n$ suffices.
\end{hint}

\begin{solution}
Define $f(x)=\sqrt[3]{x}=x^{1/3}$ and expand around $x_0=125$ (since $\sqrt[3]{125}=5$ is known exactly), with $x=129$, i.e., $x-x_0=4$.

\textbf{Derivatives:}
\[
f'(x)=\frac13 x^{-2/3},\qquad f''(x)=-\frac29 x^{-5/3}.
\]
Hence $f(125)=5$, $f'(125)=\frac13\cdot\frac{1}{25}=\frac{1}{75}$.

\textbf{Taylor polynomial of order 1:}
\[
P_1(129)=f(125)+f'(125)\cdot 4=5+\frac{4}{75}=\frac{379}{75}\approx 5.05333.
\]

\textbf{Error estimate.} By Taylor's theorem with the Lagrange remainder, there exists a point $c$ between $125$ and $129$ such that
\[
R_1(129)=\frac{f''(c)}{2!}\cdot 4^2=-\frac{2}{9}c^{-5/3}\cdot\frac{16}{2}=-\frac{16}{9\,c^{5/3}}.
\]
Since $c>125$, we have $c^{5/3}>125^{5/3}=5^5=3125$, and therefore
\[
|R_1(129)|<\frac{16}{9\cdot 3125}=\frac{16}{28125}\approx 5.7\cdot 10^{-4}<10^{-3}.
\]
Hence the Taylor polynomial of order 1 suffices.

\textbf{Answer:}
\[
\sqrt[3]{129}\approx 5+\frac{4}{75}=\frac{379}{75}\approx 5.0533,
\]
with an error smaller than $5.7\cdot10^{-4}<10^{-3}$. (The remainder is negative, i.e., the approximation is larger than the true value; indeed $\sqrt[3]{129}\approx 5.05277$.)
\end{solution}

\subsection*{Question 2a}

\begin{question}
Find the following limit:
$\displaystyle \lim_{x\to 0}\frac{e^x-1}{\sqrt{1+x}-1}$
\end{question}

\begin{solution}
As $x\to0$ the numerator and the denominator tend to $0$ (both functions are continuous at $0$ and vanish there), i.e., a limit of the form $\frac00$. The numerator and the denominator are differentiable in a neighborhood of $0$, and the derivative of the denominator $\frac{1}{2\sqrt{1+x}}\neq0$ there. By L'Hôpital's rule:
\[
\lim_{x\to0}\frac{e^x-1}{\sqrt{1+x}-1}=\lim_{x\to0}\frac{e^x}{\frac{1}{2\sqrt{1+x}}}=\lim_{x\to0}2\sqrt{1+x}\,e^x=2.
\]
(Alternative way: multiply by the conjugate: $\frac{e^x-1}{\sqrt{1+x}-1}=\frac{e^x-1}{x}\cdot(\sqrt{1+x}+1)\to 1\cdot2=2$, by the fundamental limit $\lim_{x\to0}\frac{e^x-1}{x}=1$.)

\textbf{Answer:} $2$.
\end{solution}

\subsection*{Question 2b}

\begin{question}
Find the following limit:
$\displaystyle \lim_{x\to 0}\frac{\sqrt[3]{8+3x}-2}{\sqrt[4]{16+5x}-2}$
\end{question}

\begin{solution}
At the point $x=0$: $\sqrt[3]{8}-2=0$ and $\sqrt[4]{16}-2=0$, i.e., a limit of the form $\frac00$. Both functions are differentiable in a neighborhood of $0$:
\[
\left(\sqrt[3]{8+3x}\right)'=\frac13(8+3x)^{-2/3}\cdot3=(8+3x)^{-2/3},\qquad
\left(\sqrt[4]{16+5x}\right)'=\frac54(16+5x)^{-3/4},
\]
and the derivative of the denominator is different from $0$ in a neighborhood of $0$. By L'Hôpital's rule:
\[
\lim_{x\to0}\frac{\sqrt[3]{8+3x}-2}{\sqrt[4]{16+5x}-2}
=\lim_{x\to0}\frac{(8+3x)^{-2/3}}{\frac54(16+5x)^{-3/4}}
=\frac{8^{-2/3}}{\frac54\cdot16^{-3/4}}=\frac{\frac14}{\frac54\cdot\frac18}=\frac{\frac14}{\frac{5}{32}}=\frac85.
\]
\textbf{Answer:} $\frac85$.
\end{solution}

\subsection*{Question 3}

\begin{question}
The function
\[
f(x)=\frac{2-(e^x+e^{-x})\cos(x)}{x^4}
\]
is defined for $x\neq0$. Is it possible to define $f(0)$ so that the function is continuous at the point $x=0$?
\end{question}

\begin{hint}
One must check whether $\lim_{x\to0}f(x)$ exists and is finite.
\end{hint}

\begin{hint}
Expand $e^x+e^{-x}$ and $\cos x$ into Maclaurin polynomials up to order $4$ and multiply.
\end{hint}

\begin{solution}
It is possible to define $f(0)$ so that $f$ is continuous at $0$ if and only if the limit $\lim_{x\to0}f(x)$ exists and is finite, and then $f(0)$ must be defined to be the limit.

\textbf{Maclaurin expansion.} By the Maclaurin expansions of $e^x$ and of $\cos x$:
\[
e^x+e^{-x}=2\left(1+\frac{x^2}{2}+\frac{x^4}{24}\right)+o(x^5)=2+x^2+\frac{x^4}{12}+o(x^5),
\]
\[
\cos x=1-\frac{x^2}{2}+\frac{x^4}{24}+o(x^5).
\]
Multiply and keep the terms up to order $4$:
\[
(e^x+e^{-x})\cos x=2+x^2+\frac{x^4}{12}-x^2-\frac{x^4}{2}+\frac{x^4}{12}+o(x^4)
=2+\left(\frac{1}{12}-\frac12+\frac1{12}\right)x^4+o(x^4)=2-\frac{x^4}{3}+o(x^4).
\]
Hence
\[
2-(e^x+e^{-x})\cos x=\frac{x^4}{3}+o(x^4),
\qquad\text{and therefore}\qquad
f(x)=\frac13+\frac{o(x^4)}{x^4}\xrightarrow[x\to0]{}\frac13.
\]

(Alternative way: four applications of L'Hôpital's rule to $\frac{0}{0}$; the fourth derivative of the numerator at $0$ is $8$ and that of the denominator is $24$, hence the limit is $\frac{8}{24}=\frac13$.)

\textbf{Answer:} Yes. The limit exists and equals $\frac13$, so if we define $f(0)=\frac13$ the function will be continuous at $x=0$.
\end{solution}

\subsection*{Question 4a}

\begin{question}
Analyze the function $y=\ln(e^x+e^{-2x})$ and sketch its graph.
\underline{You need to find}: the domain, asymptotes, intervals of increase and decrease, extremum points, intervals of convexity and concavity.
\end{question}

\begin{hint}
For the asymptotes: write $\ln(e^x+e^{-2x})=x+\ln(1+e^{-3x})=-2x+\ln(1+e^{3x})$.
\end{hint}

\begin{solution}
\textbf{Domain:} $e^x+e^{-2x}>0$ for every $x$, so the function is defined (and continuous and differentiable) on all of $\R$.

\textbf{Asymptotes:} There are no vertical asymptotes, since the function is continuous on all of $\R$.
For $x\to+\infty$:
\[
y=\ln\big(e^x(1+e^{-3x})\big)=x+\ln(1+e^{-3x}),\qquad y-x=\ln(1+e^{-3x})\xrightarrow[x\to+\infty]{}0,
\]
hence $y=x$ is an oblique asymptote at $+\infty$.
For $x\to-\infty$:
\[
y=\ln\big(e^{-2x}(1+e^{3x})\big)=-2x+\ln(1+e^{3x}),\qquad y+2x\xrightarrow[x\to-\infty]{}0,
\]
hence $y=-2x$ is an oblique asymptote at $-\infty$. There are no horizontal asymptotes. In addition $y-x>0$ and $y+2x>0$, i.e., the graph lies above both asymptotes.

\textbf{Increase and decrease:}
\[
y'=\frac{e^x-2e^{-2x}}{e^x+e^{-2x}}.
\]
The denominator is positive, and $e^x-2e^{-2x}>0\iff e^{3x}>2\iff x>\frac{\ln2}{3}$.
Hence $y$ is decreasing on $\left(-\infty,\frac{\ln 2}{3}\right)$ and increasing on $\left(\frac{\ln2}{3},\infty\right)$.

\textbf{Extrema:} At $x_0=\frac{\ln2}{3}$ there is a (global) minimum. There $e^{x_0}=2^{1/3}$, $e^{-2x_0}=2^{-2/3}$, and therefore
\[
y(x_0)=\ln\left(2^{1/3}+2^{-2/3}\right)=\ln\left(3\cdot2^{-2/3}\right)=\ln3-\frac23\ln2\approx0.64.
\]

\textbf{Convexity:} By the quotient rule,
\[
y''=\frac{(e^x+4e^{-2x})(e^x+e^{-2x})-(e^x-2e^{-2x})^2}{(e^x+e^{-2x})^2}
=\frac{9e^{-x}}{(e^x+e^{-2x})^2}>0,
\]
(in the numerator: $e^{2x}+5e^{-x}+4e^{-4x}-e^{2x}+4e^{-x}-4e^{-4x}=9e^{-x}$). Hence the function is convex ($\cup$) on all of $\R$, and there are no inflection points.

\textbf{Sketch:} The graph comes down from above along the asymptote $y=-2x$ (above it), reaches its minimum at the point $\left(\frac{\ln2}{3},\ \ln3-\frac23\ln2\right)\approx(0.23,\,0.64)$, and rises along the asymptote $y=x$ (above it); it is convex throughout. Intersection with the $y$-axis: $y(0)=\ln2$. There is no intersection with the $x$-axis (since $y\ge y(x_0)>0$).
\end{solution}

\subsection*{Question 4b}

\begin{question}
Compute the length of the curve $y=\frac25x\sqrt[4]{x}-\frac23\sqrt[4]{x^3}$ between its intersection points with the $x$-axis.
\end{question}

\begin{hint}
Write $y=\frac25x^{5/4}-\frac23x^{3/4}$ and check that $1+(y')^2$ is a perfect square.
\end{hint}

\begin{solution}
Write $y=\frac25x^{5/4}-\frac23x^{3/4}$, which is defined for $x\ge0$.

\textbf{Intersection with the $x$-axis:} $y=x^{3/4}\left(\frac25x^{1/2}-\frac23\right)=0\iff x=0$ or $\sqrt{x}=\frac53$, i.e., $x=0$ and $x=\frac{25}{9}$.

\textbf{Derivative:}
\[
y'=\frac25\cdot\frac54x^{1/4}-\frac23\cdot\frac34x^{-1/4}=\frac12\left(x^{1/4}-x^{-1/4}\right),
\]
\[
1+(y')^2=1+\frac14\left(x^{1/2}-2+x^{-1/2}\right)=\frac14\left(x^{1/2}+2+x^{-1/2}\right)=\left(\frac{x^{1/4}+x^{-1/4}}{2}\right)^2.
\]
\textbf{Length of the curve} (by the formula $L=\int_a^b\sqrt{1+(y')^2}\,dx$; the integral is improper at $0$ but converges, since $x^{-1/4}$ is integrable there):
\[
L=\int_0^{25/9}\frac{x^{1/4}+x^{-1/4}}{2}\,dx=\frac12\left[\frac45x^{5/4}+\frac43x^{3/4}\right]_0^{25/9}.
\]
For $x=\frac{25}{9}$: $x^{1/4}=\sqrt{\frac53}$, $x^{5/4}=\frac{25}{9}\sqrt{\frac53}$, $x^{3/4}=\frac53\sqrt{\frac53}$. Hence
\[
L=\frac12\left(\frac{20}{9}+\frac{20}{9}\right)\sqrt{\frac53}=\frac{20}{9}\sqrt{\frac53}=\frac{20\sqrt{15}}{27}\approx2.869.
\]
\end{solution}

\subsection*{Question 5a}

\begin{question}
Compute the area bounded by the lines $y=4x,\ \ y=x,\ \ y=1/x$ in the region $x\ge0,\ y\ge0$ and draw the corresponding sketch.
\end{question}

\begin{hint}
Find the intersection points of the lines, and split the integral into two intervals according to the upper curve.
\end{hint}

\begin{solution}
\textbf{Intersection points} (in the first quadrant): $4x=x\Rightarrow x=0$ (the origin); $4x=\frac1x\Rightarrow x=\frac12$ (the point $(\frac12,2)$); $x=\frac1x\Rightarrow x=1$ (the point $(1,1)$).

\textbf{Sketch:} The region is a ``curvilinear triangle'' with vertices $(0,0)$, $(\frac12,2)$, $(1,1)$: the lower boundary is the line $y=x$, and the upper boundary is the line $y=4x$ for $0\le x\le\frac12$ and the hyperbola $y=\frac1x$ for $\frac12\le x\le1$.

\textbf{Area:}
\[
S=\int_0^{1/2}(4x-x)\,dx+\int_{1/2}^1\left(\frac1x-x\right)dx
=\left[\frac{3x^2}{2}\right]_0^{1/2}+\left[\ln x-\frac{x^2}{2}\right]_{1/2}^1
=\frac38+\left(-\frac12-\ln\frac12+\frac18\right)=\frac38+\ln2-\frac38=\ln 2.
\]
\textbf{Answer:} $S=\ln2$.
\end{solution}

\subsection*{Question 5b}

\begin{question}
Sketch the curve $r=2\cos\varphi-1$ (given in polar coordinates) and compute the area of the region bounded by the curve.
\end{question}

\begin{hint}
Find for which angles $r\ge0$, and use the formula $S=\frac12\int_\alpha^\beta r^2\,d\varphi$.
\end{hint}

\begin{solution}
\textbf{Range of angles.} $r=2\cos\varphi-1\ge0\iff\cos\varphi\ge\frac12\iff -\frac\pi3\le\varphi\le\frac\pi3$ (within one period). At the endpoints $\varphi=\pm\frac\pi3$ we have $r=0$ (the curve passes through the origin), and for $\varphi=0$ we have $r=1$ (the point $(1,0)$). The curve is symmetric with respect to the $x$-axis (since $\cos$ is even).

\textbf{Sketch:} As $\varphi$ increases from $-\frac\pi3$ to $\frac\pi3$, the point leaves the origin in the direction of the angle $-\frac\pi3$, reaches $(1,0)$ and returns to the origin in the direction of the angle $\frac\pi3$: a single closed ``teardrop''-shaped loop in the half-plane $x\ge0$, symmetric with respect to the $x$-axis.

\textbf{Area} (by the formula for area in polar coordinates $S=\frac12\int_\alpha^\beta r^2\,d\varphi$):
\[
(2\cos\varphi-1)^2=4\cos^2\varphi-4\cos\varphi+1=3+2\cos2\varphi-4\cos\varphi,
\]
\[
S=\frac12\int_{-\pi/3}^{\pi/3}\left(3+2\cos2\varphi-4\cos\varphi\right)d\varphi
=\frac12\Big[3\varphi+\sin2\varphi-4\sin\varphi\Big]_{-\pi/3}^{\pi/3}
=\frac12\left(2\pi+\sqrt3-4\sqrt3\right)=\pi-\frac{3\sqrt3}{2}\approx0.544.
\]
\textbf{Answer:} $S=\pi-\frac{3\sqrt3}{2}$.

\textbf{Remark:} If negative values of $r$ are allowed (the point $(r,\varphi)$ with $r<0$ is plotted in the direction $\varphi+\pi$ at distance $|r|$), then for $\frac\pi3<\varphi<\frac{5\pi}3$ a large outer loop is obtained, and the whole curve is a limaçon (Pascal's snail) with an inner loop -- the loop computed above. In this case the area bounded by the outer loop is
\[
\frac12\int_{\pi/3}^{5\pi/3}(2\cos\varphi-1)^2\,d\varphi=2\pi+\frac{3\sqrt3}{2},
\]
and the area between the two loops is $\left(2\pi+\frac{3\sqrt3}{2}\right)-\left(\pi-\frac{3\sqrt3}{2}\right)=\pi+3\sqrt3$.
\end{solution}

\subsection*{Question 6a}

\begin{question}
Compute the following integral:
$\displaystyle\int\frac{dx}{x^2+x+1}$
\end{question}

\begin{hint}
Complete the square: $x^2+x+1=\left(x+\frac12\right)^2+\frac34$.
\end{hint}

\begin{solution}
Complete the square: $x^2+x+1=\left(x+\frac12\right)^2+\frac34$. Substitute $t=\frac{2x+1}{\sqrt3}$, i.e., $x+\frac12=\frac{\sqrt3}{2}t$, $dx=\frac{\sqrt3}{2}dt$:
\[
\int\frac{dx}{\left(x+\frac12\right)^2+\frac34}=\int\frac{\frac{\sqrt3}{2}\,dt}{\frac34(t^2+1)}=\frac{2}{\sqrt3}\int\frac{dt}{1+t^2}=\frac{2}{\sqrt3}\arctan t+C.
\]
\textbf{Answer:}
\[
\int\frac{dx}{x^2+x+1}=\frac{2}{\sqrt3}\arctan\frac{2x+1}{\sqrt3}+C.
\]
(Check: $\frac{d}{dx}\frac{2}{\sqrt3}\arctan\frac{2x+1}{\sqrt3}=\frac{2}{\sqrt3}\cdot\frac{2/\sqrt3}{1+\frac{(2x+1)^2}{3}}=\frac{4}{3+(2x+1)^2}=\frac{1}{x^2+x+1}$.)
\end{solution}

\subsection*{Question 6b}

\begin{question}
Compute the following integral:
$\displaystyle\int_1^\infty x^2e^{-2x}\,dx$
\end{question}

\begin{hint}
Integrate by parts twice (differentiating $x^2$), and then compute the limit as the upper limit of integration tends to infinity.
\end{hint}

\begin{solution}
\textbf{Antiderivative.} By integration by parts with $u=x^2$, $dv=e^{-2x}dx$, $v=-\frac12e^{-2x}$:
\[
\int x^2e^{-2x}dx=-\frac{x^2}{2}e^{-2x}+\int xe^{-2x}dx.
\]
Again by parts with $u=x$:
\[
\int xe^{-2x}dx=-\frac x2e^{-2x}+\frac12\int e^{-2x}dx=-\frac x2e^{-2x}-\frac14e^{-2x}.
\]
Hence
\[
F(x)=\int x^2e^{-2x}dx=-e^{-2x}\left(\frac{x^2}{2}+\frac x2+\frac14\right)+C.
\]
\textbf{The improper integral.} By definition,
\[
\int_1^\infty x^2e^{-2x}dx=\lim_{R\to\infty}\big(F(R)-F(1)\big).
\]
We have $\lim_{R\to\infty}e^{-2R}\left(\frac{R^2}{2}+\frac R2+\frac14\right)=0$ (the exponential dominates the polynomial; for example by L'Hôpital's rule applied twice to $\frac{R^2}{e^{2R}}$), and $F(1)=-e^{-2}\left(\frac12+\frac12+\frac14\right)=-\frac54e^{-2}$. Hence
\[
\int_1^\infty x^2e^{-2x}dx=0+\frac54e^{-2}=\frac{5}{4e^2}\approx0.169.
\]
The integral converges, and its value is $\frac{5}{4e^2}$.
\end{solution}

\section*{2016/17 Semester A Moed B}

\subsection*{Question 1}

\begin{question}
Use Taylor's formula to compute $\sin 167^\circ$ with an error not greater than $10^{-5}$. Explain how you estimated the error.
\end{question}

\begin{hint}
$\sin167^\circ=\sin(180^\circ-13^\circ)=\sin13^\circ$. Convert to radians and expand $\sin x$ around $0$.
\end{hint}

\begin{hint}
In the Maclaurin polynomial of $\sin x$ the coefficient of $x^4$ is $0$, hence $P_3=P_4$ and the remainder can be estimated using $R_4$.
\end{hint}

\begin{solution}
\textbf{Reduction.} $\sin167^\circ=\sin(180^\circ-13^\circ)=\sin13^\circ=\sin x_0$, where
\[
x_0=\frac{13\pi}{180}\approx0.226893 \quad\text{(radians)}.
\]

\textbf{Maclaurin polynomial.} For $f(x)=\sin x$: $f(0)=0$, $f'(0)=1$, $f''(0)=0$, $f'''(0)=-1$, $f^{(4)}(0)=0$, hence
\[
P_4(x)=P_3(x)=x-\frac{x^3}{6}.
\]

\textbf{Error estimate.} By Taylor's theorem with the Lagrange remainder of order $4$, there exists $c$ between $0$ and $x_0$ such that
\[
\sin x_0-P_4(x_0)=R_4(x_0)=\frac{f^{(5)}(c)}{5!}x_0^5=\frac{\cos c}{120}x_0^5 .
\]
Since $|\cos c|\le1$ and $0<x_0<0.23$:
\[
|R_4(x_0)|\le\frac{x_0^5}{120}<\frac{0.23^5}{120}\approx\frac{6.44\cdot10^{-4}}{120}\approx5.4\cdot10^{-6}<10^{-5}.
\]
(Even with the cruder estimate $x_0<0.25$ we get $\frac{0.25^5}{120}\approx8.1\cdot10^{-6}<10^{-5}$.)

\textbf{Computation:}
\[
\sin167^\circ\approx x_0-\frac{x_0^3}{6}\approx0.226893-\frac{0.011680}{6}\approx0.226893-0.001947=0.224946.
\]

\textbf{Answer:} $\sin167^\circ\approx\frac{13\pi}{180}-\frac16\left(\frac{13\pi}{180}\right)^3\approx0.22495$, with an error smaller than $10^{-5}$ (the exact value is $0.2249511\ldots$).
\end{solution}

\subsection*{Question 2a}

\begin{question}
Find the following limit:
$\displaystyle\lim_{x\to0}\frac{\sqrt{\cos x}-\sqrt[3]{\cos x}}{\sin^2x}$
\end{question}

\begin{hint}
Use $\cos x=1-\frac{x^2}{2}+o(x^2)$ and the approximation $(1+t)^\alpha=1+\alpha t+o(t)$.
\end{hint}

\begin{solution}
This is a limit of the form $\frac00$. By the Maclaurin expansion, $\cos x=1+t$ with $t=-\frac{x^2}{2}+o(x^2)\to0$. By $(1+t)^\alpha=1+\alpha t+o(t)$:
\[
\sqrt{\cos x}=1-\frac{x^2}{4}+o(x^2),\qquad \sqrt[3]{\cos x}=1-\frac{x^2}{6}+o(x^2),
\]
and therefore
\[
\sqrt{\cos x}-\sqrt[3]{\cos x}=-\frac{x^2}{4}+\frac{x^2}{6}+o(x^2)=-\frac{x^2}{12}+o(x^2).
\]
In addition, $\sin^2x=x^2+o(x^2)$ (since $\frac{\sin x}{x}\to1$). Hence
\[
\lim_{x\to0}\frac{\sqrt{\cos x}-\sqrt[3]{\cos x}}{\sin^2x}=\lim_{x\to0}\frac{-\frac{1}{12}+\frac{o(x^2)}{x^2}}{\left(\frac{\sin x}{x}\right)^2}=-\frac1{12}.
\]
(Alternative way: denote $u=\cos x$; then $\sin^2x=1-u^2$ and $u\to1$, and by L'Hôpital's rule $\lim_{u\to1}\frac{u^{1/2}-u^{1/3}}{1-u^2}$ equals $\frac{\frac12-\frac13}{-2}=-\frac1{12}$.)

\textbf{Answer:} $-\frac{1}{12}$.
\end{solution}

\subsection*{Question 2b}

\begin{question}
Find the following limit:
$\displaystyle\lim_{x\to0}\left(\frac{1+x}{2+x}\right)^{\frac{1-\sqrt{x}}{1-x}}$
\end{question}

\begin{hint}
Check separately where the base and the exponent tend. Is this an indeterminate form at all?
\end{hint}

\begin{solution}
The expression is defined only for $x\ge0$ (because of $\sqrt x$), so this is a one-sided limit $x\to0^+$.
This is \emph{not} an indeterminate form: the base $\frac{1+x}{2+x}\to\frac12>0$ and the exponent $\frac{1-\sqrt x}{1-x}\to1$. Write
\[
\left(\frac{1+x}{2+x}\right)^{\frac{1-\sqrt{x}}{1-x}}=\exp\left(\frac{1-\sqrt{x}}{1-x}\cdot\ln\frac{1+x}{2+x}\right).
\]
By the arithmetic of limits and the continuity of $\ln$, the exponent tends to $1\cdot\ln\frac12=\ln\frac12$, and by the continuity of the exponential function:
\[
\lim_{x\to0^+}\left(\frac{1+x}{2+x}\right)^{\frac{1-\sqrt{x}}{1-x}}=e^{\ln\frac12}=\frac12.
\]
\textbf{Answer:} $\frac12$.
\end{solution}

\subsection*{Question 3a}

\begin{question}
Is there a value $c$ such that
\[
f(x)=\begin{cases}\dfrac1x & x\ge1\\[1ex] 2x+c & x<1\end{cases}
\]
is continuous at $x=1$?
\end{question}

\begin{solution}
$f$ is continuous at $x=1$ if and only if $\lim_{x\to1^-}f(x)=\lim_{x\to1^+}f(x)=f(1)$.
\[
f(1)=\frac11=1,\qquad \lim_{x\to1^+}\frac1x=1,\qquad \lim_{x\to1^-}(2x+c)=2+c.
\]
Hence we need $2+c=1$, i.e., $c=-1$.

\textbf{Answer:} Yes, for $c=-1$ (and only for it).
\end{solution}

\subsection*{Question 3b}

\begin{question}
Does the equation $\cos(x)-x=0$ have a solution in the interval $(-1,1)$, and is it unique? State precisely every theorem you rely on.
\end{question}

\begin{hint}
Define $g(x)=\cos x-x$ and check the signs of $g(\pm1)$. For uniqueness, check the sign of $g'$.
\end{hint}

\begin{solution}
Define $g(x)=\cos x-x$. This is a continuous and differentiable function on all of $\R$.

\textbf{Existence.} \emph{The Intermediate Value Theorem (Bolzano):} If $g$ is continuous on a closed interval $[a,b]$ and $g(a),g(b)$ have opposite signs, then there exists $x_0\in(a,b)$ with $g(x_0)=0$.
Here $g$ is continuous on $[-1,1]$, and
\[
g(-1)=\cos1+1>0,\qquad g(1)=\cos1-1<0
\]
(since $0<\cos1<1$). Hence there exists $x_0\in(-1,1)$ with $\cos x_0-x_0=0$.

\textbf{Uniqueness.} $g'(x)=-\sin x-1$. We have $g'(x)\le0$ for every $x$, with equality $g'(x)=0$ only when $\sin x=-1$, i.e., $x=-\frac\pi2+2\pi k$; none of these points lies in the interval $(-1,1)$ (since $\frac\pi2>1$). Hence $g'(x)<0$ for every $x\in(-1,1)$.
By a corollary of Lagrange's theorem (\emph{if $g'<0$ on an interval then $g$ is strictly decreasing on it}), $g$ is strictly decreasing on $(-1,1)$ and therefore one-to-one there, so it has at most one zero.
(Equivalently: if there were two roots $x_1<x_2$, then by \emph{Rolle's theorem} -- $g$ is continuous on $[x_1,x_2]$, differentiable on $(x_1,x_2)$ and $g(x_1)=g(x_2)$ -- there would exist $\xi\in(x_1,x_2)$ with $g'(\xi)=0$, a contradiction.)

\textbf{Answer:} Yes, the equation has a solution in the interval $(-1,1)$, and it is unique.
\end{solution}

\subsection*{Question 4a}

\begin{question}
Analyze the function $f(x)=x^2e^{-x}$ and sketch its graph.
\underline{You need to find}: the domain, asymptotes; intervals of increase and decrease, extremum points, intervals of convexity and concavity.
\end{question}

\begin{solution}
\textbf{Domain:} all of $\R$; the function is continuous and differentiable at every point, and $f(x)\ge0$ with equality only at $x=0$.

\textbf{Asymptotes:} There are no vertical ones (continuous on $\R$).
As $x\to+\infty$: $\lim_{x\to\infty}\frac{x^2}{e^x}=0$ (L'Hôpital twice), hence $y=0$ is a horizontal asymptote at $+\infty$.
As $x\to-\infty$: $f(x)\to+\infty$ and $\frac{f(x)}{x}=xe^{-x}\to-\infty$, so there is no asymptote (horizontal or oblique) at $-\infty$.

\textbf{Increase and decrease:}
\[
f'(x)=2xe^{-x}-x^2e^{-x}=x(2-x)e^{-x}.
\]
$f'<0$ on $(-\infty,0)$ and on $(2,\infty)$ (decreasing), and $f'>0$ on $(0,2)$ (increasing).

\textbf{Extrema:} A (global) minimum at $(0,0)$; a local maximum at $\left(2,\frac4{e^2}\right)\approx(2,0.54)$ (not global, since $f\to\infty$ at $-\infty$).

\textbf{Convexity:}
\[
f''(x)=\big((2-2x)-(2x-x^2)\big)e^{-x}=(x^2-4x+2)e^{-x},
\]
which vanishes at $x=2\pm\sqrt2$. Hence $f$ is convex ($\cup$) on $(-\infty,2-\sqrt2)$ and on $(2+\sqrt2,\infty)$, and concave ($\cap$) on $(2-\sqrt2,2+\sqrt2)$. Inflection points at $x=2\pm\sqrt2$ (approximately $x\approx0.59$ and $x\approx3.41$).

\textbf{Sketch:} The graph decreases from $+\infty$ (at $-\infty$) to the minimum $(0,0)$, where it is tangent to the $x$-axis, rises to the maximum $(2,4e^{-2})$, and then decreases and tends to $0$ from above ($y=0$ is an asymptote at $+\infty$).
\end{solution}

\subsection*{Question 4b}

\begin{question}
Draw the graph of the function $r=f(\theta)=\theta^2$ in polar coordinates for $0\le\theta\le\pi$ and compute the length of the curve.
\end{question}

\begin{hint}
Use the formula for the arc length of a curve in polar coordinates $L=\int_\alpha^\beta\sqrt{r^2+(r')^2}\,d\theta$.
\end{hint}

\begin{solution}
\textbf{Sketch:} $r=\theta^2$ increases from $0$ to $\pi^2\approx9.87$ as $\theta$ increases from $0$ to $\pi$: a spiral that leaves the origin (tangent to the positive $x$-axis), passes through $\left(\theta=\frac\pi2,\ r=\frac{\pi^2}4\approx2.47\right)$ on the positive $y$-axis, and reaches the point $(-\pi^2,0)$ on the negative $x$-axis; the whole curve lies in the upper half-plane $y\ge0$.

\textbf{Length:} By the formula for arc length in polar coordinates, with $r'=2\theta$:
\[
L=\int_0^\pi\sqrt{r^2+(r')^2}\,d\theta=\int_0^\pi\sqrt{\theta^4+4\theta^2}\,d\theta=\int_0^\pi\theta\sqrt{\theta^2+4}\,d\theta
\]
(since $\theta\ge0$). With the substitution $u=\theta^2+4$, $du=2\theta\,d\theta$:
\[
L=\frac12\int_4^{\pi^2+4}\sqrt u\,du=\frac13\Big[u^{3/2}\Big]_4^{\pi^2+4}=\frac13\left((\pi^2+4)^{3/2}-8\right)\approx14.55.
\]
\end{solution}

\subsection*{Question 5a}

\begin{question}
Compute the integral $\displaystyle\int\frac{x^3}{\sqrt{1-x^2}}\,dx$.
\end{question}

\begin{hint}
Substitute $u=1-x^2$ (and then $x^2=1-u$).
\end{hint}

\begin{solution}
Substitute $u=1-x^2$, $du=-2x\,dx$, and $x^2=1-u$ (for $|x|<1$):
\[
\int\frac{x^3}{\sqrt{1-x^2}}dx=\int\frac{x^2\cdot x\,dx}{\sqrt{1-x^2}}=-\frac12\int\frac{1-u}{\sqrt u}\,du=-\frac12\int\left(u^{-1/2}-u^{1/2}\right)du
=-\sqrt u+\frac13u^{3/2}+C.
\]
\textbf{Answer:}
\[
\int\frac{x^3}{\sqrt{1-x^2}}dx=-\sqrt{1-x^2}+\frac13(1-x^2)^{3/2}+C=-\frac13(x^2+2)\sqrt{1-x^2}+C.
\]
(Check: differentiating $-\frac13(x^2+2)\sqrt{1-x^2}$ gives $-\frac23x\sqrt{1-x^2}+\frac{x(x^2+2)}{3\sqrt{1-x^2}}=\frac{-2x(1-x^2)+x^3+2x}{3\sqrt{1-x^2}}=\frac{x^3}{\sqrt{1-x^2}}$.)
\end{solution}

\subsection*{Question 5b}

\begin{question}
Compute the improper integral $\displaystyle\int_1^\infty\frac{\arctan(x)}{x^2}\,dx$.
\end{question}

\begin{hint}
Integrate by parts with $u=\arctan x$, $dv=\frac{dx}{x^2}$, and decompose $\frac{1}{x(1+x^2)}=\frac1x-\frac{x}{1+x^2}$.
\end{hint}

\begin{solution}
\textbf{Antiderivative.} Integration by parts with $u=\arctan x$, $dv=\frac{dx}{x^2}$, $du=\frac{dx}{1+x^2}$, $v=-\frac1x$:
\[
\int\frac{\arctan x}{x^2}dx=-\frac{\arctan x}{x}+\int\frac{dx}{x(1+x^2)}.
\]
Partial fraction decomposition: $\frac{1}{x(1+x^2)}=\frac1x-\frac{x}{1+x^2}$, hence for $x>0$
\[
\int\frac{dx}{x(1+x^2)}=\ln x-\frac12\ln(1+x^2)=\ln\frac{x}{\sqrt{1+x^2}}.
\]
\textbf{The improper integral.} By definition, $\int_1^\infty=\lim_{R\to\infty}\int_1^R$:
\[
\int_1^R\frac{\arctan x}{x^2}dx=\left[-\frac{\arctan x}{x}+\ln\frac{x}{\sqrt{1+x^2}}\right]_1^R
=-\frac{\arctan R}{R}+\ln\frac{R}{\sqrt{1+R^2}}+\frac\pi4-\ln\frac1{\sqrt2}.
\]
As $R\to\infty$: $\frac{\arctan R}{R}\to0$ (bounded numerator, denominator tending to infinity), and $\frac{R}{\sqrt{1+R^2}}\to1$, hence $\ln\frac{R}{\sqrt{1+R^2}}\to0$ (continuity of $\ln$). Therefore
\[
\int_1^\infty\frac{\arctan x}{x^2}dx=\frac\pi4+\frac12\ln2\approx1.132.
\]
\end{solution}

\subsection*{Question 6}

\begin{question}
Compute the volume of the solid of revolution obtained by rotating the region $D$ around the $x$-axis, where $D$ is the region between the graph of the function $f(x)=\dfrac{x}{x^2+6x+8}$ and the $x$-axis, for $0\le x\le2$.
\end{question}

\begin{hint}
The volume is $V=\pi\int_0^2f^2(x)\,dx$. First decompose $f$ into partial fractions: $x^2+6x+8=(x+2)(x+4)$.
\end{hint}

\begin{solution}
\textbf{The volume formula.} By the formula for the volume of a solid of revolution around the $x$-axis:
\[
V=\pi\int_0^2 f^2(x)\,dx=\pi\int_0^2\frac{x^2}{(x^2+6x+8)^2}\,dx.
\]

\textbf{Partial fractions.} $x^2+6x+8=(x+2)(x+4)$, and
\[
\frac{x}{(x+2)(x+4)}=\frac{A}{x+2}+\frac{B}{x+4},\qquad A=\frac{-2}{-2+4}=-1,\quad B=\frac{-4}{-4+2}=2.
\]
Hence $f(x)=\frac{2}{x+4}-\frac{1}{x+2}$, and
\[
f^2(x)=\frac{4}{(x+4)^2}+\frac{1}{(x+2)^2}-\frac{4}{(x+2)(x+4)},\qquad
\frac{4}{(x+2)(x+4)}=\frac{2}{x+2}-\frac{2}{x+4}.
\]

\textbf{Integration:}
\[
\int_0^2\frac{4\,dx}{(x+4)^2}=\left[-\frac4{x+4}\right]_0^2=-\frac46+1=\frac13,\qquad
\int_0^2\frac{dx}{(x+2)^2}=\left[-\frac1{x+2}\right]_0^2=-\frac14+\frac12=\frac14,
\]
\[
\int_0^2\left(\frac{2}{x+2}-\frac{2}{x+4}\right)dx=\Big[2\ln(x+2)-2\ln(x+4)\Big]_0^2=2\ln\frac46-2\ln\frac24=2\ln\frac43.
\]
Hence
\[
\int_0^2f^2(x)\,dx=\frac13+\frac14-2\ln\frac43=\frac7{12}-2\ln\frac43.
\]

\textbf{Answer:}
\[
V=\pi\left(\frac{7}{12}-2\ln\frac43\right)\approx0.0250.
\]
(The volume is very small, which is reasonable: $0\le f(x)\le f(2)=\frac1{12}$ on the interval.)
\end{solution}

\section*{2016/17 Semester B Moed A}

\subsection*{Question 1}

\begin{question}
Compute the derivative of the function $f(x)=\ln(2x)$ at the point $x$ ($x>0$), from the definition (without using known derivatives).
\end{question}

\begin{hint}
Use the laws of logarithms to write the difference quotient as $\frac{\ln(1+\frac hx)}{h}$, and then the fundamental limit $\lim_{t\to0}\frac{\ln(1+t)}{t}=1$.
\end{hint}

\begin{solution}
Let $x>0$. For $h\ne0$ with $|h|<x$ (so that $x+h>0$), by the laws of logarithms:
\[ \frac{f(x+h)-f(x)}{h}=\frac{\ln(2x+2h)-\ln(2x)}{h}=\frac{1}{h}\ln\frac{2(x+h)}{2x}=\frac1h\ln\left(1+\frac hx\right). \]
Substitute $t=\frac hx$; as $h\to0$ also $t\to0$ and $t\ne0$, so by the fundamental limit $\lim_{t\to0}\frac{\ln(1+t)}{t}=1$ (and the substitution rule for limits):
\[ \frac1h\ln\left(1+\frac hx\right)=\frac1x\cdot\frac{\ln(1+t)}{t}\xrightarrow[h\to0]{}\frac1x\cdot1. \]
Hence
\[ \boxed{f'(x)=\lim_{h\to0}\frac{f(x+h)-f(x)}{h}=\frac1x}. \]
(Note that the result is identical to the derivative of $\ln x$, since $\ln(2x)=\ln2+\ln x$.)
\end{solution}

\subsection*{Question 2}

\begin{question}
Does the function $f(x)=5-x-x^2-|x+3|$ have a maximum on the interval $[-10,10]$? If so, find it.
\end{question}

\begin{hint}
$f$ is continuous on a closed interval -- which theorem guarantees the existence of a maximum? To find it, open the absolute value and check each of the two ranges separately.
\end{hint}

\begin{solution}
\textbf{Existence:} $f$ is continuous on $[-10,10]$ (a sum of a polynomial and $-|x+3|$, which is continuous). By Weierstrass's theorem, a continuous function on a closed bounded interval attains a maximum (and a minimum) on it. Hence \textbf{yes}, there is a maximum.

\textbf{Finding it:} We open the absolute value.
\begin{itemize}
  \item For $-3\le x\le10$: $|x+3|=x+3$, hence $f(x)=5-x-x^2-x-3=2-2x-x^2=3-(x+1)^2$. This is a downward parabola with vertex at $x=-1\in[-3,10]$, so in this range $f(x)\le3$, with equality only at $x=-1$.
  \item For $-10\le x<-3$: $|x+3|=-(x+3)$, hence $f(x)=5-x-x^2+x+3=8-x^2$. Here $x^2>9$, so $f(x)<8-9=-1<3$.
\end{itemize}
(Alternatively: $f$ is differentiable except at $x=-3$; $f'(x)=-2-2x$ for $x>-3$ vanishes at $x=-1$, and $f'(x)=-2x>0$ for $x<-3$; compare the values of $f$ at the candidate points $-10,-3,-1,10$: $f(-10)=-92$, $f(-3)=-1$, $f(-1)=3$, $f(10)=-118$.)

Hence the maximum of $f$ on the interval is $\boxed{f(-1)=3}$, and it is attained at the point $x=-1$.
\end{solution}

\subsection*{Question 3}

\begin{question}
Compute:
\[ \int x^3e^{-x^2}\,dx \]
\end{question}

\begin{hint}
Substitute $t=x^2$ (or $t=-x^2$) and then integrate by parts.
\end{hint}

\begin{solution}
Substitute $t=x^2$, $dt=2x\,dx$. Then $x^3\,dx=x^2\cdot x\,dx=\frac t2\,dt$:
\[ \int x^3e^{-x^2}dx=\frac12\int te^{-t}dt. \]
Integration by parts with $u=t$, $v'=e^{-t}$ ($u'=1$, $v=-e^{-t}$):
\[ \int te^{-t}dt=-te^{-t}+\int e^{-t}dt=-te^{-t}-e^{-t}+C=-(t+1)e^{-t}+C. \]
Return to the variable $x$:
\[ \boxed{\int x^3e^{-x^2}dx=-\frac12(x^2+1)e^{-x^2}+C.} \]
\textbf{Check:} $\frac{d}{dx}\left[-\frac12(x^2+1)e^{-x^2}\right]=-xe^{-x^2}+\frac12(x^2+1)\cdot2xe^{-x^2}=x^3e^{-x^2}$.
\end{solution}

\subsection*{Question 4}

\begin{question}
Compute:
\[ \int\frac{x^4+1}{x^3-x^2-x+1}\,dx \]
\end{question}

\begin{hint}
The degree of the numerator is greater than the degree of the denominator: first perform polynomial division. Factor the denominator: $x^3-x^2-x+1=(x-1)^2(x+1)$.
\end{hint}

\begin{hint}
The double root $x=1$ contributes two partial fractions: $\frac{B}{x-1}+\frac{C}{(x-1)^2}$.
\end{hint}

\begin{solution}
\textbf{Factoring the denominator:} $x^3-x^2-x+1=x^2(x-1)-(x-1)=(x-1)(x^2-1)=(x-1)^2(x+1)$.

\textbf{Polynomial division:} The quotient is $x+1$; indeed $(x+1)(x^3-x^2-x+1)=x^4-2x^2+1$, and therefore
\[ x^4+1=(x+1)(x^3-x^2-x+1)+2x^2,\qquad \frac{x^4+1}{x^3-x^2-x+1}=x+1+\frac{2x^2}{(x-1)^2(x+1)}. \]

\textbf{Partial fractions:}
\[ \frac{2x^2}{(x-1)^2(x+1)}=\frac{A}{x+1}+\frac{B}{x-1}+\frac{C}{(x-1)^2}\iff 2x^2=A(x-1)^2+B(x-1)(x+1)+C(x+1). \]
Substituting $x=-1$: $2=4A$, i.e., $A=\frac12$. Substituting $x=1$: $2=2C$, i.e., $C=1$. Comparing the coefficients of $x^2$: $2=A+B$, hence $B=\frac32$.

\textbf{Integration:}
\[ \int\frac{x^4+1}{x^3-x^2-x+1}dx=\int\left(x+1+\frac{1/2}{x+1}+\frac{3/2}{x-1}+\frac{1}{(x-1)^2}\right)dx, \]
\[ \boxed{=\frac{x^2}{2}+x+\frac12\ln|x+1|+\frac32\ln|x-1|-\frac{1}{x-1}+C.} \]
(Check: differentiating the result gives back the integrand.)
\end{solution}

\subsection*{Question 5}

\begin{question}
Define: $f(x)=\ln(1+x^2)$. Compute $f^{(4)}(0)$.
\end{question}

\begin{hint}
Instead of differentiating four times, use the Maclaurin polynomial of $\ln(1+t)$ with $t=x^2$, and compare coefficients: the coefficient of $x^4$ is $\frac{f^{(4)}(0)}{4!}$.
\end{hint}

\begin{solution}
From the Maclaurin expansion $\ln(1+t)=t-\frac{t^2}{2}+o(t^2)$, with $t=x^2$:
\[ \ln(1+x^2)=x^2-\frac{x^4}{2}+o(x^4). \]
$f$ is infinitely differentiable, and by the uniqueness of the Taylor polynomial (the only polynomial of degree $\le4$ approximating $f$ up to $o(x^4)$ is the Taylor polynomial), the Taylor polynomial of order $4$ of $f$ around $0$ is $x^2-\frac{x^4}{2}$. The coefficient of $x^4$ in the Taylor polynomial is $\frac{f^{(4)}(0)}{4!}$, and therefore
\[ \frac{f^{(4)}(0)}{24}=-\frac12\quad\Longrightarrow\quad\boxed{f^{(4)}(0)=-12}. \]
(Check by direct differentiation: $f'(x)=\frac{2x}{1+x^2}$, $f''(x)=\frac{2-2x^2}{(1+x^2)^2}$, $f'''(x)=\frac{4x^3-12x}{(1+x^2)^3}$, $f^{(4)}(x)=\frac{-12x^4+72x^2-12}{(1+x^2)^4}$, and $f^{(4)}(0)=-12$.)
\end{solution}

\subsection*{Question 6}

\begin{question}
Let $f$ be a one-to-one differentiable function. Given: $f(1)=2$, $f(2)=4$, $f(3)=5$, $f(4)=6$, $f'(1)=0$, $f'(2)=2$, $f'(3)=0$, $f'(4)=3$.
Compute: $\left(f^{-1}\right)'(4)$
\end{question}

\begin{hint}
By the theorem on the derivative of the inverse function, $(f^{-1})'(y)=\frac{1}{f'(f^{-1}(y))}$. What is $f^{-1}(4)$?
\end{hint}

\begin{solution}
Since $f(2)=4$ and $f$ is one-to-one, $f^{-1}(4)=2$.
By the theorem on the derivative of the inverse function: if $f$ is differentiable at $x_0$, $f'(x_0)\ne0$ (and $f$ is continuous and one-to-one in a neighborhood), then $f^{-1}$ is differentiable at $y_0=f(x_0)$ and
\[ (f^{-1})'(y_0)=\frac{1}{f'(x_0)}. \]
Here $x_0=2$, $y_0=4$, and $f'(2)=2\ne0$, hence
\[ \boxed{(f^{-1})'(4)=\frac{1}{f'(2)}=\frac12}. \]
(The rest of the data -- $f'(4)=3$, for example -- are distractors: $f'$ must be evaluated at the point $f^{-1}(4)=2$ and not at the point $4$.)
\end{solution}

\subsection*{Question 7}

\begin{question}
Does the improper integral
\[ \int_0^1\frac{\sin x}{x^2}\,dx \]
converge?
\end{question}

\begin{hint}
Near $0$ we have $\sin x\approx x$. Compare (limit comparison test) with $\frac1x$.
\end{hint}

\begin{solution}
\textbf{The integral diverges.}
The integrand is continuous on $(0,1]$, and the only problem is at the point $0$. On $(0,1]$ we have $\sin x>0$ (since $0<x\le1<\pi$), so the integrand is positive and the comparison tests may be used.

Compare with $g(x)=\frac1x>0$:
\[ \lim_{x\to0^+}\frac{\sin x/x^2}{1/x}=\lim_{x\to0^+}\frac{\sin x}{x}=1\in(0,\infty). \]
By the limit comparison test, $\int_0^1\frac{\sin x}{x^2}dx$ and $\int_0^1\frac{dx}{x}$ either both converge or both diverge. But
\[ \int_\eps^1\frac{dx}{x}=-\ln\eps\xrightarrow[\eps\to0^+]{}\infty, \]
i.e., $\int_0^1\frac{dx}{x}$ diverges. Hence $\int_0^1\frac{\sin x}{x^2}dx$ also \textbf{diverges}.

(Direct comparison: since $\frac{\sin x}{x}\to1$, there exists $\delta>0$ such that $\frac{\sin x}{x}>\frac12$ on $(0,\delta)$, and hence $\frac{\sin x}{x^2}>\frac{1}{2x}$ there.)
\end{solution}

\subsection*{Question 8}

\begin{question}
Let $f(x),g(x)$ be functions defined and continuous on the ray $[0,+\infty)$. It is given that the improper integral $\int_0^\infty(f(x)+g(x))\,dx$ converges.
Prove or disprove:
The improper integral $\int_0^\infty(f(x)-g(x))\,dx$ converges.
\end{question}

\begin{hint}
Try functions whose sum is zero.
\end{hint}

\begin{solution}
\textbf{The statement is false.} Counterexample: $f(x)=1$, $g(x)=-1$ (both continuous on $[0,\infty)$).
\begin{itemize}
  \item $f+g\equiv0$, hence $\int_0^\infty(f+g)\,dx=\lim_{R\to\infty}\int_0^R0\,dx=0$ -- converges.
  \item $f-g\equiv2$, hence $\int_0^R(f-g)\,dx=2R\xrightarrow[R\to\infty]{}\infty$ -- diverges.
\end{itemize}
Hence convergence of $\int_0^\infty(f+g)$ does not imply convergence of $\int_0^\infty(f-g)$.

(Remark: if \emph{both} integrals $\int_0^\infty f$ and $\int_0^\infty g$ converged, then by linearity $\int_0^\infty(f-g)$ would converge too. Here only the convergence of the sum is given.)
\end{solution}

\section*{2016/17 Semester B Moed B}

\subsection*{Question 1}

\begin{question}
Find a rational number $x=\frac ab$ such that $|\cos(1)-x|<\frac{1}{1000}$.
\end{question}

\begin{hint}
Use the Maclaurin polynomial of $\cos$ and Lagrange's remainder formula, $|R_n(1)|\le\frac{1}{(n+1)!}$. Choose $n$ such that $(n+1)!>1000$.
\end{hint}

\begin{hint}
Note: in the expansion of $\cos$ the coefficient of $x^{2k+1}$ is zero, hence $P_{2k}=P_{2k+1}$ and the remainder of order $2k+1$ can be used.
\end{hint}

\begin{solution}
The Maclaurin polynomial of $\cos x$ has rational coefficients, so its value at $x=1$ is rational. It remains to choose an order for which the remainder is smaller than $\frac1{1000}$.

By Taylor's theorem with the Lagrange remainder, for every $n$ there exists $c$ between $0$ and $1$ such that
\[ \cos1=P_n(1)+R_n(1),\qquad R_n(1)=\frac{\cos^{(n+1)}(c)}{(n+1)!}\cdot1^{n+1}. \]
Every derivative of $\cos$ is $\pm\sin$ or $\pm\cos$, hence $|\cos^{(n+1)}(c)|\le1$ and $|R_n(1)|\le\frac1{(n+1)!}$.

Choose $n=7$: $|R_7(1)|\le\frac1{8!}=\frac1{40320}<\frac1{1000}$. Since the coefficient of $x^7$ in the expansion of $\cos$ is $0$:
\[ P_7(1)=P_6(1)=1-\frac1{2!}+\frac1{4!}-\frac1{6!}=1-\frac12+\frac1{24}-\frac1{720}=\frac{720-360+30-1}{720}=\frac{389}{720}. \]
Hence
\[ \left|\cos1-\frac{389}{720}\right|\le\frac1{40320}<\frac1{1000}, \]
and the rational number $\boxed{x=\frac{389}{720}}$ satisfies the requirement. (Numerically $\cos1\approx0.5403023$, $\frac{389}{720}\approx0.5402778$.)

\textbf{Remark:} The previous approximation $1-\frac12+\frac1{24}=\frac{13}{24}$ is \emph{not} sufficient: the bound given by Lagrange is $\frac1{6!}=\frac1{720}>\frac1{1000}$, and indeed $|\cos1-\frac{13}{24}|\approx0.00136>\frac1{1000}$.
\end{solution}

\subsection*{Question 2}

\begin{question}
For which value of the parameter $A$ is the function $f(x)$ continuous at the point $0$?
\[ f(x)=\begin{cases}\dfrac{1-\cos^2(x)}{2x^2} & x\neq0\\ A & x=0\end{cases} \]
\end{question}

\begin{solution}
$f$ is continuous at $0$ if and only if $\lim_{x\to0}f(x)=f(0)=A$. By the identity $1-\cos^2x=\sin^2x$:
\[ \lim_{x\to0}\frac{1-\cos^2x}{2x^2}=\lim_{x\to0}\frac12\left(\frac{\sin x}{x}\right)^2=\frac12\cdot1^2=\frac12, \]
by the fundamental limit $\frac{\sin x}{x}\to1$ and the arithmetic of limits. Hence $f$ is continuous at $0$ if and only if $\boxed{A=\frac12}$.
\end{solution}

\subsection*{Question 3}

\begin{question}
Compute:
\[ \int e^{\sqrt[3]{x}}\,dx \]
\end{question}

\begin{hint}
Substitute $t=\sqrt[3]x$, i.e., $x=t^3$, and then integrate by parts twice.
\end{hint}

\begin{solution}
Substitute $t=\sqrt[3]x$, i.e., $x=t^3$ and $dx=3t^2\,dt$ (the substitution is invertible and differentiable on all of $\R$):
\[ \int e^{\sqrt[3]x}dx=3\int t^2e^t\,dt. \]
Integration by parts ($u=t^2$, $v'=e^t$):
\[ \int t^2e^tdt=t^2e^t-2\int te^tdt, \]
and again by parts ($u=t$, $v'=e^t$): $\int te^tdt=te^t-e^t+C$. Hence
\[ \int t^2e^tdt=t^2e^t-2te^t+2e^t+C=e^t(t^2-2t+2)+C. \]
Return to $x$:
\[ \boxed{\int e^{\sqrt[3]x}dx=3e^{\sqrt[3]x}\left(\sqrt[3]{x^2}-2\sqrt[3]x+2\right)+C.} \]
\textbf{Check:} $\frac{d}{dt}\left[3e^t(t^2-2t+2)\right]=3e^t(t^2-2t+2+2t-2)=3t^2e^t$, which is exactly the integrand after the substitution.
\end{solution}

\subsection*{Question 4}

\begin{question}
Compute:
\[ \int\frac{x^2}{\sqrt{4-x^2}}\,dx \]
\end{question}

\begin{hint}
Trigonometric substitution: $x=2\sin t$ with $t\in(-\frac\pi2,\frac\pi2)$.
\end{hint}

\begin{solution}
The integrand is defined for $-2<x<2$. Substitute $x=2\sin t$, $t\in(-\frac\pi2,\frac\pi2)$ (an invertible substitution: $t=\arcsin\frac x2$), $dx=2\cos t\,dt$. On this interval $\cos t>0$, hence $\sqrt{4-x^2}=\sqrt{4\cos^2t}=2\cos t$. Therefore
\[ \int\frac{x^2}{\sqrt{4-x^2}}dx=\int\frac{4\sin^2t}{2\cos t}\cdot2\cos t\,dt=4\int\sin^2t\,dt=2\int(1-\cos2t)\,dt=2t-\sin2t+C. \]
Back to $x$: $t=\arcsin\frac x2$, and $\sin2t=2\sin t\cos t=2\cdot\frac x2\cdot\frac{\sqrt{4-x^2}}{2}=\frac{x\sqrt{4-x^2}}{2}$. Hence
\[ \boxed{\int\frac{x^2}{\sqrt{4-x^2}}dx=2\arcsin\frac x2-\frac{x\sqrt{4-x^2}}{2}+C.} \]
(Check: differentiating the result gives $\frac{2}{\sqrt{4-x^2}}-\frac{\sqrt{4-x^2}}{2}+\frac{x^2}{2\sqrt{4-x^2}}=\frac{4-(4-x^2)+x^2}{2\sqrt{4-x^2}}=\frac{x^2}{\sqrt{4-x^2}}$.)
\end{solution}

\subsection*{Question 5}

\begin{question}
What is the maximal area of an isosceles triangle whose perimeter is 1?
\end{question}

\begin{hint}
Denote the length of a leg by $a$ and the base by $b=1-2a$. The height to the base is $\sqrt{a^2-\frac{b^2}{4}}$. What is the allowed range of $a$?
\end{hint}

\begin{hint}
It is more convenient to maximize the square of the area.
\end{hint}

\begin{solution}
Denote the length of a leg by $a$ and the length of the base by $b$, with $2a+b=1$, i.e. $b=1-2a$. For a (non-degenerate) triangle we need $b>0$ and the triangle inequality $2a>b$, i.e. $\frac14<a<\frac12$.

The height to the base (which is also a median) is, by Pythagoras, $h=\sqrt{a^2-\frac{b^2}{4}}$, hence
\[ S=\frac12bh=\frac b4\sqrt{4a^2-b^2}=\frac{1-2a}{4}\sqrt{(2a-b)(2a+b)}=\frac{1-2a}{4}\sqrt{4a-1}, \]
since $2a+b=1$ and $2a-b=4a-1$.

We maximize $g(a)=16S^2=(1-2a)^2(4a-1)$ on $(\frac14,\frac12)$ (since $S>0$, $S$ and $S^2$ attain their maximum at the same point):
\[ g'(a)=-4(1-2a)(4a-1)+4(1-2a)^2=4(1-2a)\big[(1-2a)-(4a-1)\big]=4(1-2a)(2-6a). \]
On this interval $1-2a>0$, hence $g'(a)=0\iff a=\frac13$; $g'>0$ on $(\frac14,\frac13)$ and $g'<0$ on $(\frac13,\frac12)$. Therefore $a=\frac13$ is a point of absolute maximum on the interval (at the endpoints $g\to0$).

For $a=\frac13$: $b=\frac13$ -- the triangle is equilateral, and
\[ S_{\max}=\frac{1/3}{4}\sqrt{\frac13}=\frac{1}{12\sqrt3}=\boxed{\frac{\sqrt3}{36}}\approx0.0481. \]
\end{solution}

\subsection*{Question 6}

\begin{question}
Find the absolute minimum and maximum of the function $f(x)=\frac{\sqrt{x^2-1}}{x^2}$.
\end{question}

\begin{hint}
First find the domain, and note that $f$ is even. Also examine the limit at $\infty$ and the values of $f$ at the endpoints of the domain.
\end{hint}

\begin{solution}
\textbf{Domain:} $x^2-1\ge0$ and $x\ne0$, i.e. $|x|\ge1$. $f$ is even ($f(-x)=f(x)$), so it suffices to study $x\ge1$.

\textbf{Minimum:} $f(x)\ge0$ on the whole domain, and $f(\pm1)=0$. Hence the absolute minimum is $\boxed{0}$, attained at $x=\pm1$.

\textbf{Maximum:} For $x>1$:
\[ f'(x)=\frac{\frac{x}{\sqrt{x^2-1}}\cdot x^2-2x\sqrt{x^2-1}}{x^4}=\frac{x^2-2(x^2-1)}{x^3\sqrt{x^2-1}}=\frac{2-x^2}{x^3\sqrt{x^2-1}}. \]
Hence $f'>0$ on $(1,\sqrt2)$ and $f'<0$ on $(\sqrt2,\infty)$: $f$ is increasing on $[1,\sqrt2]$ and decreasing on $[\sqrt2,\infty)$ (using continuity at $1$). Therefore for every $x\ge1$ we have $f(x)\le f(\sqrt2)=\frac{\sqrt{2-1}}{2}=\frac12$. (In addition, $\lim_{x\to\infty}f(x)=\lim\frac{\sqrt{1-1/x^2}}{x}=0$.)
By evenness, the absolute maximum is $\boxed{\frac12}$, attained at $x=\pm\sqrt2$.
\end{solution}

\subsection*{Question 7}

\begin{question}
Does the improper integral
\[ \int_2^\infty\frac{x+7}{\sqrt{x^5-3x}}\,dx \]
converge?
\end{question}

\begin{hint}
For large $x$ the integrand behaves like $\frac{x}{x^{5/2}}=\frac{1}{x^{3/2}}$. Use the limit comparison test.
\end{hint}

\begin{solution}
\textbf{The integral converges.}
For $x\ge2$: $x^5-3x=x(x^4-3)\ge2\cdot13>0$, hence the integrand is defined, continuous and positive on $[2,\infty)$, and the only problem is at the endpoint $\infty$.

We compare with $g(x)=\frac{1}{x^{3/2}}>0$:
\[ \lim_{x\to\infty}\frac{\frac{x+7}{\sqrt{x^5-3x}}}{x^{-3/2}}=\lim_{x\to\infty}\frac{(x+7)x^{3/2}}{\sqrt{x^5-3x}}=\lim_{x\to\infty}\frac{x^{5/2}\left(1+\frac7x\right)}{x^{5/2}\sqrt{1-\frac{3}{x^4}}}=1\in(0,\infty). \]
By the limit comparison test, the two integrals $\int_2^\infty\frac{x+7}{\sqrt{x^5-3x}}dx$ and $\int_2^\infty\frac{dx}{x^{3/2}}$ either both converge or both diverge. The integral $\int_2^\infty\frac{dx}{x^p}$ converges for $p>1$, and here $p=\frac32$:
\[ \int_2^R x^{-3/2}dx=\left[-2x^{-1/2}\right]_2^R=\frac{2}{\sqrt2}-\frac{2}{\sqrt R}\xrightarrow[R\to\infty]{}\sqrt2. \]
Therefore the given integral \textbf{converges}.
\end{solution}

\subsection*{Question 8}

\begin{question}
Prove or disprove: if the function $|f(x)|$ is Riemann integrable on the interval $[a,b]$, then the function $f(x)$ is also integrable on this interval.
\end{question}

\begin{hint}
Look for a function that takes only the values $1$ and $-1$ in a ``dense'' way, similar to the Dirichlet function.
\end{hint}

\begin{solution}
\textbf{The statement is false.} (The converse is true: if $f$ is integrable then so is $|f|$.)

A counterexample (for $a<b$):
\[ f(x)=\begin{cases}1 & x\in\Q\\ -1 & x\notin\Q\end{cases}\qquad x\in[a,b]. \]
Then $|f(x)|=1$ for every $x$, a constant function, hence integrable, and $\int_a^b|f|=b-a$.

On the other hand, $f$ is not Riemann integrable: let $P$ be any partition of $[a,b]$. Every subinterval $[x_{i-1},x_i]$ (of positive length) contains both a rational number and an irrational number (density of $\Q$ and of $\R\setminus\Q$ in $\R$), hence $M_i=\sup f=1$ and $m_i=\inf f=-1$. Therefore
\[ U(f,P)=\sum M_i\Delta x_i=b-a,\qquad L(f,P)=\sum m_i\Delta x_i=-(b-a), \]
so that $U(f,P)-L(f,P)=2(b-a)$ for every partition, and it cannot be made smaller than an arbitrary $\eps$. By Riemann's criterion (or: the upper integral $b-a$ differs from the lower integral $-(b-a)$), $f$ is \textbf{not} integrable.
\end{solution}

\section*{2017/18 Semester A Moed A}

\subsection*{Question 1a}

\begin{question}
Compute the following limit:
the limit of the sequence $\displaystyle \lim_{n\to\infty}\left(\frac{n}{n+1}\right)^{3n}$.
\end{question}

\begin{hint}
Invert the fraction: $\frac{n}{n+1}=\frac{1}{1+\frac1n}$, and use the famous limit $\left(1+\frac1n\right)^n\to e$.
\end{hint}

\begin{solution}
For every $n$:
\[ \left(\frac{n}{n+1}\right)^{3n}=\frac{1}{\left(\frac{n+1}{n}\right)^{3n}}=\frac{1}{\left[\left(1+\frac1n\right)^{n}\right]^{3}}. \]
It is known that $\left(1+\frac1n\right)^n\xrightarrow[n\to\infty]{} e$. By the arithmetic of limits (a product of three convergent sequences, and a quotient where the limit of the denominator is $e^3\neq 0$) we get
\[ \lim_{n\to\infty}\left(\frac{n}{n+1}\right)^{3n}=\frac{1}{e^3}. \]
\end{solution}

\subsection*{Question 1b}

\begin{question}
Compute the following limit:
the limit of the function $\displaystyle \lim_{x\to\pi/2}(\tan x)\cdot\arctan\left(x-\frac{\pi}{2}\right)$.
\end{question}

\begin{hint}
Write $\tan x=\frac{1}{\cot x}$ and obtain an expression of the form $\frac00$.
\end{hint}

\begin{solution}
For $x$ close to $\frac\pi2$, $x\neq\frac\pi2$, we have $\tan x=\frac{1}{\cot x}$, hence
\[ (\tan x)\cdot\arctan\left(x-\frac{\pi}{2}\right)=\frac{\arctan\left(x-\frac{\pi}{2}\right)}{\cot x}. \]
As $x\to\frac\pi2$: the numerator tends to $\arctan 0=0$ and the denominator tends to $\cot\frac\pi2=0$ (both functions are continuous). This is an expression of the form $\frac00$, both functions are differentiable in a punctured neighborhood of $\frac\pi2$, and the derivative of the denominator $-\frac{1}{\sin^2x}\neq 0$ there, so L'Hôpital's rule can be applied:
\[ \lim_{x\to\pi/2}\frac{\arctan\left(x-\frac{\pi}{2}\right)}{\cot x}
   =\lim_{x\to\pi/2}\frac{\dfrac{1}{1+\left(x-\frac\pi2\right)^2}}{-\dfrac{1}{\sin^2x}}
   =\lim_{x\to\pi/2}\left(-\frac{\sin^2x}{1+\left(x-\frac\pi2\right)^2}\right)=-\frac{1}{1+0}=-1, \]
where the last limit is computed by continuity ($\sin^2\frac\pi2=1$). Since the limit of the quotient of the derivatives exists, by L'Hôpital's rule
\[ \lim_{x\to\pi/2}(\tan x)\cdot\arctan\left(x-\frac{\pi}{2}\right)=-1. \]
(Another way: $\arctan t\sim t$ as $t\to0$, and $(x-\frac\pi2)\tan x=-\frac{(x-\frac\pi2)\cos(x-\frac\pi2)}{\sin(x-\frac\pi2)}\to-1$.)
\end{solution}

\subsection*{Question 2}

\begin{question}
Find the points of discontinuity of the function
\[ f(x)=\frac{e^{\left(\frac1x-\frac{1}{x+1}\right)}}{e^{\frac1x}-2} \]
and determine their types.
\end{question}

\begin{hint}
The function is undefined at three points: where $x=0$, $x=-1$, and where $e^{1/x}=2$.
\end{hint}

\begin{hint}
For $x\to0^+$, divide the numerator and the denominator by $e^{1/x}$.
\end{hint}

\begin{solution}
\textbf{The domain.} The expression is defined when $x\neq0$, $x\neq-1$ and $e^{1/x}\neq 2$, i.e. $\frac1x\neq\ln2$, i.e. $x\neq\frac1{\ln2}$. At every other point $f$ is a composition and quotient of continuous elementary functions (with a nonzero denominator), hence continuous. The points of discontinuity are therefore $x=0,\ x=-1,\ x=\frac{1}{\ln2}$. We check each of them.

\textbf{1. $x=\frac1{\ln2}$.} The numerator is continuous at this point and its value is
$e^{\ln2-\frac{1}{\frac1{\ln2}+1}}>0$ (an exponential is always positive). The denominator $e^{1/x}-2$ tends to $0$. As $x\to\left(\frac1{\ln2}\right)^+$ we have $\frac1x<\ln2$, hence $e^{1/x}-2\to0^-$, and
\[ \lim_{x\to(1/\ln2)^+}f(x)=-\infty, \qquad \lim_{x\to(1/\ln2)^-}f(x)=+\infty \]
(from the left $\frac1x>\ln2$ and the denominator tends to $0^+$). The one-sided limits are not finite, hence this is a \textbf{discontinuity of the second kind}.

\textbf{2. $x=-1$.} As $x\to-1$, the denominator is continuous and tends to $e^{-1}-2<0$ (a nonzero negative number). We check the numerator:
\begin{itemize}
  \item As $x\to-1^+$: $x+1\to0^+$, hence $\frac1x-\frac1{x+1}\to-1-\infty=-\infty$, and therefore $e^{\left(\frac1x-\frac1{x+1}\right)}\to0$. Hence $\lim_{x\to-1^+}f(x)=\frac{0}{e^{-1}-2}=0$.
  \item As $x\to-1^-$: $x+1\to0^-$, hence $\frac1x-\frac1{x+1}\to+\infty$ and the numerator tends to $+\infty$. Since the denominator tends to a negative number, $\lim_{x\to-1^-}f(x)=-\infty$.
\end{itemize}
One of the one-sided limits is infinite, hence this is a \textbf{discontinuity of the second kind}.

\textbf{3. $x=0$.}
\begin{itemize}
  \item As $x\to0^+$: $\frac1x\to+\infty$, hence both the numerator and the denominator tend to $\infty$. Divide the numerator and the denominator by $e^{1/x}$:
  \[ f(x)=\frac{e^{\frac1x}\cdot e^{-\frac1{x+1}}}{e^{\frac1x}\left(1-2e^{-\frac1x}\right)}=\frac{e^{-\frac{1}{x+1}}}{1-2e^{-\frac1x}}\xrightarrow[x\to0^+]{}\frac{e^{-1}}{1-0}=\frac1e, \]
  since $e^{-1/x}\to0$ when $\frac1x\to+\infty$. (In the official solution the limit is computed with L'Hôpital's rule and gives the same value.)
  \item As $x\to0^-$: $\frac1x\to-\infty$, hence $e^{1/x}\to0$ and the denominator tends to $-2$; also $\frac1x-\frac1{x+1}\to-\infty-1=-\infty$, hence the numerator tends to $0$. Therefore $\lim_{x\to0^-}f(x)=\frac{0}{-2}=0$.
\end{itemize}
Both one-sided limits exist and are finite but different ($\frac1e\neq0$), hence this is a \textbf{discontinuity of the first kind (jump)}.

\textbf{Summary:} $x=0$ — discontinuity of the first kind (jump); $x=-1$ and $x=\frac1{\ln2}$ — discontinuities of the second kind.
\end{solution}

\subsection*{Question 3a}

\begin{question}
Compute the limit:
\[ \lim_{x\to0}\left(\frac{\sqrt{1-2x+x^3}-\sqrt[3]{1-3x+x^2}}{x^2}\right) \]
\underline{Hint}: you may use Maclaurin formulas.
\end{question}

\begin{hint}
Expand each root using $(1+t)^p=1+pt+\frac{p(p-1)}{2}t^2+o(t^2)$ up to order $x^2$.
\end{hint}

\begin{solution}
We use the Maclaurin expansion $(1+t)^p=1+pt+\frac{p(p-1)}{2!}t^2+o(t^2)$ as $t\to0$.

For the first root, $t=x^3-2x\to0$ and $p=\frac12$:
\begin{align*}
\sqrt{1-2x+x^3}&=1+\frac12(x^3-2x)-\frac18(x^3-2x)^2+o(x^2)\\
&=1-x+\frac{x^3}{2}-\frac18\left(4x^2-4x^4+x^6\right)+o(x^2)=1-x-\frac{x^2}{2}+o(x^2).
\end{align*}
(We used the fact that $t=O(x)$, hence $o(t^2)=o(x^2)$, and that the powers $x^3,x^4,x^6$ are $o(x^2)$.)

For the second root, $t=-3x+x^2$ and $p=\frac13$, $\frac{p(p-1)}{2}=\frac{\frac13\cdot(-\frac23)}{2}=-\frac19$:
\begin{align*}
\sqrt[3]{1-3x+x^2}&=1+\frac13(-3x+x^2)-\frac19(-3x+x^2)^2+o(x^2)\\
&=1-x+\frac{x^2}{3}-\frac19\left(9x^2-6x^3+x^4\right)+o(x^2)=1-x-\frac{2x^2}{3}+o(x^2).
\end{align*}
Hence
\[ \frac{\sqrt{1-2x+x^3}-\sqrt[3]{1-3x+x^2}}{x^2}=\frac{\left(-\frac12+\frac23\right)x^2+o(x^2)}{x^2}=\frac16+\frac{o(x^2)}{x^2}\xrightarrow[x\to0]{}\frac16. \]
\[ \boxed{\lim_{x\to0}\frac{\sqrt{1-2x+x^3}-\sqrt[3]{1-3x+x^2}}{x^2}=\frac16} \]
\end{solution}

\subsection*{Question 3b}

\begin{question}
Prove that $0<x<\arcsin x$ for every $x\in(0,1]$.
\end{question}

\begin{hint}
Consider the function $g(x)=\arcsin x-x$ and its derivative.
\end{hint}

\begin{solution}
The inequality $0<x$ is clear for $x\in(0,1]$. It remains to prove $x<\arcsin x$.

Define $g(x)=\arcsin x-x$ on the interval $[0,1]$. The function is continuous on $[0,1]$ and differentiable on $(0,1)$, with
\[ g'(x)=\frac{1}{\sqrt{1-x^2}}-1>0 \quad\text{for every } x\in(0,1), \]
since $0<\sqrt{1-x^2}<1$ there.

Let $b\in(0,1]$. By Lagrange's theorem for $g$ on the interval $[0,b]$, there exists a point $c\in(0,b)$ such that
\[ g(b)-g(0)=g'(c)\,(b-0). \]
Since $g(0)=\arcsin0-0=0$, $g'(c)>0$ and $b>0$, we get $g(b)>0$, i.e. $\arcsin b>b$.

Therefore $0<x<\arcsin x$ for every $x\in(0,1]$. $\blacksquare$

(The official solution is phrased by contradiction: assume $\arcsin b\le b$; Lagrange's theorem gives $\frac{1}{\sqrt{1-c^2}}=\frac{\arcsin b}{b}\le1$, contradicting $\frac{1}{\sqrt{1-c^2}}>1$. This is the same proof.)
\end{solution}

\subsection*{Question 4}

\begin{question}
Investigate the function $f(x)=(x+2)e^{1/x}$ and sketch its graph.
\end{question}

\begin{hint}
For the oblique asymptote: $a=\lim\frac{f(x)}{x}$, and then $b=\lim(f(x)-ax)$. Computing $b$ involves the limit $\lim_{x\to\infty}x\left(e^{1/x}-1\right)$.
\end{hint}

\begin{solution}
\textbf{1. General description.}
\begin{itemize}
  \item \textbf{Domain:} $x\neq0$.
  \item \textbf{Intersections with the axes:} there is no intersection with the $y$-axis since $0$ is not in the domain. $f(x)=0\iff x+2=0$ (since $e^{1/x}>0$), hence the intersection with the $x$-axis is $(-2,0)$.
  \item \textbf{Parity:} $f(-1)=e^{-1}$, $f(1)=3e$, hence $f$ is neither even nor odd. It is also not periodic.
  \item \textbf{Sign:} since $e^{1/x}>0$, the sign of $f$ is the sign of $x+2$: $f(x)>0$ for $x>-2$ ($x\neq0$), and $f(x)<0$ for $x<-2$.
\end{itemize}

\textbf{2. Continuity.} $f$ is elementary and hence continuous on its whole domain. At the point $x=0$:
\[ \lim_{x\to0^-}(x+2)e^{1/x}=2\cdot0=0,\qquad \lim_{x\to0^+}(x+2)e^{1/x}=2\cdot(+\infty)=+\infty, \]
since $\frac1x\to\mp\infty$. Hence at $x=0$ there is a discontinuity of the second kind.

\textbf{3. Asymptotes.}
\begin{itemize}
  \item \textbf{Vertical:} $x=0$ (from the right), since $\lim_{x\to0^+}f(x)=+\infty$. From the left the graph tends to the point $(0,0)$ (which is not on the graph).
  \item \textbf{Oblique} $y=ax+b$ as $x\to\pm\infty$:
  \[ a=\lim_{x\to\pm\infty}\frac{f(x)}{x}=\lim_{x\to\pm\infty}\left(1+\frac2x\right)e^{1/x}=1\cdot e^0=1. \]
  \[ b=\lim_{x\to\pm\infty}\left(f(x)-x\right)=\lim_{x\to\pm\infty}\left(x\left(e^{1/x}-1\right)+2e^{1/x}\right). \]
  Substitute $t=\frac1x\to0$: $x\left(e^{1/x}-1\right)=\frac{e^t-1}{t}\to1$ (a famous limit, or L'Hôpital: $\frac{e^t}{1}\to1$). Also $2e^{1/x}\to2$. Hence $b=1+2=3$.

  The oblique asymptote is $y=x+3$, both as $x\to\infty$ and as $x\to-\infty$. In particular $\lim_{x\to\pm\infty}f(x)=\pm\infty$.
\end{itemize}

\textbf{4. Increase, decrease and extrema.}
\[ f'(x)=e^{1/x}+(x+2)e^{1/x}\cdot\left(-\frac{1}{x^2}\right)=e^{1/x}\left(1-\frac{x+2}{x^2}\right)=e^{1/x}\cdot\frac{x^2-x-2}{x^2}=e^{1/x}\cdot\frac{(x-2)(x+1)}{x^2}. \]
Since $e^{1/x}>0$ and $x^2>0$, the sign of $f'$ is the sign of $(x-2)(x+1)$:
\begin{itemize}
  \item $f'>0$ on $(-\infty,-1)$ and on $(2,\infty)$ — $f$ is increasing there;
  \item $f'<0$ on $(-1,0)$ and on $(0,2)$ — $f$ is decreasing there.
\end{itemize}
Hence at $x=-1$ there is a local maximum, $f(-1)=1\cdot e^{-1}=\frac1e$, and at $x=2$ there is a local minimum, $f(2)=4e^{1/2}=4\sqrt e$.

\textbf{5. Convexity and inflection points.} Differentiate $f'(x)=e^{1/x}\cdot\frac{x^2-x-2}{x^2}=e^{1/x}\left(1-\frac1x-\frac2{x^2}\right)$:
\[ f''(x)=e^{1/x}\left(-\frac1{x^2}\right)\left(1-\frac1x-\frac2{x^2}\right)+e^{1/x}\left(\frac1{x^2}+\frac4{x^3}\right)
=e^{1/x}\cdot\frac{-x^2+x+2+x^2+4x}{x^4}=e^{1/x}\cdot\frac{5x+2}{x^4}. \]
The sign of $f''$ is the sign of $5x+2$. Hence $f$ is convex (upward, $f''>0$) on $\left(-\frac25,0\right)$ and on $(0,\infty)$, and concave (downward) on $\left(-\infty,-\frac25\right)$. At $x=-\frac25$ the second derivative changes sign and $f$ is continuous, hence this is an inflection point:
\[ \left(-\frac25,\ f\left(-\frac25\right)\right)=\left(-\frac25,\ \frac85e^{-5/2}\right). \]

\textbf{6. Sketch.} To the left of the $y$-axis: as $x\to-\infty$ the graph lies below the asymptote $y=x+3$ and approaches it (since $f(x)-(x+3)=\frac{5}{2x}+O\left(\frac1{x^2}\right)$); it is increasing and concave, crosses the $x$-axis at $(-2,0)$, reaches the local maximum $\left(-1,\frac1e\right)$, and then decreases; at the inflection point $x=-\frac25$ it becomes convex and keeps decreasing towards the point $(0,0)$ (which does not belong to the graph). To the right of the $y$-axis: the graph decreases from $+\infty$ (vertical asymptote $x=0$) to the minimum $(2,4\sqrt e)\approx(2,\,6.6)$, and then increases and approaches the asymptote $y=x+3$ from above; on this interval it is convex.
\end{solution}

\subsection*{Question 5a}

\begin{question}
Compute the integral $\displaystyle\int\frac{3x-1}{x^2-x+1}\,dx$.
\end{question}

\begin{hint}
Write the numerator as a combination of the derivative of the denominator $2x-1$ and a constant, and complete the denominator to a square.
\end{hint}

\begin{solution}
The denominator $x^2-x+1$ has no real roots (the discriminant is $1-4=-3<0$), hence $x^2-x+1>0$ for every $x$. We decompose the numerator according to the derivative of the denominator $(x^2-x+1)'=2x-1$:
\[ 3x-1=\frac32(2x-1)+\frac12. \]
Hence
\[ \int\frac{3x-1}{x^2-x+1}\,dx=\frac32\int\frac{2x-1}{x^2-x+1}\,dx+\frac12\int\frac{dx}{\left(x-\frac12\right)^2+\frac34}. \]
In the first integral substitute $u=x^2-x+1$, $du=(2x-1)\,dx$:
\[ \frac32\int\frac{du}{u}=\frac32\ln|u|=\frac32\ln\left(x^2-x+1\right). \]
In the second, substitute $v=x-\frac12$, $dv=dx$, and use $\int\frac{dv}{v^2+a^2}=\frac1a\arctan\frac va$ with $a=\frac{\sqrt3}{2}$:
\[ \frac12\int\frac{dv}{v^2+\frac34}=\frac12\cdot\frac{2}{\sqrt3}\arctan\frac{2v}{\sqrt3}=\frac{1}{\sqrt3}\arctan\frac{2x-1}{\sqrt3}. \]
In summary
\[ \boxed{\int\frac{3x-1}{x^2-x+1}\,dx=\frac32\ln\left(x^2-x+1\right)+\frac{1}{\sqrt3}\arctan\frac{2x-1}{\sqrt3}+C} \]
(No absolute value is needed inside the $\ln$ since $x^2-x+1>0$.) Check by differentiation: $\frac32\cdot\frac{2x-1}{x^2-x+1}+\frac1{\sqrt3}\cdot\frac{2/\sqrt3}{1+\frac{(2x-1)^2}{3}}=\frac{3x-\frac32}{x^2-x+1}+\frac{2}{4x^2-4x+4}=\frac{3x-1}{x^2-x+1}$.

\textbf{Remark:} in the official solution the coefficient of the $\arctan$ is written as $\frac{2}{\sqrt3}$; this is a mistake (the factor $\frac12$ was forgotten), and the correct coefficient is $\frac1{\sqrt3}$.
\end{solution}

\subsection*{Question 5b}

\begin{question}
Does the improper integral $\displaystyle\int_0^\infty\frac{x}{e^{x^2}}\,dx$ converge? If so, compute the value of the integral. If not, explain why.
\end{question}

\begin{hint}
Substitute $u=x^2$.
\end{hint}

\begin{solution}
By definition, $\int_0^\infty\frac{x}{e^{x^2}}\,dx=\lim_{b\to\infty}\int_0^b xe^{-x^2}\,dx$, if the limit exists. Substitute $u=x^2$, $du=2x\,dx$ (as $x$ runs from $0$ to $b$, $u$ runs from $0$ to $b^2$):
\[ \int_0^b xe^{-x^2}\,dx=\frac12\int_0^{b^2}e^{-u}\,du=\frac12\left[-e^{-u}\right]_0^{b^2}=\frac12\left(1-e^{-b^2}\right). \]
As $b\to\infty$, $e^{-b^2}\to0$, hence the limit exists and is finite:
\[ \boxed{\int_0^\infty\frac{x}{e^{x^2}}\,dx=\frac12} \]
that is, the integral converges and its value is $\frac12$.
\end{solution}

\subsection*{Question 6a}

\begin{question}
Compute the volume of the solid of revolution obtained by rotating the region $D$ about the $x$-axis, where $D$ is the region between the graph of the function $f(x)=\tan x$ and the $x$-axis, for $0\le x\le\frac\pi4$.
\end{question}

\begin{hint}
Use the identity $\tan^2x=\frac{1}{\cos^2x}-1$.
\end{hint}

\begin{solution}
The volume of a solid of revolution about the $x$-axis is $V=\pi\int_a^b f^2(x)\,dx$. Hence, using $\tan^2x=\frac1{\cos^2x}-1$:
\[ V=\pi\int_0^{\pi/4}\tan^2x\,dx=\pi\int_0^{\pi/4}\left(\frac1{\cos^2x}-1\right)dx=\pi\Big[\tan x-x\Big]_0^{\pi/4}=\pi\left(1-\frac\pi4\right). \]
\[ \boxed{V=\pi-\frac{\pi^2}{4}} \]
\end{solution}

\subsection*{Question 6b}

\begin{question}
Define $r=f(\theta)=\theta^2-1$ in polar coordinates. Compute the length of this curve between $\theta=1$ and $\theta=9$.

(The exam includes a figure of the curve — a spiral — on a polar grid.)
\end{question}

\begin{hint}
The arc length of a curve in polar coordinates is $\int_\alpha^\beta\sqrt{r^2+\left(\frac{dr}{d\theta}\right)^2}\,d\theta$. Check that the expression under the root is a perfect square.
\end{hint}

\begin{solution}
The arc length of a curve in polar coordinates $r=f(\theta)$, $\alpha\le\theta\le\beta$:
\[ L=\int_\alpha^\beta\sqrt{r^2+\left(\frac{dr}{d\theta}\right)^2}\,d\theta. \]
Here $\frac{dr}{d\theta}=2\theta$, hence
\[ r^2+\left(\frac{dr}{d\theta}\right)^2=(\theta^2-1)^2+4\theta^2=\theta^4-2\theta^2+1+4\theta^2=\theta^4+2\theta^2+1=(\theta^2+1)^2. \]
Since $\theta^2+1>0$, $\sqrt{(\theta^2+1)^2}=\theta^2+1$, and
\[ L=\int_1^9(\theta^2+1)\,d\theta=\left[\frac{\theta^3}{3}+\theta\right]_1^9=\left(243+9\right)-\left(\frac13+1\right)=252-\frac43=\frac{752}{3}. \]
\[ \boxed{L=\frac{752}{3}=250\tfrac23} \]
\textbf{Remark:} the official solution writes $252-\frac13+1=252\frac23$ — a sign error when substituting the lower limit; the correct value is $252-\left(\frac13+1\right)=250\frac23$.
\end{solution}

\section*{2017/18 Semester A Moed B}

\subsection*{Question 1a}

\begin{question}
Compute the following limit:
the limit of the sequence $\displaystyle\lim_{n\to\infty}\left(\sqrt{(n^2+2)(n^2-4)}-\sqrt{n^4-7}\right)$.
\end{question}

\begin{hint}
Multiply and divide by the conjugate $\sqrt{(n^2+2)(n^2-4)}+\sqrt{n^4-7}$.
\end{hint}

\begin{solution}
For every $n\ge2$ both roots are defined. Multiply and divide by the conjugate:
\[ \sqrt{(n^2+2)(n^2-4)}-\sqrt{n^4-7}=\frac{(n^2+2)(n^2-4)-(n^4-7)}{\sqrt{(n^2+2)(n^2-4)}+\sqrt{n^4-7}}. \]
In the numerator: $(n^2+2)(n^2-4)=n^4-2n^2-8$, hence the numerator is $n^4-2n^2-8-n^4+7=-2n^2-1$. Divide the numerator and the denominator by $n^2$:
\[ \frac{-2n^2-1}{\sqrt{(n^2+2)(n^2-4)}+\sqrt{n^4-7}}=\frac{-2-\frac1{n^2}}{\sqrt{\left(1+\frac2{n^2}\right)\left(1-\frac4{n^2}\right)}+\sqrt{1-\frac7{n^4}}}. \]
By the arithmetic of limits and the continuity of the square root, the numerator tends to $-2$ and the denominator to $1+1=2$. Hence
\[ \boxed{\lim_{n\to\infty}\left(\sqrt{(n^2+2)(n^2-4)}-\sqrt{n^4-7}\right)=-1} \]
\end{solution}

\subsection*{Question 1b}

\begin{question}
Compute the following limit:
the limit of the function $\displaystyle\lim_{x\to1^-}\frac{\frac\pi2-\arcsin x}{\sqrt[3]{1-x^2}}$.
\end{question}

\begin{hint}
This is an expression of the form $\frac00$; apply L'Hôpital's rule and simplify the quotient of the derivatives.
\end{hint}

\begin{solution}
As $x\to1^-$: $\arcsin x\to\frac\pi2$ (continuity), hence the numerator tends to $0$; the denominator $\sqrt[3]{1-x^2}\to0$ as well. This is an expression of the form $\frac00$. On the interval $(0,1)$ both functions are differentiable:
\[ \left(\frac\pi2-\arcsin x\right)'=-\frac{1}{\sqrt{1-x^2}},\qquad \left((1-x^2)^{1/3}\right)'=\frac13(1-x^2)^{-2/3}\cdot(-2x)=-\frac{2x}{3(1-x^2)^{2/3}}\neq0. \]
The quotient of the derivatives:
\[ \frac{-\frac{1}{(1-x^2)^{1/2}}}{-\frac{2x}{3(1-x^2)^{2/3}}}=\frac{3(1-x^2)^{2/3}}{2x(1-x^2)^{1/2}}=\frac{3(1-x^2)^{1/6}}{2x}\xrightarrow[x\to1^-]{}\frac{3\cdot0}{2}=0. \]
The limit of the quotient of the derivatives exists, hence by L'Hôpital's rule
\[ \boxed{\lim_{x\to1^-}\frac{\frac\pi2-\arcsin x}{\sqrt[3]{1-x^2}}=0} \]
(Intuition: $\frac\pi2-\arcsin x=\arccos x\approx\sqrt{2(1-x)}$, which is ``smaller'' than $\sqrt[3]{1-x^2}\approx\sqrt[3]{2(1-x)}$.)
\end{solution}

\subsection*{Question 2}

\begin{question}
Find the points of discontinuity of the function
\[ f(x)=\frac{x\ln|x|}{x^2-1} \]
and determine their types.
\end{question}

\begin{hint}
At the points $x=\pm1$ the numerator also vanishes — use the limit $\lim_{t\to0}\frac{\ln(1+t)}{t}=1$ or L'Hôpital's rule. At $x=0$ use the fact that $x\ln|x|\to0$.
\end{hint}

\begin{solution}
\textbf{The domain:} $\ln|x|$ is defined for $x\neq0$, and the denominator vanishes for $x=\pm1$. Hence $f$ is defined for $x\notin\{0,1,-1\}$, and there it is continuous as a quotient of continuous functions with a nonzero denominator. The points of discontinuity are $x=0,\,1,\,-1$.

\textbf{$x=0$:} It is known that $\lim_{x\to0}x\ln|x|=0$ (for example by L'Hôpital: $\lim_{x\to0^+}\frac{\ln x}{1/x}=\lim_{x\to0^+}\frac{1/x}{-1/x^2}=\lim_{x\to0^+}(-x)=0$, and from the left it is the same since $x\ln|x|$ is odd). The denominator tends to $-1$. Hence
\[ \lim_{x\to0}f(x)=\frac{0}{-1}=0. \]
The limit exists and is finite, but $f$ is not defined at $0$: a \textbf{removable discontinuity}.

\textbf{$x=1$:} In a neighborhood of $1$, $|x|=x$, and the expression is of the form $\frac00$. Write
\[ f(x)=\frac{x}{x+1}\cdot\frac{\ln x}{x-1}. \]
Substitute $t=x-1\to0$: $\frac{\ln x}{x-1}=\frac{\ln(1+t)}{t}\to1$ (a famous limit). Hence
\[ \lim_{x\to1}f(x)=\frac{1}{2}\cdot1=\frac12. \]
A \textbf{removable discontinuity}.

\textbf{$x=-1$:} Note that $f(-x)=\frac{-x\ln|x|}{x^2-1}=-f(x)$, i.e. $f$ is odd. Hence
\[ \lim_{x\to-1}f(x)=-\lim_{x\to1}f(x)=-\frac12. \]
(Or directly: in a neighborhood of $-1$, $f(x)=\frac{x}{x-1}\cdot\frac{\ln(-x)}{x+1}$, and with $t=-x-1\to0$: $\frac{\ln(-x)}{x+1}=\frac{\ln(1+t)}{-t}\to-1$, and $\frac{x}{x-1}\to\frac12$.)
A \textbf{removable discontinuity}.

\textbf{Summary:} the three points of discontinuity $x=0,1,-1$ are all removable; one can define $f(0)=0$, $f(1)=\frac12$, $f(-1)=-\frac12$ and obtain a function continuous on all of $\R$.
\end{solution}

\subsection*{Question 3a}

\begin{question}
Let $f(x)=x^2\ln(1+x)$. Compute $f^{(11)}(0)$.

\underline{Hint}: you may use Maclaurin formulas.
\end{question}

\begin{hint}
The coefficient of $x^{11}$ in the Maclaurin polynomial of $f$ equals $\frac{f^{(11)}(0)}{11!}$. Which term in the expansion of $\ln(1+x)$ contributes to $x^{11}$ after multiplying by $x^2$?
\end{hint}

\begin{solution}
The Maclaurin expansion of $\ln(1+x)$:
\[ \ln(1+x)=x-\frac{x^2}{2}+\frac{x^3}{3}-\dots+(-1)^{n-1}\frac{x^n}{n}+o(x^n). \]
Multiply by $x^2$ (with $n=9$):
\[ f(x)=x^2\ln(1+x)=x^3-\frac{x^4}{2}+\frac{x^5}{3}-\dots+(-1)^{8}\frac{x^{11}}{9}+o(x^{11}). \]
This is a polynomial of degree $\le 11$ plus $o(x^{11})$, hence by the uniqueness of the Taylor polynomial it is the Maclaurin polynomial of order $11$ of $f$ ($f$ is infinitely differentiable in a neighborhood of $0$). The coefficient of $x^{11}$ in the Maclaurin polynomial is $\frac{f^{(11)}(0)}{11!}$, hence
\[ \frac{f^{(11)}(0)}{11!}=\frac19\quad\Longrightarrow\quad f^{(11)}(0)=\frac{11!}{9}=\frac{39916800}{9}=4435200. \]
\[ \boxed{f^{(11)}(0)=\frac{11!}{9}=4435200} \]
\end{solution}

\subsection*{Question 3b}

\begin{question}
Prove that the polynomial $x^3-4x^2+7x+13$ has exactly one real root.
\end{question}

\begin{hint}
Existence — the Intermediate Value Theorem. Uniqueness — show that the derivative is always positive.
\end{hint}

\begin{solution}
Let $p(x)=x^3-4x^2+7x+13$.

\textbf{Existence:} $p$ is continuous (a polynomial). $p(-2)=-8-16-14+13=-25<0$ and $p(0)=13>0$. By the Intermediate Value Theorem there exists $c\in(-2,0)$ with $p(c)=0$.

\textbf{Uniqueness:} $p'(x)=3x^2-8x+7$. The discriminant is $64-84=-20<0$ and the leading coefficient is positive, hence $p'(x)>0$ for every $x\in\R$. Therefore $p$ is strictly increasing on $\R$ (a corollary of Lagrange's theorem), and hence it cannot take the value $0$ twice.

(Alternatively: if there were two roots $a<b$, then by Rolle's theorem there would be $c\in(a,b)$ with $p'(c)=0$, a contradiction.)

Therefore the polynomial has exactly one real root (lying in the interval $(-2,0)$; approximately $c\approx-1.054$). $\blacksquare$
\end{solution}

\subsection*{Question 4}

\begin{question}
Investigate the function $f(x)=x^2e^{-x^2}$ and sketch its graph.
\end{question}

\begin{hint}
The function is even — it suffices to study it for $x\ge0$.
\end{hint}

\begin{solution}
\textbf{1. General description.}
\begin{itemize}
  \item \textbf{Domain:} all of $\R$.
  \item \textbf{Intersection with the axes:} $f(x)=0\iff x=0$ (since $e^{-x^2}>0$). The only intersection point is $(0,0)$.
  \item \textbf{Parity:} $f(-x)=(-x)^2e^{-(-x)^2}=f(x)$, i.e. $f$ is even and the graph is symmetric with respect to the $y$-axis. It is not periodic.
  \item \textbf{Sign:} $f(x)\ge0$ for every $x$, with equality only at $x=0$.
\end{itemize}

\textbf{2. Continuity:} $f$ is elementary and defined on all of $\R$, hence continuous on all of $\R$ (and there are no vertical asymptotes).

\textbf{3. Behavior at infinity and asymptotes:}
\[ \lim_{x\to\pm\infty}x^2e^{-x^2}=\lim_{t\to\infty}\frac{t}{e^t}=0 \]
(we substituted $t=x^2$; the limit $\frac{t}{e^t}\to0$ — a known limit, or L'Hôpital: $\frac{1}{e^t}\to0$). Hence $y=0$ is a horizontal asymptote at $\pm\infty$.

\textbf{4. Increase, decrease and extrema:}
\[ f'(x)=2xe^{-x^2}+x^2e^{-x^2}(-2x)=2xe^{-x^2}(1-x^2)=-2x(x-1)(x+1)e^{-x^2}. \]
Critical points: $x=0,\pm1$. The sign of $f'$ is the sign of $x(1-x^2)$:
\begin{itemize}
  \item on $(-\infty,-1)$: $f'>0$, $f$ is increasing;
  \item on $(-1,0)$: $f'<0$, $f$ is decreasing;
  \item on $(0,1)$: $f'>0$, $f$ is increasing;
  \item on $(1,\infty)$: $f'<0$, $f$ is decreasing.
\end{itemize}
Hence: at $x=\pm1$ there is a local (and also absolute) maximum, $f(\pm1)=e^{-1}=\frac1e$; at $x=0$ there is a local (and also absolute) minimum, $f(0)=0$. The range of the function is $\left[0,\frac1e\right]$ (by continuity, by the Intermediate Value Theorem, and since $f\to0$ at infinity).

\textbf{5. Convexity and inflection points:}
\[ f''(x)=\left(2(x-x^3)e^{-x^2}\right)'=2(1-3x^2)e^{-x^2}+2(x-x^3)(-2x)e^{-x^2}=2e^{-x^2}\left(2x^4-5x^2+1\right). \]
$f''(x)=0\iff 2x^4-5x^2+1=0\iff x^2=\frac{5\pm\sqrt{17}}{4}$. Denote
\[ x_1=\sqrt{\tfrac{5-\sqrt{17}}{4}}\approx0.468,\qquad x_2=\sqrt{\tfrac{5+\sqrt{17}}{4}}\approx1.510. \]
The polynomial $2s^2-5s+1$ (with $s=x^2$) is positive outside its roots and negative between them, hence:
\begin{itemize}
  \item $f''>0$ (convex, ``upward'') for $|x|<x_1$ and for $|x|>x_2$;
  \item $f''<0$ (concave, ``downward'') for $x_1<|x|<x_2$.
\end{itemize}
Four inflection points: $x=\pm x_1$ with $f(\pm x_1)=\frac{5-\sqrt{17}}{4}e^{-\frac{5-\sqrt{17}}{4}}\approx0.176$, and $x=\pm x_2$ with $f(\pm x_2)=\frac{5+\sqrt{17}}{4}e^{-\frac{5+\sqrt{17}}{4}}\approx0.233$.

\textbf{6. Sketch:} a ``double hump'' graph symmetric with respect to the $y$-axis: it starts from $y=0$ at $-\infty$ (above the axis), rises to the maximum $\left(-1,\frac1e\right)$, descends to the minimum $(0,0)$ (where the graph is tangent to the $x$-axis), rises again to the maximum $\left(1,\frac1e\right)$, and then decreases and tends to $0$ as $x\to\infty$. The convexity changes at the points $\pm0.468$ and $\pm1.510$.
\end{solution}

\subsection*{Question 5a}

\begin{question}
Compute the integral $\displaystyle\int(4-x^2)^{-3/2}\,dx$. \quad [\underline{Hint}: you may substitute $x=2\sin\theta$.]
\end{question}

\begin{solution}
The integrand is defined for $|x|<2$. Substitute $x=2\sin\theta$ with $\theta\in\left(-\frac\pi2,\frac\pi2\right)$; then $dx=2\cos\theta\,d\theta$ and $\cos\theta>0$, hence
\[ 4-x^2=4-4\sin^2\theta=4\cos^2\theta,\qquad (4-x^2)^{3/2}=8\cos^3\theta. \]
\[ \int(4-x^2)^{-3/2}\,dx=\int\frac{2\cos\theta}{8\cos^3\theta}\,d\theta=\frac14\int\frac{d\theta}{\cos^2\theta}=\frac14\tan\theta+C. \]
Back to the variable $x$: $\sin\theta=\frac x2$, $\cos\theta=\sqrt{1-\frac{x^2}{4}}=\frac{\sqrt{4-x^2}}{2}$, hence $\tan\theta=\frac{x}{\sqrt{4-x^2}}$.
\[ \boxed{\int(4-x^2)^{-3/2}\,dx=\frac{x}{4\sqrt{4-x^2}}+C} \]
Check: $\left(\frac{x}{4\sqrt{4-x^2}}\right)'=\frac14\cdot\frac{\sqrt{4-x^2}+\frac{x^2}{\sqrt{4-x^2}}}{4-x^2}=\frac14\cdot\frac{4}{(4-x^2)^{3/2}}=(4-x^2)^{-3/2}$.
\end{solution}

\subsection*{Question 5b}

\begin{question}
Find the area enclosed inside the loop formed by the graph of the function $r=\tan\frac\theta2$ in polar coordinates.

(The exam includes a figure: the curve passes through the origin, forms a small loop above the origin, and its two branches extend upward and to the sides.)
\end{question}

\begin{hint}
Find for which values of $\theta$ the curve returns to the same point: compare the points obtained for $\theta=\frac\pi2$ and $\theta=-\frac\pi2$.
\end{hint}

\begin{hint}
Area in polar coordinates: $S=\frac12\int_\alpha^\beta r^2\,d\theta$, and use $\tan^2u=\frac1{\cos^2u}-1$.
\end{hint}

\begin{solution}
The curve $r=\tan\frac\theta2$ is defined for $\theta\in(-\pi,\pi)$, and $r\to\pm\infty$ as $\theta\to\pm\pi$ (these are the two branches). The Cartesian point is $(r\cos\theta,\ r\sin\theta)$.
\begin{itemize}
  \item $\theta=0$: $r=0$ — the origin.
  \item $\theta=\frac\pi2$: $r=\tan\frac\pi4=1$, the point $(0,1)$.
  \item $\theta=-\frac\pi2$: $r=\tan\left(-\frac\pi4\right)=-1$, the point $(-1\cdot\cos(-\frac\pi2),\,-1\cdot\sin(-\frac\pi2))=(0,1)$.
\end{itemize}
That is, the curve passes through the point $(0,1)$ twice, at $\theta=\pm\frac\pi2$: it intersects itself there. As $\theta$ runs from $-\frac\pi2$ to $\frac\pi2$, the curve leaves $(0,1)$, passes through the origin and returns to $(0,1)$ — this is the loop. (For $\theta\in(-\frac\pi2,0)$, $r<0$ and the points lie in the second quadrant; for $\theta\in(0,\frac\pi2)$ in the first quadrant; the loop is symmetric with respect to the $y$-axis.)

The area enclosed by the loop, by $S=\frac12\int_\alpha^\beta r^2\,d\theta$ (the formula is valid also when $r<0$, since $r^2$ is the square of the distance from the origin):
\[ S=\frac12\int_{-\pi/2}^{\pi/2}\tan^2\frac\theta2\,d\theta=\int_0^{\pi/2}\tan^2\frac\theta2\,d\theta \]
(the integrand is even). With $\tan^2u=\frac1{\cos^2u}-1$:
\[ S=\int_0^{\pi/2}\left(\frac{1}{\cos^2\frac\theta2}-1\right)d\theta=\left[2\tan\frac\theta2-\theta\right]_0^{\pi/2}=2\tan\frac\pi4-\frac\pi2=2-\frac\pi2. \]
\[ \boxed{S=2-\frac\pi2\approx0.429} \]
\end{solution}

\subsection*{Question 6a}

\begin{question}
Does the improper integral $\displaystyle\int_1^\infty\frac{e^x}{e^{e^x}}\,dx$ converge? If so, compute the value of the integral. If not, explain why.
\end{question}

\begin{hint}
Substitute $u=e^x$.
\end{hint}

\begin{solution}
By definition $\int_1^\infty\frac{e^x}{e^{e^x}}dx=\lim_{b\to\infty}\int_1^b e^xe^{-e^x}dx$. Substitute $u=e^x$, $du=e^x\,dx$; as $x$ runs from $1$ to $b$, $u$ runs from $e$ to $e^b$:
\[ \int_1^b e^xe^{-e^x}\,dx=\int_e^{e^b}e^{-u}\,du=\left[-e^{-u}\right]_e^{e^b}=e^{-e}-e^{-e^b}. \]
As $b\to\infty$: $e^b\to\infty$, hence $e^{-e^b}\to0$. The limit exists and is finite, i.e. the integral converges:
\[ \boxed{\int_1^\infty\frac{e^x}{e^{e^x}}\,dx=e^{-e}} \]
\end{solution}

\subsection*{Question 6b}

\begin{question}
Compute the volume of the solid of revolution obtained by rotating the region $D$ about the $x$-axis, where $D$ is the region between the graph of the function $f(x)=\frac{\arctan x}{\sqrt{1+x^2}}$ and the $x$-axis, for $0\le x\le1$.
\end{question}

\begin{hint}
Note that $\frac{1}{1+x^2}$ is the derivative of $\arctan x$, and substitute $u=\arctan x$.
\end{hint}

\begin{solution}
$V=\pi\int_0^1 f^2(x)\,dx=\pi\int_0^1\frac{\arctan^2x}{1+x^2}\,dx$. Substitute $u=\arctan x$, $du=\frac{dx}{1+x^2}$; as $x$ runs from $0$ to $1$, $u$ runs from $0$ to $\frac\pi4$:
\[ V=\pi\int_0^{\pi/4}u^2\,du=\pi\left[\frac{u^3}{3}\right]_0^{\pi/4}=\pi\cdot\frac{1}{3}\cdot\frac{\pi^3}{64}. \]
\[ \boxed{V=\frac{\pi^4}{192}} \]
\end{solution}

\section*{2017/18 Semester A Moed C}

\subsection*{Question 1a}

\begin{question}
Compute the following limit:
the limit of the sequence $\displaystyle\lim_{n\to\infty}\left(\frac{(n+4)!\cdot n}{(n+5)!-(n+2)!}\right)$.
\end{question}

\begin{hint}
Factor out $(n+2)!$ as a common factor in the numerator and the denominator.
\end{hint}

\begin{solution}
We use $(n+4)!=(n+2)!\,(n+3)(n+4)$ and $(n+5)!=(n+2)!\,(n+3)(n+4)(n+5)$. Cancel $(n+2)!\neq0$:
\[ \frac{(n+4)!\cdot n}{(n+5)!-(n+2)!}=\frac{n(n+3)(n+4)}{(n+3)(n+4)(n+5)-1}. \]
The numerator and the denominator are polynomials of degree $3$ with leading coefficient $1$. Divide by $n^3$:
\[ =\frac{\left(1+\frac3n\right)\left(1+\frac4n\right)}{\left(1+\frac3n\right)\left(1+\frac4n\right)\left(1+\frac5n\right)-\frac1{n^3}}\xrightarrow[n\to\infty]{}\frac{1}{1-0}=1 \]
by the arithmetic of limits.
\[ \boxed{\lim_{n\to\infty}\frac{(n+4)!\cdot n}{(n+5)!-(n+2)!}=1} \]
\end{solution}

\subsection*{Question 1b}

\begin{question}
Compute the following limit:
the limit of the function $\displaystyle\lim_{x\to0}\tan x\ln(x^2)$.
\end{question}

\begin{hint}
Write $\tan x\ln(x^2)=\frac{\tan x}{x}\cdot 2x\ln|x|$.
\end{hint}

\begin{solution}
For $0<|x|<\frac\pi2$: $\ln(x^2)=2\ln|x|$, hence
\[ \tan x\ln(x^2)=\frac{\tan x}{x}\cdot 2x\ln|x|. \]
It is known that $\lim_{x\to0}\frac{\tan x}{x}=1$ (a famous limit). Also $\lim_{x\to0}x\ln|x|=0$: for $x\to0^+$ by L'Hôpital ($\frac\infty\infty$)
\[ \lim_{x\to0^+}\frac{\ln x}{1/x}=\lim_{x\to0^+}\frac{1/x}{-1/x^2}=\lim_{x\to0^+}(-x)=0, \]
and for $x\to0^-$ it follows from the fact that $x\ln|x|$ is odd. Hence
\[ \boxed{\lim_{x\to0}\tan x\ln(x^2)=1\cdot2\cdot0=0} \]
\end{solution}

\subsection*{Question 2}

\begin{question}
Find the points of discontinuity of the function
\[ f(x)=\frac{x^2-x}{e^{x^3-x}-1} \]
and determine their types.
\end{question}

\begin{hint}
The denominator vanishes when $x^3-x=0$. Use the limit $\lim_{t\to0}\frac{e^t-1}{t}=1$ with $t=x^3-x$.
\end{hint}

\begin{solution}
\textbf{The domain:} $e^{x^3-x}-1=0\iff x^3-x=0\iff x(x-1)(x+1)=0$. Hence $f$ is defined and continuous (a quotient of continuous functions) for every $x\notin\{0,1,-1\}$, and these are the points of discontinuity.

Denote $t=t(x)=x^3-x$. For $x\notin\{0,1,-1\}$:
\[ f(x)=\frac{x^2-x}{x^3-x}\cdot\frac{x^3-x}{e^{x^3-x}-1}=\frac{x(x-1)}{x(x-1)(x+1)}\cdot\frac{t}{e^t-1}=\frac{1}{x+1}\cdot\frac{t}{e^t-1}. \]
As $x\to0$ or $x\to\pm1$, $t\to0$ and $t\neq0$, hence by the famous limit and the theorem on the limit of a composition, $\frac{t}{e^t-1}\to1$.

\textbf{$x=0$:} $\lim_{x\to0}f(x)=\frac{1}{0+1}\cdot1=1$. The limit exists and is finite — a \textbf{removable discontinuity}.

\textbf{$x=1$:} $\lim_{x\to1}f(x)=\frac{1}{2}\cdot1=\frac12$ — a \textbf{removable discontinuity}.

\textbf{$x=-1$:} Here the numerator tends to $(-1)^2-(-1)=2\neq0$ and the denominator tends to $0$. We check signs: $x^3-x=x(x-1)(x+1)$, and in a neighborhood of $-1$, $x(x-1)\approx2>0$. Hence
\begin{itemize}
  \item for $x\to-1^+$: $x+1>0$, $t\to0^+$, $e^t-1\to0^+$, hence $f(x)\to+\infty$;
  \item for $x\to-1^-$: $x+1<0$, $t\to0^-$, $e^t-1\to0^-$, hence $f(x)\to-\infty$.
\end{itemize}
The one-sided limits are infinite — a \textbf{discontinuity of the second kind}.

\textbf{Summary:} $x=0$ and $x=1$ — removable discontinuities (with limits $1$ and $\frac12$ respectively); $x=-1$ — a discontinuity of the second kind.
\end{solution}

\subsection*{Question 3a}

\begin{question}
Compute the limit:
\[ \lim_{x\to0}\left(\frac{1-\cos(x^3)}{x^2\ln(1+x^2)\sin^2x}\right) \]
\underline{Hint}: you may use Maclaurin formulas.
\end{question}

\begin{hint}
$1-\cos u=\frac{u^2}{2}+o(u^2)$, $\ln(1+u)=u+o(u)$, $\sin x=x+o(x)$. What is the order of magnitude of the numerator and of the denominator?
\end{hint}

\begin{solution}
Maclaurin expansions: $\cos u=1-\frac{u^2}{2}+o(u^2)$ with $u=x^3$ gives
\[ 1-\cos(x^3)=\frac{x^6}{2}+o(x^6). \]
Also $\ln(1+x^2)=x^2+o(x^2)$ and $\sin x=x+o(x)$, hence $\sin^2x=x^2+o(x^2)$. The denominator:
\[ x^2\ln(1+x^2)\sin^2x=x^2\left(x^2+o(x^2)\right)\left(x^2+o(x^2)\right)=x^6+o(x^6). \]
Hence
\[ \frac{1-\cos(x^3)}{x^2\ln(1+x^2)\sin^2x}=\frac{\frac12+\frac{o(x^6)}{x^6}}{1+\frac{o(x^6)}{x^6}}\xrightarrow[x\to0]{}\frac12. \]
(Equivalently: $\frac{1-\cos(x^3)}{x^6}\to\frac12$, $\frac{\ln(1+x^2)}{x^2}\to1$, $\frac{\sin^2x}{x^2}\to1$, and we split the expression into a product.)
\[ \boxed{\lim_{x\to0}\frac{1-\cos(x^3)}{x^2\ln(1+x^2)\sin^2x}=\frac12} \]
\end{solution}

\subsection*{Question 3b}

\begin{question}
Let $f$ be a function differentiable at every real number $x$, and assume that the equation $f'(x)=0$ has exactly $k$ solutions. Prove that the equation $f(x)=0$ has at most $k+1$ solutions.
\end{question}

\begin{hint}
Between any two roots of $f$ there is a root of $f'$ (Rolle's theorem). Prove by contradiction.
\end{hint}

\begin{solution}
Assume by contradiction that the equation $f(x)=0$ has at least $k+2$ distinct solutions, and denote $k+2$ of them by $x_1<x_2<\dots<x_{k+2}$.

For every $i=1,\dots,k+1$, the function $f$ is continuous on the interval $[x_i,x_{i+1}]$ and differentiable on the interval $(x_i,x_{i+1})$ (since it is differentiable on all of $\R$, and hence also continuous), and $f(x_i)=f(x_{i+1})=0$. By \textbf{Rolle's theorem} there exists a point $c_i\in(x_i,x_{i+1})$ such that $f'(c_i)=0$.

The intervals $(x_1,x_2),(x_2,x_3),\dots,(x_{k+1},x_{k+2})$ are pairwise disjoint, hence $c_1<c_2<\dots<c_{k+1}$ are $k+1$ distinct points at which $f'=0$. This contradicts the assumption that the equation $f'(x)=0$ has exactly $k$ solutions.

Therefore the equation $f(x)=0$ has at most $k+1$ solutions. $\blacksquare$
\end{solution}

\subsection*{Question 4}

\begin{question}
Investigate the function $f(x)=\sqrt{\frac{x^3}{x-1}}$ and sketch its graph. In this question there is no need to compute $f''(x)$, and no need to deal with convexity, concavity and inflection points.
\end{question}

\begin{hint}
The domain: where $\frac{x^3}{x-1}\ge0$. Note that it consists of two separate parts.
\end{hint}

\begin{hint}
For the oblique asymptote note that $\sqrt{x^2}=|x|$, hence the slopes at $+\infty$ and at $-\infty$ are different.
\end{hint}

\begin{solution}
\textbf{1. General description.}
\begin{itemize}
  \item \textbf{Domain:} we need $x\neq1$ and $\frac{x^3}{x-1}\ge0$. The sign of $x^3$ is the sign of $x$, hence the quotient is non-negative when $x\le0$ (numerator $\le0$, denominator $<0$) or $x>1$ (both positive). The domain: $(-\infty,0]\cup(1,\infty)$.
  \item \textbf{Intersection with the axes:} $f(x)=0\iff x=0$. The only intersection point (with both axes) is $(0,0)$.
  \item \textbf{Parity/periodicity:} the domain is not symmetric with respect to $0$, hence $f$ is neither even nor odd; it is not periodic.
  \item \textbf{Sign:} $f(x)\ge0$ on the whole domain, with equality only at $x=0$.
\end{itemize}

\textbf{2. Continuity:} $f$ is a composition of elementary functions, hence continuous on its whole domain (at $x=0$ — continuous from the left). The point $x=1$ is not in the domain, and as $x\to1^+$: $x^3\to1$, $x-1\to0^+$, hence $f(x)\to+\infty$.

\textbf{3. Asymptotes.}
\begin{itemize}
  \item \textbf{Vertical:} $x=1$ (from the right).
  \item \textbf{As $x\to+\infty$:} $f(x)=\sqrt{x^2\cdot\frac{x}{x-1}}=x\sqrt{\frac{x}{x-1}}$ (since $x>0$).
  \[ a=\lim_{x\to\infty}\frac{f(x)}{x}=\lim_{x\to\infty}\sqrt{\frac{x}{x-1}}=1, \]
  \[ b=\lim_{x\to\infty}\left(f(x)-x\right)=\lim_{x\to\infty}x\left(\sqrt{\tfrac{x}{x-1}}-1\right)=\lim_{x\to\infty}\frac{x\left(\frac{x}{x-1}-1\right)}{\sqrt{\frac{x}{x-1}}+1}=\lim_{x\to\infty}\frac{\frac{x}{x-1}}{\sqrt{\frac{x}{x-1}}+1}=\frac{1}{2}. \]
  Asymptote: $y=x+\frac12$.
  \item \textbf{As $x\to-\infty$:} here $\sqrt{x^2}=|x|=-x$, hence $f(x)=-x\sqrt{\frac{x}{x-1}}$.
  \[ a=\lim_{x\to-\infty}\frac{f(x)}{x}=-\lim_{x\to-\infty}\sqrt{\frac{x}{x-1}}=-1, \]
  \[ b=\lim_{x\to-\infty}\left(f(x)+x\right)=\lim_{x\to-\infty}(-x)\left(\sqrt{\tfrac{x}{x-1}}-1\right)=-\lim_{x\to-\infty}\frac{\frac{x}{x-1}}{\sqrt{\frac{x}{x-1}}+1}=-\frac12. \]
  Asymptote: $y=-x-\frac12$.
\end{itemize}

\textbf{4. Increase, decrease and extrema.} Denote $g(x)=\frac{x^3}{x-1}$. Then
\[ g'(x)=\frac{3x^2(x-1)-x^3}{(x-1)^2}=\frac{x^2(2x-3)}{(x-1)^2}, \]
and at the points where $g(x)>0$ (i.e. $x<0$ or $x>1$), by the chain rule
\[ f'(x)=\frac{g'(x)}{2\sqrt{g(x)}}=\frac{x^2(2x-3)}{2(x-1)^2\sqrt{\frac{x^3}{x-1}}}. \]
The sign of $f'$ is the sign of $2x-3$ (the other factors are positive):
\begin{itemize}
  \item On $(-\infty,0)$: $2x-3<0$, hence $f$ is (strictly) decreasing. Since $f$ is continuous at $0$ from the left, it is decreasing on all of $(-\infty,0]$.
  \item On $\left(1,\frac32\right)$: $f'<0$, $f$ is decreasing.
  \item On $\left(\frac32,\infty\right)$: $f'>0$, $f$ is increasing.
\end{itemize}
Hence:
\begin{itemize}
  \item At $x=\frac32$ there is a local minimum: $f\left(\frac32\right)=\sqrt{\frac{27/8}{1/2}}=\sqrt{\frac{27}{4}}=\frac{3\sqrt3}{2}\approx2.6$.
  \item At $x=0$ (the endpoint of the domain on the left) there is a local, and even absolute, minimum, $f(0)=0$. The one-sided derivative there: $\frac{f(x)-f(0)}{x}=\frac{|x|^{3/2}}{x\sqrt{1-x}}\to0$ as $x\to0^-$, i.e. the graph reaches $(0,0)$ horizontally (tangent to the $x$-axis).
\end{itemize}

\textbf{5. Sketch.} \emph{Left branch} ($x\le0$): the graph decreases from $+\infty$, close to and above the asymptote $y=-x-\frac12$ as $x\to-\infty$, down to the point $(0,0)$, which it reaches tangent to the $x$-axis, and stops there. \emph{Right branch} ($x>1$): the graph decreases from $+\infty$ (hugging the vertical asymptote $x=1$) to the minimum $\left(\frac32,\frac{3\sqrt3}{2}\right)$, and then increases and approaches (from above) the asymptote $y=x+\frac12$. On the interval $(0,1]$ there is no graph.
\end{solution}

\subsection*{Question 5a}

\begin{question}
Compute the integral $\displaystyle\int\frac{x^4+4x^3+4x^2+8x+1}{(x-1)(x+2)^2}\,dx$.
\end{question}

\begin{hint}
The degree of the numerator is greater than the degree of the denominator — first perform polynomial division, and then decompose into partial fractions of the form $\frac{A}{x-1}+\frac{B}{x+2}+\frac{C}{(x+2)^2}$.
\end{hint}

\begin{solution}
The denominator: $(x-1)(x+2)^2=(x-1)(x^2+4x+4)=x^3+3x^2-4$. The degree of the numerator ($4$) is greater than the degree of the denominator ($3$), so we perform polynomial division:
\[ (x+1)(x^3+3x^2-4)=x^4+4x^3+3x^2-4x-4, \]
hence
\[ x^4+4x^3+4x^2+8x+1=(x+1)(x^3+3x^2-4)+(x^2+12x+5). \]
That is,
\[ \frac{x^4+4x^3+4x^2+8x+1}{(x-1)(x+2)^2}=x+1+\frac{x^2+12x+5}{(x-1)(x+2)^2}. \]
Partial fraction decomposition:
\[ \frac{x^2+12x+5}{(x-1)(x+2)^2}=\frac{A}{x-1}+\frac{B}{x+2}+\frac{C}{(x+2)^2}, \]
\[ x^2+12x+5=A(x+2)^2+B(x-1)(x+2)+C(x-1). \]
\begin{itemize}
  \item $x=1$: $18=9A$, i.e. $A=2$.
  \item $x=-2$: $4-24+5=-15=-3C$, i.e. $C=5$.
  \item Comparing the coefficients of $x^2$: $1=A+B$, i.e. $B=-1$.
\end{itemize}
(Check, constant term: $4A-2B-C=8+2-5=5$, as required.) Hence
\[ \int\frac{x^4+4x^3+4x^2+8x+1}{(x-1)(x+2)^2}\,dx=\int\left(x+1+\frac{2}{x-1}-\frac{1}{x+2}+\frac{5}{(x+2)^2}\right)dx \]
\[ \boxed{=\frac{x^2}{2}+x+2\ln|x-1|-\ln|x+2|-\frac{5}{x+2}+C} \]
\end{solution}

\subsection*{Question 5b}

\begin{question}
Compute the length of the curve $f(x)=x^{3/2}$ between $x=0$ and $x=9$.
\end{question}

\begin{hint}
Arc length: $L=\int_a^b\sqrt{1+(f'(x))^2}\,dx$.
\end{hint}

\begin{solution}
$f'(x)=\frac32x^{1/2}$, hence $1+(f'(x))^2=1+\frac94x$. The arc length:
\[ L=\int_0^9\sqrt{1+\frac94x}\,dx. \]
Substitute $u=1+\frac94x$, $du=\frac94dx$; as $x$ runs from $0$ to $9$, $u$ runs from $1$ to $1+\frac{81}{4}=\frac{85}{4}$:
\[ L=\frac49\int_1^{85/4}u^{1/2}\,du=\frac49\cdot\frac23\left[u^{3/2}\right]_1^{85/4}=\frac{8}{27}\left(\left(\frac{85}{4}\right)^{3/2}-1\right)=\frac{8}{27}\left(\frac{85\sqrt{85}}{8}-1\right). \]
\[ \boxed{L=\frac{85\sqrt{85}-8}{27}\approx28.73} \]
\end{solution}

\subsection*{Question 6a}

\begin{question}
Does the improper integral $\displaystyle\int_0^\infty\frac{1}{e^x+e^{-x}}\,dx$ converge? If so, compute the value of the integral. If not, explain why.
\end{question}

\begin{hint}
Multiply the numerator and the denominator by $e^x$ and substitute $u=e^x$.
\end{hint}

\begin{solution}
Multiply the numerator and the denominator by $e^x>0$: $\frac{1}{e^x+e^{-x}}=\frac{e^x}{e^{2x}+1}$. By definition,
\[ \int_0^\infty\frac{dx}{e^x+e^{-x}}=\lim_{b\to\infty}\int_0^b\frac{e^x}{1+(e^x)^2}\,dx. \]
Substitute $u=e^x$, $du=e^x\,dx$; $u$ runs from $1$ to $e^b$:
\[ \int_0^b\frac{e^x\,dx}{1+e^{2x}}=\int_1^{e^b}\frac{du}{1+u^2}=\arctan(e^b)-\arctan1=\arctan(e^b)-\frac\pi4. \]
As $b\to\infty$, $e^b\to\infty$ and $\arctan(e^b)\to\frac\pi2$. The limit exists and is finite, hence the integral converges:
\[ \boxed{\int_0^\infty\frac{dx}{e^x+e^{-x}}=\frac\pi2-\frac\pi4=\frac\pi4} \]
\end{solution}

\subsection*{Question 6b}

\begin{question}
Compute the volume of the solid of revolution obtained by rotating the region $D$ about the $x$-axis, where $D$ is the region between the graph of the function $f(x)=\sqrt{\arcsin x}$ and the $x$-axis, for $0\le x\le1$.
\end{question}

\begin{hint}
We need to compute $\int\arcsin x\,dx$ — use integration by parts with $u=\arcsin x$, $v'=1$.
\end{hint}

\begin{solution}
$V=\pi\int_0^1f^2(x)\,dx=\pi\int_0^1\arcsin x\,dx$. We integrate by parts, with $u=\arcsin x$, $v'=1$, $u'=\frac{1}{\sqrt{1-x^2}}$, $v=x$:
\[ \int\arcsin x\,dx=x\arcsin x-\int\frac{x}{\sqrt{1-x^2}}\,dx=x\arcsin x+\sqrt{1-x^2}+C \]
(the last integral — by the substitution $t=1-x^2$, $dt=-2x\,dx$: $\int\frac{x\,dx}{\sqrt{1-x^2}}=-\frac12\int t^{-1/2}dt=-\sqrt{1-x^2}$.) The function $x\arcsin x+\sqrt{1-x^2}$ is continuous on $[0,1]$ (the integrand $\arcsin x$ is continuous on the closed interval), hence by the Fundamental Theorem
\[ \int_0^1\arcsin x\,dx=\left[x\arcsin x+\sqrt{1-x^2}\right]_0^1=\left(\frac\pi2+0\right)-(0+1)=\frac\pi2-1. \]
\[ \boxed{V=\pi\left(\frac\pi2-1\right)=\frac{\pi^2}{2}-\pi} \]
\end{solution}

\section*{2017/18 Semester B Moed A}

\subsection*{Question 1a}

\begin{question}
Compute the derivative of the function $f(x)=\cos x$ at the point $x$ using the definition of the derivative (an answer based on the table of derivatives will not be accepted).
\end{question}

\begin{hint}
Use the identity $\cos\alpha-\cos\beta=-2\sin\frac{\alpha+\beta}{2}\sin\frac{\alpha-\beta}{2}$ and the fundamental limit $\lim_{t\to0}\frac{\sin t}{t}=1$.
\end{hint}

\begin{solution}
By definition $f'(x)=\lim_{h\to0}\frac{f(x+h)-f(x)}{h}$. By the identity $\cos\alpha-\cos\beta=-2\sin\frac{\alpha+\beta}{2}\sin\frac{\alpha-\beta}{2}$:
\[
\frac{\cos(x+h)-\cos x}{h}=\frac{-2\sin\left(x+\frac h2\right)\sin\frac h2}{h}=-\sin\left(x+\frac h2\right)\cdot\frac{\sin\frac h2}{\frac h2}.
\]
As $h\to0$: $\sin\left(x+\frac h2\right)\to\sin x$ (continuity of the sine) and $\frac{\sin(h/2)}{h/2}\to1$ (the fundamental limit). By the arithmetic of limits:
\[
(\cos x)'=\lim_{h\to0}\frac{\cos(x+h)-\cos x}{h}=-\sin x.
\]
\end{solution}

\subsection*{Question 1b}

\begin{question}
Find all the asymptotes of the function $f(x)=\frac{\ln(2x)}{\ln(x)}$.
\end{question}

\begin{hint}
Write $\ln(2x)=\ln2+\ln x$.
\end{hint}

\begin{solution}
\textbf{Domain:} $x>0$ (for $\ln x$ and $\ln 2x$) and $\ln x\neq0$, i.e. $x\in(0,1)\cup(1,\infty)$. On this domain
\[
f(x)=\frac{\ln2+\ln x}{\ln x}=1+\frac{\ln2}{\ln x}.
\]
\textbf{Vertical:} as $x\to1^+$, $\ln x\to0^+$, hence $f(x)\to+\infty$; as $x\to1^-$, $\ln x\to0^-$, hence $f(x)\to-\infty$. Hence $x=1$ is a vertical asymptote.
As $x\to0^+$: $\ln x\to-\infty$, hence $f(x)\to1$ — a finite limit, so $x=0$ is \emph{not} a vertical asymptote.

\textbf{Horizontal/oblique:} the function is defined only for $x>0$, so we check only $x\to+\infty$: $\ln x\to\infty$, hence $f(x)\to1$. Hence $y=1$ is a horizontal asymptote at $+\infty$ (and there is no oblique one).

\textbf{Answer:} a vertical asymptote $x=1$ and a horizontal asymptote $y=1$ (as $x\to+\infty$).
\end{solution}

\subsection*{Question 1c}

\begin{question}
For which values of $a$ does the graph of the function $f(x)=x^4-ax^2$ have no inflection points?
\end{question}

\begin{hint}
An inflection point is a point where $f''$ changes sign. Compute $f''$ and check for which values of $a$ it changes sign.
\end{hint}

\begin{solution}
$f''(x)=12x^2-2a$. An inflection point is a point where $f''$ changes sign.
\begin{itemize}
  \item If $a>0$: $f''(x)=0\iff x=\pm\sqrt{a/6}$, and $f''$ (a parabola) changes sign at each of these points, hence there are two inflection points.
  \item If $a=0$: $f''(x)=12x^2\ge0$, vanishing only at $x=0$ without changing sign — there are no inflection points ($f=x^4$ is convex).
  \item If $a<0$: $f''(x)=12x^2-2a>0$ for every $x$ — there are no inflection points.
\end{itemize}
\textbf{Answer:} $a\le0$.
\end{solution}

\subsection*{Question 2a}

\begin{question}
For which values of the parameter $a$ does the equation $x^4-a\cdot x^3+27=0$ have no real roots?
\end{question}

\begin{hint}
$x=0$ is not a root. Isolate the parameter: $a=\frac{x^4+27}{x^3}=x+\frac{27}{x^3}$, and find the range of the function $g(x)=x+\frac{27}{x^3}$.
\end{hint}

\begin{solution}
$x=0$ is not a root (we get $27\neq0$). For $x\neq0$:
\[
x^4-ax^3+27=0\iff a=\frac{x^4+27}{x^3}=x+\frac{27}{x^3}=:g(x).
\]
Hence the equation has a real root if and only if $a$ belongs to the image of $g$ on $\R\setminus\{0\}$. We find the image.

\textbf{For $x>0$:} $g'(x)=1-\frac{81}{x^4}$, which vanishes at $x=3$; $g'<0$ on $(0,3)$ and $g'>0$ on $(3,\infty)$. Hence $x=3$ is an absolute minimum on $(0,\infty)$, with $g(3)=3+1=4$. In addition, $g(x)\to+\infty$ as $x\to0^+$ and as $x\to+\infty$. Since $g$ is continuous, by the Intermediate Value Theorem the image of $(0,\infty)$ is $[4,\infty)$.

\textbf{For $x<0$:} $g$ is an odd function ($g(-x)=-g(x)$), hence the image of $(-\infty,0)$ is $(-\infty,-4]$.

Hence the image of $g$ is $(-\infty,-4]\cup[4,\infty)$, and the equation has a real root if and only if $|a|\ge4$. (Check: for $a=4$, $x=3$: $81-108+27=0$.)

\textbf{Answer:} there are no real roots if and only if $-4<a<4$.
\end{solution}

\subsection*{Question 2b}

\begin{question}
Compute the volume of the solid obtained by rotating about the $x$-axis the plane region bounded by the line $y=\frac{x}{\sqrt{x^2+1}}$ and the $x$-axis on the interval $[0,1]$. Draw the corresponding sketch.
\end{question}

\begin{solution}
\textbf{Sketch:} $y=\frac{x}{\sqrt{x^2+1}}$ increases from $0$ (at $x=0$) to $\frac1{\sqrt2}\approx0.71$ (at $x=1$), and is non-negative on the interval. The region is the area between the graph and the $x$-axis for $0\le x\le1$, and rotating it about the $x$-axis gives a ``goblet''-like solid.

\textbf{Volume:} by the formula for the volume of a solid of revolution $V=\pi\int_a^by^2\,dx$:
\[
V=\pi\int_0^1\frac{x^2}{x^2+1}dx=\pi\int_0^1\left(1-\frac{1}{x^2+1}\right)dx=\pi\Big[x-\arctan x\Big]_0^1=\pi\left(1-\frac\pi4\right).
\]
\textbf{Answer:} $V=\pi\left(1-\frac\pi4\right)\approx0.674$.
\end{solution}

\subsection*{Question 3a}

\begin{question}
Compute: $\displaystyle\int_{-1}^15x^4\big(\ln(x+2)+\sin x\big)dx$
\end{question}

\begin{hint}
$5x^4\sin x$ is an odd function, hence its integral over $[-1,1]$ vanishes. For the other part, integrate by parts with $u=\ln(x+2)$, $dv=5x^4dx$.
\end{hint}

\begin{solution}
\textbf{The part with the sine.} $g(x)=5x^4\sin x$ is odd ($g(-x)=-g(x)$) and continuous, hence $\int_{-1}^1 5x^4\sin x\,dx=0$.

\textbf{The part with the logarithm.} Integration by parts with $u=\ln(x+2)$, $dv=5x^4dx$, $du=\frac{dx}{x+2}$, $v=x^5$:
\[
\int_{-1}^15x^4\ln(x+2)\,dx=\Big[x^5\ln(x+2)\Big]_{-1}^1-\int_{-1}^1\frac{x^5}{x+2}dx=\ln3-\int_{-1}^1\frac{x^5}{x+2}dx
\]
(since at $x=-1$: $(-1)^5\ln1=0$). Polynomial division:
\[
\frac{x^5}{x+2}=x^4-2x^3+4x^2-8x+16-\frac{32}{x+2}.
\]
On the symmetric interval $[-1,1]$ the integrals of the odd powers vanish, hence
\[
\int_{-1}^1\frac{x^5}{x+2}dx=\int_{-1}^1\left(x^4+4x^2+16\right)dx-32\int_{-1}^1\frac{dx}{x+2}
=\frac25+\frac83+32-32\ln3=\frac{526}{15}-32\ln3.
\]
\textbf{Answer:}
\[
\int_{-1}^15x^4\big(\ln(x+2)+\sin x\big)dx=\ln3-\frac{526}{15}+32\ln3=33\ln3-\frac{526}{15}\approx1.1875.
\]
\end{solution}

\subsection*{Question 3b}

\begin{question}
Write the Maclaurin polynomial of order 2 for the function $y(x)$ which is defined implicitly by the equation $y\cdot\sin x-x\cdot\cos x+y=2$.
\end{question}

\begin{hint}
Substitute $x=0$ to find $y(0)$, then differentiate the equation (implicit differentiation) twice and substitute $x=0$.
\end{hint}

\begin{solution}
\textbf{Value at $0$:} substituting $x=0$: $y(0)\cdot0-0+y(0)=2$, hence $y(0)=2$.

\textbf{First derivative} (implicit differentiation of both sides with respect to $x$, where $y=y(x)$):
\[
y'\sin x+y\cos x-\cos x+x\sin x+y'=0.\tag{*}
\]
Substituting $x=0$: $y(0)-1+y'(0)=0$, hence $y'(0)=1-2=-1$.

\textbf{Second derivative} (differentiating (*)):
\[
y''\sin x+y'\cos x+y'\cos x-y\sin x+\sin x+\sin x+x\cos x+y''=0.
\]
Substituting $x=0$: $2y'(0)+y''(0)=0$, hence $y''(0)=2$.

\textbf{The Maclaurin polynomial of order 2:}
\[
P_2(x)=y(0)+y'(0)x+\frac{y''(0)}{2}x^2=2-x+x^2.
\]
(Check: from the equation $y=\frac{2+x\cos x}{1+\sin x}$, and a direct expansion gives the same polynomial.)
\end{solution}

\subsection*{Question 4a}

\begin{question}
$\displaystyle\lim_{x\to3}\frac{\sqrt{x^2-2x+6}-\sqrt{x^2+2x-6}}{x^2-4x+3}$.
\end{question}

\begin{hint}
Multiply by the conjugate of the numerator.
\end{hint}

\begin{solution}
At $x=3$ the numerator and the denominator vanish ($\sqrt9-\sqrt9=0$, $9-12+3=0$). Multiply by the conjugate:
\[
\sqrt{x^2-2x+6}-\sqrt{x^2+2x-6}=\frac{(x^2-2x+6)-(x^2+2x-6)}{\sqrt{x^2-2x+6}+\sqrt{x^2+2x-6}}=\frac{-4(x-3)}{\sqrt{x^2-2x+6}+\sqrt{x^2+2x-6}},
\]
and $x^2-4x+3=(x-1)(x-3)$. Hence for $x\neq3$ (in a neighborhood of $3$):
\[
\frac{\sqrt{x^2-2x+6}-\sqrt{x^2+2x-6}}{x^2-4x+3}=\frac{-4}{(x-1)\left(\sqrt{x^2-2x+6}+\sqrt{x^2+2x-6}\right)}\xrightarrow[x\to3]{}\frac{-4}{2\cdot(3+3)}=-\frac13.
\]
\textbf{Answer:} $-\frac13$.
\end{solution}

\subsection*{Question 4b}

\begin{question}
Is the function
\[
f(x)=\begin{cases}(1+\sin2x)^{\frac{2}{3x}}+\arctan x\cdot\cos\frac1x & x\neq0\\ e^{4/3} & x=0\end{cases}
\]
continuous at the point $x=0$?
\end{question}

\begin{hint}
Treat each summand separately: the first is of the form $1^\infty$; the second is ``tends to zero times bounded''.
\end{hint}

\begin{solution}
We need to check whether $\lim_{x\to0}f(x)=f(0)=e^{4/3}$.

\textbf{The first summand.} In a neighborhood of $0$ we have $1+\sin2x>0$, hence
\[
(1+\sin2x)^{\frac{2}{3x}}=\exp\left(\frac{2}{3x}\ln(1+\sin2x)\right).
\]
We compute the exponent:
\[
\frac{2\ln(1+\sin2x)}{3x}=\frac23\cdot\frac{\ln(1+\sin2x)}{\sin2x}\cdot\frac{\sin2x}{2x}\cdot2\xrightarrow[x\to0]{}\frac23\cdot1\cdot1\cdot2=\frac43,
\]
by the fundamental limits $\lim_{t\to0}\frac{\ln(1+t)}{t}=1$ (with $t=\sin2x\to0$) and $\lim_{t\to0}\frac{\sin t}{t}=1$. By the continuity of the exponential, the first summand tends to $e^{4/3}$.

\textbf{The second summand.} $\arctan x\to0$ as $x\to0$, and $\left|\cos\frac1x\right|\le1$. A product of a function tending to $0$ and a bounded function tends to $0$, hence $\arctan x\cdot\cos\frac1x\to0$.

Hence $\lim_{x\to0}f(x)=e^{4/3}+0=e^{4/3}=f(0)$.

\textbf{Answer:} yes, $f$ is continuous at $x=0$.
\end{solution}

\subsection*{Question 4c}

\begin{question}
Does the improper integral $\displaystyle\int_0^{+\infty}\frac{x\cdot e^{-x}}{(x+1)^4}dx$ converge?
\end{question}

\begin{hint}
Compare with $e^{-x}$.
\end{hint}

\begin{solution}
The integrand $g(x)=\frac{xe^{-x}}{(x+1)^4}$ is continuous and non-negative on $[0,\infty)$, so the integral is improper only because of the infinite limit of integration, and we may use the comparison test.
For $x\ge0$: $0\le x<x+1\le(x+1)^4$, hence $\frac{x}{(x+1)^4}\le1$ and
\[
0\le g(x)\le e^{-x}.
\]
The integral $\int_0^\infty e^{-x}dx=\lim_{R\to\infty}\left(1-e^{-R}\right)=1$ converges. By the \emph{comparison test} for improper integrals of non-negative functions, $\int_0^\infty g(x)\,dx$ converges as well (and its value lies between $0$ and $1$).

\textbf{Answer:} yes, the integral converges.
\end{solution}

\section*{2017/18 Semester B Moed B}

\subsection*{Question 1a}

\begin{question}
Compute the derivative of the function $f(x)=\sin2x$ at the point $x$ using the definition of the derivative (an answer based on the table of derivatives will not be accepted).
\end{question}

\begin{hint}
Use the identity $\sin\alpha-\sin\beta=2\cos\frac{\alpha+\beta}{2}\sin\frac{\alpha-\beta}{2}$.
\end{hint}

\begin{solution}
By the definition and by the identity $\sin\alpha-\sin\beta=2\cos\frac{\alpha+\beta}{2}\sin\frac{\alpha-\beta}{2}$:
\[
\frac{\sin(2x+2h)-\sin2x}{h}=\frac{2\cos(2x+h)\sin h}{h}=2\cos(2x+h)\cdot\frac{\sin h}{h}.
\]
As $h\to0$: $\cos(2x+h)\to\cos2x$ (continuity of cosine) and $\frac{\sin h}{h}\to1$ (the fundamental limit). Therefore
\[
(\sin2x)'=\lim_{h\to0}\frac{\sin2(x+h)-\sin2x}{h}=2\cos2x.
\]
\end{solution}

\subsection*{Question 1b}

\begin{question}
Find all the points of discontinuity of the function and classify them
\[
f(x)=\frac{1}{25-5^{1/(x-2)}}.
\]
\end{question}

\begin{hint}
The suspicious points are $x=2$ (where $\frac1{x-2}$ is undefined) and the point where the denominator vanishes. At $x=2$ compute the one-sided limits separately.
\end{hint}

\begin{solution}
\textbf{Domain:} we need $x\neq2$ and also $5^{1/(x-2)}\neq25$, i.e. $\frac1{x-2}\neq2$, i.e. $x\neq\frac52$. At every other point $f$ is a composition and quotient of continuous functions, hence continuous. The points of discontinuity are $x=2$ and $x=\frac52$.

\textbf{The point $x=2$:}
\begin{itemize}
  \item $x\to2^+$: $\frac1{x-2}\to+\infty$, hence $5^{1/(x-2)}\to+\infty$ and $f(x)\to0$.
  \item $x\to2^-$: $\frac1{x-2}\to-\infty$, hence $5^{1/(x-2)}\to0$ and $f(x)\to\frac1{25}$.
\end{itemize}
Both one-sided limits exist, are finite and are different: \textbf{a discontinuity of the first kind (jump)}.

\textbf{The point $x=\frac52$:} the denominator $25-5^{1/(x-2)}$ tends to $0$. For $2<x<\frac52$: $\frac1{x-2}>2$, hence $5^{1/(x-2)}>25$ and the denominator is negative: $f(x)\to-\infty$ as $x\to\frac52^-$. For $x>\frac52$ the denominator is positive: $f(x)\to+\infty$ as $x\to\frac52^+$. The one-sided limits are infinite: \textbf{a discontinuity of the second kind} (and $x=\frac52$ is a vertical asymptote).
\end{solution}

\subsection*{Question 1c}

\begin{question}
What is the slope of the tangent to the curve $y=\int_0^x\frac{dt}{1+t^4}$ at its point with $x=1$.
\end{question}

\begin{solution}
The function $\frac1{1+t^4}$ is continuous on $\R$, so by \emph{the Fundamental Theorem of Calculus} the function $y(x)=\int_0^x\frac{dt}{1+t^4}$ is differentiable and $y'(x)=\frac1{1+x^4}$. The slope of the tangent at $x=1$:
\[
y'(1)=\frac1{1+1}=\frac12.
\]
\end{solution}

\subsection*{Question 2a}

\begin{question}
Prove that for every $x>0$ we have $3\sqrt[3]{x^2}-2x\le1$.
\end{question}

\begin{hint}
Find the maximum of $g(x)=3x^{2/3}-2x$ on the interval $(0,\infty)$.
\end{hint}

\begin{solution}
Define $g(x)=3x^{2/3}-2x$ for $x>0$; it is differentiable there and
\[
g'(x)=2x^{-1/3}-2=2\left(\frac{1}{\sqrt[3]x}-1\right).
\]
$g'(x)>0$ for $0<x<1$ and $g'(x)<0$ for $x>1$. Hence $g$ is increasing on $(0,1]$ and decreasing on $[1,\infty)$ (a corollary of Lagrange's theorem), so $x=1$ is a point of absolute maximum of $g$ on $(0,\infty)$:
\[
g(x)\le g(1)=3-2=1\qquad\text{for every }x>0.
\]
This is exactly $3\sqrt[3]{x^2}-2x\le1$, with equality only at $x=1$. $\blacksquare$
\end{solution}

\subsection*{Question 2b}

\begin{question}
Compute the area enclosed between the curves $y=x\cdot e^x$, $y=x-1$, $x=0$, $x=2$. Draw the corresponding sketch.
\end{question}

\begin{hint}
Check which of the graphs lies above the other on the interval $[0,2]$. For the integral of $xe^x$ use integration by parts.
\end{hint}

\begin{solution}
\textbf{Which is on top.} For $x\in[0,2]$: $e^x\ge1$ and $x\ge0$, hence $xe^x\ge x>x-1$. That is, the graph of $y=xe^x$ lies above the line $y=x-1$ on the whole interval (and there are no intersection points in the interval).

\textbf{Sketch:} the line $y=x-1$ goes from $(0,-1)$ to $(2,1)$; the curve $y=xe^x$ increases from $(0,0)$ to $(2,2e^2)\approx(2,14.8)$. The region is bounded on the left by $x=0$, on the right by $x=2$, from below by the line and from above by the curve.

\textbf{Area:}
\[
S=\int_0^2\big(xe^x-(x-1)\big)dx.
\]
By integration by parts ($u=x$, $dv=e^xdx$): $\int xe^xdx=xe^x-e^x=(x-1)e^x$. Therefore
\[
S=\Big[(x-1)e^x\Big]_0^2-\left[\frac{x^2}{2}-x\right]_0^2=\big(e^2-(-1)\big)-(2-2)=e^2+1.
\]
\textbf{Answer:} $S=e^2+1\approx8.39$.
\end{solution}

\subsection*{Question 3a}

\begin{question}
Compute:
\begin{enumerate}
  \item[(a)] (8 pts) $\displaystyle\int\frac{x^3+2x^2-4}{x^3-2x^2}dx$
  \item[(b)] (7 pts) $\displaystyle\int\frac{x}{\sqrt{x+1}}dx$
\end{enumerate}
\end{question}

\begin{hint}
In (a) perform polynomial division and then decompose into partial fractions: $x^3-2x^2=x^2(x-2)$. In (b) substitute $u=x+1$.
\end{hint}

\begin{solution}
\begin{enumerate}
  \item[(a)] \textbf{Polynomial division:} $x^3+2x^2-4=(x^3-2x^2)+4x^2-4$, hence
  \[
  \frac{x^3+2x^2-4}{x^3-2x^2}=1+\frac{4x^2-4}{x^2(x-2)}.
  \]
  \textbf{Partial fractions:}
  \[
  \frac{4x^2-4}{x^2(x-2)}=\frac Ax+\frac B{x^2}+\frac C{x-2}\iff 4x^2-4=Ax(x-2)+B(x-2)+Cx^2.
  \]
  Substituting $x=2$: $12=4C\Rightarrow C=3$. Substituting $x=0$: $-4=-2B\Rightarrow B=2$. Comparing the coefficients of $x^2$: $A+C=4\Rightarrow A=1$.
  (Check: $x^2-2x+2x-4+3x^2=4x^2-4$.) Therefore
  \[
  \int\frac{x^3+2x^2-4}{x^3-2x^2}dx=\int\left(1+\frac1x+\frac2{x^2}+\frac3{x-2}\right)dx=x+\ln|x|-\frac2x+3\ln|x-2|+C.
  \]

  \item[(b)] Substitute $u=x+1$, $du=dx$, $x=u-1$:
  \[
  \int\frac{x}{\sqrt{x+1}}dx=\int\frac{u-1}{\sqrt u}du=\int\left(u^{1/2}-u^{-1/2}\right)du=\frac23u^{3/2}-2u^{1/2}+C.
  \]
  \textbf{Answer:} $\displaystyle\int\frac{x}{\sqrt{x+1}}dx=\frac23(x+1)^{3/2}-2\sqrt{x+1}+C=\frac23(x-2)\sqrt{x+1}+C$.
\end{enumerate}
\end{solution}

\subsection*{Question 3b}

\begin{question}
Write the Maclaurin polynomial of order 2 for the function $f(x)=\cos(2x)-\arctan(x+1)$.
\end{question}

\begin{solution}
The Maclaurin polynomial of order 2: $P_2(x)=f(0)+f'(0)x+\frac{f''(0)}{2}x^2$.
\[
f(0)=\cos0-\arctan1=1-\frac\pi4.
\]
\[
f'(x)=-2\sin2x-\frac{1}{1+(x+1)^2},\qquad f'(0)=0-\frac12=-\frac12.
\]
\[
f''(x)=-4\cos2x+\frac{2(x+1)}{\big(1+(x+1)^2\big)^2},\qquad f''(0)=-4+\frac{2}{4}=-\frac72.
\]
\textbf{Answer:}
\[
P_2(x)=1-\frac\pi4-\frac12x-\frac74x^2.
\]
\end{solution}

\subsection*{Question 4a}

\begin{question}
For which values of the parameters $a$ and $b$ will the function
\[
f(x)=\begin{cases}\dfrac{\tan(2x)-x}{x+\sin(2x)} & -\frac\pi4<x<0\\[1ex] ax+b & 0\le x\le1\\[0.5ex] x^{1/(x-1)} & x>1\end{cases}
\]
be continuous on the domain $x>-\frac\pi4$?
\end{question}

\begin{hint}
The function is continuous inside each of the three intervals; only the junction points $x=0$ and $x=1$ need to be checked. At $x=0$ divide the numerator and denominator by $x$; at $x=1$ this is a limit of the form $1^\infty$.
\end{hint}

\begin{solution}
\textbf{Continuity inside the intervals.} On $\left(-\frac\pi4,0\right)$: $\tan2x$ is defined and continuous (since $2x\in\left(-\frac\pi2,0\right)$), and the denominator $x+\sin2x$ is strictly negative (both summands are negative), so $f$ is continuous there as a quotient of continuous functions. On $(0,1)$ $f$ is a polynomial. On $(1,\infty)$: $x^{1/(x-1)}=e^{\frac{\ln x}{x-1}}$ is continuous. It remains to check $x=0$ and $x=1$.

\textbf{The point $x=0$.} $f(0)=b$, and the limit from the right is $b$. From the left, divide the numerator and denominator by $x$:
\[
\lim_{x\to0^-}\frac{\tan2x-x}{x+\sin2x}=\lim_{x\to0^-}\frac{2\cdot\frac{\tan2x}{2x}-1}{1+2\cdot\frac{\sin2x}{2x}}=\frac{2-1}{1+2}=\frac13,
\]
by the fundamental limits $\frac{\sin t}{t}\to1$, $\frac{\tan t}{t}\to1$. Hence we need $b=\frac13$.

\textbf{The point $x=1$.} $f(1)=a+b$, and the limit from the left is $a+b$. From the right:
\[
\lim_{x\to1^+}x^{1/(x-1)}=\lim_{x\to1^+}e^{\frac{\ln x}{x-1}}=e^1=e,
\]
since $\lim_{x\to1}\frac{\ln x}{x-1}=1$ (for example by L'Hôpital's rule: $\frac{1/x}{1}\to1$, or by $\ln(1+t)/t\to1$ with $t=x-1$), and by the continuity of the exponential. Hence we need $a+b=e$, i.e. $a=e-\frac13$.

\textbf{Answer:} $a=e-\frac13$, $b=\frac13$.
\end{solution}

\subsection*{Question 4b}

\begin{question}
Compute $\displaystyle\int_1^{+\infty}\frac{\ln x}{x^4}dx$.
\end{question}

\begin{hint}
Integration by parts with $u=\ln x$, $dv=x^{-4}dx$.
\end{hint}

\begin{solution}
\textbf{Antiderivative.} Integration by parts with $u=\ln x$, $dv=x^{-4}dx$, $du=\frac{dx}x$, $v=-\frac1{3x^3}$:
\[
\int\frac{\ln x}{x^4}dx=-\frac{\ln x}{3x^3}+\frac13\int x^{-4}dx=-\frac{\ln x}{3x^3}-\frac1{9x^3}+C.
\]
\textbf{The improper integral:}
\[
\int_1^{\infty}\frac{\ln x}{x^4}dx=\lim_{R\to\infty}\left[-\frac{\ln x}{3x^3}-\frac{1}{9x^3}\right]_1^R
=\lim_{R\to\infty}\left(-\frac{\ln R}{3R^3}-\frac1{9R^3}\right)+\frac19=\frac19,
\]
since $\lim_{R\to\infty}\frac{\ln R}{R^3}=0$ (by L'Hôpital's rule: $\frac{1/R}{3R^2}=\frac1{3R^3}\to0$).

\textbf{Answer:} the integral converges and equals $\frac19$.
\end{solution}

\section*{2018/19 Semester A Moed A}

\subsection*{Question 1a}

\begin{question}
Find the following limit:
$\displaystyle\lim_{x\to0}\big(\cos x\big)^{\frac{1}{\sin^2 2x}}$
\end{question}

\begin{hint}
Write $(\cos x)^{1/\sin^22x}=e^{\frac{\ln\cos x}{\sin^22x}}$ and compute the limit of the exponent.
\end{hint}

\begin{solution}
This is a limit of the form $1^\infty$. In a small punctured neighborhood of $0$ we have $\cos x>0$ and $\sin2x\neq0$, hence
\[ (\cos x)^{\frac1{\sin^22x}}=e^{\frac{\ln\cos x}{\sin^22x}}. \]
We compute the limit of the exponent. Write
\[ \frac{\ln\cos x}{\sin^22x}=\frac{\ln\cos x}{x^2}\cdot\frac{(2x)^2}{\sin^22x}\cdot\frac14. \]
The middle factor tends to $1$ by $\lim_{t\to0}\frac{\sin t}{t}=1$. For the first factor (form $\frac00$) we use L'Hôpital's rule:
\[ \lim_{x\to0}\frac{\ln\cos x}{x^2}=\lim_{x\to0}\frac{-\tan x}{2x}=-\frac12. \]
Hence the limit of the exponent is $-\frac12\cdot1\cdot\frac14=-\frac18$, and by the continuity of the exponential function
\[ \lim_{x\to0}(\cos x)^{\frac1{\sin^22x}}=e^{-1/8}. \]
\end{solution}

\subsection*{Question 1b}

\begin{question}
Find the following limit:
$\displaystyle\lim_{x\to\infty}\frac{\sin\left(\frac{1}{\sqrt x}\right)}{\sqrt{x+1}-\sqrt x}$
\end{question}

\begin{hint}
Multiply by the conjugate $\sqrt{x+1}+\sqrt x$ and use $\lim_{t\to0}\frac{\sin t}{t}=1$ with $t=\frac1{\sqrt x}$.
\end{hint}

\begin{solution}
Multiply and divide by the conjugate:
\[ \frac{\sin\frac1{\sqrt x}}{\sqrt{x+1}-\sqrt x}=\sin\!\left(\tfrac1{\sqrt x}\right)\left(\sqrt{x+1}+\sqrt x\right)
=\frac{\sin\frac1{\sqrt x}}{\frac1{\sqrt x}}\cdot\frac{\sqrt{x+1}+\sqrt x}{\sqrt x}
=\frac{\sin\frac1{\sqrt x}}{\frac1{\sqrt x}}\cdot\left(\sqrt{1+\tfrac1x}+1\right). \]
As $x\to\infty$, $t=\frac1{\sqrt x}\to0^+$, so the first factor tends to $1$ (fundamental limit and composition of limits), and the second factor tends to $1+1=2$. Hence the limit is $\boxed{2}$.
\end{solution}

\subsection*{Question 2}

\begin{question}
To compute $a=\sqrt[3]{29}$ approximately, one writes
\[ a=3\cdot\sqrt[3]{29/27}=3\cdot\sqrt[3]{1+2/27} \]
and uses the Maclaurin polynomial of degree 2 of the function $f(x)=3\cdot(1+x)^{1/3}$. Estimate the approximation error using Taylor's formula.
\end{question}

\begin{hint}
The remainder in Lagrange form is $R_2(x)=\frac{f'''(c)}{3!}x^3$ for some $c$ between $0$ and $x$. Bound $|f'''(c)|$ for $c\in(0,\frac2{27})$.
\end{hint}

\begin{solution}
\textbf{The Maclaurin polynomial.} Differentiate $f(x)=3(1+x)^{1/3}$:
\[ f'(x)=(1+x)^{-2/3},\qquad f''(x)=-\frac23(1+x)^{-5/3},\qquad f'''(x)=\frac{10}{9}(1+x)^{-8/3}. \]
Hence $f(0)=3$, $f'(0)=1$, $f''(0)=-\frac23$ and
\[ P_2(x)=3+x-\frac13x^2. \]
\textbf{The approximation.} For $x=\frac2{27}$:
\[ a\approx P_2\!\left(\tfrac2{27}\right)=3+\frac2{27}-\frac13\cdot\frac{4}{729}=3+\frac{162}{2187}-\frac{4}{2187}=\frac{6719}{2187}\approx3.072245. \]
\textbf{Error estimate.} By Taylor's formula with the Lagrange remainder, there exists $c\in(0,\frac2{27})$ such that
\[ a-P_2\!\left(\tfrac2{27}\right)=R_2\!\left(\tfrac2{27}\right)=\frac{f'''(c)}{3!}\left(\frac2{27}\right)^3=\frac{10}{54}(1+c)^{-8/3}\cdot\frac{8}{19683}. \]
Since $c>0$, we have $0<(1+c)^{-8/3}<1$, hence
\[ 0<R_2<\frac{10}{54}\cdot\frac{8}{19683}=\frac{40}{531441}\approx7.53\cdot10^{-5}. \]
That is, the error is less than $7.6\cdot10^{-5}$ (in particular less than $10^{-4}$), and the approximation is from below (the remainder is positive). For comparison, $\sqrt[3]{29}=3.0723168\ldots$, and the actual error is about $7.17\cdot10^{-5}$.
\end{solution}

\subsection*{Question 3a}

\begin{question}
Let $f(x)$ be a function defined in a neighborhood of a point $x_0$. Complete the following definition:

\textbf{A number $a$ is the limit of $f(x)$ as $x$ tends to $x_0$ if for every \ldots}
\end{question}

\begin{solution}
(Cauchy's definition, $\eps$--$\delta$) A number $a$ is the limit of $f(x)$ as $x$ tends to $x_0$ if for every $\eps>0$ there exists $\delta>0$ such that for every $x$ satisfying $0<|x-x_0|<\delta$ we have $|f(x)-a|<\eps$.

(An equivalent definition, according to Heine: for every sequence $x_n\to x_0$ with $x_n\neq x_0$ we have $f(x_n)\to a$.)
\end{solution}

\subsection*{Question 3b}

\begin{question}
Explain why the function $f(x)=\cos\left(\dfrac{x+1}{x-1}\right)$ has no limit as $x$ tends to $1$.
\end{question}

\begin{hint}
Use Heine's definition of the limit: find two sequences tending to $1$ along which the values of the function tend to different limits.
\end{hint}

\begin{solution}
We use Heine's definition. If $L=\lim_{x\to1}f(x)$ existed, then for every sequence $x_n\to1$, $x_n\ne1$, we would have $f(x_n)\to L$.

Note that if $\frac{x+1}{x-1}=y$ then $x=\frac{y+1}{y-1}$. Choose
\[ x_n=\frac{2\pi n+1}{2\pi n-1},\qquad z_n=\frac{(2n+1)\pi+1}{(2n+1)\pi-1}. \]
Both sequences tend to $1$ (divide the numerator and denominator by $n$) and are different from $1$. We have $\frac{x_n+1}{x_n-1}=2\pi n$ and $\frac{z_n+1}{z_n-1}=(2n+1)\pi$, hence
\[ f(x_n)=\cos(2\pi n)=1\to1,\qquad f(z_n)=\cos\big((2n+1)\pi\big)=-1\to-1. \]
We obtained two sequences tending to $1$ along which the values of the function converge to different limits, contradicting the uniqueness of the limit. Hence $f$ has no limit at $1$.

(Intuitively: as $x\to1^\pm$ the expression $\frac{x+1}{x-1}\to\pm\infty$, and the cosine oscillates between $-1$ and $1$ infinitely many times.)
\end{solution}

\subsection*{Question 4}

\begin{question}
Investigate the function $y=(x+2)e^{1/x}$ and sketch its graph.

\underline{You need to investigate}: domain, continuity, asymptotes, intervals of increase and decrease, extremum points, intervals of convexity and concavity.
\end{question}

\begin{hint}
For the oblique asymptote: $m=\lim\frac{y}{x}$, $n=\lim(y-mx)$; when computing $n$ substitute $t=\frac1x$ and use $\lim_{t\to0}\frac{e^t-1}{t}=1$.
\end{hint}

\begin{hint}
Check the one-sided limits at $x=0$ separately: as $x\to0^-$, $e^{1/x}\to0$.
\end{hint}

\begin{solution}
\textbf{Domain and continuity.} $x\neq0$. On its domain the function is continuous (product and composition of continuous functions). Intersection with the $x$-axis: $x=-2$; $y>0$ for $x>-2$, $x\ne0$, and $y<0$ for $x<-2$.

\textbf{Vertical asymptotes.} As $x\to0^+$: $\frac1x\to+\infty$, $e^{1/x}\to+\infty$ and $x+2\to2$, hence $y\to+\infty$: the line $x=0$ is a vertical asymptote (from the right). As $x\to0^-$: $e^{1/x}\to0$ and hence $y\to0$ -- there is no asymptote from the left (the graph ``enters'' the origin; at $x=0$ there is an essential discontinuity).

\textbf{Oblique asymptotes.} $m=\lim_{x\to\pm\infty}\frac{(x+2)e^{1/x}}{x}=\lim\left(1+\frac2x\right)e^{1/x}=1$. Now
\[ y-x=x\big(e^{1/x}-1\big)+2e^{1/x}=\frac{e^{1/x}-1}{1/x}+2e^{1/x}\xrightarrow[x\to\pm\infty]{}1+2=3, \]
by $\lim_{t\to0}\frac{e^t-1}{t}=1$ with $t=\frac1x\to0$. Hence $y=x+3$ is an oblique asymptote in both directions.

\textbf{Increase and decrease.}
\[ y'=e^{1/x}+(x+2)e^{1/x}\cdot\left(-\frac1{x^2}\right)=e^{1/x}\,\frac{x^2-x-2}{x^2}=e^{1/x}\,\frac{(x-2)(x+1)}{x^2}. \]
The sign of $y'$ is the sign of $(x-2)(x+1)$: $y'>0$ on $(-\infty,-1)$ and on $(2,\infty)$ -- the function is increasing; $y'<0$ on $(-1,0)$ and on $(0,2)$ -- the function is decreasing.

\textbf{Extremum points.} At $x=-1$ the derivative changes sign from $+$ to $-$: a local maximum $y(-1)=e^{-1}=\frac1e$. At $x=2$ from $-$ to $+$: a local minimum $y(2)=4\sqrt e\approx6.59$.

\textbf{Convexity and concavity.} A direct computation (differentiating $e^{1/x}(1-\frac1x-\frac2{x^2})$) gives
\[ y''=e^{1/x}\left(-\frac1{x^2}\right)\left(1-\frac1x-\frac2{x^2}\right)+e^{1/x}\left(\frac1{x^2}+\frac4{x^3}\right)=e^{1/x}\,\frac{5x+2}{x^4}. \]
The sign of $y''$ is the sign of $5x+2$: $y''<0$ (concave, $\cap$) on $(-\infty,-\frac25)$, and $y''>0$ (convex, $\cup$) on $(-\frac25,0)$ and on $(0,\infty)$. Inflection point: $x=-\frac25$, $y\left(-\frac25\right)=\frac85e^{-5/2}\approx0.131$.

\textbf{Sketch.} From the expansion $y=x+3+\frac{5}{2x}+o\!\left(\frac1x\right)$ it follows that the graph lies below the asymptote as $x\to-\infty$ and above it as $x\to+\infty$. On the left: the graph approaches the line $y=x+3$ from below, increases and crosses the $x$-axis at $x=-2$, reaches the maximum $(-1,\frac1e)$, decreases through the inflection point $(-\frac25,\frac85e^{-5/2})$ and tends to $0$ as $x\to0^-$. To the right of $0$ the graph decreases from $+\infty$ (asymptote $x=0$) to the minimum $(2,4\sqrt e)$, and then increases and approaches the asymptote $y=x+3$ as $x\to+\infty$.
\end{solution}

\subsection*{Question 5a}

\begin{question}
Compute the following integral:
$\displaystyle\int e^{2x}\sin(3x)\,dx$
\end{question}

\begin{hint}
Integrate by parts twice; you will get an equation in which the desired integral appears on both sides.
\end{hint}

\begin{solution}
Denote $I=\int e^{2x}\sin3x\,dx$. Integration by parts with $u=\sin3x$, $dv=e^{2x}dx$ ($v=\frac12e^{2x}$):
\[ I=\frac12e^{2x}\sin3x-\frac32\int e^{2x}\cos3x\,dx. \]
Again by parts, $u=\cos3x$, $dv=e^{2x}dx$:
\[ \int e^{2x}\cos3x\,dx=\frac12e^{2x}\cos3x+\frac32\int e^{2x}\sin3x\,dx=\frac12e^{2x}\cos3x+\frac32I. \]
Hence
\[ I=\frac12e^{2x}\sin3x-\frac34e^{2x}\cos3x-\frac94I\;\Longrightarrow\;\frac{13}{4}I=\frac{e^{2x}}{4}\big(2\sin3x-3\cos3x\big), \]
and finally
\[ \int e^{2x}\sin3x\,dx=\frac{e^{2x}\big(2\sin3x-3\cos3x\big)}{13}+C. \]
(Check: differentiating the result gives $\frac{e^{2x}}{13}\big(4\sin3x-6\cos3x+6\cos3x+9\sin3x\big)=e^{2x}\sin3x$.)
\end{solution}

\subsection*{Question 5b}

\begin{question}
Compute the following integral:
$\displaystyle\int_0^1\frac{x^4}{x^2+2x+1}\,dx$
\end{question}

\begin{hint}
$x^2+2x+1=(x+1)^2$; substitute $u=x+1$ (or perform polynomial division).
\end{hint}

\begin{solution}
$x^2+2x+1=(x+1)^2$. Substitute $u=x+1$, $x=u-1$, $u\in[1,2]$:
\[ \int_0^1\frac{x^4}{(x+1)^2}\,dx=\int_1^2\frac{(u-1)^4}{u^2}\,du=\int_1^2\frac{u^4-4u^3+6u^2-4u+1}{u^2}\,du=\int_1^2\left(u^2-4u+6-\frac4u+\frac1{u^2}\right)du. \]
\[ =\left[\frac{u^3}3-2u^2+6u-4\ln u-\frac1u\right]_1^2=\left(\frac83-8+12-4\ln2-\frac12\right)-\left(\frac13-2+6-1\right)=\left(\frac{37}{6}-4\ln2\right)-\frac{10}{3}=\frac{17}{6}-4\ln2. \]
\textbf{Answer:} $\dfrac{17}{6}-4\ln2\approx0.0607$.
\end{solution}

\subsection*{Question 6}

\begin{question}
Does the following improper integral $\displaystyle\int_0^1\frac{\ln x}{\sqrt x}\,dx$ converge? If so, compute the integral.
\end{question}

\begin{hint}
The problematic point is $x=0$. Compute $\int_\eps^1$ by integration by parts ($u=\ln x$, $dv=x^{-1/2}dx$) and take the limit $\eps\to0^+$.
\end{hint}

\begin{solution}
The function $\frac{\ln x}{\sqrt x}$ is continuous on $(0,1]$ and unbounded near $0$ (it tends to $-\infty$), so this is an improper integral of the second kind, defined as $\lim_{\eps\to0^+}\int_\eps^1\frac{\ln x}{\sqrt x}\,dx$.

By integration by parts, $u=\ln x$, $dv=x^{-1/2}dx$, $du=\frac{dx}x$, $v=2\sqrt x$:
\[ \int_\eps^1\frac{\ln x}{\sqrt x}\,dx=\Big[2\sqrt x\ln x\Big]_\eps^1-\int_\eps^1\frac{2\sqrt x}{x}\,dx=-2\sqrt\eps\ln\eps-\Big[4\sqrt x\Big]_\eps^1=-2\sqrt\eps\ln\eps-4+4\sqrt\eps. \]
Now $\lim_{\eps\to0^+}\sqrt\eps\ln\eps=\lim_{\eps\to0^+}\frac{\ln\eps}{\eps^{-1/2}}$, and by L'Hôpital's rule (form $\frac{-\infty}{\infty}$):
\[ \lim_{\eps\to0^+}\frac{1/\eps}{-\frac12\eps^{-3/2}}=\lim_{\eps\to0^+}\left(-2\sqrt\eps\right)=0. \]
Also $4\sqrt\eps\to0$. Hence the limit exists and is finite: the integral \textbf{converges}, and
\[ \int_0^1\frac{\ln x}{\sqrt x}\,dx=-4. \]
\end{solution}

\section*{2018/19 Semester A Moed B}

\subsection*{Question 1a}

\begin{question}
Find the following limit:
$\displaystyle\lim_{x\to0}\frac{\cos x-\sqrt{1-2x^2}}{\sin^2x}$
\end{question}

\begin{hint}
Use the Maclaurin expansions $\cos x=1-\frac{x^2}2+o(x^2)$, $\sqrt{1+t}=1+\frac t2+o(t)$, and $\sin^2x\sim x^2$.
\end{hint}

\begin{solution}
By the Maclaurin expansions:
\[ \cos x=1-\frac{x^2}{2}+o(x^2),\qquad \sqrt{1-2x^2}=1+\frac12(-2x^2)+o(x^2)=1-x^2+o(x^2). \]
Hence the numerator is $\cos x-\sqrt{1-2x^2}=\frac{x^2}{2}+o(x^2)$. Also $\frac{\sin^2x}{x^2}\to1$. Therefore
\[ \lim_{x\to0}\frac{\cos x-\sqrt{1-2x^2}}{\sin^2x}=\lim_{x\to0}\frac{\frac12+\frac{o(x^2)}{x^2}}{\frac{\sin^2x}{x^2}}=\frac12. \]
(One can also use L'Hôpital's rule: after replacing $\sin^2x$ by $x^2$, $\lim\frac{-\sin x+\frac{2x}{\sqrt{1-2x^2}}}{2x}=\frac{-1+2}{2}=\frac12$.)
\end{solution}

\subsection*{Question 1b}

\begin{question}
Find the following limit:
$\displaystyle\lim_{x\to0}\left(\frac{1-4x^2}{1-3x^2}\right)^{\frac{2}{x^2}}$

Hint: $f^g=e^{g\ln f}$.
\end{question}

\begin{hint}
After passing to the form $e^{g\ln f}$, write $\ln\frac{1-4x^2}{1-3x^2}=\ln(1-4x^2)-\ln(1-3x^2)$ and use the fundamental limit $\lim_{t\to0}\frac{\ln(1+t)}{t}=1$.
\end{hint}

\begin{solution}
This is a limit of the form $1^\infty$. In a neighborhood of $0$ the base is positive, and by the hint
\[ \left(\frac{1-4x^2}{1-3x^2}\right)^{\frac2{x^2}}=\exp\!\left(\frac{2}{x^2}\ln\frac{1-4x^2}{1-3x^2}\right)=\exp\!\left(2\cdot\frac{\ln(1-4x^2)-\ln(1-3x^2)}{x^2}\right). \]
By the fundamental limit $\lim_{t\to0}\frac{\ln(1+t)}{t}=1$:
\[ \frac{\ln(1-4x^2)}{x^2}=-4\cdot\frac{\ln(1-4x^2)}{-4x^2}\to-4,\qquad \frac{\ln(1-3x^2)}{x^2}\to-3. \]
Hence the exponent tends to $2(-4+3)=-2$, and by the continuity of the exponential the limit is $e^{-2}$.
\end{solution}

\subsection*{Question 2}

\begin{question}
Let $f(x)=\ln x$. Write the Taylor polynomial of degree 3 of the function $f$ around the point $x_0=1$. Use this polynomial to compute an approximation of $\ln(1.1)$. Estimate the error of this approximation using Taylor's formula.
\end{question}

\begin{hint}
The remainder in Lagrange form is $R_3(x)=\frac{f^{(4)}(c)}{4!}(x-1)^4$ with $c$ between $1$ and $x$.
\end{hint}

\begin{solution}
\textbf{The polynomial.} $f(x)=\ln x$, $f'(x)=\frac1x$, $f''(x)=-\frac1{x^2}$, $f'''(x)=\frac2{x^3}$, $f^{(4)}(x)=-\frac6{x^4}$. At the point $1$: $f(1)=0$, $f'(1)=1$, $f''(1)=-1$, $f'''(1)=2$. Hence
\[ P_3(x)=(x-1)-\frac{(x-1)^2}{2}+\frac{(x-1)^3}{3}. \]
\textbf{The approximation.} For $x=1.1$, $x-1=0.1$:
\[ \ln(1.1)\approx P_3(1.1)=0.1-\frac{0.01}{2}+\frac{0.001}{3}=0.1-0.005+0.000\overline{3}=0.095\overline{3}. \]
\textbf{Error estimate.} By Taylor's formula with the Lagrange remainder, there exists $c\in(1,1.1)$ such that
\[ R_3(1.1)=\frac{f^{(4)}(c)}{4!}(0.1)^4=-\frac{6}{24c^4}\cdot10^{-4}=-\frac{10^{-4}}{4c^4}. \]
Since $c>1$, $c^4>1$ and hence
\[ |R_3(1.1)|<\frac{10^{-4}}{4}=2.5\cdot10^{-5}. \]
Moreover $R_3<0$, i.e. the approximation $0.09533$ is slightly larger than the true value. (For comparison, $\ln1.1=0.0953102\ldots$, and the actual error is about $2.3\cdot10^{-5}$.)
\end{solution}

\subsection*{Question 3a}

\begin{question}
Is the function
\[ f(x)=\begin{cases}\dfrac{1}{\tan x}-\dfrac1x & x\neq0\\[2mm] 0 & x=0\end{cases} \]
differentiable at the point $x=0$? Explain.
\end{question}

\begin{hint}
Compute the limit of the difference quotient $\frac{f(h)-f(0)}{h}=\frac{h-\tan h}{h^2\tan h}$, for example using the expansion $\tan h=h+\frac{h^3}{3}+o(h^3)$.
\end{hint}

\begin{solution}
(The function is defined in the punctured neighborhood $0<|x|<\frac\pi2$ of $0$, and there $\tan x\ne0$.) By the definition of the derivative:
\[ f'(0)=\lim_{h\to0}\frac{f(h)-f(0)}{h}=\lim_{h\to0}\frac{\frac1{\tan h}-\frac1h}{h}=\lim_{h\to0}\frac{h-\tan h}{h^2\tan h}. \]
By the Maclaurin expansion $\tan h=h+\frac{h^3}{3}+o(h^3)$, hence $h-\tan h=-\frac{h^3}{3}+o(h^3)$. Also $h^2\tan h=h^3\cdot\frac{\tan h}{h}$ and $\frac{\tan h}{h}\to1$. Therefore
\[ f'(0)=\lim_{h\to0}\frac{-\frac13+\frac{o(h^3)}{h^3}}{\frac{\tan h}{h}}=-\frac13. \]
(One can also use L'Hôpital's rule: $\lim\frac{h-\tan h}{h^3}=\lim\frac{1-\frac1{\cos^2h}}{3h^2}=\lim\frac{-\tan^2h}{3h^2}=-\frac13$.)

\textbf{Answer:} yes, $f$ is differentiable at $0$ and $f'(0)=-\frac13$. (In particular $f$ is continuous at $0$: $\lim_{x\to0}f(x)=\lim\frac{x-\tan x}{x\tan x}=0=f(0)$.)
\end{solution}

\subsection*{Question 3b}

\begin{question}
Prove or disprove: if a function $f(x)$ is differentiable at a point $x_0$, then it is continuous at $x_0$.
\end{question}

\begin{hint}
Write $f(x)-f(x_0)=\frac{f(x)-f(x_0)}{x-x_0}\cdot(x-x_0)$.
\end{hint}

\begin{solution}
\textbf{The statement is true.} Suppose $f$ is differentiable at $x_0$, i.e. the limit $f'(x_0)=\lim_{x\to x_0}\frac{f(x)-f(x_0)}{x-x_0}$ exists and is finite. For every $x\ne x_0$ in a neighborhood of $x_0$:
\[ f(x)-f(x_0)=\frac{f(x)-f(x_0)}{x-x_0}\cdot(x-x_0). \]
By the arithmetic of limits (product of two finite limits):
\[ \lim_{x\to x_0}\big(f(x)-f(x_0)\big)=f'(x_0)\cdot0=0, \]
i.e. $\lim_{x\to x_0}f(x)=f(x_0)$, and $f$ is continuous at $x_0$. $\blacksquare$

(The converse is false: $|x|$ is continuous at $0$ but not differentiable there.)
\end{solution}

\subsection*{Question 4}

\begin{question}
Investigate the function $f(x)=\dfrac{x^3}{x^2-4}$ and sketch its graph.

\underline{You need to investigate}: domain, continuity, asymptotes, intervals of increase and decrease, extremum points, intervals of convexity and concavity.
\end{question}

\begin{hint}
Note that the function is odd. Polynomial division gives $f(x)=x+\frac{4x}{x^2-4}$.
\end{hint}

\begin{solution}
\textbf{Domain and continuity.} $x\ne\pm2$. On its domain $f$ is continuous (a rational function). $f(-x)=-f(x)$, i.e. $f$ is odd and the graph is symmetric with respect to the origin. Intersection with the axes: only $(0,0)$.

\textbf{Asymptotes.} Vertical: at $x=\pm2$ the numerator is $\pm8\ne0$ and the denominator vanishes, hence:
\[ \lim_{x\to2^\pm}f(x)=\pm\infty,\qquad \lim_{x\to-2^\pm}f(x)=\pm\infty \]
(for example near $-2$: the numerator is $\approx-8$, and the denominator $x^2-4$ is negative as $x\to-2^+$ and positive as $x\to-2^-$). Hence $x=2$ and $x=-2$ are vertical asymptotes. Oblique: by polynomial division
\[ f(x)=x+\frac{4x}{x^2-4},\qquad \frac{4x}{x^2-4}\xrightarrow[x\to\pm\infty]{}0, \]
hence $y=x$ is an oblique asymptote in both directions.

\textbf{Increase and decrease.}
\[ f'(x)=\frac{3x^2(x^2-4)-x^3\cdot2x}{(x^2-4)^2}=\frac{x^2(x^2-12)}{(x^2-4)^2}. \]
Suspicious points: $x=0$, $x=\pm2\sqrt3$. $f'>0$ for $|x|>2\sqrt3$ -- $f$ is increasing on $(-\infty,-2\sqrt3)$ and on $(2\sqrt3,\infty)$. $f'<0$ for $|x|<2\sqrt3$, $x\ne0,\pm2$ -- $f$ is decreasing on $(-2\sqrt3,-2)$, on $(-2,2)$ (including $x=0$, where $f'=0$ but the sign does not change) and on $(2,2\sqrt3)$.

\textbf{Extrema.} $x=-2\sqrt3$: local maximum, $f(-2\sqrt3)=\frac{-24\sqrt3}{8}=-3\sqrt3$. $x=2\sqrt3$: local minimum, $f(2\sqrt3)=3\sqrt3$. At $x=0$ there is no extremum.

\textbf{Convexity and concavity.}
\[ f''(x)=\frac{8x(x^2+12)}{(x^2-4)^3}. \]
Since $x^2+12>0$, the sign of $f''$ is the sign of $\frac{x}{(x^2-4)^3}$, i.e. the sign of $x(x^2-4)$:
\begin{itemize}
  \item $x<-2$: negative -- concave ($\cap$);
  \item $-2<x<0$: positive -- convex ($\cup$);
  \item $0<x<2$: negative -- concave ($\cap$);
  \item $x>2$: positive -- convex ($\cup$).
\end{itemize}
Inflection point: $(0,0)$ (there $f''$ changes sign and $f$ is defined). At $x=\pm2$ $f$ is not defined.

\textbf{Sketch.} On $(-\infty,-2)$: the graph increases below the asymptote $y=x$ up to the maximum $(-2\sqrt3,-3\sqrt3)\approx(-3.46,-5.2)$ and then decreases to $-\infty$ near $x=-2$. On $(-2,2)$: it decreases from $+\infty$ through the inflection point $(0,0)$ (with a horizontal tangent) to $-\infty$. On $(2,\infty)$: it decreases from $+\infty$ to the minimum $(2\sqrt3,3\sqrt3)\approx(3.46,5.2)$ and then increases and approaches the asymptote $y=x$ from above.
\end{solution}

\subsection*{Question 5}

\begin{question}
Compute the following integral:
\[ \int_0^\pi x^2\sin(3x)\,dx \]
\end{question}

\begin{hint}
Integrate by parts twice, each time differentiating the power of $x$.
\end{hint}

\begin{solution}
Integrate by parts with $u=x^2$, $dv=\sin3x\,dx$ ($du=2x\,dx$, $v=-\frac13\cos3x$):
\[ \int_0^\pi x^2\sin3x\,dx=\left[-\frac{x^2\cos3x}{3}\right]_0^\pi+\frac23\int_0^\pi x\cos3x\,dx. \]
Since $\cos3\pi=-1$: $\left[-\frac{x^2\cos3x}{3}\right]_0^\pi=\frac{\pi^2}{3}$.

Again by parts, $u=x$, $dv=\cos3x\,dx$ ($v=\frac13\sin3x$):
\[ \int_0^\pi x\cos3x\,dx=\left[\frac{x\sin3x}{3}\right]_0^\pi-\frac13\int_0^\pi\sin3x\,dx=0-\frac13\left[-\frac{\cos3x}{3}\right]_0^\pi=-\frac13\cdot\frac{1+1}{3}=-\frac29, \]
(since $\sin3\pi=0$ and $-\cos3\pi+\cos0=2$). Therefore
\[ \int_0^\pi x^2\sin3x\,dx=\frac{\pi^2}{3}+\frac23\cdot\left(-\frac29\right)=\frac{\pi^2}{3}-\frac{4}{27}\approx3.142. \]
\end{solution}

\subsection*{Question 6a}

\begin{question}
Prove that the equation $e^{2x}-x-5=0$ has exactly two real solutions.
\end{question}

\begin{hint}
Check the monotonicity of $g(x)=e^{2x}-x-5$ using the derivative, and its value at the minimum point. Existence -- by the Intermediate Value Theorem; uniqueness in each interval -- by monotonicity.
\end{hint}

\begin{solution}
Define $g(x)=e^{2x}-x-5$, continuous and differentiable on all of $\R$, with $g'(x)=2e^{2x}-1$. $g'(x)=0\iff e^{2x}=\frac12\iff x_0=-\frac{\ln2}{2}$. For $x<x_0$: $g'<0$, $g$ is strictly decreasing; for $x>x_0$: $g'>0$, $g$ is strictly increasing. The minimum value:
\[ g(x_0)=\frac12+\frac{\ln2}{2}-5\approx-4.15<0. \]
\textbf{Existence:} $\lim_{x\to-\infty}g(x)=+\infty$ (since $e^{2x}\to0$ and $-x\to+\infty$), and for example $g(-6)=e^{-12}+1>0$. Also $g(2)=e^4-7>0$. By the Intermediate Value Theorem there is a root in $(-6,x_0)$ and a root in $(x_0,2)$.

\textbf{Uniqueness:} on each of the intervals $(-\infty,x_0]$ and $[x_0,\infty)$ the function is strictly monotone and hence one-to-one, so each contains at most one root (and $x_0$ itself is not a root). Hence the equation has exactly two real solutions.
\end{solution}

\subsection*{Question 6b}

\begin{question}
A curve is given in polar coordinates by $r=1+\cos\varphi$ $(0\le\varphi\le2\pi)$. Compute the length of the curve.
\end{question}

\begin{hint}
The length of a polar curve is $L=\int_\alpha^\beta\sqrt{r^2+(r')^2}\,d\varphi$; use $1+\cos\varphi=2\cos^2\frac\varphi2$.
\end{hint}

\begin{solution}
The curve (a cardioid) is $x=r\cos\varphi$, $y=r\sin\varphi$, and the arc length formula in polar coordinates is
\[ L=\int_0^{2\pi}\sqrt{r^2+\left(\frac{dr}{d\varphi}\right)^2}\,d\varphi. \]
Here $r'=-\sin\varphi$ and hence
\[ r^2+r'^2=(1+\cos\varphi)^2+\sin^2\varphi=2+2\cos\varphi=4\cos^2\frac\varphi2. \]
Therefore $\sqrt{r^2+r'^2}=2\left|\cos\frac\varphi2\right|$. For $\varphi\in[0,2\pi]$, $\frac\varphi2\in[0,\pi]$; the $\cos$ is non-negative on $[0,\pi]$ (in terms of $\varphi$) and non-positive on $[\pi,2\pi]$. By symmetry:
\[ L=2\int_0^{\pi}2\cos\frac\varphi2\,d\varphi=4\left[2\sin\frac\varphi2\right]_0^\pi=8. \]
\textbf{Answer:} the length of the curve is $8$.
\end{solution}

\section*{2018/19 Semester A Special Moed}

\subsection*{Question 1a}

\begin{question}
Find the following limit:
$\displaystyle \lim_{x\to\infty}\left(1+2x-4x^2+x^4\right)^{1/\ln(2x^2+1)}$.
\end{question}

\begin{hint}
Write the expression in the form $e^{\ln(\cdots)/\ln(2x^2+1)}$ and factor the highest power out of each logarithm.
\end{hint}

\begin{solution}
Denote $p(x)=1+2x-4x^2+x^4$. For $x$ large enough $p(x)>0$ (since $p(x)\to\infty$), hence
\[ p(x)^{1/\ln(2x^2+1)} = \exp\left(\frac{\ln p(x)}{\ln(2x^2+1)}\right). \]
Factor the leading power out of each logarithm:
\[ \ln p(x) = \ln\left(x^4\left(1+\tfrac{2}{x^3}-\tfrac{4}{x^2}+\tfrac{1}{x^4}\right)\right) = 4\ln x + \ln\left(1+\tfrac{2}{x^3}-\tfrac{4}{x^2}+\tfrac{1}{x^4}\right), \]
\[ \ln(2x^2+1) = 2\ln x + \ln\left(2+\tfrac{1}{x^2}\right). \]
As $x\to\infty$, the expressions $\ln\left(1+\tfrac{2}{x^3}-\tfrac{4}{x^2}+\tfrac{1}{x^4}\right)\to \ln 1 = 0$ and $\ln\left(2+\tfrac1{x^2}\right)\to\ln 2$ (continuity of the logarithm), whereas $\ln x\to\infty$. Divide the numerator and denominator by $\ln x$:
\[ \frac{\ln p(x)}{\ln(2x^2+1)} = \frac{4 + \frac{\ln\left(1+\frac{2}{x^3}-\frac{4}{x^2}+\frac{1}{x^4}\right)}{\ln x}}{2+\frac{\ln\left(2+\frac{1}{x^2}\right)}{\ln x}} \xrightarrow[x\to\infty]{} \frac{4+0}{2+0}=2. \]
(One can also use L'Hôpital's rule in the case $\frac{\infty}{\infty}$: $\frac{p'(x)/p(x)}{4x/(2x^2+1)}=\frac{(4x^3-8x+2)(2x^2+1)}{4x\,p(x)}\to\frac{8}{4}=2$.)

By the continuity of the exponential function we get
\[ \lim_{x\to\infty}\left(1+2x-4x^2+x^4\right)^{1/\ln(2x^2+1)} = e^{2}. \]
\textbf{Answer:} $e^2$.
\end{solution}

\subsection*{Question 1b}

\begin{question}
Find the following limit:
$\displaystyle \lim_{n\to\infty}\left(\sqrt{(n^2+1)(n^2-4)}-\sqrt{n^4-9}\right)$.
\end{question}

\begin{hint}
Multiply and divide by the conjugate $\sqrt{(n^2+1)(n^2-4)}+\sqrt{n^4-9}$.
\end{hint}

\begin{solution}
This is an expression of the form $\infty-\infty$. Multiply and divide by the conjugate (for $n\ge 2$ both roots are defined and positive):
\[ \sqrt{(n^2+1)(n^2-4)}-\sqrt{n^4-9} = \frac{(n^2+1)(n^2-4)-(n^4-9)}{\sqrt{(n^2+1)(n^2-4)}+\sqrt{n^4-9}}. \]
In the numerator: $(n^2+1)(n^2-4) = n^4-3n^2-4$, so the numerator is $n^4-3n^2-4-n^4+9 = -3n^2+5$. Divide the numerator and denominator by $n^2$:
\[ \frac{-3+\frac{5}{n^2}}{\sqrt{\left(1+\frac{1}{n^2}\right)\left(1-\frac{4}{n^2}\right)}+\sqrt{1-\frac{9}{n^4}}} \xrightarrow[n\to\infty]{} \frac{-3}{1+1} = -\frac{3}{2}, \]
by the arithmetic of limits and the continuity of the square root.

\textbf{Answer:} $-\dfrac32$.
\end{solution}

\subsection*{Question 2}

\begin{question}
Compute the following integral $\displaystyle\int_0^{\pi/4} x\tan^2 x\,dx$.
\end{question}

\begin{hint}
Use the identity $\tan^2 x = \frac{1}{\cos^2 x}-1$ and split into two integrals.
\end{hint}

\begin{hint}
Compute $\int x\cdot\frac{1}{\cos^2 x}\,dx$ by integration by parts, with $u=x$ and $v'=\frac{1}{\cos^2x}$.
\end{hint}

\begin{solution}
By the identity $1+\tan^2 x = \frac{1}{\cos^2 x}$:
\[ \int_0^{\pi/4} x\tan^2 x\,dx = \int_0^{\pi/4} \frac{x}{\cos^2 x}\,dx - \int_0^{\pi/4} x\,dx. \]
(The functions are continuous on the interval $[0,\pi/4]$, so all the integrals exist.)

\textbf{The first integral} -- integration by parts with $u=x$, $v'=\frac{1}{\cos^2x}$, i.e. $u'=1$, $v=\tan x$:
\[ \int_0^{\pi/4}\frac{x}{\cos^2x}\,dx = \Big[x\tan x\Big]_0^{\pi/4} - \int_0^{\pi/4}\tan x\,dx. \]
We have $\int \tan x\,dx = \int\frac{\sin x}{\cos x}dx = -\ln|\cos x|+C$ (substitution $t=\cos x$), hence
\[ \int_0^{\pi/4}\frac{x}{\cos^2x}\,dx = \frac{\pi}{4}\cdot 1 - 0 + \Big[\ln\cos x\Big]_0^{\pi/4} = \frac{\pi}{4} + \ln\frac{\sqrt2}{2} - \ln 1 = \frac{\pi}{4} - \frac{\ln 2}{2}. \]

\textbf{The second integral:} $\displaystyle\int_0^{\pi/4}x\,dx = \frac{x^2}{2}\Big|_0^{\pi/4} = \frac{\pi^2}{32}$.

In summary:
\[ \int_0^{\pi/4} x\tan^2 x\,dx = \frac{\pi}{4} - \frac{\ln 2}{2} - \frac{\pi^2}{32} \approx 0.1304. \]
\end{solution}

\subsection*{Question 3a}

\begin{question}
Write the Maclaurin polynomial of degree 6 for the function $y=\sin\left(\sin(x^2)\right)$.
\end{question}

\begin{hint}
Substitute $u=\sin(x^2)=x^2-\frac{x^6}{6}+o(x^6)$ into the expansion $\sin u = u-\frac{u^3}{6}+o(u^3)$.
\end{hint}

\begin{solution}
We use the known Maclaurin expansions $\sin t = t - \frac{t^3}{6} + o(t^4)$ (as $t\to0$).

First, with $t=x^2$: $\sin(x^2) = x^2 - \frac{x^6}{6} + o(x^8)$.

Denote $u=\sin(x^2)$; then $u\to 0$ as $x\to0$ and $u = O(x^2)$. By the expansion $\sin u = u - \frac{u^3}{6} + o(u^4)$, and since $o(u^4)=o(x^8)$:
\[ u^3 = \left(x^2-\tfrac{x^6}{6}+o(x^8)\right)^3 = x^6 + o(x^6), \]
hence
\[ \sin(\sin(x^2)) = x^2 - \frac{x^6}{6} - \frac{x^6}{6} + o(x^6) = x^2 - \frac{x^6}{3} + o(x^6). \]
By the uniqueness of the Taylor polynomial (a polynomial of degree $\le 6$ satisfying $f(x)-P(x)=o(x^6)$ is the Taylor polynomial), the Maclaurin polynomial of degree 6 is
\[ P_6(x) = x^2 - \frac{x^6}{3}. \]
\end{solution}

\subsection*{Question 3b}

\begin{question}
For which value of $a$ is the function
\[ f(x)=\begin{cases} \dfrac{1}{\ln x}-\dfrac{1}{x-1} & x\neq 1\\[2mm] a & x=1\end{cases} \]
continuous on its domain?
\end{question}

\begin{hint}
You need to compute $\lim_{x\to1}\left(\frac{1}{\ln x}-\frac{1}{x-1}\right)$: bring it to a common denominator and obtain the case $\frac00$.
\end{hint}

\begin{solution}
The domain of $f$ is $x>0$ ($\ln x$ must be defined; at the point $x=1$ the function is defined to be $a$). At every point $x>0$, $x\ne1$, the function is continuous as a composition, difference and quotient of continuous functions with nonzero denominators ($\ln x\ne0$ and $x-1\ne0$). Hence $f$ is continuous on its domain if and only if it is continuous at $x=1$, i.e. if and only if $a=\lim_{x\to1}f(x)$.

We compute: $\displaystyle \frac{1}{\ln x}-\frac{1}{x-1} = \frac{x-1-\ln x}{(x-1)\ln x}$, the case $\frac{0}{0}$. By L'Hôpital's rule (the numerator and denominator are differentiable in a neighborhood of 1):
\[ \lim_{x\to1}\frac{x-1-\ln x}{(x-1)\ln x} = \lim_{x\to1}\frac{1-\frac1x}{\ln x+\frac{x-1}{x}} = \lim_{x\to1}\frac{x-1}{x\ln x + x-1}, \]
again the case $\frac00$, and by L'Hôpital's rule once more:
\[ \lim_{x\to1}\frac{1}{\ln x + 1 + 1} = \frac{1}{2}. \]
(Alternatively, with $x=1+h$: $\ln(1+h)=h-\frac{h^2}{2}+o(h^2)$, and the numerator is $\frac{h^2}{2}+o(h^2)$ and the denominator $h^2+o(h^2)$.)

\textbf{Answer:} $a=\dfrac12$.
\end{solution}

\subsection*{Question 4}

\begin{question}
Investigate the function $f(x)=\dfrac{e^{2x-1}}{x^2}$ and sketch its graph.
You need to find: domain, asymptotes, intervals of increase and decrease, extremum points, intersection points of the graph with the axes, intervals of convexity and concavity.
\end{question}

\begin{hint}
$f'(x)=\frac{2e^{2x-1}(x-1)}{x^3}$. Note that the sign of $x^3$ changes at $x=0$.
\end{hint}

\begin{solution}
\textbf{Domain:} $x\neq 0$, i.e. $(-\infty,0)\cup(0,\infty)$.

\textbf{Intersection with the axes:} $x=0$ is not in the domain, so there is no intersection with the $y$-axis. Since $e^{2x-1}>0$ and $x^2>0$, we have $f(x)>0$ for every $x$ in the domain, so there is no intersection with the $x$-axis.

\textbf{Asymptotes:}
\begin{itemize}
  \item \emph{Vertical:} $\lim_{x\to0^\pm}\frac{e^{2x-1}}{x^2} = \frac{e^{-1}}{0^+} = +\infty$, hence $x=0$ is a vertical asymptote (the graph rises to $+\infty$ on both sides).
  \item \emph{At $-\infty$:} $e^{2x-1}\to0$ and $\frac1{x^2}\to0$, hence $\lim_{x\to-\infty}f(x)=0$, and $y=0$ is a horizontal asymptote on the left.
  \item \emph{At $+\infty$:} $\frac{f(x)}{x}=\frac{e^{2x-1}}{x^3}\to\infty$ (the exponential dominates every power, for example by applying L'Hôpital's rule three times), so there is no horizontal or oblique asymptote on the right, and $f(x)\to+\infty$.
\end{itemize}

\textbf{First derivative:} by the quotient rule,
\[ f'(x) = \frac{2e^{2x-1}x^2 - e^{2x-1}\cdot 2x}{x^4} = \frac{2e^{2x-1}(x-1)}{x^3}. \]
$f'(x)=0 \iff x=1$. The sign of $f'$ is determined by $\frac{x-1}{x^3}$:
\begin{itemize}
  \item $x<0$: $x-1<0$, $x^3<0$ and hence $f'>0$ -- $f$ is \textbf{increasing} on $(-\infty,0)$.
  \item $0<x<1$: $x-1<0$, $x^3>0$ and hence $f'<0$ -- $f$ is \textbf{decreasing} on $(0,1)$.
  \item $x>1$: $f'>0$ -- $f$ is \textbf{increasing} on $(1,\infty)$.
\end{itemize}
Hence at $x=1$ there is a \textbf{local minimum}: $f(1)=\frac{e^{1}}{1}=e$, i.e. the point $(1,e)$. There are no other extremum points ($x=0$ is not in the domain).

\textbf{Second derivative:} differentiate $f'(x)=2e^{2x-1}\cdot\frac{x-1}{x^3}$:
\[ \left(\frac{x-1}{x^3}\right)' = \frac{x^3-3x^2(x-1)}{x^6} = \frac{-2x+3}{x^4}, \]
\[ f''(x) = 2e^{2x-1}\left(\frac{2(x-1)}{x^3}+\frac{3-2x}{x^4}\right) = \frac{2e^{2x-1}\left(2x^2-4x+3\right)}{x^4}. \]
The quadratic $2x^2-4x+3$ has discriminant $16-24<0$ and a positive leading coefficient, so it is positive for every $x$. Also $e^{2x-1}>0$ and $x^4>0$. Hence $f''(x)>0$ for every $x\ne0$:
$f$ is \textbf{convex} (concave up, $\cup$) on each of the intervals $(-\infty,0)$ and $(0,\infty)$, and there are no inflection points.

\textbf{Sketch:} on the left the graph starts close to the $x$-axis (above it), increases, is convex, and tends to $+\infty$ as $x\to0^-$. To the right of $0$ the graph decreases from $+\infty$ to the minimum $(1,e)$ and then increases at an exponential rate to $+\infty$. The whole graph lies above the $x$-axis.
\end{solution}

\subsection*{Question 5}

\begin{question}
The equation of the curve in polar coordinates is: $r=2(1-\sin\varphi)$.
Sketch the curve and compute the area of the region bounded by the curve.
\end{question}

\begin{hint}
Area of a region in polar coordinates: $S=\frac12\int_\alpha^\beta r^2(\varphi)\,d\varphi$. Here $r\ge0$ for every $\varphi$, so the curve closes up as $\varphi$ runs over $[0,2\pi]$.
\end{hint}

\begin{solution}
\textbf{Sketch:} we have $0\le r=2(1-\sin\varphi)\le 4$ for every $\varphi$, and $r$ is periodic with period $2\pi$, so the curve is closed and is obtained for $\varphi\in[0,2\pi]$. A few values:
\[ \varphi=0:\ r=2;\quad \varphi=\tfrac{\pi}{2}:\ r=0;\quad \varphi=\pi:\ r=2;\quad \varphi=\tfrac{3\pi}{2}:\ r=4. \]
That is, the curve passes through the points $(2,0)$, $(0,0)$, $(-2,0)$, $(0,-4)$ (in Cartesian coordinates). Since $r(\pi-\varphi)=r(\varphi)$, the curve is symmetric with respect to the $y$-axis. This is a \textbf{cardioid} (``heart'') whose cusp is at the origin (when $\varphi=\frac\pi2$) and which ``opens'' downwards: the farthest point is $(0,-4)$.

\textbf{Area:} by the formula for area in polar coordinates, $S=\frac12\int_0^{2\pi}r^2\,d\varphi$:
\[ S = \frac12\int_0^{2\pi}4(1-\sin\varphi)^2\,d\varphi = 2\int_0^{2\pi}\left(1-2\sin\varphi+\sin^2\varphi\right)d\varphi. \]
We compute each term separately:
\[ \int_0^{2\pi}1\,d\varphi=2\pi,\qquad \int_0^{2\pi}\sin\varphi\,d\varphi = \Big[-\cos\varphi\Big]_0^{2\pi}=0, \]
\[ \int_0^{2\pi}\sin^2\varphi\,d\varphi = \int_0^{2\pi}\frac{1-\cos2\varphi}{2}\,d\varphi = \Big[\frac{\varphi}{2}-\frac{\sin2\varphi}{4}\Big]_0^{2\pi} = \pi. \]
Hence
\[ S = 2\left(2\pi - 0 + \pi\right) = 6\pi. \]
\textbf{Answer:} the area of the region is $6\pi$.
\end{solution}

\subsection*{Question 6a}

\begin{question}
Does the following improper integral $\displaystyle\int_{-\infty}^{-2}\frac{dx}{x^2-2x}$ converge? If so, find it.
\end{question}

\begin{hint}
Decompose into partial fractions: $\frac{1}{x(x-2)}=\frac12\left(\frac{1}{x-2}-\frac1x\right)$.
\end{hint}

\begin{solution}
On the interval $(-\infty,-2]$ the denominator $x^2-2x=x(x-2)>0$ (both factors are negative), so the function is continuous and positive there, and the only problem is the limit $-\infty$. By definition,
\[ \int_{-\infty}^{-2}\frac{dx}{x^2-2x} = \lim_{R\to-\infty}\int_R^{-2}\frac{dx}{x(x-2)}. \]
Partial fraction decomposition: $\frac{1}{x(x-2)}=\frac{A}{x}+\frac{B}{x-2}$, and from $1=A(x-2)+Bx$ we get ($x=0$) $A=-\frac12$ and ($x=2$) $B=\frac12$. Hence
\[ \int\frac{dx}{x(x-2)} = \frac12\ln|x-2|-\frac12\ln|x| + C = \frac12\ln\left|\frac{x-2}{x}\right| + C. \]
Therefore
\[ \int_R^{-2}\frac{dx}{x(x-2)} = \frac12\ln\frac{-4}{-2} - \frac12\ln\left|\frac{R-2}{R}\right| = \frac12\ln2-\frac12\ln\left|1-\frac2R\right|. \]
As $R\to-\infty$, $1-\frac2R\to1$ and by the continuity of the logarithm $\ln\left|1-\frac2R\right|\to0$. Hence the limit exists and is finite:
\[ \int_{-\infty}^{-2}\frac{dx}{x^2-2x} = \frac{\ln2}{2}. \]
\textbf{Answer:} the integral converges and its value is $\frac12\ln2$.
(The convergence can also be seen from the limit comparison test with $\frac1{x^2}$, but the direct computation also gives the value.)
\end{solution}

\subsection*{Question 6b}

\begin{question}
Compute the following integral: $\displaystyle\int_{\ln2}^{\ln3}\frac{dx}{e^x-e^{-x}}$.
\end{question}

\begin{hint}
Multiply the numerator and denominator by $e^x$ and substitute $t=e^x$.
\end{hint}

\begin{solution}
On the interval $[\ln2,\ln3]$ we have $e^x>e^{-x}$, so the integrand is continuous and the integral is proper. Multiply the numerator and denominator by $e^x$:
\[ \int_{\ln2}^{\ln3}\frac{dx}{e^x-e^{-x}} = \int_{\ln2}^{\ln3}\frac{e^x\,dx}{e^{2x}-1}. \]
Substitute $t=e^x$, $dt=e^x\,dx$; the limits: $x=\ln2\mapsto t=2$, $x=\ln3\mapsto t=3$:
\[ = \int_2^3\frac{dt}{t^2-1} = \int_2^3\frac12\left(\frac{1}{t-1}-\frac{1}{t+1}\right)dt = \frac12\Big[\ln\frac{t-1}{t+1}\Big]_2^3 = \frac12\left(\ln\frac12-\ln\frac13\right) = \frac12\ln\frac32. \]
\textbf{Answer:} $\dfrac12\ln\dfrac32\approx0.2027$.
\end{solution}

\section*{2018/19 Semester B Moed A}

\subsection*{Question 1a}

\begin{question}
Compute the derivative of the function $f(x)=\sqrt{5x+3}$ at the point $x$ using the definition of the derivative (an answer based on the table of derivatives will not be accepted).
\end{question}

\begin{hint}
Multiply the numerator and denominator by the conjugate $\sqrt{5(x+h)+3}+\sqrt{5x+3}$.
\end{hint}

\begin{solution}
The function is defined for $x\ge -\tfrac35$, and we compute the derivative at a point $x>-\tfrac35$. By the definition, multiplying by the conjugate:
\begin{align*}
f'(x)&=\lim_{h\to0}\frac{\sqrt{5(x+h)+3}-\sqrt{5x+3}}{h}
=\lim_{h\to0}\frac{\big(5(x+h)+3\big)-(5x+3)}{h\left(\sqrt{5x+5h+3}+\sqrt{5x+3}\right)}\\
&=\lim_{h\to0}\frac{5}{\sqrt{5x+5h+3}+\sqrt{5x+3}}=\frac{5}{2\sqrt{5x+3}},
\end{align*}
where in the last step we used the continuity of the square root function ($\sqrt{5x+5h+3}\to\sqrt{5x+3}>0$). Hence
\[ f'(x)=\frac{5}{2\sqrt{5x+3}},\qquad x>-\tfrac35. \]
(At the point $x=-\tfrac35$ the quotient $\frac{\sqrt{5h}}{h}\to+\infty$, so the function is not differentiable there.)
\end{solution}

\subsection*{Question 1b}

\begin{question}
Find all the asymptotes of the function: $f(x)=\dfrac{x^2+x-6}{|x|-2}$.
\end{question}

\begin{hint}
Factor the numerator $x^2+x-6=(x+3)(x-2)$ and treat the cases $x>0$ and $x<0$ separately.
\end{hint}

\begin{solution}
Domain: $|x|\neq2$, i.e. $x\neq\pm2$. Factor $x^2+x-6=(x+3)(x-2)$.

\textbf{For $x\ge0$, $x\ne 2$:} $|x|-2=x-2$ and hence $f(x)=\dfrac{(x+3)(x-2)}{x-2}=x+3$.

\textbf{For $x<0$, $x\ne-2$:} $|x|-2=-x-2=-(x+2)$ and hence
\[ f(x)=-\frac{(x+3)(x-2)}{x+2}=-\frac{x^2+x-6}{x+2}=-x+1+\frac{4}{x+2}, \]
(polynomial division: $x^2+x-6=(x+2)(x-1)-4$).

\textbf{Vertical asymptotes:} at $x=2$: $\lim_{x\to2}f(x)=\lim_{x\to2}(x+3)=5$ is finite, so this is a removable discontinuity and there is no asymptote there. At $x=-2$: $\lim_{x\to-2^\pm}\left(-x+1+\frac{4}{x+2}\right)=\pm\infty$, hence $x=-2$ is a vertical asymptote.

\textbf{Oblique/horizontal asymptotes:} as $x\to+\infty$ we have exactly $f(x)=x+3$, hence $y=x+3$ is an oblique asymptote at $+\infty$. As $x\to-\infty$:
\[ f(x)-(-x+1)=\frac{4}{x+2}\xrightarrow[x\to-\infty]{}0, \]
hence $y=-x+1$ is an oblique asymptote at $-\infty$. (There are no horizontal asymptotes.)

\textbf{Answer:} $x=-2$ (vertical), $y=x+3$ (at $+\infty$), $y=-x+1$ (at $-\infty$).
\end{solution}

\subsection*{Question 1c}

\begin{question}
What is the slope of the tangent to the curve $y=\int_0^{x+1}e^{-2t^2}\,dt$ at its point with $x=0$.
\end{question}

\begin{hint}
Use the Fundamental Theorem of Calculus together with the chain rule.
\end{hint}

\begin{solution}
Denote $F(u)=\int_0^u e^{-2t^2}\,dt$. The function $e^{-2t^2}$ is continuous, so by the Fundamental Theorem $F'(u)=e^{-2u^2}$. Hence $y(x)=F(x+1)$ and by the chain rule
\[ y'(x)=F'(x+1)\cdot 1=e^{-2(x+1)^2}. \]
The slope of the tangent at the point $x=0$ is $y'(0)=e^{-2}=\dfrac{1}{e^2}$.
\end{solution}

\subsection*{Question 2a}

\begin{question}
Prove that the equation $x^2-\ln x=0$ has no real solutions.
\end{question}

\begin{hint}
Find the minimum value of $g(x)=x^2-\ln x$ on the domain $x>0$ and show that it is positive.
\end{hint}

\begin{solution}
Define $g(x)=x^2-\ln x$, which is defined and differentiable on the domain $x>0$ (so the equation has no solutions with $x\le0$). Differentiate:
\[ g'(x)=2x-\frac1x=\frac{2x^2-1}{x}. \]
For $x>0$: $g'(x)<0$ when $0<x<\frac{1}{\sqrt2}$ and $g'(x)>0$ when $x>\frac1{\sqrt2}$. Hence $g$ is strictly decreasing on $(0,\frac1{\sqrt2}]$ and strictly increasing on $[\frac1{\sqrt2},\infty)$, and $x=\frac1{\sqrt2}$ is a point of absolute minimum of $g$ on the domain:
\[ g\!\left(\tfrac1{\sqrt2}\right)=\frac12-\ln\frac1{\sqrt2}=\frac12+\frac{\ln2}{2}>0. \]
Hence $g(x)\ge\frac{1+\ln 2}{2}>0$ for every $x>0$, and in particular $g(x)\ne0$: the equation has no real solutions.
\end{solution}

\subsection*{Question 2b}

\begin{question}
Compute the area enclosed between the curves $y=x\cdot e^x$, $y=x-1$, $x=0$, $x=2$. Draw the corresponding sketch.
\end{question}

\begin{hint}
Check which of the graphs lies above the other on the interval $[0,2]$, and compute $\int xe^x\,dx$ by integration by parts.
\end{hint}

\begin{solution}
For every $x\in[0,2]$ we have $e^x\ge1$ and hence $xe^x\ge x>x-1$. That is, on the interval the graph of $y=xe^x$ lies above the line $y=x-1$ (the sketch: the line $y=x-1$ passes through $(0,-1)$ and $(2,1)$, and the curve $y=xe^x$ passes through $(0,0)$ and $(2,2e^2)$; the region is bounded on the left and right by the lines $x=0$, $x=2$). The area:
\[ S=\int_0^2\big(xe^x-(x-1)\big)\,dx. \]
By integration by parts ($u=x$, $dv=e^xdx$): $\int xe^x\,dx=xe^x-\int e^x\,dx=(x-1)e^x+C$. Hence
\[ \int_0^2 xe^x\,dx=\big[(x-1)e^x\big]_0^2=e^2-(-1)=e^2+1, \qquad \int_0^2(x-1)\,dx=\Big[\tfrac{x^2}{2}-x\Big]_0^2=0. \]
\textbf{Answer:} $S=e^2+1$.
\end{solution}

\subsection*{Question 3a}

\begin{question}
Compute:
\begin{enumerate}
  \item[a.] (8 pts) $\displaystyle\int\frac{(x+1)^2}{x(x+2)^2}\,dx$
  \item[b.] (7 pts) $\displaystyle\int_{-\pi/4}^{\pi/4}\left(\frac{1}{\cos^2x\cdot(2+\tan x)^2}+x\cos^3x\right)dx$
\end{enumerate}
\end{question}

\begin{hint}
In part (a) decompose into partial fractions of the form $\frac{A}{x}+\frac{B}{x+2}+\frac{C}{(x+2)^2}$.
\end{hint}

\begin{hint}
In part (b), the integral of an odd function over a symmetric interval is $0$; in the second part substitute $u=\tan x$.
\end{hint}

\begin{solution}
\begin{enumerate}
  \item[a.] The degree of the numerator (2) is less than the degree of the denominator (3), so we decompose directly into partial fractions:
  \[ \frac{(x+1)^2}{x(x+2)^2}=\frac{A}{x}+\frac{B}{x+2}+\frac{C}{(x+2)^2}
  \;\Longrightarrow\; (x+1)^2=A(x+2)^2+Bx(x+2)+Cx. \]
  Substituting $x=0$: $1=4A$, i.e. $A=\frac14$. Substituting $x=-2$: $1=-2C$, i.e. $C=-\frac12$. Comparing the coefficients of $x^2$: $1=A+B$, hence $B=\frac34$. Therefore
  \[ \int\frac{(x+1)^2}{x(x+2)^2}\,dx=\frac14\ln|x|+\frac34\ln|x+2|+\frac{1}{2(x+2)}+C. \]
  \item[b.] Split into two integrals. The function $x\cos^3x$ is odd (the product of the odd function $x$ and the even function $\cos^3x$) and continuous, so its integral over the symmetric interval $[-\frac\pi4,\frac\pi4]$ equals $0$.

  In the first integral substitute $u=\tan x$, $du=\frac{dx}{\cos^2x}$; as $x$ runs from $-\frac\pi4$ to $\frac\pi4$, $u$ runs from $-1$ to $1$ (and on this interval $2+\tan x\ge1>0$):
  \[ \int_{-\pi/4}^{\pi/4}\frac{dx}{\cos^2x\,(2+\tan x)^2}=\int_{-1}^{1}\frac{du}{(2+u)^2}=\left[-\frac{1}{2+u}\right]_{-1}^{1}=-\frac13+1=\frac23. \]
  \textbf{Answer:} $\frac23$.
\end{enumerate}
\end{solution}

\subsection*{Question 3b}

\begin{question}
Write the Maclaurin polynomial of order 2 for the function $f(x)=\cos(2x)-\arctan(x+1)$.
\end{question}

\begin{hint}
Expand each summand separately: for $\cos(2x)$ use the known expansion of $\cos t$, and for $\arctan(x+1)$ compute $g(0),g'(0),g''(0)$.
\end{hint}

\begin{solution}
The Maclaurin polynomial of order 2: $P_2(x)=f(0)+f'(0)x+\frac{f''(0)}{2}x^2$. We compute separately for each summand.

For $\cos(2x)$: by the expansion $\cos t=1-\frac{t^2}{2}+o(t^2)$ we get $\cos 2x=1-2x^2+o(x^2)$.

For $g(x)=\arctan(x+1)$:
\[ g(0)=\arctan1=\frac\pi4,\quad g'(x)=\frac{1}{1+(x+1)^2}\Rightarrow g'(0)=\frac12,\quad g''(x)=-\frac{2(x+1)}{\big(1+(x+1)^2\big)^2}\Rightarrow g''(0)=-\frac24=-\frac12. \]
Hence $\arctan(x+1)=\frac\pi4+\frac x2-\frac{x^2}{4}+o(x^2)$. Subtracting:
\[ P_2(x)=\left(1-\frac\pi4\right)-\frac12x-\frac74x^2. \]
\end{solution}

\subsection*{Question 4a}

\begin{question}
For which values of the parameters $a$ and $b$ will the function
\[ f(x)=\begin{cases}\dfrac{1-\cos x}{x\ln(1+x)} & -1<x<0\\[2mm] ax+b & 0\le x\le1\\[1mm] x^{1/(x-1)} & x>1\end{cases} \]
be continuous on the domain $x>-1$?
\end{question}

\begin{hint}
On each of the open intervals the function is continuous as a composition of elementary functions; only $x=0$ and $x=1$ need to be checked.
\end{hint}

\begin{hint}
Compute the limit $\lim_{x\to1^+}x^{1/(x-1)}$ using $x^{1/(x-1)}=e^{\frac{\ln x}{x-1}}$.
\end{hint}

\begin{solution}
On the intervals $(-1,0)$, $(0,1)$, $(1,\infty)$ the function is continuous (quotient/composition of continuous functions, and the denominator $x\ln(1+x)\ne0$ on $(-1,0)$; on $(1,\infty)$: $x^{1/(x-1)}=e^{\ln x/(x-1)}$). It remains to check the junction points.

\textbf{The point $x=0$:} $f(0)=b=\lim_{x\to0^+}f(x)$. From the left, using the known limits $\lim_{x\to0}\frac{1-\cos x}{x^2}=\frac12$ and $\lim_{x\to0}\frac{\ln(1+x)}{x}=1$:
\[ \lim_{x\to0^-}\frac{1-\cos x}{x\ln(1+x)}=\lim_{x\to0^-}\frac{1-\cos x}{x^2}\cdot\frac{x}{\ln(1+x)}=\frac12\cdot1=\frac12. \]
Hence continuity at $0$ is equivalent to $b=\frac12$.

\textbf{The point $x=1$:} $f(1)=a+b=\lim_{x\to1^-}f(x)$. From the right:
\[ \lim_{x\to1^+}x^{1/(x-1)}=\lim_{x\to1^+}e^{\frac{\ln x}{x-1}}. \]
The limit $\lim_{x\to1}\frac{\ln x}{x-1}$ is of the form $\frac00$, and by L'Hôpital's rule $=\lim_{x\to1}\frac{1/x}{1}=1$ (this is also the derivative of $\ln$ at $1$). By the continuity of the exponential the limit is $e^1=e$. Hence continuity at $1$ is equivalent to $a+b=e$.

\textbf{Answer:} $b=\frac12$, $a=e-\frac12$.
\end{solution}

\subsection*{Question 4b}

\begin{question}
Compute $\displaystyle\int_1^{+\infty}\frac{\ln x}{x^4}\,dx$.
\end{question}

\begin{hint}
Integrate by parts with $u=\ln x$, $dv=x^{-4}dx$ and then let the upper limit tend to infinity.
\end{hint}

\begin{solution}
By the definition of the improper integral, $\int_1^\infty=\lim_{R\to\infty}\int_1^R$. By integration by parts with $u=\ln x$, $dv=x^{-4}dx$ ($du=\frac{dx}{x}$, $v=-\frac{1}{3x^3}$):
\[ \int_1^R\frac{\ln x}{x^4}\,dx=\left[-\frac{\ln x}{3x^3}\right]_1^R+\frac13\int_1^R\frac{dx}{x^4}=-\frac{\ln R}{3R^3}+\frac13\left[-\frac{1}{3x^3}\right]_1^R=-\frac{\ln R}{3R^3}+\frac19\left(1-\frac{1}{R^3}\right). \]
As $R\to\infty$: $\frac{\ln R}{R^3}\to0$ (by L'Hôpital's rule: $\lim\frac{1/R}{3R^2}=0$) and $\frac1{R^3}\to0$. Hence the integral converges and
\[ \int_1^{+\infty}\frac{\ln x}{x^4}\,dx=\frac19. \]
\end{solution}

\section*{2018/19 Semester B Moed B}

\subsection*{Question 1a}

\begin{question}
Compute the derivative of the function $f(x)=\sin(3x+4)$ at the point $x$ using the definition of the derivative (an answer based on the table of derivatives will not be accepted).
\end{question}

\begin{hint}
Use the identity $\sin\alpha-\sin\beta=2\cos\frac{\alpha+\beta}{2}\sin\frac{\alpha-\beta}{2}$ and the limit $\lim_{t\to0}\frac{\sin t}{t}=1$.
\end{hint}

\begin{solution}
By the definition and using the identity $\sin\alpha-\sin\beta=2\cos\frac{\alpha+\beta}{2}\sin\frac{\alpha-\beta}{2}$:
\begin{align*}
f'(x)&=\lim_{h\to0}\frac{\sin(3x+3h+4)-\sin(3x+4)}{h}
=\lim_{h\to0}\frac{2\cos\!\left(3x+4+\frac{3h}{2}\right)\sin\frac{3h}{2}}{h}\\
&=\lim_{h\to0}3\cos\!\left(3x+4+\tfrac{3h}{2}\right)\cdot\frac{\sin\frac{3h}{2}}{\frac{3h}{2}}=3\cos(3x+4)\cdot1,
\end{align*}
where we used the continuity of $\cos$ and the fundamental limit $\lim_{t\to0}\frac{\sin t}{t}=1$. Hence $f'(x)=3\cos(3x+4)$.
\end{solution}

\subsection*{Question 1b}

\begin{question}
Find all the asymptotes of the function: $f(x)=\dfrac{x\sqrt{x^2-1}}{2x^2-1}$.
\end{question}

\begin{hint}
Pay attention to the domain $|x|\ge1$ and to the fact that $\sqrt{x^2}=|x|$.
\end{hint}

\begin{solution}
\textbf{Domain:} we need $x^2-1\ge0$, i.e. $|x|\ge1$, and also $2x^2-1\ne0$, i.e. $x\ne\pm\frac1{\sqrt2}$ -- but these points are not in the domain anyway. Hence the domain is $(-\infty,-1]\cup[1,\infty)$.

\textbf{Vertical asymptotes:} on the domain the denominator satisfies $2x^2-1\ge1>0$, so $f$ is continuous on its whole domain (also at the endpoints $x=\pm1$, where $f(\pm1)=0$). There is no point at which $f$ tends to infinity, hence there are no vertical asymptotes.

\textbf{Horizontal asymptotes:} since $\sqrt{x^2-1}=|x|\sqrt{1-\frac1{x^2}}$,
\[ f(x)=\frac{x|x|\sqrt{1-\frac1{x^2}}}{x^2\left(2-\frac1{x^2}\right)}=\frac{|x|}{x}\cdot\frac{\sqrt{1-\frac1{x^2}}}{2-\frac1{x^2}}. \]
As $x\to+\infty$: $\frac{|x|}{x}=1$ and hence $f(x)\to\frac12$. As $x\to-\infty$: $\frac{|x|}{x}=-1$ and hence $f(x)\to-\frac12$. Since the limits are finite, there are no additional oblique asymptotes.

\textbf{Answer:} $y=\frac12$ (at $+\infty$) and $y=-\frac12$ (at $-\infty$); there are no vertical asymptotes.
\end{solution}

\subsection*{Question 1c}

\begin{question}
For which values of $a$ does the curve $y=x^2+a\sin x$ have inflection points?
\end{question}

\begin{hint}
Check when $y''$ changes sign.
\end{hint}

\begin{solution}
$y'=2x+a\cos x$, $y''=2-a\sin x$. An inflection point is a point at which $y''$ changes sign.
\begin{itemize}
  \item If $|a|<2$: $|a\sin x|\le|a|<2$ and hence $y''>0$ for every $x$ -- there are no inflection points.
  \item If $a=2$: $y''=2(1-\sin x)\ge0$, and it vanishes only at isolated points; the sign does not change -- no inflection. Similarly for $a=-2$: $y''=2(1+\sin x)\ge0$.
  \item If $|a|>2$: the equation $\sin x=\frac2a$ with $0<\left|\frac2a\right|<1$ has solutions $x_0$, and at them $\cos x_0\ne0$. Hence $(y'')'(x_0)=-a\cos x_0\ne0$, i.e. $y''$ passes through $0$ in a strictly monotone way and changes sign at $x_0$ -- an inflection point.
\end{itemize}
\textbf{Answer:} $|a|>2$.
\end{solution}

\subsection*{Question 2a}

\begin{question}
Prove that the equation $\dfrac{1}{x+1}+\dfrac1x+\dfrac{1}{x-1}=1$ has at least two solutions.
\end{question}

\begin{hint}
Define $g(x)=\frac{1}{x+1}+\frac1x+\frac1{x-1}-1$ and examine the one-sided limits of $g$ at the endpoints of the intervals $(-1,0)$ and $(0,1)$. Use the Intermediate Value Theorem.
\end{hint}

\begin{solution}
Define $g(x)=\frac{1}{x+1}+\frac1x+\frac1{x-1}-1$, which is continuous on every interval not containing $-1,0,1$.

\textbf{The interval $(-1,0)$:} as $x\to-1^+$, the term $\frac1{x+1}\to+\infty$ while the other terms are bounded, hence $g(x)\to+\infty$. As $x\to0^-$, $\frac1x\to-\infty$ and hence $g(x)\to-\infty$. Therefore there exist points $-1<p<q<0$ with $g(p)>0$ and $g(q)<0$. $g$ is continuous on $[p,q]$, so by the Intermediate Value Theorem there exists $x_1\in(p,q)$ with $g(x_1)=0$.

\textbf{The interval $(0,1)$:} as $x\to0^+$, $\frac1x\to+\infty$ and hence $g\to+\infty$; as $x\to1^-$, $\frac1{x-1}\to-\infty$ and hence $g\to-\infty$. In the same way, by the Intermediate Value Theorem there exists $x_2\in(0,1)$ with $g(x_2)=0$.

The intervals are disjoint, hence $x_1\ne x_2$ and the equation has at least two solutions. (In the same way there is a third solution in $(1,\infty)$: $g\to+\infty$ as $x\to1^+$ and $g\to-1$ as $x\to\infty$.)
\end{solution}

\subsection*{Question 2b}

\begin{question}
Compute the volume of the solid obtained by revolving the plane region
\[ 0\le x\le2,\quad 0\le y\le\sqrt{\frac{x}{x+2}} \]
about the $x$-axis. Draw the corresponding sketch.
\end{question}

\begin{hint}
The volume of a solid of revolution about the $x$-axis is $V=\pi\int_a^b y^2\,dx$.
\end{hint}

\begin{solution}
The sketch: the function $y=\sqrt{\frac{x}{x+2}}$ increases from $y(0)=0$ to $y(2)=\sqrt{\frac12}$; the region is the area between the graph and the $x$-axis over the interval $[0,2]$, and the solid is generated by revolving it about the $x$-axis. By the formula for the volume of a solid of revolution:
\[ V=\pi\int_0^2\frac{x}{x+2}\,dx=\pi\int_0^2\left(1-\frac{2}{x+2}\right)dx=\pi\Big[x-2\ln(x+2)\Big]_0^2=\pi\big(2-2\ln4+2\ln2\big). \]
\textbf{Answer:} $V=\pi(2-2\ln2)=2\pi(1-\ln2)\approx1.928$.
\end{solution}

\subsection*{Question 3a}

\begin{question}
Compute:
\begin{enumerate}
  \item[a.] (7 pts) $\displaystyle\int\frac{x^2-8}{x^3+4x^2}\,dx$
  \item[b.] (8 pts) $\displaystyle\int_0^1\ln(1+\sqrt x)\,dx$
\end{enumerate}
\end{question}

\begin{hint}
In part (a) factor $x^3+4x^2=x^2(x+4)$ and decompose into partial fractions. In part (b) substitute $t=\sqrt x$ and then integrate by parts.
\end{hint}

\begin{solution}
\begin{enumerate}
  \item[a.] $x^3+4x^2=x^2(x+4)$, and the degree of the numerator is less than the degree of the denominator. We decompose:
  \[ \frac{x^2-8}{x^2(x+4)}=\frac Ax+\frac B{x^2}+\frac C{x+4}\;\Longrightarrow\;x^2-8=Ax(x+4)+B(x+4)+Cx^2. \]
  Substituting $x=0$: $-8=4B$, $B=-2$. Substituting $x=-4$: $8=16C$, $C=\frac12$. Coefficients of $x^2$: $1=A+C$, $A=\frac12$. Hence
  \[ \int\frac{x^2-8}{x^3+4x^2}\,dx=\frac12\ln|x|+\frac2x+\frac12\ln|x+4|+C. \]
  \item[b.] Substitute $x=t^2$, $dx=2t\,dt$, $t\in[0,1]$:
  \[ \int_0^1\ln(1+\sqrt x)\,dx=\int_0^1 2t\ln(1+t)\,dt. \]
  Integrating by parts ($u=\ln(1+t)$, $dv=2t\,dt$, $v=t^2$):
  \[ =\Big[t^2\ln(1+t)\Big]_0^1-\int_0^1\frac{t^2}{1+t}\,dt=\ln2-\int_0^1\left(t-1+\frac1{1+t}\right)dt=\ln2-\left(\frac12-1+\ln2\right)=\frac12. \]
  \textbf{Answer:} $\frac12$.
\end{enumerate}
\end{solution}

\subsection*{Question 3b}

\begin{question}
Write the Taylor polynomial of order 2 for the function
\[ f(x)=\int_0^x\frac{1-t}{1+t^2}\,dt \]
about the point $x=1$.
\end{question}

\begin{hint}
By the Fundamental Theorem, $f'(x)=\frac{1-x}{1+x^2}$; compute $f(1)$ directly from the integral.
\end{hint}

\begin{solution}
The Taylor polynomial of order 2 about $1$: $P_2(x)=f(1)+f'(1)(x-1)+\frac{f''(1)}{2}(x-1)^2$.

The function $\frac{1-t}{1+t^2}$ is continuous on all of $\R$, hence by the Fundamental Theorem of Calculus $f'(x)=\frac{1-x}{1+x^2}$, and therefore $f'(1)=0$. Differentiating again:
\[ f''(x)=\frac{-(1+x^2)-(1-x)\cdot2x}{(1+x^2)^2}=\frac{x^2-2x-1}{(1+x^2)^2},\qquad f''(1)=\frac{-2}{4}=-\frac12. \]
Finally,
\[ f(1)=\int_0^1\frac{dt}{1+t^2}-\int_0^1\frac{t\,dt}{1+t^2}=\Big[\arctan t\Big]_0^1-\Big[\tfrac12\ln(1+t^2)\Big]_0^1=\frac\pi4-\frac{\ln2}{2}. \]
\textbf{Answer:}
\[ P_2(x)=\frac\pi4-\frac{\ln2}{2}-\frac14(x-1)^2. \]
\end{solution}

\subsection*{Question 4a}

\begin{question}
For which values of the parameters $a$ and $b$ will the function
\[ f(x)=\begin{cases}\dfrac{\sin(3x)}{\ln(1-x)} & x<0\\[2mm] ax+b & 0\le x\le1\\[1mm] \big(\cos(x-1)\big)^{1/(x-1)} & x>1\end{cases} \]
be continuous for every real $x$?
\end{question}

\begin{hint}
Continuity needs to be checked only at the junction points $x=0$ and $x=1$. Write the right-hand limit at $1$ in the form $e^{\frac{\ln\cos(x-1)}{x-1}}$.
\end{hint}

\begin{solution}
For $x<0$ we have $1-x>1$ and hence $\ln(1-x)>0$, so on the interval $(-\infty,0)$ the function is continuous as a quotient of continuous functions with a nonzero denominator. On $(0,1)$ it is a polynomial. For $x>1$ the expression $\big(\cos(x-1)\big)^{1/(x-1)}=e^{\frac{\ln\cos(x-1)}{x-1}}$ is defined and continuous when $\cos(x-1)>0$, in particular on the interval $1<x<1+\frac\pi2$, which is the neighborhood relevant to continuity at $1$. (Remark: for $x>1$ with $\cos(x-1)\le0$ the expression is not defined in the real sense, so the requirement ``continuous for every $x$'' should be understood as continuity on the domain.) We check the junction points.

\textbf{The point $x=0$:} $f(0)=b=\lim_{x\to0^+}f(x)$. From the left:
\[ \lim_{x\to0^-}\frac{\sin3x}{\ln(1-x)}=\lim_{x\to0^-}\frac{\sin3x}{3x}\cdot\frac{-x}{\ln(1-x)}\cdot(-3)=1\cdot1\cdot(-3)=-3, \]
by $\lim_{t\to0}\frac{\sin t}{t}=1$ and $\lim_{t\to0}\frac{\ln(1+t)}{t}=1$ (with $t=-x$). Hence $b=-3$.

\textbf{The point $x=1$:} $f(1)=a+b=\lim_{x\to1^-}f(x)$. From the right, let $u=x-1\to0^+$:
\[ \big(\cos u\big)^{1/u}=e^{\frac{\ln\cos u}{u}},\qquad \lim_{u\to0^+}\frac{\ln\cos u}{u}\overset{\text{L'H}}{=}\lim_{u\to0^+}\frac{-\tan u}{1}=0. \]
By continuity of the exponential, the limit is $e^0=1$. Hence $a+b=1$.

\textbf{Answer:} $b=-3$, $a=4$.
\end{solution}

\subsection*{Question 4b}

\begin{question}
Compute $\displaystyle\int_{0.25}^{+\infty}\frac{dx}{(4x+1)\sqrt x}$.
\end{question}

\begin{hint}
Substitute $t=\sqrt x$.
\end{hint}

\begin{solution}
Substitute $t=\sqrt x$, $x=t^2$, $dx=2t\,dt$; when $x=0.25$, $t=\frac12$, and when $x\to\infty$, $t\to\infty$:
\[ \int_{0.25}^{R}\frac{dx}{(4x+1)\sqrt x}=\int_{1/2}^{\sqrt R}\frac{2t\,dt}{(4t^2+1)t}=\int_{1/2}^{\sqrt R}\frac{2\,dt}{1+(2t)^2}=\Big[\arctan(2t)\Big]_{1/2}^{\sqrt R}=\arctan(2\sqrt R)-\frac\pi4. \]
As $R\to\infty$, $\arctan(2\sqrt R)\to\frac\pi2$. Hence the integral converges and
\[ \int_{0.25}^{+\infty}\frac{dx}{(4x+1)\sqrt x}=\frac\pi2-\frac\pi4=\frac\pi4. \]
\end{solution}

\section*{2019/20 Semester A Quiz}

\subsection*{Question 1}

\begin{question}
The limit $\displaystyle\lim_{x\to 0}\frac{\sin ax}{\sin bx}\,\frac{x}{\ln(1+cx)}$ is equal to
\begin{enumerate}
  \item $\frac{ac}{b}$
  \item $\frac{a}{bc}$
  \item $\frac{bc}{a}$
  \item $\frac{b}{ac}$
  \item the limit does not exist
\end{enumerate}
\end{question}

\begin{hint}
Break the expression into a product of basic limits: $\frac{\sin t}{t}\to 1$ and $\frac{\ln(1+t)}{t}\to 1$ as $t\to 0$.
\end{hint}

\begin{solution}
\textbf{The correct answer: (b)} $\frac{a}{bc}$ (under the implicit assumption in the question that $a,b,c\neq 0$; otherwise the expression is not defined or the limit is different).

We multiply and divide so that basic limits appear. For $x$ close to 0 and different from it:
\[
\frac{\sin ax}{\sin bx}\cdot\frac{x}{\ln(1+cx)}
=\frac{\sin ax}{ax}\cdot\frac{bx}{\sin bx}\cdot\frac{cx}{\ln(1+cx)}\cdot\frac{ax\cdot x}{bx\cdot cx}
=\frac{\sin ax}{ax}\cdot\frac{bx}{\sin bx}\cdot\frac{cx}{\ln(1+cx)}\cdot\frac{a}{bc}.
\]
As $x\to 0$ also $ax,bx,cx\to 0$, hence by the basic limit $\lim_{t\to 0}\frac{\sin t}{t}=1$ (with the substitution $t=ax$, $t=bx$) we have
$\frac{\sin ax}{ax}\to 1$ and $\frac{bx}{\sin bx}\to 1$, and by the basic limit $\lim_{t\to 0}\frac{\ln(1+t)}{t}=1$ (from the formula sheet, $\lim_{t\to0}\frac{\log_a(1+t)}{t}=\frac{1}{\ln a}$ with base $e$) we have $\frac{cx}{\ln(1+cx)}\to 1$.
By the arithmetic of limits:
\[
\lim_{x\to 0}\frac{\sin ax}{\sin bx}\,\frac{x}{\ln(1+cx)}=1\cdot 1\cdot 1\cdot\frac{a}{bc}=\boxed{\frac{a}{bc}}.
\]

\textbf{Why the others are wrong:} (a), (c), (d) are other combinations of $a,b,c$, which differ from $\frac{a}{bc}$ in general: for example, for $a=2,\ b=3,\ c=5$ the limit is $\frac{2}{15}$, whereas (a) gives $\frac{10}{3}$, (c) gives $\frac{15}{2}$ and (d) gives $\frac{3}{10}$. (e) is wrong because we showed that the limit exists.
\end{solution}

\subsection*{Question 2}

\begin{question}
The limit $\displaystyle\lim_{x\to\infty}\left(\frac{x^2-1}{x^2+3x}\right)^{\frac{x}{3}}$ is equal to
\begin{enumerate}
  \item $e^3$
  \item $e$
  \item $e^{\frac13}$
  \item $e^{-1}$
  \item $\infty$
\end{enumerate}
\end{question}

\begin{hint}
This is a limit of the form $1^\infty$. Write the base as $1+t(x)$ with $t(x)\to 0$ and use $\lim_{y\to\infty}\left(1+\frac1y\right)^y=e$.
\end{hint}

\begin{solution}
\textbf{The correct answer: (d)} $e^{-1}$.

The base tends to 1 (dividing by $x^2$: $\frac{1-1/x^2}{1+3/x}\to 1$) and the exponent tends to $\infty$, i.e. this is a limit of the form $1^\infty$. We write
\[
\frac{x^2-1}{x^2+3x}=1+\frac{x^2-1-x^2-3x}{x^2+3x}=1+t(x),\qquad t(x)=\frac{-3x-1}{x^2+3x}\xrightarrow[x\to\infty]{}0 .
\]
For large $x$, $t(x)\neq 0$, hence
\[
\left(1+t(x)\right)^{\frac x3}=\Bigl[\left(1+t(x)\right)^{\frac{1}{t(x)}}\Bigr]^{t(x)\cdot\frac{x}{3}} .
\]
By the basic limit $\lim_{s\to 0}(1+s)^{1/s}=e$ (substituting $s=t(x)\to 0$) the inner base tends to $e$. The exponent:
\[
t(x)\cdot\frac x3=\frac{-3x-1}{x^2+3x}\cdot\frac{x}{3}=\frac{-3x^2-x}{3x^2+9x}=\frac{-3-\frac1x}{3+\frac9x}\xrightarrow[x\to\infty]{}-1 .
\]
Since $u^v=e^{v\ln u}$ and the functions $\ln$ and $\exp$ are continuous, if $u\to e$ and $v\to -1$ then $u^v\to e^{-1}$. Hence
\[
\lim_{x\to\infty}\left(\frac{x^2-1}{x^2+3x}\right)^{\frac{x}{3}}=\boxed{e^{-1}} .
\]

\textbf{Why the others are wrong:} (a), (b), (c) result from an error in the sign or in the coefficient of the exponent (for example, ignoring the fact that $t(x)$ is negative). (e) is wrong: the base is less than 1 for large $x$ (the numerator is smaller than the denominator), hence the power is even bounded by 1 and cannot tend to $\infty$.
\end{solution}

\subsection*{Question 3}

\begin{question}
The derivative of $e^{2x}\ln\sqrt{1+x}$ is equal to
\begin{enumerate}
  \item $e^{2x}\ln\sqrt{1+x}+e^{2x}\frac{1}{2(1+x)}$
  \item $2e^{2x}\ln\sqrt{1+x}+e^{2x}\frac{1}{2(1+x)}$
  \item $2e^{2x}\ln\sqrt{1+x}+e^{2x}\frac{1}{(1+x)}$
  \item $e^{2x}\ln\sqrt{1+x}+e^{2x}\frac{1}{(1+x)}$
  \item $2e^{2x}\sqrt{1+x}+2e^{2x}\frac{1}{(1+x)}$
\end{enumerate}
\end{question}

\begin{solution}
\textbf{The correct answer: (b)}.

By the product rule $(fg)'=f'g+fg'$ with $f(x)=e^{2x}$, $g(x)=\ln\sqrt{1+x}$ (defined for $x>-1$).

By the chain rule: $f'(x)=e^{2x}\cdot 2=2e^{2x}$.

For $g$ we use the fact that $\ln\sqrt{1+x}=\frac12\ln(1+x)$, hence $g'(x)=\frac12\cdot\frac{1}{1+x}=\frac{1}{2(1+x)}$.
(Alternatively, by the chain rule: $g'(x)=\frac{1}{\sqrt{1+x}}\cdot\frac{1}{2\sqrt{1+x}}=\frac{1}{2(1+x)}$.)

Hence
\[
\left(e^{2x}\ln\sqrt{1+x}\right)'=\boxed{2e^{2x}\ln\sqrt{1+x}+e^{2x}\frac{1}{2(1+x)}} .
\]

\textbf{Why the others are wrong:} (a) and (d) forget the factor 2 coming from the derivative of the exponent $2x$ (chain rule). (c) and (d) forget the factor $\frac12$ that comes from the square root. (e) replaces $\ln\sqrt{1+x}$ by $\sqrt{1+x}$ and the coefficient of the second term by 2, and this is not the derivative.
\end{solution}

\subsection*{Question 4}

\begin{question}
The derivative $y'(0)$ of the implicit function $y(x)$ defined by the equation
$\cos\left(\frac{\pi}{2}+x+y\right)-y=0$ in a neighborhood of the point $(0,0)$ is equal to
\begin{enumerate}
  \item $\pi/2$
  \item $-\frac12$
  \item $1+\pi/2$
  \item $\frac12$
  \item $0$
\end{enumerate}
\end{question}

\begin{hint}
Differentiate both sides of the equation with respect to $x$, where $y=y(x)$ (chain rule), and then substitute $x=0,\ y=0$.
\end{hint}

\begin{solution}
\textbf{The correct answer: (b)} $-\frac12$.

First, the point $(0,0)$ satisfies the equation: $\cos\frac{\pi}{2}-0=0$.
Differentiate both sides of the equation $\cos\left(\frac{\pi}{2}+x+y(x)\right)-y(x)=0$ with respect to $x$, using the chain rule:
\[
-\sin\left(\frac{\pi}{2}+x+y\right)\cdot\left(1+y'\right)-y'=0 .
\]
Substitute $x=0$, $y(0)=0$: $\sin\frac{\pi}{2}=1$, hence
\[
-(1+y'(0))-y'(0)=0\quad\Longrightarrow\quad -1-2y'(0)=0\quad\Longrightarrow\quad y'(0)=\boxed{-\frac12}.
\]
(Check: by the identity $\cos(\frac{\pi}{2}+\theta)=-\sin\theta$ the equation is $-\sin(x+y)=y$; differentiating gives $-\cos(x+y)(1+y')=y'$, and at the point $(0,0)$ again $y'(0)=-\frac12$.)

\textbf{Why the others are wrong:} (d) results from a sign error in the derivative of $\cos$. (e) results from forgetting to differentiate the $y$ inside the cosine or the term $-y$. (a) and (c) result from confusing the value of the angle $\frac{\pi}{2}$ with the derivative; the derivative is a single value which we computed, $-\frac12$.
\end{solution}

\subsection*{Question 5}

\begin{question}
The function $f(x)$ is continuous at the point $x=x_0$. Which statement is true?
\begin{enumerate}
  \item $f(x)$ is differentiable at the point $x=x_0$
  \item the limit $\displaystyle\lim_{x\to x_0}\frac{f(x)-f(x_0)}{x-x_0}$ exists
  \item the limit $\displaystyle\lim_{x\to x_0}\left(f(x)-f(x_0)\right)$ exists
  \item the limit $\displaystyle\lim_{h\to 0}\frac{f(x_0+h)-f(x_0)}{h}$ exists
  \item $f'(x)$ is continuous at the point $x=x_0$
\end{enumerate}
\end{question}

\begin{hint}
Continuity does not imply differentiability. Think of $f(x)=|x|$ at the point $x_0=0$.
\end{hint}

\begin{solution}
\textbf{The correct answer: (c)}.

By definition, $f$ is continuous at $x_0$ if and only if $\lim_{x\to x_0}f(x)=f(x_0)$. By the arithmetic of limits (the limit of a constant is the constant itself) this is equivalent to
\[
\lim_{x\to x_0}\bigl(f(x)-f(x_0)\bigr)=f(x_0)-f(x_0)=0 ,
\]
and in particular the limit in statement (c) exists (and equals 0).

\textbf{Why the others are wrong:} consider $f(x)=|x|$ and $x_0=0$. The function is continuous at 0, but
\[
\frac{f(0+h)-f(0)}{h}=\frac{|h|}{h}=\begin{cases}1 & h>0\\ -1 & h<0\end{cases}
\]
hence the one-sided limits of the difference quotient are $1$ and $-1$ and the limit does not exist. Therefore:
\begin{itemize}
  \item (b) and (d) are wrong: these are two ways of writing the definition of the derivative $f'(x_0)$, and we showed that the limit does not exist.
  \item (a) is wrong: differentiability at $x_0$ means exactly the existence of the limit in (b)/(d).
  \item (e) is wrong: $f'(0)$ is not even defined, hence $f'$ is not continuous at 0. (Moreover, even a differentiable function can have a discontinuous derivative, for example $x^2\sin\frac1x$ with $f(0)=0$.)
\end{itemize}
\end{solution}

\section*{2019/20 Semester A Moed A}

\subsection*{Part A, Question 1}

\begin{question}
Investigate the function $y=\frac{e^{2x}}{x}$ according to the following items:
\begin{enumerate}
  \item Domain, intersection points with the axes, evenness, oddness
  \item Continuity, points of discontinuity
  \item Asymptotes
  \item Intervals of increase and decrease, extremum points
  \item Intervals of convexity and concavity, inflection points
  \item Graphical description.
\end{enumerate}
\end{question}

\begin{hint}
$y'=\frac{e^{2x}(2x-1)}{x^2}$ and $y''=\frac{2e^{2x}(2x^2-2x+1)}{x^3}$. Check the sign of each factor.
\end{hint}

\begin{solution}
\begin{enumerate}
  \item \textbf{Domain:} $x\neq0$, i.e. $(-\infty,0)\cup(0,\infty)$.

  \textbf{Intersection with the axes:} $x=0$ is not in the domain --- there is no intersection with the $y$-axis. $e^{2x}>0$ always, hence $y\ne0$ --- there is no intersection with the $x$-axis. (The sign: $y<0$ for $x<0$ and $y>0$ for $x>0$.)

  \textbf{Evenness:} $y(-x)=\frac{e^{-2x}}{-x}$. For example $y(1)=e^2$, $y(-1)=-e^{-2}$, hence $y(-1)\ne y(1)$ and also $y(-1)\ne -y(1)$. The function is \textbf{neither even nor odd}.

  \item The function is a quotient of continuous functions ($e^{2x}$ and $x$) and hence is continuous at every point where the denominator is nonzero, i.e. on its entire domain. The only problematic point is $x=0$ (where the function is not defined):
  \[ \lim_{x\to0^-}\frac{e^{2x}}{x}=\frac{1}{0^-}=-\infty,\qquad \lim_{x\to0^+}\frac{e^{2x}}{x}=\frac{1}{0^+}=+\infty, \]
  hence $x=0$ is a discontinuity of the second kind (essential).

  \item \textbf{Vertical asymptote:} $x=0$ (by the limits in the previous item).

  \textbf{At $-\infty$:} $e^{2x}\to0$ and $\frac1x\to0$, hence $\lim_{x\to-\infty}y=0$ --- a horizontal asymptote $y=0$ on the left (the graph is below the axis).

  \textbf{At $+\infty$:} $\frac{y}{x}=\frac{e^{2x}}{x^2}\to\infty$ (by L'Hôpital twice: $\lim\frac{e^{2x}}{x^2}=\lim\frac{2e^{2x}}{2x}=\lim\frac{4e^{2x}}{2}=\infty$), hence there is no horizontal or oblique asymptote on the right.

  \item By the quotient rule:
  \[ y'=\frac{2e^{2x}\cdot x - e^{2x}}{x^2} = \frac{e^{2x}(2x-1)}{x^2}. \]
  Since $e^{2x}>0$ and $x^2>0$, the sign of $y'$ is the sign of $2x-1$:
  \begin{itemize}
    \item $y'<0$ for $x<0$ and for $0<x<\frac12$: the function is \textbf{decreasing} on $(-\infty,0)$ and on $(0,\frac12)$.
    \item $y'>0$ for $x>\frac12$: the function is \textbf{increasing} on $(\frac12,\infty)$.
  \end{itemize}
  At $x=\frac12$ the derivative changes sign from $-$ to $+$, hence this is a \textbf{local minimum} point: $y(\tfrac12)=\frac{e}{1/2}=2e$, i.e. $(\frac12,2e)$. There are no other extremum points.

  \item Differentiate $y'=e^{2x}\cdot\frac{2x-1}{x^2}$:
  \[ \left(\frac{2x-1}{x^2}\right)'=\frac{2x^2-2x(2x-1)}{x^4}=\frac{2-2x}{x^3}, \]
  \[ y''=e^{2x}\left(\frac{2(2x-1)}{x^2}+\frac{2-2x}{x^3}\right)=\frac{e^{2x}\left(4x^2-2x+2-2x\right)}{x^3}=\frac{2e^{2x}\left(2x^2-2x+1\right)}{x^3}. \]
  The trinomial $2x^2-2x+1$ has discriminant $4-8<0$ and a positive leading coefficient, hence it is always positive. Therefore the sign of $y''$ is the sign of $x^3$, i.e. of $x$:
  \begin{itemize}
    \item $x<0$: $y''<0$ --- the function is \textbf{concave} ($\cap$) on $(-\infty,0)$.
    \item $x>0$: $y''>0$ --- the function is \textbf{convex} ($\cup$) on $(0,\infty)$.
  \end{itemize}
  The sign of $y''$ changes only at $x=0$, which is not in the domain, hence there are \textbf{no inflection points}.

  \item \textbf{Graphical description:} to the left of the $y$-axis the graph is below the $x$-axis: it starts close to $y=0$ at $-\infty$, decreases (and is concave) and tends to $-\infty$ as $x\to0^-$. To the right of the $y$-axis the graph is above the $x$-axis: it decreases from $+\infty$ (near $x=0^+$) to the minimum point $(\frac12,2e)\approx(0.5,\,5.44)$, and then increases (convexly) to $+\infty$ at an exponential rate.
\end{enumerate}
\end{solution}

\subsection*{Part A, Question 2}

\begin{question}
Find the largest value and the smallest value of the function $y=\sqrt[3]{x}(1+2x)$ on the interval $[-8,1]$.
\end{question}

\begin{hint}
The candidates are the endpoints of the interval, points where $y'=0$, and points where $y$ is not differentiable (pay attention to $x=0$).
\end{hint}

\begin{solution}
The function $y=x^{1/3}+2x^{4/3}$ is continuous on the closed interval $[-8,1]$, hence by \textbf{Weierstrass's theorem} it attains a maximum and a minimum there. These points are at the endpoints of the interval or at interior critical points (where $y'=0$ or $y'$ does not exist), by Fermat's theorem.

\textbf{Derivative} (for $x\ne0$):
\[ y' = \frac13x^{-2/3}+\frac83x^{1/3} = \frac{1+8x}{3\sqrt[3]{x^2}}. \]
\begin{itemize}
  \item $y'=0\iff x=-\frac18\in(-8,1)$.
  \item At $x=0$ the derivative does not exist: $\frac{y(h)-y(0)}{h}=\frac{h^{1/3}(1+2h)}{h}=\frac{1+2h}{h^{2/3}}\to+\infty$. Hence $x=0$ is a critical point.
\end{itemize}

\textbf{Values of the function at the candidates:}
\begin{align*}
y(-8) &= \sqrt[3]{-8}\,(1-16) = (-2)(-15) = 30,\\
y\left(-\tfrac18\right) &= \sqrt[3]{-\tfrac18}\left(1-\tfrac14\right) = -\tfrac12\cdot\tfrac34 = -\tfrac38,\\
y(0) &= 0,\\
y(1) &= 1\cdot 3 = 3.
\end{align*}

\textbf{Answer:} the largest value is $30$ (attained at $x=-8$), and the smallest value is $-\frac38$ (attained at $x=-\frac18$).
\end{solution}

\subsection*{Part A, Question 3}

\begin{question}
Find the integral
\[ \int\left(\frac{x^2-x-4}{(x-1)(x^2+2x+1)}+e^{2x}x^2\right)dx. \]
\end{question}

\begin{hint}
Note that $x^2+2x+1=(x+1)^2$, and decompose into partial fractions of the form $\frac{A}{x-1}+\frac{B}{x+1}+\frac{C}{(x+1)^2}$.
\end{hint}

\begin{hint}
Compute $\int x^2e^{2x}dx$ by integrating by parts twice (differentiating the polynomial).
\end{hint}

\begin{solution}
By linearity of the integral we compute each term separately.

\textbf{The first term.} $x^2+2x+1=(x+1)^2$. The numerator has degree 2 and the denominator has degree 3, hence we decompose into partial fractions:
\[ \frac{x^2-x-4}{(x-1)(x+1)^2} = \frac{A}{x-1}+\frac{B}{x+1}+\frac{C}{(x+1)^2}, \]
i.e. $x^2-x-4 = A(x+1)^2+B(x-1)(x+1)+C(x-1)$.
\begin{itemize}
  \item $x=1$: $1-1-4=4A$, hence $A=-1$.
  \item $x=-1$: $1+1-4=-2C$, hence $C=1$.
  \item Comparing coefficients of $x^2$: $1=A+B$, hence $B=2$.
\end{itemize}
(Check, constant term: $A-B-C=-1-2-1=-4$ $\checkmark$.) Hence
\[ \int\frac{x^2-x-4}{(x-1)(x+1)^2}dx = -\ln|x-1| + 2\ln|x+1| - \frac{1}{x+1}+C_1. \]

\textbf{The second term.} Integration by parts with $u=x^2$, $v'=e^{2x}$ ($v=\frac12e^{2x}$):
\[ \int x^2e^{2x}dx = \frac{x^2}{2}e^{2x} - \int xe^{2x}dx. \]
Again by parts with $u=x$, $v'=e^{2x}$:
\[ \int xe^{2x}dx = \frac{x}{2}e^{2x}-\frac12\int e^{2x}dx = \frac x2e^{2x}-\frac14e^{2x}+C_2. \]
Hence
\[ \int x^2e^{2x}dx = e^{2x}\left(\frac{x^2}{2}-\frac{x}{2}+\frac14\right)+C_3. \]

\textbf{Answer:}
\[ \int\left(\frac{x^2-x-4}{(x-1)(x^2+2x+1)}+e^{2x}x^2\right)dx = -\ln|x-1|+2\ln|x+1|-\frac{1}{x+1}+e^{2x}\left(\frac{x^2}{2}-\frac{x}{2}+\frac14\right)+C. \]
(Check: differentiating the result gives back the integrand.)
\end{solution}

\subsection*{Part B, Question 1}

\begin{question}
Let $f(x)$ be continuous on the interval $[a,b]$. Which of the following must be true?
\begin{enumerate}
  \item[(a)] $f(x)$ is not differentiable on the interval $(a,b)$
  \item[(b)] $f(x)$ satisfies the conditions of Lagrange's theorem on the interval $[a,b]$
  \item[(c)] $f(x)$ satisfies the conditions of Rolle's theorem on the interval $[a,b]$
  \item[(d)] $f(x)$ is differentiable on the interval $(a,b)$
  \item[(e)] there exists $M$ such that $f(x)<M$ for every $x$ belonging to $[a,b]$
\end{enumerate}
\end{question}

\begin{hint}
Which theorem guarantees something for every continuous function on a closed interval?
\end{hint}

\begin{solution}
\textbf{The correct answer: (e).}

By \textbf{Weierstrass's first theorem}, a function continuous on a closed interval $[a,b]$ is bounded there. Hence there exists $K$ such that $f(x)\le K$ for every $x\in[a,b]$, and then $M=K+1$ satisfies $f(x)<M$ for every $x\in[a,b]$.

The other answers need not hold:
\begin{itemize}
  \item (a) is false: for example $f(x)=x$ is continuous and differentiable on $(a,b)$.
  \item (b) and (d) are false: Lagrange's theorem also requires differentiability on $(a,b)$, and a continuous function need not be differentiable. For example, on the interval $[-1,1]$ the function $f(x)=|x|$ is continuous but not differentiable at $0$.
  \item (c) is false: Rolle's theorem requires both differentiability and $f(a)=f(b)$; for example $f(x)=x$ is continuous but $f(a)\ne f(b)$.
\end{itemize}
\end{solution}

\subsection*{Part B, Question 2}

\begin{question}
Let $f(x)$ be differentiable on the interval $(a,b)$. Which of the following must be true?
\begin{enumerate}
  \item[(a)] $f'(x)$ is continuous on $(a,b)$
  \item[(b)] $f(x)$ is monotone on $(a,b)$
  \item[(c)] $f(x)+\sin(x)$ is continuous on $(a,b)$
  \item[(d)] $f(x)$ is monotone increasing on $(a,b)$
  \item[(e)] there exists $c\in(a,b)$ such that $f'(c)=0$
\end{enumerate}
\end{question}

\begin{hint}
What does differentiability at a point imply?
\end{hint}

\begin{solution}
\textbf{The correct answer: (c).}

A function differentiable at a point is continuous there, hence $f$ is continuous on $(a,b)$. Also $\sin x$ is continuous, and a sum of continuous functions is continuous; hence $f(x)+\sin(x)$ is continuous on $(a,b)$.

The other answers need not hold:
\begin{itemize}
  \item (a) is false: the derivative of a differentiable function need not be continuous. The classical example: $f(x)=x^2\sin\frac1x$ for $x\neq0$ and $f(0)=0$, on the interval $(-1,1)$. It is differentiable at every point ($f'(0)=\lim_{h\to0}h\sin\frac1h=0$), but $f'(x)=2x\sin\frac1x-\cos\frac1x$ for $x\ne0$, and $f'$ has no limit at $0$.
  \item (b), (d) are false: for example $f(x)=x^2$ on the interval $(-1,1)$ is differentiable but not monotone.
  \item (e) is false: for example $f(x)=x$ is differentiable and $f'(x)=1\ne0$ for every $x$. (Rolle's theorem also requires continuity on a closed interval and $f(a)=f(b)$.)
\end{itemize}
\end{solution}

\subsection*{Part B, Question 3}

\begin{question}
The derivative $y'(0)$ of the implicit function $y(x)$ defined by the equation $\cos\left(\frac{\pi}{2}+x+y\right)-y=0$ in a neighborhood of the point $(0,0)$ is equal to
\begin{enumerate}
  \item[(a)] $\frac12$
  \item[(b)] $0$
  \item[(c)] $\frac{\pi}{2}$
  \item[(d)] $1+\frac{\pi}{2}$
  \item[(e)] $-\frac12$
\end{enumerate}
\end{question}

\begin{hint}
Differentiate both sides of the equation with respect to $x$, where $y=y(x)$ (chain rule), and substitute $x=0$, $y=0$.
\end{hint}

\begin{solution}
\textbf{The correct answer: (e).}

First, the point $(0,0)$ satisfies the equation: $\cos\frac{\pi}{2}-0=0$. Differentiate both sides of the equation with respect to $x$, where $y=y(x)$, using the chain rule:
\[ -\sin\left(\frac{\pi}{2}+x+y\right)\cdot\left(1+y'\right) - y' = 0. \]
Substitute $x=0$, $y(0)=0$; then $\sin\frac{\pi}{2}=1$:
\[ -(1+y'(0)) - y'(0) = 0 \quad\Longrightarrow\quad -1-2y'(0)=0 \quad\Longrightarrow\quad y'(0) = -\frac12. \]
(Alternatively: $\cos\left(\frac\pi2+t\right)=-\sin t$, hence the equation is $\sin(x+y)+y=0$; differentiating gives $\cos(x+y)(1+y')+y'=0$, and at $(0,0)$: $1+2y'=0$.)
\end{solution}

\subsection*{Part B, Question 4}

\begin{question}
If $f(x)=\frac{1-e^{-\frac{x^2}{2}}}{x}$, then in a neighborhood of $0$ we have
\begin{enumerate}
  \item[(a)] $f(x)=\frac{x}{2}+\frac{x^3}{4}+o(x^3)$
  \item[(b)] $f(x)=\frac{x}{2}+\frac{x^3}{8}+o(x^3)$
  \item[(c)] $f(x)=\frac{x}{2}-\frac{x^3}{8}+o(x^3)$
  \item[(d)] $f(x)=\frac{x}{2}-\frac{x^3}{2}+o(x^3)$
  \item[(e)] $f(x)=\frac{x}{2}-\frac{x^3}{4}+o(x^3)$
\end{enumerate}
\end{question}

\begin{hint}
Substitute $t=-\frac{x^2}{2}$ into the expansion $e^t=1+t+\frac{t^2}{2}+o(t^2)$.
\end{hint}

\begin{solution}
\textbf{The correct answer: (c).}

By the Maclaurin expansion $e^t=1+t+\frac{t^2}{2}+o(t^2)$ as $t\to0$, with $t=-\frac{x^2}{2}$ (and then $o(t^2)=o(x^4)$):
\[ e^{-x^2/2} = 1-\frac{x^2}{2}+\frac{x^4}{8}+o(x^4). \]
Hence
\[ 1-e^{-x^2/2} = \frac{x^2}{2}-\frac{x^4}{8}+o(x^4), \]
and dividing by $x$ (for $x\neq0$; $\frac{o(x^4)}{x}=o(x^3)$):
\[ f(x)=\frac{x}{2}-\frac{x^3}{8}+o(x^3). \]
\end{solution}

\subsection*{Part B, Question 5}

\begin{question}
The area bounded by the curves $y(x)=5x-x^2$ and $y(x)=2x$ is equal to
\begin{enumerate}
  \item[(a)] $\frac92$
  \item[(b)] $1$
  \item[(c)] $\frac{45}{2}$
  \item[(d)] $\frac{25}{6}$
  \item[(e)] $\frac{27}{2}$
\end{enumerate}
\end{question}

\begin{hint}
Find the intersection points of the curves, and determine which of them is on top between the intersection points.
\end{hint}

\begin{solution}
\textbf{The correct answer: (a).}

Intersection points: $5x-x^2=2x\iff x^2-3x=0\iff x=0$ or $x=3$.
On the interval $(0,3)$ we have $(5x-x^2)-2x = 3x-x^2 = x(3-x)>0$, i.e. the parabola is above the line. Hence
\[ S=\int_0^3\left[(5x-x^2)-2x\right]dx = \int_0^3(3x-x^2)\,dx = \left[\frac{3x^2}{2}-\frac{x^3}{3}\right]_0^3 = \frac{27}{2}-9 = \frac92. \]
\end{solution}

\section*{2019/20 Semester A Moed B}

\subsection*{Part A, Question 1}

\begin{question}
Investigate the function $y=\frac{x^3-1}{4x^2}$ according to the following items:
\begin{enumerate}
  \item Domain, intersection points with the axes, evenness, oddness
  \item Continuity, points of discontinuity
  \item Asymptotes
  \item Intervals of increase and decrease, extremum points
  \item Intervals of convexity and concavity, inflection points
  \item Graphical description.
\end{enumerate}
\end{question}

\begin{hint}
It is convenient to write the function in the form $y=\frac{x}{4}-\frac{1}{4x^2}$: from it one immediately sees the oblique asymptote, and differentiation is simple.
\end{hint}

\begin{solution}
Denote $f(x)=\frac{x^3-1}{4x^2}=\frac{x}{4}-\frac{1}{4x^2}$.
\begin{enumerate}
  \item \textbf{Domain:} the denominator vanishes only at $x=0$, hence $D=\{x\in\R : x\neq 0\}$.

  \textbf{Intersection with the axes:} $x=0$ is not in the domain, hence there is no intersection with the $y$-axis. Intersection with the $x$-axis: $x^3-1=0 \iff x=1$, i.e. the point $(1,0)$.

  \textbf{Evenness:} $f(-x)=\frac{-x^3-1}{4x^2}$. For example $f(1)=0$ whereas $f(-1)=-\frac12$, so that $f(-1)\neq f(1)$ and also $f(-1)\neq -f(1)$. Hence the function is neither even nor odd.

  \item $f$ is a quotient of polynomials, hence continuous on its entire domain. The only point where $f$ is not defined is $x=0$. There the numerator tends to $-1$ and the denominator $4x^2\to 0^+$, hence
  \[ \lim_{x\to 0^-} f(x)=\lim_{x\to 0^+} f(x)=-\infty , \]
  i.e. at $x=0$ there is a discontinuity of the second kind (infinite).

  \item \textbf{Vertical asymptote:} by the previous item, $x=0$ is a vertical asymptote (from both sides the function tends to $-\infty$). There are no other vertical asymptotes since $f$ is continuous at every other point.

  \textbf{Oblique asymptote:}
  \[ m=\lim_{x\to\pm\infty}\frac{f(x)}{x}=\lim_{x\to\pm\infty}\frac{x^3-1}{4x^3}=\frac14,\qquad
     n=\lim_{x\to\pm\infty}\Big(f(x)-\frac{x}{4}\Big)=\lim_{x\to\pm\infty}\Big(-\frac{1}{4x^2}\Big)=0 . \]
  Hence $y=\frac{x}{4}$ is an oblique asymptote both at $+\infty$ and at $-\infty$. Since $f(x)-\frac{x}{4}=-\frac{1}{4x^2}<0$, the graph lies below the asymptote. There is no horizontal asymptote (since $f(x)\to\pm\infty$ as $x\to\pm\infty$).

  \item Differentiating:
  \[ f'(x)=\frac14+\frac{1}{2x^3}=\frac{x^3+2}{4x^3}. \]
  $f'(x)=0 \iff x^3=-2 \iff x=-\sqrt[3]{2}$. Sign table (the numerator $x^3+2$ is positive for $x>-\sqrt[3]{2}$, the denominator $4x^3$ is positive for $x>0$):
  \begin{itemize}
    \item $x<-\sqrt[3]{2}$: numerator negative, denominator negative, $f'>0$ --- $f$ is increasing.
    \item $-\sqrt[3]{2}<x<0$: numerator positive, denominator negative, $f'<0$ --- $f$ is decreasing.
    \item $x>0$: numerator positive, denominator positive, $f'>0$ --- $f$ is increasing.
  \end{itemize}
  Hence $f$ is increasing on $(-\infty,-\sqrt[3]{2})$ and on $(0,\infty)$, and decreasing on $(-\sqrt[3]{2},0)$.
  At $x=-\sqrt[3]{2}$ the derivative changes from positive to negative, hence this is a local maximum point:
  \[ f(-\sqrt[3]{2})=\frac{-2-1}{4\sqrt[3]{4}}=-\frac{3}{4\sqrt[3]{4}}\approx -0.47 . \]
  There are no other extremum points (at $x=0$ the function is not defined).

  \item
  \[ f''(x)=\Big(\frac14+\frac12x^{-3}\Big)'=-\frac{3}{2x^4}. \]
  For every $x\neq 0$ we have $f''(x)<0$. Hence $f$ is concave ($\cap$) on each of the intervals $(-\infty,0)$ and $(0,\infty)$, and is not convex ($\cup$) on any interval. The sign of $f''$ does not change within the domain, hence there are no inflection points.

  \item \textbf{Graphical description:} on the left side the graph lies below the asymptote $y=\frac{x}{4}$ and hugs it as $x\to-\infty$. It increases up to the maximum $\big(-\sqrt[3]{2},-\frac{3}{4\sqrt[3]{4}}\big)\approx(-1.26,-0.47)$, and then decreases to $-\infty$ as $x\to0^-$. On the right side it increases from $-\infty$ (as $x\to0^+$), crosses the $x$-axis at $(1,0)$, and approaches the asymptote $y=\frac{x}{4}$ from below as $x\to\infty$. The graph is concave everywhere.
\end{enumerate}
\end{solution}

\subsection*{Part A, Question 2}

\begin{question}
Compute the limit
\[ \lim_{x\to 0}\left(\frac{\sqrt{x+9}-3}{\sqrt{(x+2)^2+5}-x-3}+(\cos(x))^{\frac{1}{\sin^2(2x)}}\right) \]
\end{question}

\begin{hint}
Compute each term separately. In the first term multiply the numerator and the denominator by the conjugates (to get rid of the roots).
\end{hint}

\begin{hint}
In the second term write $(\cos x)^{\frac{1}{\sin^2(2x)}}=e^{\frac{\ln\cos x}{\sin^2(2x)}}$, and compute the limit of the exponent.
\end{hint}

\begin{solution}
We compute each term separately. If both have a finite limit, the limit of the sum equals the sum of the limits (arithmetic of limits).

\textbf{The first term.} For $x=0$ we get $\frac{0}{0}$. We multiply by the conjugates. In the numerator:
\[ \sqrt{x+9}-3=\frac{(x+9)-9}{\sqrt{x+9}+3}=\frac{x}{\sqrt{x+9}+3}. \]
In the denominator, $(x+2)^2+5=x^2+4x+9$, hence
\[ \sqrt{x^2+4x+9}-(x+3)=\frac{x^2+4x+9-(x+3)^2}{\sqrt{x^2+4x+9}+x+3}=\frac{-2x}{\sqrt{x^2+4x+9}+x+3}. \]
Hence, for small $x\neq0$,
\[ \frac{\sqrt{x+9}-3}{\sqrt{(x+2)^2+5}-x-3}=\frac{x}{\sqrt{x+9}+3}\cdot\frac{\sqrt{x^2+4x+9}+x+3}{-2x}
=-\frac{\sqrt{x^2+4x+9}+x+3}{2(\sqrt{x+9}+3)}\xrightarrow[x\to0]{}-\frac{3+3}{2\cdot 6}=-\frac12 , \]
where we used the continuity of the square root.

\textbf{The second term.} Near $0$ we have $\cos x>0$, hence
\[ (\cos x)^{\frac{1}{\sin^2(2x)}}=\exp\Big(\frac{\ln(\cos x)}{\sin^2(2x)}\Big). \]
We compute the limit of the exponent using Taylor expansions about $0$:
$\cos x=1-\frac{x^2}{2}+o(x^2)$, hence $\ln(\cos x)=\ln\big(1+(-\frac{x^2}{2}+o(x^2))\big)=-\frac{x^2}{2}+o(x^2)$.
Also $\sin(2x)=2x+o(x)$, hence $\sin^2(2x)=4x^2+o(x^2)$. Therefore
\[ \lim_{x\to0}\frac{\ln(\cos x)}{\sin^2(2x)}=\lim_{x\to0}\frac{-\frac{x^2}{2}+o(x^2)}{4x^2+o(x^2)}=-\frac18 . \]
(This can also be obtained by L'Hôpital: $\frac{\ln\cos x}{\sin^2 2x}$ is of the form $\frac00$, and the ratio of the derivatives $\frac{-\tan x}{4\sin(2x)\cos(2x)}=\frac{-\tan x}{2\sin(4x)}\to -\frac{1}{8}$, since $\frac{\tan x}{x}\to1$ and $\frac{\sin 4x}{4x}\to 1$.)
The exponential function is continuous, hence the second term tends to $e^{-1/8}$.

\textbf{Answer:}
\[ \lim_{x\to 0}\left(\frac{\sqrt{x+9}-3}{\sqrt{(x+2)^2+5}-x-3}+(\cos x)^{\frac{1}{\sin^2(2x)}}\right)=\boxed{e^{-1/8}-\frac12}\approx 0.3825 . \]
\end{solution}

\subsection*{Part A, Question 3}

\begin{question}
What is the area enclosed between the curves $y=\frac{4}{x}$ and $y=5-x$?
What is the volume obtained by revolving this area about the $x$-axis?
\end{question}

\begin{hint}
First find the intersection points of the curves, and check which of them lies above the other between the intersection points.
\end{hint}

\begin{hint}
Revolving about the $x$-axis gives a solid with a hole (the disk/washer method): $V=\pi\int_a^b\big(y_{\text{upper}}^2-y_{\text{lower}}^2\big)\,dx$.
\end{hint}

\begin{solution}
\textbf{Intersection points:} $\frac4x=5-x \iff 4=5x-x^2 \iff x^2-5x+4=0 \iff (x-1)(x-4)=0$, i.e. $x=1$ and $x=4$ (the points $(1,4)$ and $(4,1)$).

\textbf{Which is above which:} on the interval $(1,4)$, for example at $x=2$: $5-2=3>2=\frac42$. And in general, for $x>0$,
$5-x-\frac4x=\frac{-(x^2-5x+4)}{x}=\frac{-(x-1)(x-4)}{x}>0$ when $1<x<4$. Hence the line lies above the hyperbola, and both functions are positive on the interval.

\textbf{The area:} by the formula for the area between two graphs,
\[ S=\int_1^4\Big(5-x-\frac4x\Big)dx=\Big[5x-\frac{x^2}{2}-4\ln x\Big]_1^4
=\Big(20-8-4\ln4\Big)-\Big(5-\frac12\Big)=\frac{15}{2}-8\ln2 . \]
That is, $\boxed{S=\frac{15}{2}-8\ln 2\approx 1.955}$.

\textbf{The volume:} when revolving about the $x$-axis, every vertical cross-section at a point $x\in[1,4]$ is a washer with outer radius $5-x$ and inner radius $\frac4x$ (both radii are positive). Hence
\[ V=\pi\int_1^4\Big((5-x)^2-\frac{16}{x^2}\Big)dx=\pi\Big[-\frac{(5-x)^3}{3}+\frac{16}{x}\Big]_1^4
=\pi\Big[\Big(-\frac13+4\Big)-\Big(-\frac{64}{3}+16\Big)\Big]=\pi\Big(\frac{11}{3}+\frac{16}{3}\Big)=9\pi . \]
That is, $\boxed{V=9\pi}$.
\end{solution}

\subsection*{Part B, Question 1}

\begin{question}
Let $f(x)$ be a function such that $f(x)+\sin(x)$ is differentiable on the open interval $(a,b)$. Which of the following must be true?
\begin{enumerate}
  \item $f(x)+\cos(x)$ is differentiable on the open interval $(a,b)$.
  \item $f(x)+\sin(x)$ is continuous on the closed interval $[a,b]$.
  \item $f(x)$ is continuous on the closed interval $[a,b]$.
  \item $f(x)$ is bounded on the open interval $(a,b)$.
  \item $f(x)+\sin(x)$ is bounded on the open interval $(a,b)$.
\end{enumerate}
\end{question}

\begin{hint}
Write $f+\cos$ as the sum of $f+\sin$ and a known differentiable function.
\end{hint}

\begin{solution}
\textbf{The correct answer: a.}

Denote $g(x)=f(x)+\sin x$, differentiable on $(a,b)$. Then
\[ f(x)+\cos x=g(x)-\sin x+\cos x , \]
and this is a sum of functions differentiable on $(a,b)$ ($\sin$ and $\cos$ are differentiable on all of $\R$), hence it is differentiable on $(a,b)$. That is, (a) is always true.

The other statements need not hold. Counterexample: $f(x)=\frac{1}{x-a}-\sin x$ on $(a,b)$. Then $f(x)+\sin x=\frac{1}{x-a}$ is differentiable on $(a,b)$, but:
\begin{itemize}
  \item $f+\sin$ and $f$ are not defined (and certainly not continuous) at $x=a$ --- (b) and (c) fail. (Even if we assign them some value at $a$, they are unbounded near $a$ and hence not continuous there.)
  \item $\frac{1}{x-a}\to+\infty$ as $x\to a^+$, hence $f+\sin$ is unbounded on $(a,b)$ --- (e) fails; and also $f=\frac{1}{x-a}-\sin x$ is unbounded (since $\sin$ is bounded) --- (d) fails.
\end{itemize}
\end{solution}

\subsection*{Part B, Question 2}

\begin{question}
The given graph is a graph in polar coordinates of the function:

(In the figure: two equal closed loops, circle-shaped, one above the $x$-axis and one below it, touching each other at the origin; the loops are symmetric about the $y$-axis and stretched along it.)
\begin{enumerate}
  \item $r(\phi)=\sin^2\phi$
  \item $r(\phi)=1+\cos\phi$
  \item $r(\phi)=1-\cos\phi$
  \item $r(\phi)=\cos^2\phi$
  \item $r(\phi)=\cos\phi$
\end{enumerate}
\end{question}

\begin{hint}
For each function, check at which angles $r=0$ (there the curve passes through the origin) and at which angles $r$ is maximal (the points farthest from the origin).
\end{hint}

\begin{solution}
\textbf{The correct answer: a.}

In the graph the curve passes through the origin in the direction of the $x$-axis (where the two loops touch), and the points farthest from the origin lie on the $y$-axis, above and below, at the same distance. Also, there are two loops.

\begin{itemize}
  \item $r=\sin^2\phi$: we have $r\ge0$ for every $\phi$, $r=0$ exactly for $\phi=0,\pi$ (the direction of the $x$-axis), and $r=1$ is maximal for $\phi=\frac{\pi}{2},\frac{3\pi}{2}$, i.e. at the points $(0,1)$ and $(0,-1)$. As $\phi$ goes from $0$ to $\pi$ a loop is formed in the upper half, and as $\phi$ goes from $\pi$ to $2\pi$ a loop is formed in the lower half. This is exactly the given graph.
  \item $r=1\pm\cos\phi$: cardioids --- only one loop (with a ``dent'' at the origin), symmetric about the $x$-axis. Does not fit.
  \item $r=\cos^2\phi$: two loops, but along the $x$-axis (maximum at $\phi=0,\pi$). Does not fit.
  \item $r=\cos\phi$: a single circle, $x^2+y^2=x$, centered at $(\frac12,0)$. Does not fit.
\end{itemize}
\end{solution}

\subsection*{Part B, Question 3}

\begin{question}
The derivative $y'(0)$ of the implicit function $y(x)$ defined by the equation $\sin(x-y)-y=0$ in a neighborhood of the point $(0,0)$ is equal to
\begin{enumerate}
  \item $\frac12$
  \item $-\frac12$
  \item $\frac{\pi}{2}$
  \item $1+\frac{\pi}{2}$
  \item $0$
\end{enumerate}
\end{question}

\begin{solution}
\textbf{The correct answer: a.}

The point $(0,0)$ satisfies the equation: $\sin(0)-0=0$. Differentiate both sides with respect to $x$, where $y=y(x)$ (chain rule):
\[ \cos(x-y)\cdot(1-y')-y'=0 . \]
Substitute $x=0,\ y=0$: $\cos 0\cdot(1-y'(0))-y'(0)=0$, i.e. $1-2y'(0)=0$, hence $y'(0)=\frac12$.

(In general: $y'=\frac{\cos(x-y)}{1+\cos(x-y)}$, and at $(0,0)$ this is $\frac{1}{2}$.)
\end{solution}

\subsection*{Part B, Question 4}

\begin{question}
If $f(x)=\sin^2(x)$, then in a neighborhood of $0$ we have
\begin{enumerate}
  \item $f(x)=x^2-\frac{x^4}{3}+o(x^4)$
  \item $f(x)=x^2+\frac{x^4}{3}+o(x^4)$
  \item $f(x)=x^2-\frac{x^4}{6}+o(x^4)$
  \item $f(x)=x^2+\frac{x^4}{6}+o(x^4)$
  \item $f(x)=x^2+\frac{x^4}{36}+o(x^4)$
\end{enumerate}
\end{question}

\begin{hint}
Use the identity $\sin^2x=\frac{1-\cos(2x)}{2}$ and the familiar expansion of $\cos$.
\end{hint}

\begin{solution}
\textbf{The correct answer: a.}

By the identity $\sin^2x=\frac{1-\cos 2x}{2}$ and the Maclaurin expansion $\cos t=1-\frac{t^2}{2}+\frac{t^4}{24}+o(t^4)$ with $t=2x$:
\[ \cos 2x=1-2x^2+\frac{16x^4}{24}+o(x^4)=1-2x^2+\frac{2x^4}{3}+o(x^4), \]
hence
\[ \sin^2x=\frac{1-\cos2x}{2}=x^2-\frac{x^4}{3}+o(x^4). \]
(Check in another way: $\sin x=x-\frac{x^3}{6}+o(x^4)$, hence $\sin^2x=x^2-2\cdot x\cdot\frac{x^3}{6}+o(x^4)=x^2-\frac{x^4}{3}+o(x^4)$.)
\end{solution}

\subsection*{Part B, Question 5}

\begin{question}
If $f(x)$ is a real function whose derivative is $f'(x)=(x-1)^2(x+2)(x-5)$, then:
\begin{enumerate}
  \item $f(x)$ has a maximum at $x=-2$ and a minimum at $x=5$, and no extremum point at $x=1$.
  \item $f(x)$ has a maximum at $x=5$ and a minimum at $x=-2$, and no extremum point at $x=1$.
  \item $f(x)$ has a maximum at $x=-2$ and at $x=1$ and a minimum at $x=5$.
  \item $f(x)$ has a maximum at $x=5$ and a minimum at $x=-2$ and at $x=1$.
  \item $f(x)$ has no extremum points.
\end{enumerate}
\end{question}

\begin{hint}
Check the sign of the derivative on both sides of each critical point. Note that the factor $(x-1)^2$ does not change sign.
\end{hint}

\begin{solution}
\textbf{The correct answer: a.}

$f$ is differentiable on all of $\R$, hence extremum points are possible only where $f'=0$ (Fermat's theorem): $x=-2,\,1,\,5$. We check the sign of $f'$. The factor $(x-1)^2\ge0$ does not change sign, hence the sign is determined by $(x+2)(x-5)$:
\begin{itemize}
  \item $x<-2$: $(x+2)(x-5)>0$, hence $f'>0$ --- $f$ is increasing.
  \item $-2<x<5$, $x\neq 1$: $(x+2)(x-5)<0$, hence $f'<0$ --- $f$ is decreasing.
  \item $x>5$: $f'>0$ --- $f$ is increasing.
\end{itemize}
Hence at $x=-2$ the derivative changes from positive to negative --- a local maximum; at $x=5$ it changes from negative to positive --- a local minimum; and at $x=1$ the derivative is negative on both sides, so $f$ is strictly decreasing in a neighborhood of $1$ and there is no extremum there.
\end{solution}

\section*{2019/20 Semester A Moed C}

\subsection*{Part A, Question 1}

\begin{question}
Investigate the function $y=\frac{2x-1}{(x-1)^2}$ according to the following items:
\begin{enumerate}
  \item The domain of the function, intersection points with the axes, evenness, oddness.
  \item Continuity, points of discontinuity.
  \item Asymptotes
  \item Intervals of increase and decrease, extremum points.
  \item Intervals of convexity and concavity, inflection points.
  \item Graphical description.
\end{enumerate}
\end{question}

\begin{hint}
When differentiating by the quotient rule, it is convenient to cancel one factor $(x-1)$ from the numerator and the denominator before simplifying.
\end{hint}

\begin{solution}
Denote $f(x)=\frac{2x-1}{(x-1)^2}$.
\begin{enumerate}
  \item \textbf{Domain:} $D=\{x\in\R: x\neq1\}$.

  \textbf{Intersection with the axes:} $f(0)=\frac{-1}{1}=-1$, i.e. $(0,-1)$ on the $y$-axis. $f(x)=0\iff 2x-1=0\iff x=\frac12$, i.e. $(\frac12,0)$ on the $x$-axis.

  \textbf{Evenness:} the domain is not symmetric with respect to $0$ ($-1\in D$ but $1\notin D$), hence the function is neither even nor odd.

  \item $f$ is a quotient of polynomials, hence continuous on its entire domain. At $x=1$: the numerator tends to $1>0$ and the denominator $(x-1)^2\to0^+$, hence $\lim_{x\to1^\pm}f(x)=+\infty$ --- a discontinuity of the second kind (infinite).

  \item \textbf{Vertical:} $x=1$ (from both sides $f\to+\infty$).
  \textbf{Horizontal:} $\lim_{x\to\pm\infty}\frac{2x-1}{(x-1)^2}=\lim_{x\to\pm\infty}\frac{\frac2x-\frac1{x^2}}{(1-\frac1x)^2}=0$, hence $y=0$ is a horizontal asymptote at $\pm\infty$ (and hence there is no oblique asymptote).

  \item By the quotient rule:
  \[ f'(x)=\frac{2(x-1)^2-(2x-1)\cdot2(x-1)}{(x-1)^4}=\frac{2(x-1)-2(2x-1)}{(x-1)^3}=\frac{-2x}{(x-1)^3}. \]
  $f'(x)=0\iff x=0$. Signs:
  \begin{itemize}
    \item $x<0$: numerator $-2x>0$, denominator $<0$, hence $f'<0$ --- decreasing.
    \item $0<x<1$: numerator $<0$, denominator $<0$, hence $f'>0$ --- increasing.
    \item $x>1$: numerator $<0$, denominator $>0$, hence $f'<0$ --- decreasing.
  \end{itemize}
  Hence $f$ is decreasing on $(-\infty,0)$ and on $(1,\infty)$ and increasing on $(0,1)$. At $x=0$ the derivative changes from negative to positive: a \textbf{local minimum} $(0,-1)$. (This is also an absolute minimum: for $x<\frac12$ the function decreases to $-1$ and then increases, and for $x\ge\frac12$, $f\ge0$.)

  \item Differentiate $f'(x)=-2x(x-1)^{-3}$:
  \[ f''(x)=-2(x-1)^{-3}+6x(x-1)^{-4}=\frac{-2(x-1)+6x}{(x-1)^4}=\frac{2(2x+1)}{(x-1)^4}. \]
  The denominator is positive for every $x\neq1$, hence the sign is that of $2x+1$:
  $f''<0$ for $x<-\frac12$ --- $f$ is concave ($\cap$) on $(-\infty,-\frac12)$;
  $f''>0$ for $-\frac12<x<1$ and $x>1$ --- $f$ is convex ($\cup$) on $(-\frac12,1)$ and on $(1,\infty)$.
  At $x=-\frac12$ the sign changes and the function is continuous, hence this is an \textbf{inflection point}: $f(-\frac12)=\frac{-2}{9/4}=-\frac89$, i.e. $(-\frac12,-\frac89)$.

  \item \textbf{Graphical description:} on the left the graph starts below the $x$-axis, close to the asymptote $y=0$ as $x\to-\infty$, decreases (concave) to the inflection point $(-\frac12,-\frac89)$, continues to decrease (convex) to the minimum $(0,-1)$, increases and crosses the $x$-axis at $(\frac12,0)$ and tends to $+\infty$ as $x\to1^-$. To the right of $x=1$ it decreases from $+\infty$ and approaches $y=0$ from above as $x\to\infty$ (convex).
\end{enumerate}
\end{solution}

\subsection*{Part A, Question 2}

\begin{question}
A curve is given in polar form: $r=a(1+\cos\varphi)$, where $a$ is a real positive parameter, $0\le\varphi\le2\pi$.
Compute the \textbf{length} of the curve. \textbf{Compute the area} bounded by the curve
\end{question}

\begin{hint}
Length of a polar curve: $L=\int_\alpha^\beta\sqrt{r^2+(r')^2}\,d\varphi$. Area: $S=\frac12\int_\alpha^\beta r^2\,d\varphi$.
\end{hint}

\begin{hint}
Use the identity $1+\cos\varphi=2\cos^2\frac{\varphi}{2}$, and pay attention to the sign of $\cos\frac\varphi2$ on the interval.
\end{hint}

\begin{solution}
The curve (a cardioid) is $r(\varphi)=a(1+\cos\varphi)\ge0$, and $r'(\varphi)=-a\sin\varphi$.

\textbf{Length:} by the formula for arc length in polar coordinates,
\[ r^2+(r')^2=a^2\big(1+2\cos\varphi+\cos^2\varphi+\sin^2\varphi\big)=2a^2(1+\cos\varphi)=4a^2\cos^2\frac{\varphi}{2}. \]
Hence
\[ L=\int_0^{2\pi}\sqrt{r^2+(r')^2}\,d\varphi=\int_0^{2\pi}2a\Big|\cos\frac{\varphi}{2}\Big|\,d\varphi . \]
For $0\le\varphi\le\pi$, $\cos\frac\varphi2\ge0$, and for $\pi\le\varphi\le2\pi$, $\cos\frac\varphi2\le0$. By symmetry (or directly):
\[ L=2\int_0^{\pi}2a\cos\frac{\varphi}{2}\,d\varphi=4a\Big[2\sin\frac\varphi2\Big]_0^{\pi}=8a . \]
That is, $\boxed{L=8a}$.

\textbf{Area:} by the formula for area in polar coordinates,
\[ S=\frac12\int_0^{2\pi}a^2(1+\cos\varphi)^2d\varphi=\frac{a^2}{2}\int_0^{2\pi}\Big(1+2\cos\varphi+\frac{1+\cos2\varphi}{2}\Big)d\varphi
=\frac{a^2}{2}\Big[\frac32\varphi+2\sin\varphi+\frac{\sin2\varphi}{4}\Big]_0^{2\pi}=\frac{a^2}{2}\cdot3\pi . \]
That is, $\boxed{S=\frac{3\pi a^2}{2}}$.
\end{solution}

\subsection*{Part A, Question 3}

\begin{question}
Compute the following limit:
\[ \lim_{x\to0^+}\left((x^2+x+1))^{\frac{1}{x^2}}+x\cdot\ln x\right) \]
\end{question}

\begin{hint}
Compute each term separately. In the first term write $(x^2+x+1)^{1/x^2}=e^{\frac{\ln(1+x+x^2)}{x^2}}$.
\end{hint}

\begin{solution}
(The original wording contains a superfluous parenthesis; the first term is $(x^2+x+1)^{1/x^2}$.)

\textbf{The second term:} $x\ln x=\frac{\ln x}{1/x}$, of the form $\frac{-\infty}{\infty}$. By L'Hôpital's rule:
\[ \lim_{x\to0^+}\frac{\ln x}{1/x}=\lim_{x\to0^+}\frac{1/x}{-1/x^2}=\lim_{x\to0^+}(-x)=0 . \]

\textbf{The first term:} for $x>0$, $x^2+x+1>1$, hence
\[ (x^2+x+1)^{\frac{1}{x^2}}=\exp\Big(\frac{\ln(1+x+x^2)}{x^2}\Big). \]
The exponent: $\frac{\ln(1+x+x^2)}{x^2}=\frac{\ln(1+x+x^2)}{x+x^2}\cdot\frac{x+x^2}{x^2}=\frac{\ln(1+t)}{t}\Big|_{t=x+x^2}\cdot\Big(\frac1x+1\Big)$.
As $x\to0^+$: $t=x+x^2\to0^+$ and hence $\frac{\ln(1+t)}{t}\to1$, whereas $\frac1x+1\to+\infty$. Hence the exponent tends to $+\infty$, and therefore
\[ \lim_{x\to0^+}(x^2+x+1)^{\frac{1}{x^2}}=+\infty . \]
(Intuitively: this is a limit of the form $1^\infty$, but the base approaches $1$ ``slowly'' --- like $1+x$ --- while the exponent grows like $\frac{1}{x^2}$, hence $\big(1+x\big)^{1/x^2}\approx e^{1/x}\to\infty$.)

\textbf{Summary:} the sum of a function tending to $+\infty$ and a function tending to a finite limit ($0$) tends to $+\infty$. Hence
\[ \lim_{x\to0^+}\left((x^2+x+1)^{\frac{1}{x^2}}+x\ln x\right)=\boxed{+\infty} . \]
\end{solution}

\subsection*{Part B, Question 1}

\begin{question}
Given $h(x)=\begin{cases}x^2+a & x<-1\\ x^3-8 & x\ge-1\end{cases}$. For which value of the parameter $a$ is the function continuous for every real $x$,
\begin{enumerate}
  \item $-1$
  \item $-9$
  \item $-10$
  \item $-8$
  \item $-2$
\end{enumerate}
\end{question}

\begin{solution}
\textbf{The correct answer: c ($a=-10$).}

On each of the intervals $x<-1$ and $x>-1$ the function is a polynomial and hence continuous. It remains to check $x=-1$:
\[ h(-1)=\lim_{x\to-1^+}h(x)=(-1)^3-8=-9,\qquad \lim_{x\to-1^-}h(x)=(-1)^2+a=1+a . \]
Continuity at $-1$ holds if and only if $1+a=-9$, i.e. $a=-10$.
\end{solution}

\subsection*{Part B, Question 2}

\begin{question}
Let $f(x)$ be a function differentiable on an open interval $(a,b)$, and it is known that its derivative is positive, i.e. $f'(x)>0$. Which of the following must be true?
\begin{enumerate}
  \item $f(x)$ is a bounded function on the open interval $(a,b)$.
  \item $f(x)$ is a function satisfying $f(x)>0$ on the open interval $(a,b)$.
  \item $f(x)$ is a continuous function on the closed interval $[a,b]$.
  \item $f(x)$ is a monotone increasing function on the open interval $(a,b)$.
  \item $f(x)$ is a monotone increasing function on the closed interval $[a,b]$.
\end{enumerate}
\end{question}

\begin{hint}
What does Lagrange's theorem say about $f(x_2)-f(x_1)$ when $a<x_1<x_2<b$?
\end{hint}

\begin{solution}
\textbf{The correct answer: d.}

Let $a<x_1<x_2<b$. $f$ is differentiable and hence continuous on $[x_1,x_2]$ and differentiable on $(x_1,x_2)$, so by Lagrange's theorem there exists $c\in(x_1,x_2)$ such that
\[ f(x_2)-f(x_1)=f'(c)(x_2-x_1)>0 . \]
Hence $f$ is (strictly) increasing on $(a,b)$.

Counterexamples to the other statements, for example on the interval $(0,1)$:
\begin{itemize}
  \item $f(x)=-\frac1x$: $f'(x)=\frac1{x^2}>0$, but $f$ is not bounded on $(0,1)$ ((a) fails), is not defined at $0$ and in particular is not continuous on $[0,1]$ ((c) fails), and is not monotone on the closed interval ((e) fails, $f$ is not even defined there).
  \item $f(x)=x-5$: $f'=1>0$ but $f<0$ on all of $(0,1)$ ((b) fails).
\end{itemize}
\end{solution}

\subsection*{Part B, Question 3}

\begin{question}
Given the function $F(x)=\int_0^x\sqrt{t^2+1}\,dt$. Compute $F'(\sqrt2)$.
\begin{enumerate}
  \item $\sqrt3$
  \item $\frac{1}{\sqrt3}$
  \item $\sqrt2$
  \item $3$
  \item $\sqrt2+1$
\end{enumerate}
\end{question}

\begin{solution}
\textbf{The correct answer: a.}

The function $t\mapsto\sqrt{t^2+1}$ is continuous on all of $\R$, hence by the Fundamental Theorem of Calculus $F$ is differentiable and $F'(x)=\sqrt{x^2+1}$. Hence
\[ F'(\sqrt2)=\sqrt{2+1}=\sqrt3 . \]
\end{solution}

\subsection*{Part B, Question 4}

\begin{question}
Find the smallest value on the domain of the function: $f(x)=x^4+4x$
\begin{enumerate}
  \item $-4$
  \item $1$
  \item $-3$
  \item $-5$
  \item $-2$
\end{enumerate}
\end{question}

\begin{solution}
\textbf{The correct answer: c ($-3$).}

$f$ is a polynomial, defined and differentiable on all of $\R$. $f'(x)=4x^3+4=4(x+1)(x^2-x+1)$, and the factor $x^2-x+1$ is always positive (negative discriminant). Hence $f'<0$ for $x<-1$ and $f'>0$ for $x>-1$: $f$ is decreasing on $(-\infty,-1]$ and increasing on $[-1,\infty)$. Therefore $x=-1$ is an absolute minimum point, and the smallest value is
\[ f(-1)=1-4=-3 . \]
\end{solution}

\subsection*{Part B, Question 5}

\begin{question}
Let $f(x)$ be a function that is \textbf{twice} differentiable on an open interval $(a,b)$. Which of the following must be true?
\begin{enumerate}
  \item $f'(x)$ is continuous on the open interval $(a,b)$.
  \item $f(x)$ has an inflection point in the open interval $(a,b)$.
  \item $f(x)$ is a monotone function on the open interval $(a,b)$.
  \item $f''(x)$ is a continuous function on the open interval $(a,b)$.
  \item $f(x)$ is a bounded function on the open interval $(a,b)$.
\end{enumerate}
\end{question}

\begin{hint}
A function differentiable at a point is continuous there. Apply this to $f'$.
\end{hint}

\begin{solution}
\textbf{The correct answer: a.}

$f$ is twice differentiable on $(a,b)$, i.e. $f'$ is differentiable at every point of $(a,b)$. A function differentiable at a point is continuous there, hence $f'$ is continuous on $(a,b)$.

Counterexamples to the other statements, on the interval $(0,1)$ (unless stated otherwise):
\begin{itemize}
  \item $f(x)=x^2$ (on $(0,1)$): $f''=2>0$, there is no inflection point ((b) fails).
  \item $f(x)=\left(x-\frac12\right)^2$: not monotone on $(0,1)$ ((c) fails).
  \item $f(x)=x^4\sin\frac1x$ ($f(0)=0$) on the interval $(-1,1)$: twice differentiable, but $f''(x)=12x^2\sin\frac1x-6x\cos\frac1x-\sin\frac1x$ for $x\neq0$ and $f''(0)=0$, and $f''$ is not continuous at $0$ ((d) fails).
  \item $f(x)=\frac1x$ on $(0,1)$: twice differentiable but not bounded ((e) fails).
\end{itemize}
\end{solution}

\section*{2019/20 Semester B Moed B}

\subsection*{Question 1a}

\begin{question}
Compute the derivative of the function $f(x)=\sin(2x+1)$ at the point $x$ using the definition of the derivative (an answer based on the table of derivatives will not be accepted).
\end{question}

\begin{hint}
Use the identity $\sin\alpha-\sin\beta=2\sin\frac{\alpha-\beta}{2}\cos\frac{\alpha+\beta}{2}$ and the limit $\lim_{t\to0}\frac{\sin t}{t}=1$.
\end{hint}

\begin{solution}
By definition, $f'(x)=\lim_{h\to0}\frac{f(x+h)-f(x)}{h}$. By the sum-to-product identity,
\[ \sin(2x+2h+1)-\sin(2x+1)=2\sin(h)\cos(2x+1+h). \]
Hence
\[ f'(x)=\lim_{h\to0}\frac{2\sin h\cos(2x+1+h)}{h}=2\lim_{h\to0}\frac{\sin h}{h}\cdot\lim_{h\to0}\cos(2x+1+h)=2\cdot1\cdot\cos(2x+1), \]
where we used the basic limit $\frac{\sin h}{h}\to1$ and the continuity of $\cos$. That is, $\boxed{f'(x)=2\cos(2x+1)}$.
\end{solution}

\subsection*{Question 1b}

\begin{question}
Compute the following limits:
\begin{enumerate}
  \item (5 pts) $\displaystyle\lim_{x\to\infty}\frac{x+2\sin x}{3x-\cos x}$,
  \item (5 pts) $\displaystyle\lim_{x\to+\infty}\left(\sqrt{x+\sqrt{x}}-\sqrt{x-\sqrt{x}}\right)$
\end{enumerate}
\end{question}

\begin{hint}
In the first limit divide the numerator and the denominator by $x$ and use the fact that $\sin x,\cos x$ are bounded.
\end{hint}

\begin{hint}
In the second limit multiply and divide by the conjugate.
\end{hint}

\begin{solution}
\begin{enumerate}
  \item Divide the numerator and the denominator by $x$ (for $x>0$):
  \[ \frac{x+2\sin x}{3x-\cos x}=\frac{1+2\frac{\sin x}{x}}{3-\frac{\cos x}{x}} . \]
  Since $|\sin x|,|\cos x|\le1$ and $\frac1x\to0$, by ``bounded times tending to zero'' $\frac{\sin x}{x}\to0$ and $\frac{\cos x}{x}\to0$. Hence the limit is $\boxed{\frac13}$.
  \item Multiply and divide by the conjugate (for $x>1$ both roots are defined):
  \[ \sqrt{x+\sqrt x}-\sqrt{x-\sqrt x}=\frac{(x+\sqrt x)-(x-\sqrt x)}{\sqrt{x+\sqrt x}+\sqrt{x-\sqrt x}}=\frac{2\sqrt x}{\sqrt{x+\sqrt x}+\sqrt{x-\sqrt x}}
  =\frac{2}{\sqrt{1+\frac{1}{\sqrt x}}+\sqrt{1-\frac{1}{\sqrt x}}} , \]
  where we divided the numerator and the denominator by $\sqrt x$. As $x\to+\infty$, $\frac{1}{\sqrt x}\to0$, hence (by continuity of the square root) the limit is $\frac{2}{1+1}=\boxed{1}$.
\end{enumerate}
\end{solution}

\subsection*{Question 1c}

\begin{question}
For which values of the parameter $a$ does the graph of the function $f(x)=\frac{x^2-9}{x+a}$ have no vertical asymptote?
\end{question}

\begin{hint}
When is the root of the denominator also a root of the numerator?
\end{hint}

\begin{solution}
The function is continuous at every point where $x\neq -a$, hence a vertical asymptote can occur only at $x=-a$. As $x\to-a$ the denominator tends to $0$, hence:
\begin{itemize}
  \item If the numerator does not vanish at $-a$, i.e. $a^2-9\neq0$, then $|f(x)|\to\infty$ as $x\to-a$, and $x=-a$ is a vertical asymptote.
  \item If $a^2=9$, i.e. $a=3$ or $a=-3$: for $a=3$, $f(x)=\frac{(x-3)(x+3)}{x+3}=x-3$ for every $x\neq-3$, hence $\lim_{x\to-3}f(x)=-6$ is finite (a removable discontinuity), and there is no vertical asymptote. Similarly for $a=-3$, $f(x)=x+3$ for every $x\neq3$ and $\lim_{x\to3}f(x)=6$.
\end{itemize}
Hence there is no vertical asymptote exactly when $\boxed{a=3 \text{ or } a=-3}$.
\end{solution}

\subsection*{Question 2a}

\begin{question}
For which values of the parameter $a$ does the equation $(x-1)^2e^x-a=0$ have exactly one real solution?
\end{question}

\begin{hint}
Investigate the function $g(x)=(x-1)^2e^x$ (intervals of monotonicity, extrema, limits at $\pm\infty$), and count how many times the horizontal line $y=a$ intersects the graph.
\end{hint}

\begin{solution}
The equation is equivalent to $g(x)=a$ for $g(x)=(x-1)^2e^x$. We investigate $g$, which is differentiable on all of $\R$:
\[ g'(x)=2(x-1)e^x+(x-1)^2e^x=e^x(x-1)(x+1). \]
$e^x>0$, hence $g'>0$ on $(-\infty,-1)$, $g'<0$ on $(-1,1)$ and $g'>0$ on $(1,\infty)$. That is, $g$ is strictly increasing on $(-\infty,-1]$, strictly decreasing on $[-1,1]$ and strictly increasing on $[1,\infty)$, with a local maximum $g(-1)=\frac{4}{e}$ and a local minimum $g(1)=0$.

Limits: $\lim_{x\to\infty}g(x)=\infty$; and as $x\to-\infty$, $(x-1)^2e^x=\frac{(x-1)^2}{e^{-x}}\to0$ (the exponential dominates a polynomial, e.g. by L'Hôpital twice). Also $g(x)\ge0$ for every $x$, with equality only at $x=1$.

By the Intermediate Value Theorem and strict monotonicity on each interval, on each interval of monotonicity every value in the range of that interval is attained exactly once. The ranges of the intervals: $(-\infty,-1]\mapsto(0,\frac4e]$, $[-1,1]\mapsto[0,\frac4e]$, $[1,\infty)\mapsto[0,\infty)$. Hence the number of solutions:
\begin{itemize}
  \item $a<0$: no solutions ($g\ge0$).
  \item $a=0$: a unique solution, $x=1$.
  \item $0<a<\frac4e$: three solutions (one in each interval of monotonicity).
  \item $a=\frac4e$: two solutions ($x=-1$, and one more in $(1,\infty)$).
  \item $a>\frac4e$: a unique solution (in $(1,\infty)$; on the other two intervals $g\le\frac4e<a$).
\end{itemize}
\textbf{Answer:} the equation has a unique real solution for $\boxed{a=0 \text{ or } a>\frac{4}{e}}$.
\end{solution}

\subsection*{Question 2b}

\begin{question}
Find the volume of the solid generated by revolving about the $x$-axis the region bounded by the lines $y=xe^{-x}$, $y=0$, $x=2$. Draw the corresponding sketch.
\end{question}

\begin{hint}
The limits of the integral are $x=0$ (where $xe^{-x}=0$) and $x=2$. The integral $\int x^2e^{-2x}dx$ is computed by integrating by parts twice.
\end{hint}

\begin{solution}
\textbf{Sketch:} $y=xe^{-x}$ passes through the origin, is positive for $x>0$, increases up to the maximum $(1,\frac1e)$ ($y'=(1-x)e^{-x}$) and then decreases, and at $x=2$ its value is $2e^{-2}\approx0.27$. The bounded region is the area between the graph and the $x$-axis, for $0\le x\le 2$ (the graph meets $y=0$ at $x=0$).

\textbf{The volume:} by the formula for the volume of a solid of revolution about the $x$-axis,
\[ V=\pi\int_0^2\big(xe^{-x}\big)^2dx=\pi\int_0^2x^2e^{-2x}dx . \]
We compute by integrating by parts twice:
\[ \int x^2e^{-2x}dx=-\frac{x^2}{2}e^{-2x}+\int xe^{-2x}dx=-\frac{x^2}{2}e^{-2x}-\frac{x}{2}e^{-2x}+\frac12\int e^{-2x}dx
=-e^{-2x}\Big(\frac{x^2}{2}+\frac{x}{2}+\frac14\Big)+C . \]
Hence
\[ V=\pi\Big[-e^{-2x}\Big(\frac{x^2}{2}+\frac{x}{2}+\frac14\Big)\Big]_0^2=\pi\Big(-\frac{13}{4}e^{-4}+\frac14\Big)=\boxed{\frac{\pi}{4}\big(1-13e^{-4}\big)}\approx0.4357 . \]
\end{solution}

\subsection*{Question 3a}

\begin{question}
Compute:
\begin{enumerate}
  \item (7 pts) $\displaystyle\int\frac{5x+2}{x^4+x^3+2x^2}\,dx$
  \item (8 pts) $\displaystyle\int_0^1\ln(1+\sqrt{x})\,dx$
\end{enumerate}
\end{question}

\begin{hint}
In the first integral, decompose into partial fractions: $x^4+x^3+2x^2=x^2(x^2+x+2)$, and the quadratic factor is irreducible.
\end{hint}

\begin{hint}
In the second integral, substitute $t=\sqrt x$ and then integrate by parts.
\end{hint}

\begin{solution}
\begin{enumerate}
  \item The denominator: $x^4+x^3+2x^2=x^2(x^2+x+2)$, and $x^2+x+2$ is irreducible (the discriminant is $1-8<0$). We decompose into partial fractions:
  \[ \frac{5x+2}{x^2(x^2+x+2)}=\frac{A}{x}+\frac{B}{x^2}+\frac{Cx+D}{x^2+x+2}, \]
  i.e. $5x+2=Ax(x^2+x+2)+B(x^2+x+2)+(Cx+D)x^2$. Comparing coefficients: constant term $2B=2\Rightarrow B=1$; coefficient of $x$: $2A+B=5\Rightarrow A=2$; coefficient of $x^3$: $A+C=0\Rightarrow C=-2$; coefficient of $x^2$: $A+B+D=0\Rightarrow D=-3$. Hence
  \[ \frac{5x+2}{x^4+x^3+2x^2}=\frac{2}{x}+\frac{1}{x^2}-\frac{2x+3}{x^2+x+2}. \]
  We split the last term: $2x+3=(2x+1)+2$, and $x^2+x+2=\big(x+\frac12\big)^2+\frac74$. Hence
  \[ \int\frac{2x+3}{x^2+x+2}dx=\ln(x^2+x+2)+2\int\frac{dx}{(x+\frac12)^2+\frac74}=\ln(x^2+x+2)+\frac{4}{\sqrt7}\arctan\frac{2x+1}{\sqrt7}. \]
  \textbf{Answer:}
  \[ \int\frac{5x+2}{x^4+x^3+2x^2}dx=2\ln|x|-\frac1x-\ln(x^2+x+2)-\frac{4}{\sqrt7}\arctan\frac{2x+1}{\sqrt7}+C . \]
  \item Substitute $t=\sqrt x$, i.e. $x=t^2$, $dx=2t\,dt$, and the limits are $0\to0$, $1\to1$:
  \[ \int_0^1\ln(1+\sqrt x)\,dx=\int_0^12t\ln(1+t)\,dt . \]
  Integration by parts with $u=\ln(1+t)$, $v'=2t$ ($v=t^2$):
  \[ =\Big[t^2\ln(1+t)\Big]_0^1-\int_0^1\frac{t^2}{1+t}dt=\ln2-\int_0^1\Big(t-1+\frac{1}{1+t}\Big)dt=\ln2-\Big(\frac12-1+\ln2\Big)=\boxed{\frac12} . \]
  (We used the polynomial division $\frac{t^2}{1+t}=t-1+\frac{1}{1+t}$.)
\end{enumerate}
\end{solution}

\subsection*{Question 3b}

\begin{question}
Write the Maclaurin polynomial of order $n=3$ for the function:
\[ f(x)=\sin(\sin(x)) \]
\end{question}

\begin{hint}
Compose the expansion of $\sin$ with itself.
\end{hint}

\begin{solution}
By the Maclaurin expansion, $\sin u=u-\frac{u^3}{6}+o(u^3)$ and $\sin x=x-\frac{x^3}{6}+o(x^3)$. Substitute $u=\sin x$; since $u=x+o(x)$, we have $u^3=x^3+o(x^3)$ and $o(u^3)=o(x^3)$. Hence
\[ \sin(\sin x)=\Big(x-\frac{x^3}{6}\Big)-\frac{x^3}{6}+o(x^3)=x-\frac{x^3}{3}+o(x^3). \]
By the uniqueness of the Taylor polynomial, the Maclaurin polynomial of order 3 is
\[ \boxed{P_3(x)=x-\frac{x^3}{3}} . \]
(Direct check: $f'(x)=\cos(\sin x)\cos x$, $f'(0)=1$; $f$ is odd, hence $f(0)=f''(0)=0$; and a computation gives $f'''(0)=-2$, and indeed $\frac{-2}{3!}=-\frac13$.)
\end{solution}

\subsection*{Question 4a}

\begin{question}
For which values of the parameters $a$ and $b$ will the function
\[ f(x)=\begin{cases} \dfrac{\sin(3x)}{\ln(1-x)} & x<0\\[2mm] ax+b & 0\le x\le1\\[2mm] (\cos(x-1))^{1/(x-1)} & 1<x<\dfrac{\pi}{2}+1\end{cases} \]
be continuous for every real $x$ less than $\frac{\pi}{2}+1$?
\end{question}

\begin{hint}
Require the one-sided limits at the junction points $x=0$ and $x=1$ to equal the value of the function. For the limit at $x=1^+$ (of the form $1^\infty$), pass to $e^{\ln(\cdot)}$.
\end{hint}

\begin{solution}
On each of the open intervals $(-\infty,0)$, $(0,1)$, $(1,\frac\pi2+1)$ the function is continuous as a composition and quotient of elementary functions: for $x<0$, $1-x>1$ and hence $\ln(1-x)>0$ (the denominator does not vanish); for $1<x<\frac\pi2+1$, $0<x-1<\frac\pi2$ and hence $\cos(x-1)>0$ and the power is defined. It remains to check the junction points.

\textbf{At $x=0$:} $f(0)=b=\lim_{x\to0^+}f(x)$. From the left, using $\frac{\sin 3x}{3x}\to1$ and $\frac{\ln(1-x)}{-x}\to1$:
\[ \lim_{x\to0^-}\frac{\sin3x}{\ln(1-x)}=\lim_{x\to0^-}\frac{\sin3x}{3x}\cdot\frac{-x}{\ln(1-x)}\cdot\frac{3x}{-x}=1\cdot1\cdot(-3)=-3 . \]
Hence continuity at $0$ requires $b=-3$.

\textbf{At $x=1$:} $f(1)=a+b=\lim_{x\to1^-}f(x)$. From the right, let $t=x-1\to0^+$:
\[ (\cos t)^{1/t}=\exp\Big(\frac{\ln\cos t}{t}\Big),\qquad \lim_{t\to0^+}\frac{\ln\cos t}{t}\overset{\text{L'H}}{=}\lim_{t\to0^+}\frac{-\tan t}{1}=0 , \]
(the limit is of the form $\frac00$, and the derivatives exist), and hence by continuity of the exponential $\lim_{x\to1^+}f(x)=e^0=1$. Continuity at $1$ requires $a+b=1$, i.e. $a=4$.

\textbf{Answer:} $\boxed{a=4,\ b=-3}$.
\end{solution}

\subsection*{Question 4b}

\begin{question}
Find the equation of the tangent line to the function $y(x)$, given implicitly by $e^{x\cdot y}+y^3=9$, at the point $(0,2)$.
\end{question}

\begin{hint}
Differentiate both sides of the equation with respect to $x$, where $y=y(x)$, and substitute the point.
\end{hint}

\begin{solution}
The point is on the curve: $e^{0}+2^3=1+8=9$. Differentiate the equation with respect to $x$ (implicit differentiation, $y=y(x)$):
\[ e^{xy}\,(y+xy')+3y^2y'=0 . \]
Substitute $x=0$, $y=2$: $1\cdot(2+0)+12y'=0$, hence $y'(0)=-\frac16$. The equation of the tangent line:
\[ \boxed{y=2-\frac{x}{6}} . \]
\end{solution}

\subsection*{Question 4c}

\begin{question}
Compute $\displaystyle\int_{0.25}^{+\infty}\frac{dx}{(4x+1)\sqrt{x}}$.
\end{question}

\begin{hint}
Substitute $t=\sqrt x$.
\end{hint}

\begin{solution}
This is an improper integral (infinite upper limit). Substitute $t=\sqrt x$, $x=t^2$, $dx=2t\,dt$; when $x=0.25$, $t=\frac12$, and when $x\to\infty$, $t\to\infty$:
\[ \int_{0.25}^{R}\frac{dx}{(4x+1)\sqrt x}=\int_{1/2}^{\sqrt R}\frac{2t\,dt}{(4t^2+1)t}=\int_{1/2}^{\sqrt R}\frac{2\,dt}{1+(2t)^2}=\Big[\arctan(2t)\Big]_{1/2}^{\sqrt R}=\arctan(2\sqrt R)-\frac{\pi}{4}. \]
As $R\to\infty$, $\arctan(2\sqrt R)\to\frac\pi2$. Hence the integral converges and
\[ \int_{0.25}^{+\infty}\frac{dx}{(4x+1)\sqrt{x}}=\frac\pi2-\frac\pi4=\boxed{\frac{\pi}{4}} . \]
\end{solution}

\section*{2021/22 Semester A Moed A}

\subsection*{Question 1}

\begin{question}
Define a function by
\[ f(x)=\begin{cases} x^3 \sin(\frac{5}{x}) & x\neq 0\\ 0 & x=0 \end{cases} \]
\begin{enumerate}
  \item (6 pts) Is $f$ continuous at the point $x=0$ ? Justify.
  \item (6 pts) Is $f$ differentiable at the point $x=0$ ? Justify.
  \item (5 pts) Is the derivative $f'(x)$ continuous at the point $x=0$ ? Justify.
\end{enumerate}
\end{question}

\begin{hint}
``Tends to zero times bounded'': $\left|\sin(\frac{5}{x})\right|\le 1$ for every $x\neq 0$.
\end{hint}

\begin{hint}
In part (b), compute the derivative at $0$ by definition (limit of the difference quotient). In part (c), compute $f'(x)$ for $x\neq0$ using the differentiation rules and check the limit as $x\to 0$.
\end{hint}

\begin{solution}
\begin{enumerate}
  \item \textbf{Yes.} For every $x\neq 0$ we have
  \[ 0\le |f(x)| = |x|^3\left|\sin\left(\tfrac{5}{x}\right)\right| \le |x|^3 \xrightarrow[x\to0]{} 0, \]
  and hence by the Squeeze theorem $\lim_{x\to0} f(x)=0=f(0)$. That is, $f$ is continuous at $0$.
  \item \textbf{Yes.} By the definition of the derivative:
  \[ f'(0)=\lim_{h\to0}\frac{f(h)-f(0)}{h}=\lim_{h\to0}\frac{h^3\sin(5/h)}{h}=\lim_{h\to0} h^2\sin\left(\tfrac{5}{h}\right). \]
  Here $h^2\to0$ and $\sin(5/h)$ is bounded, and hence (tends to zero times bounded) the limit equals $0$. Hence $f$ is differentiable at $0$ and $f'(0)=0$.
  \item \textbf{Yes.} For $x\neq0$, by the product rule and the chain rule:
  \[ f'(x)=3x^2\sin\left(\tfrac{5}{x}\right)+x^3\cos\left(\tfrac{5}{x}\right)\cdot\left(-\tfrac{5}{x^2}\right)=3x^2\sin\left(\tfrac{5}{x}\right)-5x\cos\left(\tfrac{5}{x}\right). \]
  Both terms are of the form ``tends to zero times bounded'' ($3x^2\to0$, $5x\to0$ and $|\sin|,|\cos|\le1$), and hence
  \[ \lim_{x\to0}f'(x)=0=f'(0). \]
  It follows that $f'$ is continuous at the point $x=0$.
\end{enumerate}
\end{solution}

\subsection*{Question 2a}

\begin{question}
Find all the asymptotes of the function
\[ f(x)=\frac{2x^2-3x+1}{x} \]
\end{question}

\begin{hint}
Divide the numerator by $x$: $f(x)=2x-3+\frac1x$.
\end{hint}

\begin{solution}
Domain: $x\neq0$. For every $x\neq 0$:
\[ f(x)=2x-3+\frac{1}{x}. \]
\textbf{Vertical asymptote:} the numerator at $x=0$ equals $1\neq0$, and hence
\[ \lim_{x\to0^+}f(x)=+\infty,\qquad \lim_{x\to0^-}f(x)=-\infty, \]
i.e. $x=0$ is a vertical asymptote. This is the only point outside the domain, and $f$ is continuous at every other point, so there are no other vertical asymptotes.

\textbf{Oblique/horizontal asymptote:} we compute
\[ m=\lim_{x\to\pm\infty}\frac{f(x)}{x}=\lim_{x\to\pm\infty}\left(2-\frac3x+\frac1{x^2}\right)=2, \]
\[ n=\lim_{x\to\pm\infty}\big(f(x)-2x\big)=\lim_{x\to\pm\infty}\left(-3+\frac1x\right)=-3. \]
Hence $y=2x-3$ is an oblique asymptote both at $+\infty$ and at $-\infty$. Since $f(x)\to\pm\infty$ as $x\to\pm\infty$, there is no horizontal asymptote.

\textbf{Answer:} $x=0$ (vertical) and $y=2x-3$ (oblique in both directions).
\end{solution}

\subsection*{Question 2b}

\begin{question}
Find the intervals of convexity and the inflection points of the function
\[ f(x)=x^2\cdot e^{-x} \]
\end{question}

\begin{hint}
Compute $f''$ and check its sign; note that $e^{-x}>0$ always.
\end{hint}

\begin{solution}
$f$ is defined and twice differentiable on all of $\R$.
\[ f'(x)=2xe^{-x}-x^2e^{-x}=(2x-x^2)e^{-x}, \]
\[ f''(x)=(2-2x)e^{-x}-(2x-x^2)e^{-x}=(x^2-4x+2)e^{-x}. \]
Since $e^{-x}>0$, the sign of $f''$ is the sign of $x^2-4x+2$, whose zeros are $x=2\pm\sqrt2$. This is a quadratic with a positive leading coefficient, and hence:
\begin{itemize}
  \item $f''>0$ on $(-\infty,\,2-\sqrt2)$ and on $(2+\sqrt2,\,\infty)$: there $f$ is \textbf{convex} (concave up).
  \item $f''<0$ on $(2-\sqrt2,\,2+\sqrt2)$: there $f$ is \textbf{concave} (concave down).
\end{itemize}
At the points $x=2\pm\sqrt2$ the function is continuous and $f''$ changes sign, and hence these are \textbf{inflection points}:
\[ \left(2-\sqrt2,\ (2-\sqrt2)^2e^{-(2-\sqrt2)}\right)=\left(2-\sqrt2,\ (6-4\sqrt2)e^{\sqrt2-2}\right), \]
\[ \left(2+\sqrt2,\ (6+4\sqrt2)e^{-2-\sqrt2}\right). \]
\end{solution}

\subsection*{Question 3}

\begin{question}
Estimate $\sqrt[4]{18}$ using a linear approximation (a Taylor polynomial of degree one) and give a bound for the error.
\end{question}

\begin{hint}
Expand $f(x)=\sqrt[4]{x}$ around a point close to $18$ at which the fourth root is known.
\end{hint}

\begin{hint}
Bound the error using the Lagrange remainder: $R_1(x)=\frac{f''(c)}{2}(x-a)^2$ for some $c$ between $a$ and $x$.
\end{hint}

\begin{solution}
Take $f(x)=x^{1/4}$ and expand around $a=16$ (since $\sqrt[4]{16}=2$).
\[ f'(x)=\tfrac14x^{-3/4},\qquad f''(x)=-\tfrac{3}{16}x^{-7/4}. \]
Hence $f(16)=2$, $f'(16)=\frac14\cdot16^{-3/4}=\frac14\cdot\frac18=\frac1{32}$, and the first-order Taylor polynomial is
\[ P_1(x)=2+\frac{1}{32}(x-16). \]
The approximation:
\[ \sqrt[4]{18}\approx P_1(18)=2+\frac{2}{32}=\frac{33}{16}=2.0625. \]

\textbf{Error bound.} By Taylor's theorem with the Lagrange remainder ($f$ is twice differentiable on $(0,\infty)$), there exists $c\in(16,18)$ such that
\[ R_1(18)=f(18)-P_1(18)=\frac{f''(c)}{2!}(18-16)^2=-\frac{3}{16}c^{-7/4}\cdot\frac{4}{2}=-\frac{3}{8}c^{-7/4}. \]
The function $c^{-7/4}$ is decreasing, and hence for $c>16$:
\[ |R_1(18)|=\frac38c^{-7/4}<\frac38\cdot16^{-7/4}=\frac38\cdot\frac{1}{2^7}=\frac{3}{1024}\approx0.0029. \]

\textbf{Answer:} $\sqrt[4]{18}\approx2.0625$ with an error less than $\frac{3}{1024}\approx 0.003$. Moreover, since $R_1<0$, the approximation is larger than the true value:
$2.0625-\frac{3}{1024}<\sqrt[4]{18}<2.0625$.
(Check: $\sqrt[4]{18}\approx2.05977$.)
\end{solution}

\subsection*{Question 4a}

\begin{question}
Compute the antiderivative of
\[ f(x)=\frac{e^x}{e^{2x}+1} \]
\end{question}

\begin{hint}
Substitute $t=e^x$.
\end{hint}

\begin{solution}
Substitute $t=e^x$, $dt=e^x\,dx$:
\[ \int\frac{e^x}{e^{2x}+1}\,dx=\int\frac{dt}{t^2+1}=\arctan t+C=\boxed{\arctan(e^x)+C}. \]
Check: $\left(\arctan e^x\right)'=\frac{e^x}{1+e^{2x}}$.
\end{solution}

\subsection*{Question 4b}

\begin{question}
Find the general solution of the differential equation
\[ (x^2+4)y'+2xy^2=0 \]
\end{question}

\begin{hint}
This is a separable equation: move all the $y$ terms to one side and all the $x$ terms to the other. Do not forget the solution $y\equiv0$.
\end{hint}

\begin{solution}
The equation is equivalent to
\[ y'=-\frac{2x}{x^2+4}\,y^2 \]
(since $x^2+4>0$), and this is a separable equation.

\textbf{Constant solution:} $y\equiv0$ satisfies the equation.

\textbf{Solutions with $y\neq0$:} divide by $y^2$:
\[ \frac{y'}{y^2}=-\frac{2x}{x^2+4}\ \Longrightarrow\ \int\frac{dy}{y^2}=-\int\frac{2x}{x^2+4}\,dx \ \Longrightarrow\ -\frac1y=-\ln(x^2+4)+C. \]
(We used $\int\frac{2x}{x^2+4}dx=\ln(x^2+4)$, since the numerator is the derivative of the denominator.) Hence
\[ \boxed{y=\frac{1}{\ln(x^2+4)+C}},\quad C\in\R, \]
on every interval on which the denominator does not vanish, and in addition the solution $\boxed{y\equiv0}$.

Check: for $y=\frac1{\ln(x^2+4)+C}$ we have $y'=-\frac{2x}{x^2+4}\cdot y^2$, and hence $(x^2+4)y'+2xy^2=0$.
\end{solution}

\subsection*{Question 5}

\begin{question}
Compute the integral
\[ \int\frac{6x^2+13x+8}{x^3+2x^2+2x}dx \]
\end{question}

\begin{hint}
Factor the denominator: $x^3+2x^2+2x=x(x^2+2x+2)$, and the quadratic factor has no real roots. Use partial fraction decomposition.
\end{hint}

\begin{hint}
Split $\frac{2x+5}{x^2+2x+2}$ into $\frac{2x+2}{x^2+2x+2}+\frac{3}{(x+1)^2+1}$.
\end{hint}

\begin{solution}
\textbf{Factoring the denominator.} $x^3+2x^2+2x=x(x^2+2x+2)$, and the discriminant of $x^2+2x+2$ is $4-8<0$, hence it is irreducible over $\R$. The degree of the numerator is less than the degree of the denominator, so we look for a partial fraction decomposition:
\[ \frac{6x^2+13x+8}{x(x^2+2x+2)}=\frac{A}{x}+\frac{Bx+C}{x^2+2x+2}. \]
Multiplying by the denominator: $6x^2+13x+8=A(x^2+2x+2)+(Bx+C)x$.
\begin{itemize}
  \item Substituting $x=0$: $8=2A$, hence $A=4$.
  \item Then $6x^2+13x+8-4(x^2+2x+2)=2x^2+5x=Bx^2+Cx$, hence $B=2$, $C=5$.
\end{itemize}
\[ \frac{6x^2+13x+8}{x^3+2x^2+2x}=\frac4x+\frac{2x+5}{x^2+2x+2}=\frac4x+\frac{2x+2}{x^2+2x+2}+\frac{3}{(x+1)^2+1}. \]

\textbf{Integration.}
\begin{itemize}
  \item $\int\frac4x\,dx=4\ln|x|$.
  \item $\int\frac{2x+2}{x^2+2x+2}\,dx=\ln(x^2+2x+2)$ (the numerator is the derivative of the denominator, and the denominator is positive).
  \item $\int\frac{3}{(x+1)^2+1}\,dx=3\arctan(x+1)$ (substitution $t=x+1$).
\end{itemize}

\textbf{Answer:}
\[ \int\frac{6x^2+13x+8}{x^3+2x^2+2x}dx=\boxed{4\ln|x|+\ln(x^2+2x+2)+3\arctan(x+1)+C}. \]
(Check: differentiating the result gives back the integrand.)
\end{solution}

\subsection*{Question 6a}

\begin{question}
Compute the finite area bounded between the function $f(x)=x^2$ and the line $y=-x+2$.
\end{question}

\begin{hint}
Find the intersection points, and determine which of the two functions is on top.
\end{hint}

\begin{solution}
\textbf{Intersection points:} $x^2=-x+2\iff x^2+x-2=0\iff (x+2)(x-1)=0$, i.e. $x=-2$ and $x=1$.
On the interval $(-2,1)$ we have $-x+2-x^2=-(x+2)(x-1)>0$, hence the line is above the parabola. The area:
\[ S=\int_{-2}^{1}\left(-x+2-x^2\right)dx=\left[-\frac{x^2}{2}+2x-\frac{x^3}{3}\right]_{-2}^{1}
=\left(-\frac12+2-\frac13\right)-\left(-2-4+\frac83\right)=\frac76+\frac{10}{3}=\boxed{\frac92}. \]
\end{solution}

\subsection*{Question 6b}

\begin{question}
Compute the improper integral $\int_0^\infty xe^{-x}dx$, or show that it does not converge.
\end{question}

\begin{hint}
Compute $\int_0^R xe^{-x}dx$ by integration by parts, and then let $R\to\infty$.
\end{hint}

\begin{solution}
By definition, $\int_0^\infty xe^{-x}dx=\lim_{R\to\infty}\int_0^R xe^{-x}dx$.
By integration by parts ($u=x$, $v'=e^{-x}$, i.e. $u'=1$, $v=-e^{-x}$):
\[ \int_0^R xe^{-x}dx=\Big[-xe^{-x}\Big]_0^R+\int_0^R e^{-x}dx=-Re^{-R}+\Big[-e^{-x}\Big]_0^R=-Re^{-R}-e^{-R}+1. \]
As $R\to\infty$: $e^{-R}\to0$, and also $Re^{-R}=\frac{R}{e^R}\to0$ (for example by L'Hôpital's rule: $\lim\frac{R}{e^R}=\lim\frac1{e^R}=0$). Hence
\[ \int_0^\infty xe^{-x}dx=\lim_{R\to\infty}\left(1-Re^{-R}-e^{-R}\right)=\boxed{1}, \]
and the integral converges.
\end{solution}

\section*{2021/22 Semester A Moed B}

\subsection*{Question 1}

\begin{question}
For arbitrary parameters $a,b\in\R$ define a function by
\[ f(x)=\begin{cases} \sin(ax)\sin(\frac{1}{x}) & x<0\\ b+3 & x=0\\ x\ln(x) & x>0 \end{cases} \]
For which values of the parameters $a,b$ is the function continuous at $0$? Justify.
\end{question}

\begin{hint}
Compute the one-sided limits at $0$ separately. From the left: ``tends to zero times bounded''. From the right: write $x\ln x=\frac{\ln x}{1/x}$ and use L'Hôpital's rule.
\end{hint}

\begin{solution}
$f$ is continuous at $0$ if and only if both one-sided limits exist and equal $f(0)=b+3$.

\textbf{Limit from the left.} For every $a\in\R$: $\sin(ax)\to\sin0=0$ as $x\to0^-$ (continuity of sine), and $|\sin(1/x)|\le1$. By ``tends to zero times bounded'':
\[ \lim_{x\to0^-}\sin(ax)\sin\left(\tfrac1x\right)=0. \]
(Explicitly: $0\le|\sin(ax)\sin(1/x)|\le|\sin(ax)|\le|a||x|\to0$.) The limit equals $0$ \textbf{for every} value of $a$.

\textbf{Limit from the right.} This is a limit of the form $0\cdot(-\infty)$. We write it as a quotient of the form $\frac{-\infty}{\infty}$ and use L'Hôpital's rule:
\[ \lim_{x\to0^+}x\ln x=\lim_{x\to0^+}\frac{\ln x}{1/x}\overset{L}{=}\lim_{x\to0^+}\frac{1/x}{-1/x^2}=\lim_{x\to0^+}(-x)=0. \]

\textbf{Conclusion.} Both one-sided limits equal $0$, and hence $\lim_{x\to0}f(x)=0$. Continuity at $0$ is equivalent to $b+3=0$.

\textbf{Answer:} $f$ is continuous at $0$ if and only if $\boxed{b=-3}$, with $a\in\R$ arbitrary.
\end{solution}

\subsection*{Question 2a}

\begin{question}
Find the local extremum points and the intervals of increase and decrease of the function
\[ f(x)=\frac{2x}{x^2+1} \]
\end{question}

\begin{hint}
Compute $f'$ by the quotient rule and check its sign; the denominator of $f'$ is always positive.
\end{hint}

\begin{solution}
$f$ is defined and differentiable on all of $\R$ (the denominator is positive). By the quotient rule:
\[ f'(x)=\frac{2(x^2+1)-2x\cdot2x}{(x^2+1)^2}=\frac{2(1-x^2)}{(x^2+1)^2}. \]
$f'(x)=0\iff x=\pm1$. The sign of $f'$ is the sign of $1-x^2$:
\begin{itemize}
  \item $f'<0$ on $(-\infty,-1)$: $f$ is \textbf{decreasing}.
  \item $f'>0$ on $(-1,1)$: $f$ is \textbf{increasing}.
  \item $f'<0$ on $(1,\infty)$: $f$ is \textbf{decreasing}.
\end{itemize}
Hence at $x=-1$ there is a \textbf{local minimum} $f(-1)=-1$, and at $x=1$ there is a \textbf{local maximum} $f(1)=1$.
(In fact these are also global extrema: $|2x|\le x^2+1$ since $(|x|-1)^2\ge0$, and hence $-1\le f\le 1$.)
\end{solution}

\subsection*{Question 2b}

\begin{question}
Compute the absolute minimum and maximum values of the following function on its domain:
\[ f(x)=\frac{\sqrt{x^2-1}}{x^2} \]
or explain why they do not exist.
\end{question}

\begin{hint}
First find the domain, and note that the function is even and non-negative. Also check the limit at infinity.
\end{hint}

\begin{solution}
\textbf{Domain:} $x^2-1\ge0$ and $x\neq0$, i.e. $D=(-\infty,-1]\cup[1,\infty)$.
$f$ is even, so it suffices to study $[1,\infty)$.

\textbf{Absolute minimum:} $f(x)\ge0$ for every $x\in D$, and $f(\pm1)=0$. Hence the absolute minimum is $\boxed{0}$, attained at $x=\pm1$.

\textbf{Absolute maximum:} for $x>1$:
\[ f'(x)=\frac{\frac{x}{\sqrt{x^2-1}}\cdot x^2-2x\sqrt{x^2-1}}{x^4}=\frac{x^2-2(x^2-1)}{x^3\sqrt{x^2-1}}=\frac{2-x^2}{x^3\sqrt{x^2-1}}. \]
Hence $f'>0$ on $(1,\sqrt2)$ and $f'<0$ on $(\sqrt2,\infty)$: $f$ is increasing on $[1,\sqrt2]$ and decreasing on $[\sqrt2,\infty)$ (it is continuous at the endpoint $1$). It follows that $f(x)\le f(\sqrt2)$ for every $x\ge1$, and
\[ f(\sqrt2)=\frac{\sqrt{2-1}}{2}=\frac12. \]
(Note also that $\lim_{x\to\infty}f(x)=\lim\frac{\sqrt{1-1/x^2}}{x}=0$, so nothing ``escapes'' at infinity.) By evenness, the same holds on $(-\infty,-1]$.

\textbf{Answer:} the absolute maximum is $\boxed{\tfrac12}$, attained at $x=\pm\sqrt2$; the absolute minimum is $\boxed{0}$, attained at $x=\pm1$.
\end{solution}

\subsection*{Question 3}

\begin{question}
Approximate the value of $\sqrt[4]{e}$ with a remainder less than $\frac{1}{1000}$. You must justify your answer carefully.
\end{question}

\begin{hint}
$\sqrt[4]{e}=e^{1/4}$. Use the Maclaurin polynomial of $e^x$ and the Lagrange remainder, and find the smallest order $n$ that guarantees an error less than $\frac{1}{1000}$.
\end{hint}

\begin{hint}
To bound $e^c$ for $0<c<\frac14$ you can use $e^c<e<3$.
\end{hint}

\begin{solution}
Write $\sqrt[4]{e}=e^{1/4}$ and use the Maclaurin polynomial of $f(x)=e^x$. All the derivatives are $f^{(k)}(x)=e^x$, hence $f^{(k)}(0)=1$ and
\[ P_n(x)=\sum_{k=0}^n\frac{x^k}{k!}. \]
By Taylor's theorem with the Lagrange remainder, for every $n$ there exists $c\in(0,\tfrac14)$ such that
\[ R_n\left(\tfrac14\right)=\frac{e^c}{(n+1)!}\left(\tfrac14\right)^{n+1}. \]
Since $e^x$ is increasing, $e^c<e^{1/4}<e<3$, and hence
\[ 0<R_n\left(\tfrac14\right)<\frac{3}{(n+1)!\,4^{n+1}}. \]
\begin{itemize}
  \item $n=2$: the bound is $\frac{3}{6\cdot64}=\frac{1}{128}$ --- not enough.
  \item $n=3$: the bound is $\frac{3}{24\cdot256}=\frac{1}{2048}<\frac{1}{1000}$ --- enough.
\end{itemize}
(And indeed $n=2$ really is not enough, not just its bound: $R_2=\frac{e^c}{6\cdot64}>\frac{1}{384}>\frac1{1000}$, since $e^c>1$.)

Hence we choose $n=3$:
\[ \sqrt[4]{e}\approx P_3\left(\tfrac14\right)=1+\frac14+\frac{1}{2\cdot16}+\frac{1}{6\cdot64}=1+\frac14+\frac1{32}+\frac1{384}=\frac{493}{384}\approx1.28385, \]
and the error satisfies $0<\sqrt[4]{e}-\frac{493}{384}<\frac{1}{2048}<\frac{1}{1000}$.

\textbf{Answer:} $\boxed{\sqrt[4]{e}\approx\frac{493}{384}\approx1.2839}$ (the true value is $\approx1.28403$).
\end{solution}

\subsection*{Question 4a}

\begin{question}
Compute the antiderivative of
\[ f(x)=\cos^2x\sin^3x \]
\end{question}

\begin{hint}
Write $\sin^3x=(1-\cos^2x)\sin x$ and substitute $t=\cos x$.
\end{hint}

\begin{solution}
Write $\cos^2x\sin^3x=\cos^2x(1-\cos^2x)\sin x$ and substitute $t=\cos x$, $dt=-\sin x\,dx$:
\[ \int\cos^2x\sin^3x\,dx=-\int t^2(1-t^2)\,dt=-\frac{t^3}{3}+\frac{t^5}{5}+C=\boxed{\frac{\cos^5x}{5}-\frac{\cos^3x}{3}+C}. \]
Check: $\left(\frac{\cos^5x}{5}-\frac{\cos^3x}{3}\right)'=-\cos^4x\sin x+\cos^2x\sin x=\cos^2x\sin x(1-\cos^2x)=\cos^2x\sin^3x$.
\end{solution}

\subsection*{Question 4b}

\begin{question}
Solve the differential equation
\[ xe^y=(x^2+1)y' \]
\end{question}

\begin{hint}
This is a separable equation: $e^{-y}y'=\frac{x}{x^2+1}$.
\end{hint}

\begin{solution}
Since $x^2+1>0$ and $e^y>0$, we can write
\[ e^{-y}\,y'=\frac{x}{x^2+1}. \]
(There are no constant solutions: for $y\equiv c$ we get $xe^c=0$ for every $x$, which is impossible.) Integrate both sides:
\[ \int e^{-y}dy=\int\frac{x}{x^2+1}dx\ \Longrightarrow\ -e^{-y}=\frac12\ln(x^2+1)-C. \]
Hence $e^{-y}=C-\frac12\ln(x^2+1)$, and therefore
\[ \boxed{y=-\ln\left(C-\tfrac12\ln(x^2+1)\right)},\qquad C\in\R, \]
on the domain where $C-\frac12\ln(x^2+1)>0$ (in particular $C>0$ is required; the domain is $|x|<\sqrt{e^{2C}-1}$).

Check: $y'=\frac{\frac{x}{x^2+1}}{C-\frac12\ln(x^2+1)}=\frac{x}{x^2+1}e^{y}$, and hence $(x^2+1)y'=xe^y$.
\end{solution}

\subsection*{Question 5}

\begin{question}
Compute the integral
\[ \int\frac{5x^3+4x^2+8x+12}{x^4+4x^2}dx \]
\end{question}

\begin{hint}
$x^4+4x^2=x^2(x^2+4)$. The partial fraction decomposition is of the form $\frac Ax+\frac B{x^2}+\frac{Cx+D}{x^2+4}$.
\end{hint}

\begin{solution}
\textbf{Partial fraction decomposition.} $x^4+4x^2=x^2(x^2+4)$, and $x^2+4$ is irreducible over $\R$. The degree of the numerator is less than the degree of the denominator, and hence
\[ \frac{5x^3+4x^2+8x+12}{x^2(x^2+4)}=\frac Ax+\frac B{x^2}+\frac{Cx+D}{x^2+4}, \]
\[ 5x^3+4x^2+8x+12=Ax(x^2+4)+B(x^2+4)+(Cx+D)x^2=(A+C)x^3+(B+D)x^2+4Ax+4B. \]
Comparing coefficients:
\[ 4B=12\Rightarrow B=3,\quad 4A=8\Rightarrow A=2,\quad A+C=5\Rightarrow C=3,\quad B+D=4\Rightarrow D=1. \]
Hence
\[ \frac{5x^3+4x^2+8x+12}{x^4+4x^2}=\frac2x+\frac3{x^2}+\frac{3x}{x^2+4}+\frac{1}{x^2+4}. \]

\textbf{Integration.}
\begin{itemize}
  \item $\int\frac2x\,dx=2\ln|x|$, \quad $\int\frac3{x^2}\,dx=-\frac3x$.
  \item $\int\frac{3x}{x^2+4}\,dx=\frac32\int\frac{2x}{x^2+4}\,dx=\frac32\ln(x^2+4)$.
  \item $\int\frac{dx}{x^2+4}=\frac14\int\frac{dx}{(x/2)^2+1}=\frac12\arctan\frac x2$ (substitution $t=\frac x2$).
\end{itemize}

\textbf{Answer:}
\[ \int\frac{5x^3+4x^2+8x+12}{x^4+4x^2}dx=\boxed{2\ln|x|-\frac3x+\frac32\ln(x^2+4)+\frac12\arctan\frac x2+C}. \]
\end{solution}

\subsection*{Question 6a}

\begin{question}
Compute the volume of the solid of revolution of the function $f(x)=\sin x$ around the $x$-axis on the interval $[0,\pi]$.
\end{question}

\begin{hint}
Use the formula $V=\pi\int_a^b f^2(x)\,dx$ and the identity $\sin^2x=\frac{1-\cos2x}{2}$.
\end{hint}

\begin{solution}
The volume of the solid of revolution around the $x$-axis:
\[ V=\pi\int_0^\pi\sin^2x\,dx=\pi\int_0^\pi\frac{1-\cos2x}{2}\,dx=\pi\left[\frac x2-\frac{\sin2x}{4}\right]_0^\pi=\pi\cdot\frac\pi2=\boxed{\frac{\pi^2}{2}}. \]
\end{solution}

\subsection*{Question 6b}

\begin{question}
Determine the convergence of the improper integral
\[ \int_1^\infty\frac{\sin(\frac1x)}{2+x\sqrt{x}}dx \]
\end{question}

\begin{hint}
The integrand is positive on the interval. Use the comparison test, with $0<\sin t\le t$ for $0<t\le1$.
\end{hint}

\begin{solution}
For $x\ge1$ we have $0<\frac1x\le1<\pi$, hence $\sin\frac1x>0$, and the integrand is positive and continuous on $[1,\infty)$ (so the only problem is at infinity). From the inequality $\sin t\le t$ for $t\ge0$:
\[ 0<\frac{\sin(1/x)}{2+x\sqrt x}\le\frac{1/x}{x\sqrt x}=\frac{1}{x^{5/2}}. \]
The integral $\int_1^\infty\frac{dx}{x^{p}}$ converges for $p>1$, and in particular for $p=\frac52$. By the \textbf{comparison test} for improper integrals of non-negative functions, the integral
\[ \int_1^\infty\frac{\sin(1/x)}{2+x\sqrt x}dx \]
\textbf{converges}.

(Alternative method: the limit comparison test with $g(x)=x^{-5/2}$: $\lim_{x\to\infty}\frac{f(x)}{g(x)}=\lim\frac{\sin(1/x)}{1/x}\cdot\frac{x^{3/2}}{2+x^{3/2}}=1$.)
\end{solution}

\section*{2021/22 Semester A Special Moed}

\subsection*{Question 1}

\begin{question}
For arbitrary parameters $a,b\in\R$ define a function by
\[ f(x)=\begin{cases} e^{\frac1x} & x<0\\ ax+b-6 & 0\le x\le 9\\ \frac{x^2-8x-9}{3-\sqrt{x}} & x>9 \end{cases} \]
For which values of the parameters $a,b$ is the function continuous on $\R$? Justify.
\end{question}

\begin{hint}
On each of the open intervals, $f$ is a composition/quotient of continuous functions. It remains to check the junction points $x=0$ and $x=9$.
\end{hint}

\begin{hint}
At $x=9$: $x^2-8x-9=(x-9)(x+1)$ and $x-9=(\sqrt x-3)(\sqrt x+3)$.
\end{hint}

\begin{solution}
\textbf{Continuity inside the intervals.} On $(-\infty,0)$ the function $e^{1/x}$ is continuous (a composition of continuous functions). On $(0,9)$ $f$ is a polynomial. On $(9,\infty)$ $f$ is a quotient of continuous functions whose denominator $3-\sqrt x<0$ does not vanish. Hence $f$ is continuous at every $x\neq0,9$ for all $a,b$, and we only need to check $x=0$ and $x=9$.

\textbf{The point $x=0$.} $f(0)=b-6$, and the limit from the right is $\lim_{x\to0^+}(ax+b-6)=b-6$. From the left: as $x\to0^-$ we have $\frac1x\to-\infty$, and hence
\[ \lim_{x\to0^-}e^{1/x}=0. \]
Continuity at $0$ is equivalent to $b-6=0$, i.e. $b=6$.

\textbf{The point $x=9$.} $f(9)=9a+b-6$, and this is also the limit from the left. From the right, for $x>9$:
\[ \frac{x^2-8x-9}{3-\sqrt x}=\frac{(x-9)(x+1)}{3-\sqrt x}=\frac{(\sqrt x-3)(\sqrt x+3)(x+1)}{-(\sqrt x-3)}=-(\sqrt x+3)(x+1), \]
and hence
\[ \lim_{x\to9^+}f(x)=-(3+3)(9+1)=-60. \]
Continuity at $9$ is equivalent to $9a+b-6=-60$.

\textbf{Solving the system.} From $b=6$ we get $9a=-60$, i.e. $a=-\frac{20}{3}$.

\textbf{Answer:} $f$ is continuous on all of $\R$ if and only if $\boxed{a=-\tfrac{20}{3},\ b=6}$.
\end{solution}

\subsection*{Question 2a}

\begin{question}
Find all the local and absolute (global) extremum points on the whole line of the function
\[ f(x)=x^2\cdot e^{-x} \]
\end{question}

\begin{hint}
Check the limits at $\pm\infty$ to decide whether there is an absolute extremum.
\end{hint}

\begin{solution}
$f$ is differentiable on all of $\R$, and
\[ f'(x)=2xe^{-x}-x^2e^{-x}=x(2-x)e^{-x}. \]
Since $e^{-x}>0$, $f'(x)=0\iff x=0$ or $x=2$, and the sign of $f'$ is the sign of $x(2-x)$:
\begin{itemize}
  \item $f'<0$ on $(-\infty,0)$: $f$ is decreasing.
  \item $f'>0$ on $(0,2)$: $f$ is increasing.
  \item $f'<0$ on $(2,\infty)$: $f$ is decreasing.
\end{itemize}
Hence at $x=0$ there is a \textbf{local minimum} $f(0)=0$, and at $x=2$ there is a \textbf{local maximum} $f(2)=\frac{4}{e^2}$.

\textbf{Absolute extrema.}
\begin{itemize}
  \item $f(x)=x^2e^{-x}\ge0=f(0)$ for every $x$, and hence $x=0$ is an \textbf{absolute minimum}, with value $0$.
  \item $\lim_{x\to-\infty}x^2e^{-x}=+\infty$ (both factors tend to $+\infty$), hence $f$ is not bounded above and \textbf{there is no absolute maximum}. In particular, the local maximum at $x=2$ is not absolute (for example $f(-2)=4e^2>\frac4{e^2}$).
\end{itemize}
(For completeness: $\lim_{x\to+\infty}\frac{x^2}{e^x}=0$ by L'Hôpital's rule twice.)
\end{solution}

\subsection*{Question 2b}

\begin{question}
Prove that the function $f(x)=\frac15x^5+\frac23x^3+x+7$ has a unique real root.
\end{question}

\begin{hint}
Existence follows from the Intermediate Value Theorem, and uniqueness from monotonicity: $f'(x)=x^4+2x^2+1$ (or from Rolle's theorem).
\end{hint}

\begin{solution}
\textbf{Existence.} $f$ is a polynomial and hence continuous on all of $\R$. We have $f(0)=7>0$ and
\[ f(-2)=-\frac{32}{5}-\frac{16}{3}-2+7=-\frac{101}{15}<0. \]
By the \textbf{Intermediate Value Theorem} there exists $x_0\in(-2,0)$ with $f(x_0)=0$.

\textbf{Uniqueness.}
\[ f'(x)=x^4+2x^2+1=(x^2+1)^2>0\quad\text{for every } x, \]
and hence $f$ is strictly increasing on $\R$, and in particular one-to-one --- it attains the value $0$ at most once.
(Alternatively, by Rolle's theorem: if there were two roots $x_1<x_2$, there would exist $c\in(x_1,x_2)$ with $f'(c)=0$, contradicting $f'>0$.)

Hence $f$ has a unique real root, and it lies in the interval $(-2,0)$.
\end{solution}

\subsection*{Question 3}

\begin{question}
Estimate $\arctan(0.2)$ using a linear approximation (a first-order Taylor polynomial) and give a bound for the error.

You must justify your answer carefully.
\end{question}

\begin{hint}
Expand $f(x)=\arctan x$ around $a=0$ and bound the error using the Lagrange remainder $R_1=\frac{f''(c)}{2}x^2$.
\end{hint}

\begin{solution}
Expand $f(x)=\arctan x$ around $a=0$:
\[ f'(x)=\frac{1}{1+x^2},\qquad f''(x)=-\frac{2x}{(1+x^2)^2}. \]
$f(0)=0$, $f'(0)=1$, hence $P_1(x)=x$ and the approximation is
\[ \arctan(0.2)\approx P_1(0.2)=\boxed{0.2}. \]

\textbf{Error bound.} By Taylor's theorem with the Lagrange remainder, there exists $c\in(0,0.2)$ such that
\[ R_1(0.2)=\frac{f''(c)}{2}(0.2)^2=-\frac{c}{(1+c^2)^2}\cdot0.04. \]
For $0<c<0.2$ we have $(1+c^2)^2>1$, hence $\frac{c}{(1+c^2)^2}<c<0.2$, and therefore
\[ |R_1(0.2)|<0.2\cdot0.04=0.008. \]

\textbf{Answer:} $\arctan(0.2)\approx0.2$ with an error less than $0.008$. Moreover $R_1<0$, and hence $0.192<\arctan(0.2)<0.2$.

\textbf{Remark (a better bound).} Since $f''(0)=0$, we have $P_2=P_1$, so we can use the second-order remainder: $R_2(0.2)=\frac{f'''(c)}{6}(0.2)^3$ with $f'''(x)=\frac{6x^2-2}{(1+x^2)^3}$. For $0<c<0.2$ we have $|6c^2-2|\le 2$ and $(1+c^2)^3>1$, hence $|f'''(c)|<2$ and
\[ |R_2(0.2)|<\frac{2}{6}\cdot0.008=\frac{0.008}{3}\approx0.0027. \]
(Check: $\arctan(0.2)\approx0.19740$, and the actual error is $\approx0.0026$.)
\end{solution}

\subsection*{Question 4a}

\begin{question}
Compute the antiderivative of
\[ f(x)=\cos(3x)\cos(7x) \]
\end{question}

\begin{hint}
Use the identity $\cos\alpha\cos\beta=\frac12\left[\cos(\alpha-\beta)+\cos(\alpha+\beta)\right]$.
\end{hint}

\begin{solution}
By the identity $\cos\alpha\cos\beta=\frac12\left[\cos(\alpha-\beta)+\cos(\alpha+\beta)\right]$:
\[ \cos(3x)\cos(7x)=\frac12\left[\cos(4x)+\cos(10x)\right]. \]
Hence
\[ \int\cos(3x)\cos(7x)\,dx=\frac12\left[\frac{\sin(4x)}{4}+\frac{\sin(10x)}{10}\right]+C=\boxed{\frac{\sin(4x)}{8}+\frac{\sin(10x)}{20}+C}. \]
\end{solution}

\subsection*{Question 4b}

\begin{question}
Solve the differential equation
\[ y'=e^{x-y} \]
\end{question}

\begin{hint}
$e^{x-y}=e^x\cdot e^{-y}$, so this is a separable equation.
\end{hint}

\begin{solution}
$y'=e^x e^{-y}$, hence $e^{y}y'=e^x$ (there are no constant solutions, since $e^{x-y}>0$). Integrating both sides:
\[ \int e^y\,dy=\int e^x\,dx\ \Longrightarrow\ e^y=e^x+C. \]
Hence
\[ \boxed{y=\ln\left(e^x+C\right)},\qquad C\in\R, \]
on the domain where $e^x+C>0$ (for every $x$ if $C\ge0$; for $C<0$, on the domain $x>\ln(-C)$).

Check: $y'=\frac{e^x}{e^x+C}=\frac{e^x}{e^y}=e^{x-y}$.
\end{solution}

\subsection*{Question 5}

\begin{question}
Compute the integral
\[ \int\frac{5x^2+2x+2}{x^3-1}dx \]
\end{question}

\begin{hint}
$x^3-1=(x-1)(x^2+x+1)$, and the quadratic factor is irreducible. Decompose into partial fractions.
\end{hint}

\begin{solution}
\textbf{Partial fraction decomposition.} $x^3-1=(x-1)(x^2+x+1)$, and the factor $x^2+x+1$ has discriminant $1-4<0$, hence it is irreducible. The degree of the numerator is less than the degree of the denominator:
\[ \frac{5x^2+2x+2}{(x-1)(x^2+x+1)}=\frac{A}{x-1}+\frac{Bx+C}{x^2+x+1}, \]
\[ 5x^2+2x+2=A(x^2+x+1)+(Bx+C)(x-1). \]
\begin{itemize}
  \item Substituting $x=1$: $9=3A$, hence $A=3$.
  \item Then $5x^2+2x+2-3(x^2+x+1)=2x^2-x-1=(x-1)(2x+1)$, hence $Bx+C=2x+1$, i.e. $B=2$, $C=1$.
\end{itemize}
\[ \frac{5x^2+2x+2}{x^3-1}=\frac{3}{x-1}+\frac{2x+1}{x^2+x+1}. \]

\textbf{Integration.} $\int\frac{3}{x-1}dx=3\ln|x-1|$, and the numerator $2x+1$ is exactly the derivative of the denominator $x^2+x+1>0$, hence $\int\frac{2x+1}{x^2+x+1}dx=\ln(x^2+x+1)$.

\textbf{Answer:}
\[ \int\frac{5x^2+2x+2}{x^3-1}dx=\boxed{3\ln|x-1|+\ln(x^2+x+1)+C}. \]
\end{solution}

\subsection*{Question 6a}

\begin{question}
Compute the finite area bounded between the graph of the function $f(x)=x^3-x$ and the $x$-axis.
\end{question}

\begin{hint}
Note that the function changes sign between its intersection points with the axis --- the area of each part must be computed separately (or use oddness).
\end{hint}

\begin{solution}
\textbf{Intersection with the $x$-axis:} $x^3-x=x(x-1)(x+1)=0\iff x\in\{-1,0,1\}$.
On the interval $(-1,0)$ we have $f>0$, and on the interval $(0,1)$ we have $f<0$. Outside $[-1,1]$ the region between the graph and the axis is unbounded, so the finite bounded area consists of the two ``leaves'' between $-1$ and $1$:
\[ S=\int_{-1}^0(x^3-x)\,dx+\int_0^1(x-x^3)\,dx. \]
\[ \int_0^1(x-x^3)\,dx=\left[\frac{x^2}{2}-\frac{x^4}{4}\right]_0^1=\frac14,\qquad \int_{-1}^0(x^3-x)\,dx=\left[\frac{x^4}{4}-\frac{x^2}{2}\right]_{-1}^0=-\left(\frac14-\frac12\right)=\frac14. \]
(The equality also follows from the fact that $f$ is odd.) Hence
\[ S=\frac14+\frac14=\boxed{\frac12}. \]
(Note: $\int_{-1}^1(x^3-x)\,dx=0$ --- this is not the area.)
\end{solution}

\subsection*{Question 6b}

\begin{question}
Determine the convergence of the improper integral
\[ \int_1^\infty\frac{e^x}{e^{2x}+1}dx \]
\end{question}

\begin{hint}
The antiderivative is $\arctan(e^x)$ (substitution $t=e^x$).
\end{hint}

\begin{solution}
Substitute $t=e^x$, $dt=e^xdx$: $\int\frac{e^x}{e^{2x}+1}dx=\int\frac{dt}{t^2+1}=\arctan(e^x)+C$. Hence
\[ \int_1^R\frac{e^x}{e^{2x}+1}dx=\arctan(e^R)-\arctan(e). \]
As $R\to\infty$: $e^R\to\infty$ and $\lim_{t\to\infty}\arctan t=\frac\pi2$, and hence
\[ \int_1^\infty\frac{e^x}{e^{2x}+1}dx=\lim_{R\to\infty}\left(\arctan(e^R)-\arctan e\right)=\boxed{\frac\pi2-\arctan e}\approx0.3466. \]
The integral \textbf{converges}.

(Alternatively, without computation: $0<\frac{e^x}{e^{2x}+1}<\frac{e^x}{e^{2x}}=e^{-x}$, and $\int_1^\infty e^{-x}dx=e^{-1}$ converges, hence by the comparison test the integral converges.)
\end{solution}

\section*{2022/23 Semester A Moed A}

\subsection*{Question 1a}

\begin{question}
Check whether an inverse function exists for the function $f(x)=|x|+x$ on its domain. If it exists, find it.
\end{question}

\begin{hint}
Open the absolute value: what is the value of $f$ for $x<0$? A function has an inverse only if it is one-to-one.
\end{hint}

\begin{solution}
The function is defined for every $x\in\R$. Open the absolute value:
\[ f(x)=\begin{cases} 2x & x\ge 0\\ 0 & x<0.\end{cases} \]
A function has an inverse function (on its domain) if and only if it is one-to-one. But, for example, $f(-1)=f(-2)=0$ while $-1\neq -2$, hence $f$ is not one-to-one on $\R$, and \textbf{it has no inverse function} on its domain.

(Remark: if we restrict $f$ to the ray $[0,\infty)$, there it is $f(x)=2x$, one-to-one and onto $[0,\infty)$, and the inverse of the restriction is $f^{-1}(y)=\frac{y}{2}$, $y\ge 0$. But this is no longer $f$ on its whole domain.)
\end{solution}

\subsection*{Question 1b}

\begin{question}
For which values of the parameter $a$ does the following function have a limit at $x=0$?
\[ f(x)=\begin{cases} \dfrac{\sin(2x)}{x} & x<0 \\[2mm] \dfrac{x+3}{4x+a} & x\ge 0 \end{cases} \]
\end{question}

\begin{hint}
Compute the one-sided limits separately and require them to be equal. Pay attention to the case $a=0$.
\end{hint}

\begin{solution}
The function has a limit at $x=0$ if and only if both one-sided limits exist (are finite) and are equal.

\textbf{Limit from the left:} by the fundamental limit $\lim_{t\to 0}\frac{\sin t}{t}=1$,
\[ \lim_{x\to 0^-}\frac{\sin(2x)}{x}=\lim_{x\to 0^-}2\cdot\frac{\sin(2x)}{2x}=2. \]

\textbf{Limit from the right:} if $a\neq 0$, the numerator and the denominator are continuous at $0$ and the denominator is different from $0$ there, hence
\[ \lim_{x\to 0^+}\frac{x+3}{4x+a}=\frac{3}{a}. \]
If $a=0$, then $\frac{x+3}{4x}\to +\infty$ as $x\to 0^+$ (the numerator tends to $3$ and the denominator to $0^+$), and the limit from the right is not finite, so there is no limit.

Hence for $a\ne 0$ we need $\frac{3}{a}=2$, i.e. $a=\frac32$. (In this case the denominator $4x+\frac32>0$ for every $x\ge 0$, so the function is well defined in a right neighborhood of $0$.)

\textbf{Answer:} the function has a limit at $x=0$ only for $a=\frac{3}{2}$, and the value of the limit is $2$.
\end{solution}

\subsection*{Question 2a}

\begin{question}
Use a linear approximation to estimate the expression $\sqrt[3]{28}$.
\end{question}

\begin{hint}
Choose a point close to $28$ at which the cube root is known.
\end{hint}

\begin{solution}
Define $g(x)=\sqrt[3]{x}=x^{1/3}$ and the base point $x_0=27$, at which $g(27)=3$. The derivative:
\[ g'(x)=\frac{1}{3}x^{-2/3}=\frac{1}{3\sqrt[3]{x^2}},\qquad g'(27)=\frac{1}{3\cdot 9}=\frac{1}{27}. \]
The linear approximation (the tangent line) near $x_0$: $g(x)\approx g(x_0)+g'(x_0)(x-x_0)$. With $x=28$, $x-x_0=1$:
\[ \sqrt[3]{28}\approx 3+\frac{1}{27}\approx 3.037. \]
(The true value is $3.0366\ldots$.)
\end{solution}

\subsection*{Question 2b}

\begin{question}
Let $f(x)$ be a function continuous on the whole line, and assume that $|f(x)|\le 7$ for every $x$. Prove that the equation $2x+f(x)=3$ has at least one solution.
\end{question}

\begin{hint}
Define $g(x)=2x+f(x)-3$ and find a point where $g$ is negative and a point where $g$ is positive, using the bound $|f(x)|\le 7$.
\end{hint}

\begin{solution}
Define $g(x)=2x+f(x)-3$. The function $g$ is continuous on all of $\R$ as a sum of continuous functions. By assumption $-7\le f(x)\le 7$ for every $x$, and hence:
\[ g(-10)=-20+f(-10)-3\le -23+7=-16<0,\qquad g(10)=20+f(10)-3\ge 17-7=10>0. \]
$g$ is continuous on the closed interval $[-10,10]$ and takes values of opposite signs at its endpoints, hence by the \textbf{Intermediate Value Theorem} (Bolzano's theorem) there exists $c\in(-10,10)$ with $g(c)=0$, i.e. $2c+f(c)=3$. Hence the equation has at least one solution. $\blacksquare$
\end{solution}

\subsection*{Question 3a}

\begin{question}
Find the intervals of convexity and the inflection points of the function $f(x)=\ln(1+4x^2)$.
\end{question}

\begin{solution}
Domain: $1+4x^2>0$ for every $x$, hence $D=\R$. Differentiate:
\[ f'(x)=\frac{8x}{1+4x^2},\qquad
f''(x)=\frac{8(1+4x^2)-8x\cdot 8x}{(1+4x^2)^2}=\frac{8-32x^2}{(1+4x^2)^2}=\frac{8(1-2x)(1+2x)}{(1+4x^2)^2}. \]
The denominator is always positive, hence the sign of $f''$ is the sign of $1-4x^2$:
\begin{itemize}
  \item $f''(x)>0$ for $-\frac12<x<\frac12$: the function is \textbf{convex} (concave up, $\cup$) on the interval $\left(-\frac12,\frac12\right)$.
  \item $f''(x)<0$ for $|x|>\frac12$: the function is \textbf{concave} ($\cap$) on $\left(-\infty,-\frac12\right)$ and $\left(\frac12,\infty\right)$.
\end{itemize}
At $x=\pm\frac12$ the second derivative changes sign and the function is continuous (and differentiable) there, hence these are inflection points. $f(\pm\tfrac12)=\ln(1+1)=\ln 2$.

\textbf{Answer:} convex on $\left(-\frac12,\frac12\right)$, concave on $\left(-\infty,-\frac12\right)\cup\left(\frac12,\infty\right)$; inflection points $\left(-\frac12,\ln 2\right)$ and $\left(\frac12,\ln 2\right)$.
\end{solution}

\subsection*{Question 3b}

\begin{question}
Prove that for every $x,y$ we have $|\arctan x-\arctan y|\le|x-y|$.
\end{question}

\begin{hint}
Apply Lagrange's theorem to $\arctan$ on the interval between $x$ and $y$.
\end{hint}

\begin{solution}
If $x=y$ both sides equal $0$ and there is nothing to prove. Assume $x\neq y$, and without loss of generality $x<y$. The function $g(t)=\arctan t$ is continuous on $[x,y]$ and differentiable on $(x,y)$, hence by \textbf{Lagrange's theorem} there exists $c\in(x,y)$ such that
\[ \arctan y-\arctan x=g'(c)(y-x)=\frac{1}{1+c^2}(y-x). \]
Since $1+c^2\ge 1$, we have $0<\frac{1}{1+c^2}\le 1$, and hence
\[ |\arctan x-\arctan y|=\frac{1}{1+c^2}\,|x-y|\le|x-y|. \qquad\blacksquare \]
\end{solution}

\subsection*{Question 4a}

\begin{question}
Compute the limit $\displaystyle\lim_{x\to 1}\frac{x^2+\ln(x)-1}{e^x-e}$.
\end{question}

\begin{solution}
Substituting $x=1$: the numerator is $1+0-1=0$ and the denominator is $e-e=0$, i.e. an expression of the form $\frac00$. The numerator and the denominator are differentiable near $1$, and the derivative of the denominator is $e^x\neq 0$. By \textbf{L'Hôpital's rule}:
\[ \lim_{x\to1}\frac{x^2+\ln x-1}{e^x-e}=\lim_{x\to1}\frac{2x+\frac1x}{e^x}=\frac{2+1}{e}=\frac{3}{e}, \]
where the last limit is computed by substitution (continuity), and hence the original limit also exists and equals $\frac{3}{e}$.
\end{solution}

\subsection*{Question 4b}

\begin{question}
Compute $\sin 0.1$ with an error not exceeding $0.001$.
\end{question}

\begin{hint}
Use the Maclaurin polynomial of $\sin x$ and estimate the remainder in Lagrange form.
\end{hint}

\begin{solution}
We use the Maclaurin polynomial of $f(x)=\sin x$. We have $f(0)=0$, $f'(0)=\cos 0=1$, $f''(0)=-\sin 0=0$, and hence
\[ P_2(x)=x. \]
By Taylor's formula with the Lagrange remainder, there exists $c$ between $0$ and $x$ such that
\[ \sin x=P_2(x)+R_2(x),\qquad R_2(x)=\frac{f'''(c)}{3!}x^3=\frac{-\cos c}{6}x^3. \]
For $x=0.1$, and since $|\cos c|\le 1$:
\[ |R_2(0.1)|\le\frac{(0.1)^3}{6}=\frac{0.001}{6}\approx 0.000167<0.001. \]
Hence $\sin 0.1\approx 0.1$, with an error less than $0.001$, as required.

(For a more accurate approximation one can take $P_4(x)=x-\frac{x^3}{6}$, and then $\sin 0.1\approx 0.0998333$ with an error $\le\frac{(0.1)^5}{120}<10^{-7}$; the true value is $0.0998334\ldots$.)
\end{solution}

\subsection*{Question 5a}

\begin{question}
Compute the integral $\displaystyle\int\sqrt{2x+3}\,dx$.
\end{question}

\begin{solution}
Substitute $t=2x+3$; then $dt=2\,dx$, i.e. $dx=\frac{dt}{2}$:
\[ \int\sqrt{2x+3}\,dx=\frac12\int t^{1/2}\,dt=\frac12\cdot\frac{t^{3/2}}{3/2}+C=\frac13(2x+3)^{3/2}+C. \]
Check: $\left(\frac13(2x+3)^{3/2}\right)'=\frac13\cdot\frac32(2x+3)^{1/2}\cdot 2=\sqrt{2x+3}$.
\end{solution}

\subsection*{Question 5b}

\begin{question}
Compute the integral $\displaystyle\int_{-1}^{1}\tan x\,dx$.
\end{question}

\begin{hint}
Check whether the function $\tan x$ is even or odd, and whether it is continuous on the interval $[-1,1]$.
\end{hint}

\begin{solution}
Since $1<\frac{\pi}{2}$, the function $\tan x=\frac{\sin x}{\cos x}$ is continuous on the interval $[-1,1]$ (there $\cos x>0$), and hence the definite integral exists. It is odd: $\tan(-x)=-\tan x$. The integral of a continuous odd function over a symmetric interval $[-a,a]$ equals $0$. Directly: an antiderivative is $-\ln|\cos x|$ (since $(-\ln\cos x)'=\frac{\sin x}{\cos x}$), and by \textbf{the Fundamental Theorem} (Newton--Leibniz):
\[ \int_{-1}^{1}\tan x\,dx=\Big[-\ln(\cos x)\Big]_{-1}^{1}=-\ln(\cos 1)+\ln(\cos(-1))=0, \]
since $\cos$ is even. \textbf{Answer:} $0$.
\end{solution}

\subsection*{Question 6a}

\begin{question}
Compute the area bounded between the graphs of the functions $f(x)=\sin^2x$ and $g(x)=\frac14$ on the interval $[0,\pi]$.
\end{question}

\begin{hint}
Find the intersection points $\sin^2x=\frac14$ in the interval, and split the interval according to which function is larger. Use the identity $\sin^2x=\frac{1-\cos 2x}{2}$.
\end{hint}

\begin{solution}
\textbf{Intersection points:} $\sin^2x=\frac14\iff\sin x=\pm\frac12$. On the interval $[0,\pi]$ we have $\sin x\ge 0$, hence $\sin x=\frac12$, i.e. $x=\frac{\pi}{6}$ or $x=\frac{5\pi}{6}$.

On the intervals $[0,\frac{\pi}{6})$ and $(\frac{5\pi}{6},\pi]$ we have $0\le\sin x<\frac12$, hence $\sin^2x<\frac14$; on the interval $(\frac{\pi}{6},\frac{5\pi}{6})$ we have $\sin^2x>\frac14$.

\textbf{Antiderivative:} by $\sin^2x=\frac{1-\cos2x}{2}$,
\[ \int\left(\sin^2x-\frac14\right)dx=\int\left(\frac14-\frac{\cos 2x}{2}\right)dx=\frac{x}{4}-\frac{\sin 2x}{4}=:F(x). \]
Values: $F(0)=0$, $F(\frac{\pi}{6})=\frac{\pi}{24}-\frac{\sqrt3}{8}$, $F(\frac{5\pi}{6})=\frac{5\pi}{24}+\frac{\sqrt3}{8}$, $F(\pi)=\frac{\pi}{4}$.

\textbf{The area} is $\int_0^\pi|f-g|\,dx$:
\begin{align*}
S_1&=\int_0^{\pi/6}\left(\tfrac14-\sin^2x\right)dx=-\big(F(\tfrac{\pi}{6})-F(0)\big)=\frac{\sqrt3}{8}-\frac{\pi}{24},\\
S_2&=\int_{\pi/6}^{5\pi/6}\left(\sin^2x-\tfrac14\right)dx=F(\tfrac{5\pi}{6})-F(\tfrac{\pi}{6})=\frac{\pi}{6}+\frac{\sqrt3}{4},\\
S_3&=\int_{5\pi/6}^{\pi}\left(\tfrac14-\sin^2x\right)dx=-\big(F(\pi)-F(\tfrac{5\pi}{6})\big)=\frac{\sqrt3}{8}-\frac{\pi}{24}.
\end{align*}
In total:
\[ S=S_1+S_2+S_3=\frac{\sqrt3}{2}+\frac{\pi}{12}\approx 1.128. \]
(If only the closed region between the two intersection points is meant, its area is $S_2=\frac{\pi}{6}+\frac{\sqrt3}{4}\approx 0.957$.)
\end{solution}

\subsection*{Question 6b}

\begin{question}
Compute the improper integral $\displaystyle\int_0^{\infty}e^{-2x}\sin x\,dx$ or show that it does not converge.
\end{question}

\begin{hint}
Compute the integral up to $R$ using integration by parts twice (the integral ``repeats itself''), and then let $R\to\infty$.
\end{hint}

\begin{solution}
\textbf{Convergence:} $|e^{-2x}\sin x|\le e^{-2x}$ and $\int_0^\infty e^{-2x}dx=\frac12$ converges, hence by the \textbf{comparison test} the integral converges (absolutely).

\textbf{Computation:} let $I(R)=\int_0^R e^{-2x}\sin x\,dx$, and find an antiderivative $J=\int e^{-2x}\sin x\,dx$ by integrating by parts twice ($u=e^{-2x}$):
\begin{align*}
J&=-e^{-2x}\cos x-\int 2e^{-2x}\cos x\,dx
 =-e^{-2x}\cos x-2\left(e^{-2x}\sin x+\int 2e^{-2x}\sin x\,dx\right)\\
 &=-e^{-2x}\cos x-2e^{-2x}\sin x-4J.
\end{align*}
Hence $5J=-e^{-2x}(\cos x+2\sin x)$, i.e.
\[ \int e^{-2x}\sin x\,dx=-\frac{e^{-2x}(\cos x+2\sin x)}{5}+C. \]
Therefore
\[ I(R)=\frac{1}{5}-\frac{e^{-2R}(\cos R+2\sin R)}{5}. \]
As $R\to\infty$: $|\cos R+2\sin R|\le 3$ and $e^{-2R}\to 0$, hence (bounded times tending to zero) the second term tends to $0$. Hence
\[ \int_0^\infty e^{-2x}\sin x\,dx=\frac15. \]
\end{solution}

\section*{2022/23 Semester A Moed B}

\subsection*{Question 1a}

\begin{question}
Compute the one-sided limits of the function $f(x)=2^{\frac{1}{x}}$ at the given point $x=0$. Determine whether it has a limit $\lim_{x\to 0}f(x)$.
\end{question}

\begin{solution}
Let $t=\frac1x$. As $x\to 0^+$ we have $t\to+\infty$, and as $x\to 0^-$ we have $t\to-\infty$. Since $2>1$, $\lim_{t\to+\infty}2^t=+\infty$ and $\lim_{t\to-\infty}2^t=0$. Hence (limit of a composition)
\[ \lim_{x\to0^+}2^{1/x}=+\infty,\qquad \lim_{x\to0^-}2^{1/x}=0. \]
The one-sided limits are different (and the right one is not even finite), and hence \textbf{the limit $\lim_{x\to0}2^{1/x}$ does not exist}.
\end{solution}

\subsection*{Question 1b}

\begin{question}
Classify the points of discontinuity of the function
\[ f(x)=\begin{cases}\dfrac{x^3-8}{x-2} & x\neq 2\\[2mm] 1 & x=2\end{cases} \]
\end{question}

\begin{solution}
For every $x\neq 2$ the function is a quotient of polynomials with a nonzero denominator, and hence it is continuous there. We check $x=2$. By the difference of cubes formula $x^3-8=(x-2)(x^2+2x+4)$, and hence for $x\neq2$:
\[ f(x)=x^2+2x+4\ \Longrightarrow\ \lim_{x\to2}f(x)=4+4+4=12. \]
The limit exists and is finite, but $f(2)=1\neq 12$. Hence $x=2$ is a \textbf{removable} discontinuity (of the first kind, removable): if we redefine $f(2)=12$, the function becomes continuous. This is the only point of discontinuity.
\end{solution}

\subsection*{Question 2a}

\begin{question}
Compute the equation of the tangent line to the circle $x^2+y^2=25$ at the point $(3,4)$.
\end{question}

\begin{hint}
Differentiate the equation implicitly.
\end{hint}

\begin{solution}
The point is on the circle: $9+16=25$. Near $(3,4)$ ($y>0$) the circle is the graph of a differentiable function $y(x)$. Differentiate the equation implicitly with respect to $x$:
\[ 2x+2y\,y'=0\ \Longrightarrow\ y'=-\frac{x}{y},\qquad y'(3)=-\frac34. \]
The equation of the tangent line: $y-4=-\frac34(x-3)$, i.e.
\[ y=-\frac34x+\frac{25}{4}\qquad\Longleftrightarrow\qquad 3x+4y=25. \]
(In agreement with the geometry: the tangent is perpendicular to the radius in the direction $(3,4)$.)
\end{solution}

\subsection*{Question 2b}

\begin{question}
Use a linear approximation to estimate the expression $\arcsin(0.54)$.
\end{question}

\begin{hint}
Use the point $0.5$, at which $\arcsin$ is known.
\end{hint}

\begin{solution}
Define $f(x)=\arcsin x$ and the base point $x_0=0.5$: $f(0.5)=\frac{\pi}{6}$, and
\[ f'(x)=\frac{1}{\sqrt{1-x^2}},\qquad f'(0.5)=\frac{1}{\sqrt{3/4}}=\frac{2}{\sqrt3}. \]
Linear approximation: $f(x)\approx f(x_0)+f'(x_0)(x-x_0)$, with $x-x_0=0.04$:
\[ \arcsin(0.54)\approx\frac{\pi}{6}+\frac{2}{\sqrt3}\cdot0.04=\frac{\pi}{6}+\frac{0.08}{\sqrt3}\approx 0.5236+0.0462=0.5698. \]
(The true value is $0.5704\ldots$.)
\end{solution}

\subsection*{Question 3a}

\begin{question}
Compute the absolute extremum points of the function $f(x)=\frac{x^2}{x-3}$ on the interval $[-2,2]$.
\end{question}

\begin{solution}
The point $x=3$ is not in the interval, hence $f$ is continuous on the closed interval $[-2,2]$, and by \textbf{Weierstrass's theorem} it attains a maximum and a minimum there. They are attained at interior critical points or at the endpoints.
\[ f'(x)=\frac{2x(x-3)-x^2}{(x-3)^2}=\frac{x^2-6x}{(x-3)^2}=\frac{x(x-6)}{(x-3)^2}. \]
$f'(x)=0\iff x=0$ or $x=6$; in the interval $(-2,2)$ only $x=0$. Compare the values:
\[ f(-2)=\frac{4}{-5}=-\frac45,\qquad f(0)=0,\qquad f(2)=\frac{4}{-1}=-4. \]
\textbf{Answer:} absolute maximum $0$ at the point $x=0$; absolute minimum $-4$ at the point $x=2$.
\end{solution}

\subsection*{Question 3b}

\begin{question}
For the equation $e^x=x+1$, prove that a solution exists or prove that no solution exists in the domain $x\neq0$.
\end{question}

\begin{hint}
Study the function $g(x)=e^x-x-1$: find its minimum point.
\end{hint}

\begin{solution}
We prove that there is \textbf{no} solution with $x\neq 0$. Define $g(x)=e^x-x-1$, differentiable on $\R$, with $g(0)=0$ and
\[ g'(x)=e^x-1. \]
For $x<0$: $g'(x)<0$, hence $g$ is strictly decreasing on $(-\infty,0]$; for $x>0$: $g'(x)>0$, hence $g$ is strictly increasing on $[0,\infty)$. (Strict monotonicity follows from Lagrange's theorem.) Hence:
\begin{itemize}
  \item for every $x<0$: $g(x)>g(0)=0$;
  \item for every $x>0$: $g(x)>g(0)=0$.
\end{itemize}
That is, $e^x>x+1$ for every $x\neq0$, and the equation $e^x=x+1$ has no solution in the domain $x\ne 0$ (its only solution is $x=0$). $\blacksquare$
\end{solution}

\subsection*{Question 4a}

\begin{question}
Compute the limit $\displaystyle\lim_{x\to0}\frac{x-\sin(x)}{x^3}$.
\end{question}

\begin{solution}
An expression of the form $\frac00$. We apply \textbf{L'Hôpital's rule} (the conditions hold: differentiability near $0$, the derivative of the denominator is different from $0$ near $0$):
\[ \lim_{x\to0}\frac{x-\sin x}{x^3}=\lim_{x\to0}\frac{1-\cos x}{3x^2}=\lim_{x\to0}\frac{\sin x}{6x}=\frac16, \]
where in the second step we applied L'Hôpital's rule again (again $\frac00$), and in the last we used $\lim\frac{\sin x}{x}=1$.
(Alternatively: $\sin x=x-\frac{x^3}{6}+o(x^3)$, and hence $\frac{x-\sin x}{x^3}=\frac16+o(1)$.) \textbf{Answer:} $\frac16$.
\end{solution}

\subsection*{Question 4b}

\begin{question}
Use a second-order Taylor polynomial to estimate the expression $\sqrt[4]{82}$, and estimate the error.
\end{question}

\begin{hint}
Expand $f(x)=x^{1/4}$ around $x_0=81$ and use the Lagrange remainder.
\end{hint}

\begin{solution}
$f(x)=x^{1/4}$, $x_0=81$, $f(81)=3$.
\[ f'(x)=\tfrac14x^{-3/4},\quad f''(x)=-\tfrac{3}{16}x^{-7/4},\quad f'''(x)=\tfrac{21}{64}x^{-11/4}. \]
\[ f'(81)=\frac{1}{4\cdot27}=\frac{1}{108},\qquad f''(81)=-\frac{3}{16\cdot 3^7}=-\frac{1}{16\cdot729}=-\frac{1}{11664}. \]
The Taylor polynomial of order 2 around $81$, at the point $x=82$ ($x-x_0=1$):
\[ \sqrt[4]{82}\approx P_2(82)=3+\frac{1}{108}-\frac{1}{2\cdot11664}=3+\frac{1}{108}-\frac{1}{23328}=\frac{70199}{23328}\approx 3.0092164. \]
\textbf{Error estimate} (Lagrange remainder): there exists $c\in(81,82)$ such that
\[ R_2=\frac{f'''(c)}{3!}(82-81)^3=\frac{21}{64\cdot 6}\,c^{-11/4}. \]
$c^{-11/4}$ is decreasing, hence $c^{-11/4}<81^{-11/4}=3^{-11}$, and therefore
\[ 0<R_2<\frac{21}{384\cdot 3^{11}}=\frac{21}{384\cdot177147}\approx 3.1\cdot10^{-7}. \]
That is, $\sqrt[4]{82}\approx3.0092164$ with an error less than $3.1\cdot 10^{-7}$ (and the approximation is slightly smaller than the true value, $3.0092167\ldots$).
\end{solution}

\subsection*{Question 5a}

\begin{question}
Compute the integral $\displaystyle\int\frac{x}{\sqrt{1+x^2}}\,dx$.
\end{question}

\begin{solution}
Substitute $t=1+x^2$, $dt=2x\,dx$:
\[ \int\frac{x}{\sqrt{1+x^2}}\,dx=\frac12\int t^{-1/2}\,dt=t^{1/2}+C=\sqrt{1+x^2}+C. \]
Check: $\left(\sqrt{1+x^2}\right)'=\frac{2x}{2\sqrt{1+x^2}}=\frac{x}{\sqrt{1+x^2}}$.
\end{solution}

\subsection*{Question 5b}

\begin{question}
Solve the equation $e^x\,dx-y\,dy=0$ with the initial condition $y(0)=1$.
\end{question}

\begin{hint}
This is a separable equation: move one term to the other side and integrate each side separately.
\end{hint}

\begin{solution}
The equation is separable: $y\,dy=e^x\,dx$ (i.e. $y\,y'=e^x$). Integrate both sides:
\[ \int y\,dy=\int e^x\,dx\ \Longrightarrow\ \frac{y^2}{2}=e^x+C. \]
From the initial condition $y(0)=1$: $\frac12=1+C$, hence $C=-\frac12$, and
\[ y^2=2e^x-1. \]
Since $y(0)=1>0$ and a continuous solution cannot pass to negative values without passing through $0$, we choose the positive branch:
\[ y(x)=\sqrt{2e^x-1}, \]
which is defined when $2e^x-1>0$, i.e. $x>-\ln 2$.
Check: $y'=\frac{e^x}{\sqrt{2e^x-1}}$, hence $y\,y'=e^x$ ✓, and $y(0)=\sqrt{1}=1$ ✓.
\end{solution}

\subsection*{Question 6a}

\begin{question}
Compute the finite area bounded between the graph of the function $f(x)=x^3+x^2-2x$ and the $x$-axis.
\end{question}

\begin{hint}
Factor the polynomial to find the intersection points with the $x$-axis and the sign on each interval.
\end{hint}

\begin{solution}
Factoring: $f(x)=x(x^2+x-2)=x(x+2)(x-1)$. The intersection points with the $x$-axis: $x=-2,0,1$. Sign: on $(-2,0)$ $f>0$ (for example $f(-1)=2$), and on $(0,1)$ $f<0$ (for example $f(\frac12)=-\frac58$). The finite region bounded between the graph and the axis lies over $[-2,1]$.

Antiderivative: $F(x)=\frac{x^4}{4}+\frac{x^3}{3}-x^2$. Then $F(-2)=4-\frac83-4=-\frac83$, $F(0)=0$, $F(1)=\frac14+\frac13-1=-\frac{5}{12}$.
\[ S=\int_{-2}^0 f\,dx-\int_0^1 f\,dx=\big(F(0)-F(-2)\big)-\big(F(1)-F(0)\big)=\frac83+\frac{5}{12}=\frac{37}{12}. \]
\textbf{Answer:} $S=\frac{37}{12}$.
\end{solution}

\subsection*{Question 6b}

\begin{question}
Determine the convergence of the integral $\displaystyle\int_1^{\infty}\frac{\sqrt{x^5}\,dx}{2x^2+3x+40}$.
\end{question}

\begin{hint}
Check how the function behaves as $x\to\infty$ and compare with a suitable $x^{p}$ (the limit comparison test).
\end{hint}

\begin{solution}
The function $g(x)=\frac{x^{5/2}}{2x^2+3x+40}$ is continuous and positive on $[1,\infty)$. Compare with $h(x)=x^{1/2}$:
\[ \lim_{x\to\infty}\frac{g(x)}{h(x)}=\lim_{x\to\infty}\frac{x^2}{2x^2+3x+40}=\frac12\in(0,\infty). \]
By the \textbf{limit comparison test}, $\int_1^\infty g$ and $\int_1^\infty x^{1/2}dx$ converge or diverge together. $\int_1^\infty x^{p}dx$ converges only for $p<-1$, and here $p=\frac12$, hence $\int_1^\infty x^{1/2}dx=\infty$.
(One can also note that even $g(x)\to\infty$, so the integral certainly diverges.)

\textbf{Answer:} the integral \textbf{diverges} (equals $+\infty$).
\end{solution}

\section*{2022/23 Semester A Moed C}

\subsection*{Question 1a}

\begin{question}
Check whether an inverse function exists for the function $f(x)=|x|+2x$ on its domain. If it exists, find it.
\end{question}

\begin{hint}
Open the absolute value: write $f$ separately for $x\ge 0$ and for $x<0$, and check that it is one-to-one and onto.
\end{hint}

\begin{solution}
The function is defined on all of $\R$. Open the absolute value:
\[ f(x)=\begin{cases} 3x & x\ge 0\\ x & x<0. \end{cases} \]
\textbf{One-to-one:} on each of the two domains $f$ is strictly increasing (slope $3$ and slope $1$), and in addition $f(x)\ge 0$ for $x\ge0$ and $f(x)<0$ for $x<0$, so values from the two domains cannot coincide. Hence $f$ is strictly increasing on all of $\R$, and in particular one-to-one.

\textbf{Onto:} the image of $[0,\infty)$ under $3x$ is $[0,\infty)$, and the image of $(-\infty,0)$ under $x$ is $(-\infty,0)$. Hence the image is all of $\R$.

It follows that $f:\R\to\R$ is invertible. We find the inverse: if $y\ge 0$ then $y=3x$ with $x\ge 0$, i.e. $x=y/3$; if $y<0$ then $y=x$. Hence
\[ \boxed{f^{-1}(y)=\begin{cases} \dfrac{y}{3} & y\ge 0\\[2mm] y & y<0. \end{cases}} \]
(Check: $f^{-1}(f(x))=x$ on both domains.)
\end{solution}

\subsection*{Question 1b}

\begin{question}
Classify the points of discontinuity of the function
\[ f(x)=\begin{cases} \dfrac{\sin(2x)}{x} & x\neq 0\\[2mm] 1 & x=0 \end{cases} \]
\end{question}

\begin{solution}
For every $x\neq 0$ the function is a quotient of continuous functions ($\sin(2x)$ and $x$) whose denominator does not vanish, hence it is continuous. We check $x=0$: by the fundamental limit $\lim_{t\to0}\frac{\sin t}{t}=1$,
\[ \lim_{x\to0}\frac{\sin(2x)}{x}=\lim_{x\to0}2\cdot\frac{\sin(2x)}{2x}=2. \]
The limit exists and is finite, but $f(0)=1\neq 2$. Hence $x=0$ is a \textbf{removable discontinuity} (it can be fixed by redefining $f(0)=2$), and it is the only point of discontinuity.
\end{solution}

\subsection*{Question 2a}

\begin{question}
Use a linear approximation to estimate the expression $\log_2 10$.
\end{question}

\begin{hint}
Choose a point close to $10$ at which $\log_2$ is known, for example $x_0=8$.
\end{hint}

\begin{solution}
Define $f(x)=\log_2 x=\frac{\ln x}{\ln 2}$; then $f'(x)=\frac{1}{x\ln 2}$. The linear approximation around $x_0$ is
\[ f(x)\approx f(x_0)+f'(x_0)(x-x_0). \]
Take $x_0=8$ (where $\log_2 8=3$) and $x=10$:
\[ \log_2 10\approx 3+\frac{1}{8\ln 2}\cdot 2=3+\frac{1}{4\ln 2}\approx 3+0.3607=\boxed{3.3607}. \]
(The exact value is $\log_2 10\approx 3.3219$; the approximation is slightly larger than the true value, as expected, since $\log_2$ is concave and therefore the tangent line lies above the graph.)
\end{solution}

\subsection*{Question 2b}

\begin{question}
Compute the derivative at the point $x=0$ of the function $y(x)$ given implicitly by the equation
\[ y+x^3y^3-3xy=x^5 \]
\end{question}

\begin{hint}
First find $y(0)$ from the equation, then differentiate both sides with respect to $x$ (chain rule and product rule).
\end{hint}

\begin{solution}
\textbf{Value of the function:} substitute $x=0$ in the equation: $y+0-0=0$, hence $y(0)=0$.

\textbf{Implicit differentiation:} differentiate both sides with respect to $x$ (where $y=y(x)$):
\[ y'+3x^2y^3+3x^3y^2y'-3y-3xy'=5x^4. \]
Substitute $x=0,\ y=0$:
\[ y'(0)+0+0-3\cdot 0-0=0 \quad\Longrightarrow\quad \boxed{y'(0)=0}. \]
(In general $y'=\dfrac{5x^4-3x^2y^3+3y}{1+3x^3y^2-3x}$, and the denominator equals $1\neq 0$ at the point $(0,0)$.)
\end{solution}

\subsection*{Question 3a}

\begin{question}
Find the points on the curve $y=2x^3-3x^2-12x-5$ at which the tangent line is parallel to the $x$-axis.
\end{question}

\begin{solution}
The tangent line is parallel to the $x$-axis exactly when its slope is $0$, i.e. $y'=0$:
\[ y'=6x^2-6x-12=6(x^2-x-2)=6(x-2)(x+1)=0\iff x=2\ \text{or}\ x=-1. \]
We compute the values of $y$: $y(2)=16-12-24-5=-25$ and $y(-1)=-2-3+12-5=2$. The points are
\[ \boxed{(-1,\,2)\quad\text{and}\quad(2,\,-25)}. \]
\end{solution}

\subsection*{Question 3b}

\begin{question}
Find the intervals of convexity and the inflection points of the function $f(x)=e^{-x^2+2x}$
\end{question}

\begin{solution}
$f$ is defined and twice differentiable on all of $\R$. By the chain rule
\[ f'(x)=(2-2x)e^{-x^2+2x}, \]
\[ f''(x)=-2e^{-x^2+2x}+(2-2x)^2e^{-x^2+2x}=(4x^2-8x+2)e^{-x^2+2x}=2(2x^2-4x+1)e^{-x^2+2x}. \]
Since $e^{-x^2+2x}>0$, the sign of $f''$ is the sign of $2x^2-4x+1$, whose zeros are
\[ x_{1,2}=\frac{4\pm\sqrt{16-8}}{4}=1\pm\frac{\sqrt2}{2}. \]
This is a parabola with a positive leading coefficient, hence:
\begin{itemize}
  \item $f''>0$ ($f$ is \textbf{convex}) for $x<1-\frac{\sqrt2}{2}$ and for $x>1+\frac{\sqrt2}{2}$;
  \item $f''<0$ ($f$ is \textbf{concave}) for $1-\frac{\sqrt2}{2}<x<1+\frac{\sqrt2}{2}$.
\end{itemize}
At both points $x=1\pm\frac{\sqrt2}{2}$ the second derivative changes sign and the function is continuous, hence these are inflection points. Since $-x^2+2x=1-(x-1)^2=1-\frac12=\frac12$ at these points, the inflection points are
\[ \boxed{\left(1-\tfrac{\sqrt2}{2},\ \sqrt e\right),\quad \left(1+\tfrac{\sqrt2}{2},\ \sqrt e\right)}. \]
\end{solution}

\subsection*{Question 4a}

\begin{question}
Compute the limit $\displaystyle\lim_{x\to0}\frac{e^{2x}-e^x}{x}$
\end{question}

\begin{solution}
This is a limit of the form $\frac00$, and the numerator and denominator are differentiable, so by L'Hôpital's rule
\[ \lim_{x\to0}\frac{e^{2x}-e^x}{x}=\lim_{x\to0}\frac{2e^{2x}-e^x}{1}=2-1=\boxed{1}. \]
(Alternatively: $\frac{e^{2x}-e^x}{x}=2\cdot\frac{e^{2x}-1}{2x}-\frac{e^x-1}{x}\to 2-1=1$ by the fundamental limit $\lim_{t\to0}\frac{e^t-1}{t}=1$.)
\end{solution}

\subsection*{Question 4b}

\begin{question}
Use the second-order Maclaurin polynomial to estimate the expression $e^{0.1}$ and estimate the error
\end{question}

\begin{hint}
Use the Lagrange remainder $R_2(x)=\frac{f'''(c)}{3!}x^3$ with $c$ between $0$ and $0.1$, and bound $e^c$.
\end{hint}

\begin{solution}
For $f(x)=e^x$ all derivatives are $e^x$ and equal $1$ at $0$, so the second-order Maclaurin polynomial is
\[ P_2(x)=1+x+\frac{x^2}{2}, \qquad e^{0.1}\approx P_2(0.1)=1+0.1+0.005=\boxed{1.105}. \]
\textbf{Error estimate:} by Taylor's theorem with the Lagrange remainder, there exists $c\in(0,0.1)$ such that
\[ R_2(0.1)=e^{0.1}-P_2(0.1)=\frac{e^{c}}{3!}(0.1)^3. \]
Since $e^x$ is increasing, $e^c<e^{0.1}<e^1<3$, and therefore
\[ 0<R_2(0.1)<\frac{3}{6}\cdot 0.001=0.0005. \]
(A tighter bound: $e^{0.1}<1.2$ gives $R_2<0.0002$.) The error is positive, i.e. the approximation is an underestimate. The true value is $e^{0.1}\approx 1.10517$, and the actual error is about $0.00017$.
\end{solution}

\subsection*{Question 5a}

\begin{question}
Compute the integral $\displaystyle\int x\cdot\sqrt{1-x^2}\,dx$
\end{question}

\begin{solution}
Substitute $t=1-x^2$, $dt=-2x\,dx$, i.e. $x\,dx=-\frac12dt$:
\[ \int x\sqrt{1-x^2}\,dx=-\frac12\int t^{1/2}\,dt=-\frac12\cdot\frac{2}{3}t^{3/2}+C=\boxed{-\frac13(1-x^2)^{3/2}+C}. \]
Check by differentiation: $\left(-\frac13(1-x^2)^{3/2}\right)'=-\frac13\cdot\frac32(1-x^2)^{1/2}\cdot(-2x)=x\sqrt{1-x^2}$.
\end{solution}

\subsection*{Question 5b}

\begin{question}
Compute the integral $\displaystyle\int_2^4\frac{dx}{x+4}$
\end{question}

\begin{solution}
The function $\frac{1}{x+4}$ is continuous on $[2,4]$, and an antiderivative of it is $\ln|x+4|$. By the Fundamental Theorem (Newton--Leibniz):
\[ \int_2^4\frac{dx}{x+4}=\ln(x+4)\Big|_2^4=\ln8-\ln6=\boxed{\ln\frac43}\approx 0.2877. \]
\end{solution}

\subsection*{Question 6a}

\begin{question}
Determine whether the integral $\displaystyle\int_0^\infty e^{-x^2}dx$ converges
\end{question}

\begin{hint}
The integral cannot be computed in terms of elementary functions. Split at $x=1$ and compare $e^{-x^2}$ with a simpler function for $x\ge1$.
\end{hint}

\begin{solution}
Split: $\int_0^\infty e^{-x^2}dx=\int_0^1 e^{-x^2}dx+\int_1^\infty e^{-x^2}dx$.
The first integral is an ordinary definite integral of a continuous function on a closed interval, hence it exists and is finite.

For $x\ge 1$ we have $x^2\ge x$, and therefore $0<e^{-x^2}\le e^{-x}$. The integral
\[ \int_1^\infty e^{-x}dx=\lim_{R\to\infty}\left(e^{-1}-e^{-R}\right)=e^{-1} \]
converges. By the \textbf{comparison test} for improper integrals of non-negative functions, $\int_1^\infty e^{-x^2}dx$ converges as well. Hence $\boxed{\int_0^\infty e^{-x^2}dx\ \text{converges}}$
(and its value, as is well known, is $\frac{\sqrt\pi}{2}$, but this is not required).
\end{solution}

\subsection*{Question 6b}

\begin{question}
Compute the improper integral $\displaystyle\int_1^\infty\frac{dx}{x^3}$ or show that it does not converge.
\end{question}

\begin{solution}
By definition:
\[ \int_1^\infty\frac{dx}{x^3}=\lim_{R\to\infty}\int_1^R x^{-3}dx=\lim_{R\to\infty}\left[-\frac{1}{2x^2}\right]_1^R=\lim_{R\to\infty}\left(\frac12-\frac{1}{2R^2}\right)=\boxed{\frac12}. \]
The integral converges (in accordance with the fact that $\int_1^\infty\frac{dx}{x^p}$ converges for $p>1$).
\end{solution}

\section*{2022/23 Semester B Moed A (Dr. Peter Samovol)}

\subsection*{Question 1}

\begin{question}
(Mandatory question)
\begin{enumerate}
  \item (22 pts) Fully analyze the function: $\displaystyle f(x)=\frac{e^{2x+1}}{(x+1)^2}$ (domain, intersections with the axes, sign of the function, asymptotes, intervals of increase and decrease, extremum points, intervals of convexity and inflection points) and draw a sketch of the graph of the function.
  \item (4 pts) Sketch the graph of the function $f(-|x|)$ ?
\end{enumerate}
\end{question}

\begin{hint}
When differentiating, factor out a common factor: $f'(x)=\frac{e^{2x+1}\left(2(x+1)-2\right)}{(x+1)^3}$.
\end{hint}

\begin{hint}
In part (b), note that $f(-|x|)$ is an even function, and for $x\ge0$ it equals $f(-x)$, i.e. the reflection of the part of the graph of $f$ to the left of the $y$-axis.
\end{hint}

\begin{solution}
\begin{enumerate}
  \item \textbf{Domain:} $x\neq -1$, i.e. $D=(-\infty,-1)\cup(-1,\infty)$.

  \textbf{Intersections with the axes and sign:} $e^{2x+1}>0$ and $(x+1)^2>0$ in the domain, hence $f(x)>0$ for every $x\in D$; there is no intersection with the $x$-axis. Intersection with the $y$-axis: $f(0)=e$, the point $(0,e)$.

  \textbf{Asymptotes:}
  \begin{itemize}
    \item Vertical: $\lim_{x\to-1^\pm}f(x)=\frac{e^{-1}}{0^+}=+\infty$, hence $x=-1$ is a vertical asymptote (from both sides the function tends to $+\infty$).
    \item At $-\infty$: $e^{2x+1}\to 0$ and $(x+1)^2\to\infty$, so $f(x)\to0$: $y=0$ is a horizontal asymptote on the left.
    \item At $+\infty$: by L'Hôpital's rule (twice) $\lim_{x\to\infty}\frac{e^{2x+1}}{(x+1)^2}=\lim\frac{4e^{2x+1}}{2}=\infty$, and similarly $\frac{f(x)}{x}\to\infty$, so there is no horizontal or oblique asymptote on the right.
  \end{itemize}

  \textbf{Monotonicity and extrema:} by the quotient rule,
  \[ f'(x)=\frac{2e^{2x+1}(x+1)^2-e^{2x+1}\cdot2(x+1)}{(x+1)^4}=\frac{2e^{2x+1}\,x}{(x+1)^3}. \]
  $f'(x)=0\iff x=0$. The sign of $f'$ is the sign of $\frac{x}{(x+1)^3}$:
  \begin{itemize}
    \item $x<-1$: negative numerator, negative denominator $\Rightarrow f'>0$, $f$ is \textbf{increasing}.
    \item $-1<x<0$: negative numerator, positive denominator $\Rightarrow f'<0$, $f$ is \textbf{decreasing}.
    \item $x>0$: $f'>0$, $f$ is \textbf{increasing}.
  \end{itemize}
  Hence at $x=0$ there is a \textbf{local minimum} $(0,e)$. There is no local maximum (at $x=-1$ the function is not defined).

  \textbf{Convexity:} differentiating again gives
  \[ f''(x)=\frac{2e^{2x+1}(2x^2+1)}{(x+1)^4}. \]
  (Differentiate $f'=2e^{2x+1}\cdot x(x+1)^{-3}$: $f''=2e^{2x+1}\left[2x(x+1)^{-3}+(x+1)^{-3}-3x(x+1)^{-4}\right]=\frac{2e^{2x+1}\left[(2x+1)(x+1)-3x\right]}{(x+1)^4}$ and $(2x+1)(x+1)-3x=2x^2+1$.)
  Since $2x^2+1>0$, we have $f''>0$ on the whole domain: $f$ is \textbf{convex} on $(-\infty,-1)$ and on $(-1,\infty)$, and there are \textbf{no inflection points}.

  \textbf{Sketch (described in words):} on the left the graph starts just above the asymptote $y=0$ (close to $0$ at $-\infty$), increases in a convex manner and tends to $+\infty$ as $x\to-1^-$. To the right of the asymptote $x=-1$ the graph decreases from $+\infty$ to the minimum $(0,e)$ and then increases rapidly (exponentially) to $+\infty$. The whole graph lies above the $x$-axis.

  \item Denote $g(x)=f(-|x|)$. We have $g(-x)=g(x)$, i.e. $g$ is \textbf{even} -- the graph is symmetric with respect to the $y$-axis. For $x\ge0$: $g(x)=f(-x)$, i.e. the graph of $g$ for $x\ge0$ is the reflection with respect to the $y$-axis of the graph of $f$ for $x\le 0$, and the graph for $x<0$ is simply the graph of $f$ for $x<0$.

  In particular: $g$ is defined for $x\neq\pm1$, $g(0)=e$, and it has vertical asymptotes $x=\pm1$ (where $g\to+\infty$) and the horizontal asymptote $y=0$ at $\pm\infty$. On $(-1,0)$, $g=f$ is decreasing, and on $(0,1)$, $g$ is increasing; at $x=0$ there is a local minimum $(0,e)$ (and since $f'(0)=0$ there is no ``corner'' there: the one-sided derivatives of $g$ are both $0$). On $(-\infty,-1)$, $g$ increases from $0$ to $+\infty$, and on $(1,\infty)$ it decreases from $+\infty$ to $0$. $g$ is convex on each of the intervals.
\end{enumerate}
\end{solution}

\subsection*{Question 2a}

\begin{question}
(Mandatory question)
Solve the following integral:
$\displaystyle\int\frac{(x^2+1)\cdot dx}{x^3-x}$
\end{question}

\begin{hint}
Decompose into partial fractions: $x^3-x=x(x-1)(x+1)$.
\end{hint}

\begin{solution}
$x^3-x=x(x-1)(x+1)$, and the degree of the numerator is less than the degree of the denominator. Decompose:
\[ \frac{x^2+1}{x(x-1)(x+1)}=\frac{A}{x}+\frac{B}{x-1}+\frac{C}{x+1}. \]
By the cover-up method: $A=\frac{0+1}{(0-1)(0+1)}=-1$, $B=\frac{1+1}{1\cdot 2}=1$, $C=\frac{1+1}{(-1)(-2)}=1$. Hence
\[ \int\frac{x^2+1}{x^3-x}dx=-\ln|x|+\ln|x-1|+\ln|x+1|+C=\boxed{\ln\left|\frac{x^2-1}{x}\right|+C}. \]
\end{solution}

\subsection*{Question 2b}

\begin{question}
(Mandatory question)
Solve the following integral:
$\displaystyle\int\sin^2 4x\cos^2 x\,dx$
\end{question}

\begin{hint}
Use the power-reduction formulas $\sin^2\alpha=\frac{1-\cos2\alpha}{2}$, $\cos^2\alpha=\frac{1+\cos2\alpha}{2}$.
\end{hint}

\begin{solution}
By the power-reduction formulas:
\[ \sin^2 4x\cos^2x=\frac{1-\cos8x}{2}\cdot\frac{1+\cos2x}{2}=\frac14\left(1+\cos2x-\cos8x-\cos8x\cos2x\right). \]
By the product-to-sum formula $\cos8x\cos2x=\frac12(\cos10x+\cos6x)$, and therefore
\[ \sin^2 4x\cos^2x=\frac14+\frac14\cos2x-\frac14\cos8x-\frac18\cos10x-\frac18\cos6x. \]
Integrate term by term:
\[ \boxed{\int\sin^2 4x\cos^2x\,dx=\frac{x}{4}+\frac{\sin2x}{8}-\frac{\sin8x}{32}-\frac{\sin6x}{48}-\frac{\sin10x}{80}+C}. \]
\end{solution}

\subsection*{Question 2c}

\begin{question}
(Mandatory question)
Solve the following integral:
$\displaystyle\int_{-3}^{-1}\operatorname{arctg}(x+2)\,dx$
\end{question}

\begin{hint}
Substitute $t=x+2$ and pay attention to evenness/oddness.
\end{hint}

\begin{solution}
Substitute $t=x+2$, $dt=dx$; the limits: $x=-3\mapsto t=-1$, $x=-1\mapsto t=1$:
\[ \int_{-3}^{-1}\operatorname{arctg}(x+2)\,dx=\int_{-1}^{1}\arctan t\,dt. \]
$\arctan$ is an odd continuous function, and the interval is symmetric about $0$, hence the integral $\boxed{=0}$.

(Directly: by integration by parts $\int\arctan t\,dt=t\arctan t-\frac12\ln(1+t^2)+C$, and between $-1$ and $1$: $\left(\frac\pi4-\frac12\ln2\right)-\left(\frac\pi4-\frac12\ln2\right)=0$.)
\end{solution}

\subsection*{Question 2d}

\begin{question}
(Mandatory question)
Solve the following integral:
$\displaystyle\int_2^5\frac{2x-1}{\sqrt{x^2-x}}\,dx$
\end{question}

\begin{hint}
The numerator is exactly the derivative of the expression under the square root.
\end{hint}

\begin{solution}
On $[2,5]$ we have $x^2-x>0$, so the integrand is continuous. Substitute $u=x^2-x$, $du=(2x-1)dx$; the limits are $u(2)=2$, $u(5)=20$:
\[ \int_2^5\frac{2x-1}{\sqrt{x^2-x}}dx=\int_2^{20}u^{-1/2}du=2\sqrt u\Big|_2^{20}=2\sqrt{20}-2\sqrt2=\boxed{4\sqrt5-2\sqrt2}\approx 6.116. \]
\end{solution}

\subsection*{Question 3a}

\begin{question}
Compute the area of the region bounded by the lines $y=\min\left(e^x,3\right),\ x=4,\ x=0,\ y=0$
\end{question}

\begin{hint}
Find where $e^x=3$ and split the integral at that point.
\end{hint}

\begin{solution}
$e^x\le 3\iff x\le\ln3$ (since $e^x$ is increasing). Hence on $[0,4]$:
\[ \min(e^x,3)=\begin{cases} e^x & 0\le x\le\ln 3\\ 3 & \ln3\le x\le 4\end{cases} \]
and note that $\ln3\approx1.0986\in[0,4]$. The function is positive, hence the area is
\[ S=\int_0^{\ln3}e^x\,dx+\int_{\ln3}^{4}3\,dx=\left(e^{\ln3}-e^0\right)+3(4-\ln3)=2+12-3\ln3=\boxed{14-3\ln3}\approx 10.70. \]
\end{solution}

\subsection*{Question 3b}

\begin{question}
Find the values of the parameter $b$ for which the graph of the function $\displaystyle f(x)=\frac{1}{x^2+b}$ has no inflection points
\end{question}

\begin{hint}
Compute $f''$ and split into the cases $b>0$, $b=0$, $b<0$. Remember that an inflection point must lie in the domain.
\end{hint}

\begin{solution}
Differentiate: $f'(x)=-\frac{2x}{(x^2+b)^2}$, and
\[ f''(x)=\frac{-2(x^2+b)^2+2x\cdot2(x^2+b)\cdot2x}{(x^2+b)^4}=\frac{6x^2-2b}{(x^2+b)^3}. \]
\begin{itemize}
  \item $b>0$: $f$ is defined on all of $\R$, the denominator is positive, and the numerator $6x^2-2b$ changes sign at $x=\pm\sqrt{b/3}$. Hence there are two inflection points.
  \item $b=0$: $f=\frac1{x^2}$, $x\neq0$, and $f''=\frac{6}{x^4}>0$: there are no inflection points.
  \item $b<0$: $f$ is defined for $x\neq\pm\sqrt{-b}$. The numerator $6x^2-2b=6x^2+2|b|>0$ always, hence $f''$ does not vanish. The sign of $f''$ changes only when passing through $x=\pm\sqrt{-b}$, which are not in the domain (vertical asymptotes), and hence are not inflection points. There are no inflection points.
\end{itemize}
\textbf{Answer:} the graph has no inflection points exactly for $\boxed{b\le 0}$.
\end{solution}

\subsection*{Question 4a}

\begin{question}
Compute the following limits
\[ (1)\ \lim_{x\to0}\frac{\ln(1+2x)}{\sin2x-\sin6x}\qquad\qquad (2)\ \lim_{x\to0}\left(\sqrt{1+2x}\right)^{\frac{1}{(-x)}} \]
\end{question}

\begin{hint}
In limit (2), write the expression in the form $e^{g(x)}$ and compute the limit of the exponent.
\end{hint}

\begin{solution}
\textbf{(1)} A limit of the form $\frac00$; by L'Hôpital's rule:
\[ \lim_{x\to0}\frac{\ln(1+2x)}{\sin2x-\sin6x}=\lim_{x\to0}\frac{\frac{2}{1+2x}}{2\cos2x-6\cos6x}=\frac{2}{2-6}=\boxed{-\frac12}. \]
(Alternatively, using $\ln(1+2x)\sim2x$, $\sin2x\sim2x$, $\sin6x\sim6x$: $\frac{2x}{2x-6x}=-\frac12$.)

\textbf{(2)} For $x$ close to $0$ ($x>-\frac12$, $x\neq0$) the base is positive, and
\[ \left(\sqrt{1+2x}\right)^{-1/x}=e^{-\frac{\ln(1+2x)}{2x}}. \]
By the fundamental limit (or L'Hôpital's rule) $\lim_{x\to0}\frac{\ln(1+2x)}{2x}=1$. By continuity of the exponential function,
\[ \lim_{x\to0}\left(\sqrt{1+2x}\right)^{\frac{1}{-x}}=e^{-1}=\boxed{\frac1e}. \]
\end{solution}

\subsection*{Question 4b}

\begin{question}
For which values of $a$ is the function $f(x)=\min\left(x^2+x,a\right)$ differentiable for every real $x$? Justify your answer.
\end{question}

\begin{hint}
Compare $a$ with the minimum value of $x^2+x$. If the line $y=a$ intersects the parabola, what happens to the one-sided derivatives at the intersection point?
\end{hint}

\begin{solution}
The parabola $p(x)=x^2+x=\left(x+\frac12\right)^2-\frac14$ attains its minimum value $-\frac14$ at $x=-\frac12$.
\begin{itemize}
  \item If $a\le-\frac14$: $p(x)\ge-\frac14\ge a$ for every $x$, hence $f(x)\equiv a$ is constant, and differentiable on all of $\R$.
  \item If $a>-\frac14$: the equation $x^2+x=a$ has two roots $x_1<x_2$, $x_{1,2}=\frac{-1\mp\sqrt{1+4a}}{2}$. Then $f(x)=x^2+x$ on $[x_1,x_2]$ and $f(x)=a$ outside it. At the point $x_1$ (where $f$ is continuous) the left derivative is $0$ (a constant function), and the right derivative is $p'(x_1)=2x_1+1=-\sqrt{1+4a}\neq0$. The one-sided derivatives are different, hence $f$ is not differentiable at $x_1$ (and similarly at $x_2$, where the left derivative is $\sqrt{1+4a}$ and the right derivative is $0$).
\end{itemize}
\textbf{Answer:} $f$ is differentiable for every real $x$ if and only if $\boxed{a\le-\frac14}$.
\end{solution}

\subsection*{Question 5a}

\begin{question}
Find the equation of the tangent line to the graph of the function-
\[ \begin{cases} x=\ln\left(3t^2+e\cdot t-3\right)\\ y=t^3-3t\end{cases} \]
at the point $M_0(1,-2)$.
\end{question}

\begin{hint}
First find the value of the parameter $t_0$ corresponding to the point $M_0$ (both equations must hold simultaneously), and then use $\frac{dy}{dx}=\frac{y'(t)}{x'(t)}$.
\end{hint}

\begin{solution}
\textbf{Finding the parameter:} from $y=-2$: $t^3-3t+2=0\iff(t-1)^2(t+2)=0$, i.e. $t=1$ or $t=-2$.
From $x=1$: $3t^2+et-3=e$. For $t=1$: $3+e-3=e$ \checkmark. For $t=-2$: $12-2e-3=9-2e\neq e$ (since $e\neq3$). Hence $t_0=1$.

\textbf{The derivative:} by the formula for the derivative of a parametrically given function,
\[ x'(t)=\frac{6t+e}{3t^2+et-3},\qquad y'(t)=3t^2-3. \]
At $t_0=1$: $x'(1)=\frac{6+e}{e}\neq 0$ and $y'(1)=0$, hence
\[ \frac{dy}{dx}\Big|_{M_0}=\frac{y'(1)}{x'(1)}=0. \]
\textbf{Equation of the tangent line:} $y-(-2)=0\cdot(x-1)$, i.e. $\boxed{y=-2}$ (a horizontal tangent).
\end{solution}

\subsection*{Question 5b}

\begin{question}
Give the definition of an antiderivative of a function $f(x)$. Include an example.
\end{question}

\begin{solution}
\textbf{Definition:} let $f$ be defined on an interval $I$. A function $F$ is called an \textbf{antiderivative} of $f$ on $I$ if $F$ is differentiable on $I$ and $F'(x)=f(x)$ for every $x\in I$.
(If $F$ is an antiderivative of $f$ on an interval, then all antiderivatives are of the form $F+C$, $C$ a constant.)

\textbf{Example:} $F(x)=\frac{x^3}{3}$ is an antiderivative of $f(x)=x^2$ on $\R$, since $\left(\frac{x^3}{3}\right)'=x^2$; $\frac{x^3}{3}+5$ is also an antiderivative of it. Another example: $-\cos x$ is an antiderivative of $\sin x$.
\end{solution}

\subsection*{Question 6a}

\begin{question}
Write the Maclaurin polynomial of order 2 of the function $y=f(x)$ defined implicitly by: $x^2+x+e^x+2y=0$.
\end{question}

\begin{solution}
We find $y(0)$, $y'(0)$, $y''(0)$ by implicit differentiation.
\begin{itemize}
  \item $x=0$: $0+0+1+2y(0)=0\Rightarrow y(0)=-\frac12$.
  \item Differentiating: $2x+1+e^x+2y'=0\Rightarrow y'=-\frac{2x+1+e^x}{2}$, hence $y'(0)=-\frac{2}{2}=-1$.
  \item Differentiating again: $2+e^x+2y''=0\Rightarrow y''=-\frac{2+e^x}{2}$, hence $y''(0)=-\frac32$.
\end{itemize}
The Maclaurin polynomial of order 2:
\[ P_2(x)=y(0)+y'(0)x+\frac{y''(0)}{2}x^2=\boxed{-\frac12-x-\frac34x^2}. \]
(Check: here one can also isolate $y=-\frac12(x^2+x+e^x)$ and substitute $e^x=1+x+\frac{x^2}{2}+\dots$, which gives the same polynomial.)
\end{solution}

\subsection*{Question 6b}

\begin{question}
Is it true that the equation $e^x-\sin x=0$ has no real solutions? Justify
\end{question}

\begin{hint}
Check what happens for $x<0$: $e^x$ is very small, and $\sin x$ reaches $1$. Try the Intermediate Value Theorem on a suitable interval.
\end{hint}

\begin{solution}
\textbf{No} -- the equation has real solutions. Define $g(x)=e^x-\sin x$, continuous on $\R$.
\[ g(-\pi)=e^{-\pi}-0>0,\qquad g\left(-\tfrac{3\pi}{2}\right)=e^{-3\pi/2}-\sin\left(-\tfrac{3\pi}{2}\right)=e^{-3\pi/2}-1<0, \]
since $e^{-3\pi/2}<1$. By the \textbf{Intermediate Value Theorem} there exists $c\in\left(-\frac{3\pi}2,-\pi\right)$ with $g(c)=0$, i.e. $e^c=\sin c$ (numerically $c\approx-3.183$). In the same way, every interval $\left(-\frac{3\pi}{2}-2\pi k,\,-\pi-2\pi k\right)$ contains a solution, so there are infinitely many solutions (all of them negative: for $x\ge0$, $e^x\ge1\ge\sin x$ with equality impossible simultaneously, hence there are no non-negative solutions).
\end{solution}

\section*{2023/24 Semester A Moed A}

\subsection*{Question 1}

\begin{question}
Compute the following limit:
\[ \lim_{x\to\infty}\left(\frac{x^2+x+2}{x^2+x+1}\right)^{2x^2+4} \]
\end{question}

\begin{hint}
This is a limit of the form $1^\infty$. Write the base in the form $1+\frac{1}{x^2+x+1}$ and use the limit $\lim_{t\to\infty}\left(1+\frac1t\right)^t=e$.
\end{hint}

\begin{solution}
As $x\to\infty$ the base tends to $1$ and the exponent tends to $\infty$, i.e. this is an indeterminate form $1^\infty$.
Write the base in the form
\[ \frac{x^2+x+2}{x^2+x+1} = 1+\frac{1}{x^2+x+1}. \]
Denote $t=x^2+x+1$; as $x\to\infty$ also $t\to\infty$. Then
\[ \left(1+\frac{1}{x^2+x+1}\right)^{2x^2+4} = \left[\left(1+\frac1t\right)^{t}\right]^{\frac{2x^2+4}{x^2+x+1}}. \]
The expression in square brackets tends to $e$ (the famous limit $\lim_{t\to\infty}(1+\frac1t)^t=e$), and the exponent satisfies
\[ \lim_{x\to\infty}\frac{2x^2+4}{x^2+x+1} = \lim_{x\to\infty}\frac{2+\frac{4}{x^2}}{1+\frac1x+\frac1{x^2}} = 2. \]
By continuity of the function $(u,v)\mapsto u^v=e^{v\ln u}$ for $u>0$ (or: via $e^{v\ln u}$ and the continuity of $\ln$ and $\exp$) we get
\[ \lim_{x\to\infty}\left(\frac{x^2+x+2}{x^2+x+1}\right)^{2x^2+4} = e^{2}. \]

\textbf{Another method:} write the expression as $e^{(2x^2+4)\ln\left(1+\frac{1}{x^2+x+1}\right)}$. Since $\ln(1+s)\sim s$ as $s\to0$,
\[ (2x^2+4)\ln\left(1+\frac{1}{x^2+x+1}\right) \sim \frac{2x^2+4}{x^2+x+1}\xrightarrow[x\to\infty]{}2, \]
and again, by continuity of the exponential, the limit is $\boxed{e^2}$.
\end{solution}

\subsection*{Question 2}

\begin{question}
Let $a,b\in\R$. Define a function by
\[ f(x)=\begin{cases} \frac{\sin(8x)}{x} & x<0\\[2pt] ax+b & 0\le x\le 4\\[2pt] \frac{x^2-3x-4}{2-\sqrt{x}} & x>4 \end{cases} \]
For which values of $a,b$ is the function continuous on $\R$? Justify.
\end{question}

\begin{hint}
In each of the three open intervals the function is continuous as a composition/quotient of continuous functions. It remains to check the junction points $x=0$ and $x=4$.
\end{hint}

\begin{hint}
For the limit at $x=4$, factor the numerator $x^2-3x-4=(x-4)(x+1)$ and write $x-4=(\sqrt x-2)(\sqrt x+2)$.
\end{hint}

\begin{solution}
\textbf{Continuity inside each interval.}
\begin{itemize}
  \item On $x<0$: $f(x)=\frac{\sin 8x}{x}$ is a quotient of continuous functions whose denominator does not vanish, hence it is continuous.
  \item On $0<x<4$: $f(x)=ax+b$ is a polynomial, continuous.
  \item On $x>4$: $f(x)=\frac{x^2-3x-4}{2-\sqrt x}$ is a quotient of continuous functions, and the denominator vanishes only at $x=4$, which is not in this interval. Hence it is continuous.
\end{itemize}
Therefore $f$ is continuous on $\R$ if and only if it is continuous at the points $x=0$ and $x=4$.

\textbf{The point $x=0$.} $f(0)=b$, and from the right $\lim_{x\to0^+}f(x)=\lim_{x\to0^+}(ax+b)=b$. From the left, by the fundamental limit $\lim_{t\to0}\frac{\sin t}{t}=1$:
\[ \lim_{x\to0^-}\frac{\sin(8x)}{x} = \lim_{x\to0^-}8\cdot\frac{\sin(8x)}{8x} = 8. \]
Hence continuity at $0$ holds if and only if $b=8$.

\textbf{The point $x=4$.} $f(4)=4a+b$, and from the left $\lim_{x\to4^-}(ax+b)=4a+b$. From the right, for $x>4$:
\[ \frac{x^2-3x-4}{2-\sqrt x} = \frac{(x-4)(x+1)}{2-\sqrt x} = \frac{(\sqrt x-2)(\sqrt x+2)(x+1)}{-(\sqrt x-2)} = -(\sqrt x+2)(x+1), \]
and therefore
\[ \lim_{x\to4^+}f(x) = -(2+2)(4+1) = -20. \]
Continuity at $4$ holds if and only if $4a+b=-20$.

\textbf{Conclusion.} With $b=8$ we get $4a=-28$, i.e.
\[ \boxed{a=-7,\quad b=8}, \]
and only for these values is $f$ continuous on all of $\R$.
\end{solution}

\subsection*{Question 3}

\begin{question}
Find the intervals of increase and decrease and the absolute extrema of the function
\[ f(x)=\sqrt{2x-x^2} \]
on its domain.
\end{question}

\begin{solution}
\textbf{Domain.} We must require $2x-x^2=x(2-x)\ge0$, i.e. $0\le x\le 2$. The domain is the closed interval $[0,2]$.

\textbf{Derivative.} On the open interval $(0,2)$ the expression under the root is positive, so by the chain rule
\[ f'(x)=\frac{2-2x}{2\sqrt{2x-x^2}}=\frac{1-x}{\sqrt{2x-x^2}},\qquad 0<x<2. \]
(At the points $x=0,2$ the function is not differentiable -- the one-sided derivative is infinite -- but they are endpoints.)

\textbf{Sign of the derivative.} The denominator is positive, so the sign of $f'$ is the sign of $1-x$:
\begin{itemize}
  \item For $0<x<1$: $f'(x)>0$, hence $f$ is strictly increasing on $[0,1]$.
  \item For $1<x<2$: $f'(x)<0$, hence $f$ is strictly decreasing on $[1,2]$.
\end{itemize}
(Passing to the closed intervals is justified since $f$ is continuous on them -- a consequence of Lagrange's theorem.)

\textbf{Absolute extrema.} $f$ is continuous on the closed interval $[0,2]$, so by Weierstrass's theorem it attains a maximum and a minimum there. The candidate points are the critical point $x=1$ and the endpoints:
\[ f(0)=0,\qquad f(1)=\sqrt{2-1}=1,\qquad f(2)=0. \]
Hence
\[ \boxed{\max_{[0,2]}f=f(1)=1,\qquad \min_{[0,2]}f=f(0)=f(2)=0.} \]
(Indeed $f\ge0$ always, being a square root, so the minimum $0$ is also clear directly. In addition, $y=\sqrt{2x-x^2}$ is equivalent to $(x-1)^2+y^2=1,\ y\ge0$ -- the upper half of the circle with center $(1,0)$ and radius $1$, which confirms the results.)
\end{solution}

\subsection*{Question 4}

\begin{question}
Estimate $\sqrt[3]{30}$ using a linear approximation, and give a bound for the error.
\end{question}

\begin{hint}
Use $f(x)=\sqrt[3]{x}$ around the point $x_0=27$, where the root is known.
\end{hint}

\begin{hint}
For the error bound, use the Lagrange remainder of the Taylor polynomial of order 1: $R_1(x)=\frac{f''(c)}{2}(x-x_0)^2$.
\end{hint}

\begin{solution}
Choose $f(x)=\sqrt[3]{x}=x^{1/3}$ and $x_0=27$ (a point close to $30$ where $f(27)=3$).

\textbf{The linear approximation.} $f'(x)=\frac13x^{-2/3}$, hence $f'(27)=\frac{1}{3\cdot 9}=\frac1{27}$. The linear approximation (Taylor polynomial of order 1):
\[ f(x)\approx f(27)+f'(27)(x-27)=3+\frac{x-27}{27}, \]
and in particular
\[ \sqrt[3]{30}\approx 3+\frac{3}{27}=3+\frac19=\frac{28}{9}\approx 3.1111. \]

\textbf{Error bound.} By Taylor's theorem with the Lagrange remainder, there exists $c\in(27,30)$ such that
\[ \sqrt[3]{30}-\frac{28}{9} = R_1(30) = \frac{f''(c)}{2!}(30-27)^2. \]
$f''(x)=-\frac29x^{-5/3}$, and therefore
\[ |R_1(30)| = \frac{2}{9}c^{-5/3}\cdot\frac{9}{2} = \frac{1}{c^{5/3}} < \frac{1}{27^{5/3}}=\frac{1}{3^5}=\frac1{243}\approx 0.0041, \]
since $c>27$ and the function $c^{-5/3}$ is decreasing.

In addition, $f''<0$ and therefore $R_1<0$: the approximation is an overestimate. To summarize
\[ \boxed{\sqrt[3]{30}\approx\frac{28}{9}\approx3.111,\qquad 0<\frac{28}{9}-\sqrt[3]{30}<\frac{1}{243}\approx0.0041.} \]
(For verification: $\sqrt[3]{30}\approx3.1072$, and the actual error is about $0.0039$.)
\end{solution}

\subsection*{Question 5}

\begin{question}
Compute the antiderivative of
\[ f(x)=\frac{4x^2+9x+10}{x^3+2x^2+5x} \]
\end{question}

\begin{hint}
Factor the denominator: $x^3+2x^2+5x=x(x^2+2x+5)$, where the quadratic factor has no real roots. Decompose into partial fractions of the form $\frac{A}{x}+\frac{Bx+C}{x^2+2x+5}$.
\end{hint}

\begin{hint}
In the integral of $\frac{Bx+C}{x^2+2x+5}$, split into a numerator that is the derivative of the denominator plus a constant, and complete the square: $x^2+2x+5=(x+1)^2+4$.
\end{hint}

\begin{solution}
\textbf{Factoring the denominator.} $x^3+2x^2+5x=x(x^2+2x+5)$, and the discriminant of $x^2+2x+5$ is $4-20<0$, so the quadratic factor is irreducible. The degree of the numerator is less than the degree of the denominator, so we can decompose directly into partial fractions:
\[ \frac{4x^2+9x+10}{x(x^2+2x+5)}=\frac{A}{x}+\frac{Bx+C}{x^2+2x+5}. \]
Multiplying by the denominator: $4x^2+9x+10=A(x^2+2x+5)+(Bx+C)x$.
\begin{itemize}
  \item Substituting $x=0$: $10=5A$, i.e. $A=2$.
  \item Comparing the coefficients of $x^2$: $4=A+B$, i.e. $B=2$.
  \item Comparing the coefficients of $x$: $9=2A+C$, i.e. $C=5$.
\end{itemize}
Hence
\[ f(x)=\frac{2}{x}+\frac{2x+5}{x^2+2x+5}. \]

\textbf{Integration.} $\int\frac{2}{x}\,dx=2\ln|x|$. In the second fraction we split $2x+5=(2x+2)+3$, where $2x+2$ is the derivative of the denominator:
\[ \int\frac{2x+5}{x^2+2x+5}\,dx=\int\frac{2x+2}{x^2+2x+5}\,dx+3\int\frac{dx}{(x+1)^2+4}. \]
The first integral equals $\ln(x^2+2x+5)$ (substitution $u=x^2+2x+5$; no absolute value since the expression is positive). In the second we substitute $x+1=2t$:
\[ 3\int\frac{dx}{(x+1)^2+4}=\frac32\arctan\frac{x+1}{2}. \]

\textbf{Answer.}
\[ \boxed{\int f(x)\,dx = 2\ln|x|+\ln(x^2+2x+5)+\frac32\arctan\frac{x+1}{2}+C} \]
(on each of the intervals $x>0$ and $x<0$ separately; one can check by differentiation that $f$ is obtained.)
\end{solution}

\subsection*{Question 6}

\begin{question}
Compute the definite integral
\[ \int_0^{\pi}x\sin^2x\,dx \]
\end{question}

\begin{hint}
Use the identity $\sin^2x=\frac{1-\cos 2x}{2}$ and then integration by parts for $\int x\cos2x\,dx$.
\end{hint}

\begin{solution}
By the identity $\sin^2x=\frac{1-\cos2x}{2}$:
\[ \int_0^{\pi}x\sin^2x\,dx=\frac12\int_0^{\pi}x\,dx-\frac12\int_0^{\pi}x\cos2x\,dx. \]
\textbf{The first integral:} $\frac12\int_0^\pi x\,dx=\frac12\cdot\frac{\pi^2}{2}=\frac{\pi^2}{4}$.

\textbf{The second integral} -- integration by parts with $u=x$, $dv=\cos2x\,dx$, i.e. $du=dx$, $v=\frac12\sin2x$:
\[ \int_0^{\pi}x\cos2x\,dx=\left[\frac{x\sin2x}{2}\right]_0^{\pi}-\frac12\int_0^{\pi}\sin2x\,dx = 0+\left[\frac{\cos2x}{4}\right]_0^{\pi}=\frac14-\frac14=0. \]
Hence
\[ \boxed{\int_0^{\pi}x\sin^2x\,dx=\frac{\pi^2}{4}} \]

\textbf{Check (another method):} the substitution $x=\pi-t$ gives $I=\int_0^\pi(\pi-t)\sin^2t\,dt=\pi\int_0^\pi\sin^2t\,dt-I$, hence $I=\frac\pi2\cdot\frac\pi2=\frac{\pi^2}{4}$.
\end{solution}

\subsection*{Question 7}

\begin{question}
Compute the volume of the solid of revolution of the following function, around the $x$-axis, on the interval $[0,\pi]$:
\[ f(x)=\sqrt{\sin x}\cdot\cos^2x \]
\end{question}

\begin{hint}
The volume of a solid of revolution around the $x$-axis is $V=\pi\int_a^b f^2(x)\,dx$. Here $f^2(x)=\sin x\cos^4x$, and the substitution $u=\cos x$ is suitable.
\end{hint}

\begin{solution}
On the interval $[0,\pi]$ we have $\sin x\ge0$, so $f$ is defined and continuous on the whole interval. The formula for the volume of a solid of revolution around the $x$-axis:
\[ V=\pi\int_0^{\pi}f^2(x)\,dx=\pi\int_0^{\pi}\sin x\cos^4x\,dx. \]
Substitute $u=\cos x$, $du=-\sin x\,dx$; when $x=0$: $u=1$, and when $x=\pi$: $u=-1$. Hence
\[ \int_0^{\pi}\sin x\cos^4x\,dx=\int_{1}^{-1}u^4\,(-du)=\int_{-1}^{1}u^4\,du=\left[\frac{u^5}{5}\right]_{-1}^{1}=\frac25. \]
And therefore
\[ \boxed{V=\frac{2\pi}{5}} \]
\end{solution}

\section*{2023/24 Semester A Moed B}

\subsection*{Question 1}

\begin{question}
Compute the following limit:
\[ \lim_{x\to 1}\frac{(x^2+\ln x-1)\cdot\sin(x-1)}{(x-1)(e^x-e)} \]
\end{question}

\begin{hint}
Split the expression into a product of $\frac{\sin(x-1)}{x-1}$ and $\frac{x^2+\ln x-1}{e^x-e}$, and compute each limit separately.
\end{hint}

\begin{solution}
As $x\to1$ the numerator and the denominator both tend to $0$. Write the expression as a product (for $x$ close to $1$, $x\neq1$, $x>0$):
\[ \frac{(x^2+\ln x-1)\sin(x-1)}{(x-1)(e^x-e)}=\frac{\sin(x-1)}{x-1}\cdot\frac{x^2+\ln x-1}{e^x-e}. \]
\textbf{The first factor:} substitute $t=x-1\to0$; from the fundamental limit $\lim_{t\to0}\frac{\sin t}{t}=1$ we get $\lim_{x\to1}\frac{\sin(x-1)}{x-1}=1$.

\textbf{The second factor:} the numerator $x^2+\ln x-1\to 1+0-1=0$ and the denominator $e^x-e\to0$, i.e. this is a $\frac00$ case. Both functions are differentiable in a neighborhood of $1$ and the derivative of the denominator $e^x\neq0$, hence by L'Hôpital's rule
\[ \lim_{x\to1}\frac{x^2+\ln x-1}{e^x-e}=\lim_{x\to1}\frac{2x+\frac1x}{e^x}=\frac{2+1}{e}=\frac3e, \]
provided that the limit on the right exists -- and it indeed exists (the function $\frac{2x+1/x}{e^x}$ is continuous at $1$).

\textbf{Summary:} by the arithmetic of limits (product of two finite limits),
\[ \lim_{x\to 1}\frac{(x^2+\ln x-1)\sin(x-1)}{(x-1)(e^x-e)}=1\cdot\frac3e=\boxed{\frac3e}. \]
\end{solution}

\subsection*{Question 2}

\begin{question}
Define a function by
\[ f(x)=\begin{cases}\frac{\ln(1+x^2)}{\sin^2 x} & x<0\\ 0 & 0\le x\le 1\\ \frac{x}{x^2-1} & x>1\end{cases} \]
Find all the points of discontinuity of the function, and classify them as removable, jump or essential discontinuities. Justify.
\end{question}

\begin{hint}
Inside each of the intervals the function is a composition/quotient of elementary functions. Check where a denominator vanishes, and separately the junction points $x=0$ and $x=1$.
\end{hint}

\begin{hint}
Near $0$: $\ln(1+x^2)\approx x^2$ and $\sin^2x\approx x^2$.
\end{hint}

\begin{solution}
\textbf{Inside the open intervals.}
\begin{itemize}
  \item On the interval $(0,1)$ the function is constant and hence continuous.
  \item On the interval $(1,\infty)$: $f(x)=\frac{x}{x^2-1}$ is a quotient of polynomials whose denominator does not vanish there ($x^2-1>0$), hence it is continuous.
  \item On the interval $(-\infty,0)$: $f(x)=\frac{\ln(1+x^2)}{\sin^2x}$ is a quotient of continuous functions, hence continuous at every point where $\sin x\neq0$. The denominator vanishes at the points $x=-k\pi$, $k=1,2,3,\dots$, and there the function is not defined. At these points the numerator tends to $\ln(1+k^2\pi^2)>0$ and the denominator tends to $0^+$ (since $\sin^2x>0$ in a punctured neighborhood), and therefore
  \[ \lim_{x\to-k\pi^\pm}f(x)=+\infty. \]
  The one-sided limits are not finite, hence $x=-k\pi$ ($k\in\N$) are \textbf{essential} discontinuities (of the second kind).
\end{itemize}

\textbf{The point $x=0$.} $f(0)=0$, and on the right $f\equiv0$, hence $\lim_{x\to0^+}f(x)=0$. From the left:
\[ \lim_{x\to0^-}\frac{\ln(1+x^2)}{\sin^2x}=\lim_{x\to0^-}\frac{\ln(1+x^2)}{x^2}\cdot\left(\frac{x}{\sin x}\right)^2=1\cdot1^2=1, \]
where we used the fundamental limit $\lim_{t\to0}\frac{\ln(1+t)}{t}=1$ (with $t=x^2\to0$) and $\lim_{x\to0}\frac{\sin x}{x}=1$. The two one-sided limits are finite and different ($1\neq0$), hence at $x=0$ there is a \textbf{jump} discontinuity.

\textbf{The point $x=1$.} $f(1)=0$ and $\lim_{x\to1^-}f(x)=0$. From the right, the numerator tends to $1$ and the denominator $x^2-1\to0^+$, hence
\[ \lim_{x\to1^+}\frac{x}{x^2-1}=+\infty. \]
The right-hand limit is not finite, hence at $x=1$ there is an \textbf{essential} discontinuity.

\textbf{Answer:} at $x=0$ -- a jump discontinuity; at $x=1$ -- an essential discontinuity; at the points $x=-k\pi$, $k=1,2,\dots$ (where $f$ is not defined) -- essential discontinuities (the limit there is $+\infty$). At all other points $f$ is continuous.
\end{solution}

\subsection*{Question 3}

\begin{question}
Find the intervals where the function is concave up and concave down (convexity and concavity), and the inflection points of the function
\[ f(x)=\ln\left(\frac{e^x}{1+e^x}\right) \]
on its domain.
\end{question}

\begin{hint}
Simplify first: $\ln\frac{e^x}{1+e^x}=x-\ln(1+e^x)$.
\end{hint}

\begin{solution}
\textbf{Domain.} For every real $x$, $\frac{e^x}{1+e^x}>0$, hence $f$ is defined on all of $\R$. By the laws of logarithms
\[ f(x)=\ln e^x-\ln(1+e^x)=x-\ln(1+e^x). \]

\textbf{First derivative.}
\[ f'(x)=1-\frac{e^x}{1+e^x}=\frac{1}{1+e^x}. \]

\textbf{Second derivative.} By the chain rule,
\[ f''(x)=-\frac{e^x}{(1+e^x)^2}. \]
For every $x\in\R$ we have $e^x>0$ and $(1+e^x)^2>0$, hence $f''(x)<0$ for every $x$.

\textbf{Conclusion.} Since $f''<0$ on all of $\R$, the function is concave down (concave, $\cap$) on the whole line $(-\infty,\infty)$, and there is no interval on which it is concave up. An inflection point is a point where the direction of convexity changes; since $f''$ does not change sign (and does not even vanish), \textbf{the function has no inflection points}.
\end{solution}

\subsection*{Question 4}

\begin{question}
Estimate $\sqrt[3]{e}$ using a second-order Taylor approximation, and give a bound for the error.
\end{question}

\begin{hint}
Use the Maclaurin polynomial of $e^x$ of order $2$ at the point $x=\frac13$, and the remainder in Lagrange form.
\end{hint}

\begin{solution}
We use the function $f(x)=e^x$ around $x_0=0$, and substitute $x=\frac13$, since $\sqrt[3]{e}=e^{1/3}$.

\textbf{Taylor polynomial.} For every $n$, $f^{(n)}(x)=e^x$, hence $f^{(n)}(0)=1$. The Maclaurin polynomial of order $2$:
\[ P_2(x)=1+x+\frac{x^2}{2}. \]
Hence
\[ \sqrt[3]{e}\approx P_2\!\left(\tfrac13\right)=1+\frac13+\frac1{18}=\frac{25}{18}\approx1.3889. \]

\textbf{Error bound.} By Taylor's theorem with the Lagrange remainder, there exists $c$ between $0$ and $\frac13$ such that
\[ R_2\!\left(\tfrac13\right)=e^{1/3}-P_2\!\left(\tfrac13\right)=\frac{f'''(c)}{3!}\left(\frac13\right)^3=\frac{e^c}{6\cdot27}=\frac{e^c}{162}. \]
Since $e^x$ is increasing and $0<c<\frac13$, we have $e^c<e^{1/3}$. Since $e<3<\frac{27}{8}=1.5^3$, we get $e^{1/3}<1.5$. Hence
\[ 0<R_2\!\left(\tfrac13\right)<\frac{1.5}{162}=\frac1{108}\approx0.0093. \]
(One can also use the cruder bound $e^c<e<3$ and get an error smaller than $\frac{3}{162}=\frac1{54}$.)

\textbf{Answer:} $\sqrt[3]{e}\approx\frac{25}{18}\approx1.3889$, and the error is positive and smaller than $\frac1{108}$. (For verification: $\sqrt[3]{e}=1.39561\ldots$, and the actual error is about $0.0067$.)
\end{solution}

\subsection*{Question 5}

\begin{question}
Compute the antiderivative of
\[ f(x)=\frac{5x^3+x^2+x-2}{x^4+x^2} \]
\end{question}

\begin{hint}
Factor the denominator: $x^4+x^2=x^2(x^2+1)$, and use a partial fraction decomposition of the form $\frac Ax+\frac B{x^2}+\frac{Cx+D}{x^2+1}$.
\end{hint}

\begin{solution}
The degree of the numerator ($3$) is less than the degree of the denominator ($4$), so there is no need for polynomial division. Factor the denominator: $x^4+x^2=x^2(x^2+1)$, where $x^2+1$ is irreducible over $\R$. We look for a partial fraction decomposition:
\[ \frac{5x^3+x^2+x-2}{x^2(x^2+1)}=\frac Ax+\frac B{x^2}+\frac{Cx+D}{x^2+1}. \]
Multiply by $x^2(x^2+1)$:
\[ 5x^3+x^2+x-2=Ax(x^2+1)+B(x^2+1)+(Cx+D)x^2=(A+C)x^3+(B+D)x^2+Ax+B. \]
Comparing coefficients: $B=-2$, $A=1$, $A+C=5\Rightarrow C=4$, $B+D=1\Rightarrow D=3$. Hence
\[ f(x)=\frac1x-\frac2{x^2}+\frac{4x+3}{x^2+1}=\frac1x-\frac2{x^2}+2\cdot\frac{2x}{x^2+1}+\frac{3}{x^2+1}. \]
Integrate term by term using the basic integrals $\int\frac{dx}x=\ln|x|$, $\int x^{-2}dx=-x^{-1}$, $\int\frac{2x}{x^2+1}dx=\ln(x^2+1)$ (substitution $t=x^2+1$) and $\int\frac{dx}{x^2+1}=\arctan x$:
\[ \boxed{\int f(x)\,dx=\ln|x|+\frac2x+2\ln(x^2+1)+3\arctan x+C} \]
(on each of the intervals $x>0$ and $x<0$ separately; the constant may be different on each interval).

\textbf{Check:} differentiation gives $\frac1x-\frac2{x^2}+\frac{4x}{x^2+1}+\frac3{x^2+1}=f(x)$.
\end{solution}

\subsection*{Question 6}

\begin{question}
Compute the definite integral
\[ \int_0^1 (x+1)^2\cdot e^x\,dx \]
\end{question}

\begin{hint}
Integration by parts twice, each time differentiating the polynomial.
\end{hint}

\begin{solution}
We use integration by parts: $\int u\,v'=uv-\int u'v$.

\textbf{First step:} $u=(x+1)^2$, $v'=e^x$, hence $u'=2(x+1)$, $v=e^x$:
\[ \int_0^1(x+1)^2e^x\,dx=\Big[(x+1)^2e^x\Big]_0^1-2\int_0^1(x+1)e^x\,dx=(4e-1)-2\int_0^1(x+1)e^x\,dx. \]

\textbf{Second step:} $u=x+1$, $v'=e^x$, hence $u'=1$, $v=e^x$:
\[ \int_0^1(x+1)e^x\,dx=\Big[(x+1)e^x\Big]_0^1-\int_0^1e^x\,dx=(2e-1)-(e-1)=e. \]

\textbf{Summary:}
\[ \int_0^1(x+1)^2e^x\,dx=4e-1-2e=\boxed{2e-1}\approx 4.4366. \]
(Alternatively: an antiderivative is $F(x)=e^x\big((x+1)^2-2(x+1)+2\big)=e^x(x^2+1)$, and by the Fundamental Theorem $F(1)-F(0)=2e-1$.)
\end{solution}

\subsection*{Question 7}

\begin{question}
Determine whether the following improper integral converges:
\[ \int_2^\infty \frac{\sqrt[5]{x}}{3x^2+x+100}\,dx \]
\end{question}

\begin{hint}
At infinity, the function behaves like $\frac{x^{1/5}}{3x^2}$. Compare (using the limit comparison test) with $\frac{1}{x^{9/5}}$.
\end{hint}

\begin{solution}
The function $f(x)=\frac{\sqrt[5]{x}}{3x^2+x+100}$ is continuous and positive on $[2,\infty)$ (the denominator is positive), hence the integral is improper only because of the upper limit, and we may use the comparison tests for non-negative functions.

Compare with $g(x)=\frac{1}{x^{9/5}}=\frac{x^{1/5}}{x^2}$, which is also positive on $[2,\infty)$:
\[ \lim_{x\to\infty}\frac{f(x)}{g(x)}=\lim_{x\to\infty}\frac{x^{1/5}\cdot x^{2}}{x^{1/5}(3x^2+x+100)}=\lim_{x\to\infty}\frac{1}{3+\frac1x+\frac{100}{x^2}}=\frac13. \]
The limit is finite and positive, hence by the \textbf{limit comparison test} the integrals $\int_2^\infty f$ and $\int_2^\infty g$ either both converge or both diverge.

The integral $\int_2^\infty\frac{dx}{x^p}$ converges if and only if $p>1$; here $p=\frac95>1$, and explicitly
\[ \int_2^\infty x^{-9/5}dx=\lim_{R\to\infty}\left[-\frac54x^{-4/5}\right]_2^R=\frac54\cdot2^{-4/5}<\infty. \]

\textbf{Answer:} the integral $\int_2^\infty\frac{\sqrt[5]{x}}{3x^2+x+100}\,dx$ \textbf{converges}.
\end{solution}

\section*{2023/24 Semester A Special Moed}

\subsection*{Question 1}

\begin{question}
Compute the following limit:
\[ \lim_{x\to0^+}\frac{\ln(e^x-1)}{3\cdot\ln x} \]
\end{question}

\begin{hint}
This is an $\frac{\infty}{\infty}$ case; use L'Hôpital's rule.
\end{hint}

\begin{solution}
As $x\to0^+$: $e^x-1\to0^+$, hence $\ln(e^x-1)\to-\infty$, and also $3\ln x\to-\infty$. This is an $\frac{\infty}{\infty}$ case. The numerator and the denominator are differentiable on $(0,\infty)$ and the derivative of the denominator $\frac3x\neq0$, so by L'Hôpital's rule (provided that the limit on the right exists):
\[ \lim_{x\to0^+}\frac{\ln(e^x-1)}{3\ln x}=\lim_{x\to0^+}\frac{\frac{e^x}{e^x-1}}{\frac3x}=\lim_{x\to0^+}\frac{e^x}{3}\cdot\frac{x}{e^x-1}. \]
We have $e^x\to1$, and from the fundamental limit $\lim_{x\to0}\frac{e^x-1}{x}=1$ we get $\frac{x}{e^x-1}\to1$. By the arithmetic of limits the limit on the right exists and equals $\frac13\cdot1=\frac13$, and therefore
\[ \lim_{x\to0^+}\frac{\ln(e^x-1)}{3\ln x}=\boxed{\frac13}. \]
\end{solution}

\subsection*{Question 2}

\begin{question}
Define a function by
\[ f(x)=\begin{cases}2^{\frac{\sin x}{x}} & x<0\\ ax+b & 0\le x\le 1\\ \frac{x^2+x-2}{\sqrt{x}-1} & x>1\end{cases} \]
For which values of the parameters $a,b$ is the function continuous on all of $\R$? Justify.
\end{question}

\begin{hint}
Inside each interval the function is continuous; it remains to require continuity at the junction points $x=0$ and $x=1$.
\end{hint}

\begin{hint}
For $x>1$: $x^2+x-2=(x-1)(x+2)$ and $x-1=(\sqrt x-1)(\sqrt x+1)$.
\end{hint}

\begin{solution}
\textbf{Inside the intervals.} On $(-\infty,0)$ the function $\frac{\sin x}{x}$ is continuous (a quotient of continuous functions, denominator different from $0$), and $2^t$ is continuous, hence the composition is continuous. On $(0,1)$ it is a linear function, continuous. On $(1,\infty)$ it is a quotient of continuous functions whose denominator is $\sqrt x-1>0$, hence continuous. Therefore, for all $a,b$ the function is continuous at every point except possibly $x=0$ and $x=1$.

\textbf{Continuity at $x=0$.} $f(0)=b$, and from the right $\lim_{x\to0^+}(ax+b)=b$. From the left, since $\lim_{x\to0}\frac{\sin x}{x}=1$ and $2^t$ is continuous at $t=1$, by the theorem on the limit of a composite function
\[ \lim_{x\to0^-}2^{\frac{\sin x}{x}}=2^1=2. \]
Hence $f$ is continuous at $0$ if and only if $b=2$.

\textbf{Continuity at $x=1$.} $f(1)=a+b$ and $\lim_{x\to1^-}f(x)=a+b$. From the right, for $x>1$:
\[ \frac{x^2+x-2}{\sqrt x-1}=\frac{(x-1)(x+2)}{\sqrt x-1}=\frac{(\sqrt x-1)(\sqrt x+1)(x+2)}{\sqrt x-1}=(\sqrt x+1)(x+2), \]
and therefore $\lim_{x\to1^+}f(x)=(1+1)(1+2)=6$. Hence $f$ is continuous at $1$ if and only if $a+b=6$.

\textbf{Answer:} $f$ is continuous on all of $\R$ if and only if $\boxed{b=2,\ a=4}$.
\end{solution}

\subsection*{Question 3}

\begin{question}
Find the absolute minimum and maximum of the function
\[ f(x)=(x+2)\cdot e^x \]
on its domain, or explain why they do not exist.
\end{question}

\begin{hint}
Find the intervals of increase and decrease using $f'$, and check the behavior of the function as $x\to\pm\infty$.
\end{hint}

\begin{solution}
\textbf{Domain:} all of $\R$; $f$ is differentiable at every point.

\textbf{Derivative:} by the product rule
\[ f'(x)=e^x+(x+2)e^x=(x+3)e^x. \]
Since $e^x>0$, the sign of $f'$ is the sign of $x+3$: $f'<0$ on $(-\infty,-3)$ and $f'>0$ on $(-3,\infty)$. Hence (by a corollary of Lagrange's theorem) $f$ is strictly decreasing on $(-\infty,-3]$ and strictly increasing on $[-3,\infty)$.

\textbf{Absolute minimum:} by monotonicity, for every $x\le-3$ we have $f(x)\ge f(-3)$, and for every $x\ge-3$ we have $f(x)\ge f(-3)$. Hence $x=-3$ is a point of absolute minimum, and the minimum value is
\[ f(-3)=(-3+2)e^{-3}=-\frac1{e^3}\approx-0.0498. \]

\textbf{Absolute maximum:} $\lim_{x\to\infty}(x+2)e^x=+\infty$, hence the function is not bounded above and has no absolute maximum. (For completeness: $\lim_{x\to-\infty}(x+2)e^x=\lim_{x\to-\infty}\frac{x+2}{e^{-x}}=0$ by L'Hôpital's rule, so the function tends to $0^-$ on the left, but on the right it is unbounded.)

\textbf{Answer:} the absolute minimum is $-e^{-3}$, attained at $x=-3$; an absolute maximum does not exist since $f(x)\to+\infty$ as $x\to\infty$.
\end{solution}

\subsection*{Question 4}

\begin{question}
Prove that the following equation has a unique positive real solution:
\[ 2x^3+\arctan x+\sqrt{x}=7 \]
\end{question}

\begin{hint}
Existence: the Intermediate Value Theorem for $g(x)=2x^3+\arctan x+\sqrt x-7$ on a suitable interval. Uniqueness: show that $g$ is strictly increasing.
\end{hint}

\begin{solution}
Define $g(x)=2x^3+\arctan x+\sqrt x-7$ on $[0,\infty)$. $g$ is continuous there as a sum of continuous functions, and differentiable on $(0,\infty)$.

\textbf{Existence.} $g(0)=-7<0$, and
\[ g(2)=16+\arctan2+\sqrt2-7=9+\arctan 2+\sqrt2>0. \]
By the Intermediate Value Theorem there exists $x_0\in(0,2)$ with $g(x_0)=0$, i.e. a positive solution of the equation.

\textbf{Uniqueness.} For every $x>0$:
\[ g'(x)=6x^2+\frac1{1+x^2}+\frac1{2\sqrt x}>0. \]
Hence $g$ is strictly increasing on $(0,\infty)$ (a corollary of Lagrange's theorem: if $0<x_1<x_2$ then $g(x_2)-g(x_1)=g'(c)(x_2-x_1)>0$). A strictly increasing function attains each value at most once, hence it has at most one zero in $(0,\infty)$.

(Alternatively: if there were two positive solutions $x_1<x_2$, then by Rolle's theorem there would be $c\in(x_1,x_2)$ with $g'(c)=0$, contradicting $g'>0$.)

\textbf{Conclusion:} the equation has a unique positive real solution (numerically $x_0\approx1.3486$).
\end{solution}

\subsection*{Question 5}

\begin{question}
Compute the antiderivative of
\[ f(x)=\sin^2x\cdot\cos^5x \]
\end{question}

\begin{hint}
The power of $\cos x$ is odd: separate one factor $\cos x$, write $\cos^4x=(1-\sin^2x)^2$ and substitute $t=\sin x$.
\end{hint}

\begin{solution}
Write $\cos^5x=\cos^4x\cdot\cos x=(1-\sin^2x)^2\cos x$. Hence
\[ \int\sin^2x\cos^5x\,dx=\int\sin^2x(1-\sin^2x)^2\cos x\,dx. \]
Substitute $t=\sin x$, $dt=\cos x\,dx$:
\[ \int t^2(1-t^2)^2\,dt=\int\left(t^2-2t^4+t^6\right)dt=\frac{t^3}3-\frac{2t^5}5+\frac{t^7}7+C. \]
Return to the variable $x$:
\[ \boxed{\int\sin^2x\cos^5x\,dx=\frac{\sin^3x}{3}-\frac{2\sin^5x}{5}+\frac{\sin^7x}{7}+C}. \]
\textbf{Check:} the derivative is $\cos x\left(\sin^2x-2\sin^4x+\sin^6x\right)=\cos x\sin^2x(1-\sin^2x)^2=\sin^2x\cos^5x$.
\end{solution}

\subsection*{Question 6}

\begin{question}
Compute the volume of the solid of revolution around the $x$-axis of the function
\[ f(x)=\frac{1}{\sqrt{x}(x-1)} \]
on the interval $[4,9]$.
\end{question}

\begin{hint}
Volume of a solid of revolution: $V=\pi\int_a^b f^2(x)\,dx$. Here $f^2(x)=\frac1{x(x-1)^2}$ -- use partial fractions.
\end{hint}

\begin{solution}
$f$ is continuous on $[4,9]$, and the volume of the solid of revolution around the $x$-axis is
\[ V=\pi\int_4^9f^2(x)\,dx=\pi\int_4^9\frac{dx}{x(x-1)^2}. \]
\textbf{Partial fractions:}
\[ \frac1{x(x-1)^2}=\frac Ax+\frac B{x-1}+\frac C{(x-1)^2}\ \Longrightarrow\ 1=A(x-1)^2+Bx(x-1)+Cx. \]
Substituting $x=0$ gives $A=1$; substituting $x=1$ gives $C=1$; comparing the coefficient of $x^2$: $A+B=0$, hence $B=-1$:
\[ \frac1{x(x-1)^2}=\frac1x-\frac1{x-1}+\frac1{(x-1)^2}. \]
\textbf{Integration:} an antiderivative (for $x>1$) is $F(x)=\ln x-\ln(x-1)-\frac1{x-1}$, and by the Fundamental Theorem
\begin{align*}
\int_4^9\frac{dx}{x(x-1)^2}&=\left[\ln\frac{x}{x-1}-\frac1{x-1}\right]_4^9=\left(\ln\frac98-\frac18\right)-\left(\ln\frac43-\frac13\right)\\
&=\ln\left(\frac98\cdot\frac34\right)+\frac13-\frac18=\ln\frac{27}{32}+\frac5{24}.
\end{align*}
\textbf{Answer:}
\[ V=\pi\left(\frac5{24}+\ln\frac{27}{32}\right)=\pi\left(\frac5{24}-\ln\frac{32}{27}\right)\approx0.1207. \]
\end{solution}

\subsection*{Question 7}

\begin{question}
Compute the following improper integral, or show that it diverges:
\[ \int_2^\infty x\cdot e^{-x}\,dx \]
\end{question}

\begin{hint}
Integration by parts with $u=x$, $v'=e^{-x}$, and then a limit as the upper limit tends to infinity.
\end{hint}

\begin{solution}
By definition, $\int_2^\infty xe^{-x}\,dx=\lim_{R\to\infty}\int_2^Rxe^{-x}\,dx$, if the limit exists.

\textbf{Antiderivative:} integration by parts with $u=x$, $v'=e^{-x}$ ($u'=1$, $v=-e^{-x}$):
\[ \int xe^{-x}\,dx=-xe^{-x}+\int e^{-x}\,dx=-xe^{-x}-e^{-x}+C=-(x+1)e^{-x}+C. \]
Hence
\[ \int_2^Rxe^{-x}\,dx=\Big[-(x+1)e^{-x}\Big]_2^R=3e^{-2}-\frac{R+1}{e^R}. \]
\textbf{Passing to the limit:} $\lim_{R\to\infty}\frac{R+1}{e^R}$ is an $\frac\infty\infty$ case, and by L'Hôpital's rule $\lim_{R\to\infty}\frac{1}{e^R}=0$. Hence
\[ \int_2^\infty xe^{-x}\,dx=\boxed{\frac{3}{e^2}}\approx0.406, \]
and the integral converges.
\end{solution}

\section*{2023/24 Semester B Moed A}

\subsection*{Question 1a}

\begin{question}
Compute
\[ \lim_{x\to\infty} x\left(\ln(x+1)-\ln(x)\right) \]
\end{question}

\begin{hint}
Write $\ln(x+1)-\ln x=\ln\left(1+\frac1x\right)$ and substitute $t=\frac1x$.
\end{hint}

\begin{solution}
By the laws of logarithms, $\ln(x+1)-\ln(x)=\ln\frac{x+1}{x}=\ln\left(1+\frac1x\right)$. Substitute $t=\frac1x$; as $x\to\infty$ we have $t\to0^+$, and therefore
\[ \lim_{x\to\infty} x\ln\left(1+\frac1x\right)=\lim_{t\to0^+}\frac{\ln(1+t)}{t}. \]
This is a limit of the form $\frac00$; by L'Hôpital's rule (or the fundamental limit $\frac{\ln(1+t)}{t}\to1$):
\[ \lim_{t\to0^+}\frac{\ln(1+t)}{t}=\lim_{t\to0^+}\frac{1/(1+t)}{1}=1. \]
The answer: $\boxed{1}$.
\end{solution}

\subsection*{Question 1b}

\begin{question}
Compute
\[ \lim_{x\to 0} \frac{x^2+\sin(x)\cos(x)}{1-e^x} \]
\end{question}

\begin{solution}
As $x\to0$ the numerator tends to $0+\sin0\cos0=0$ and the denominator to $1-e^0=0$, i.e. a limit of the form $\frac00$. The numerator and the denominator are differentiable, and the derivative of the denominator $-e^x\neq0$ in a neighborhood of $0$, so by L'Hôpital's rule:
\[ \lim_{x\to0}\frac{x^2+\sin x\cos x}{1-e^x}=\lim_{x\to0}\frac{2x+\cos^2x-\sin^2x}{-e^x}=\frac{0+1-0}{-1}=-1, \]
where we used the continuity at $0$ of the resulting expression. (Alternatively: $\sin x\cos x=\frac12\sin2x$, and dividing the numerator and denominator by $x$: $\frac{x+\frac{\sin 2x}{2x}}{\frac{1-e^x}{x}}\to\frac{0+1}{-1}$.)
The answer: $\boxed{-1}$.
\end{solution}

\subsection*{Question 2}

\begin{question}
Define a function as follows:
\[ f(x)=\begin{cases} x^2\sin\left(\frac{3}{x}\right) & x\neq0\\ 0 & x=0\end{cases} \]
\begin{enumerate}
  \item Is $f(x)$ continuous at the point $x=0$?
  \item Is $f(x)$ differentiable at the point $x=0$?
\end{enumerate}
\end{question}

\begin{hint}
$\sin\left(\frac3x\right)$ is bounded; use the Squeeze theorem (or ``tends to zero times bounded'').
\end{hint}

\begin{solution}
\begin{enumerate}
  \item For every $x\neq0$ we have $\left|\sin\frac3x\right|\le1$, and therefore
  \[ 0\le |f(x)|=x^2\left|\sin\frac3x\right|\le x^2. \]
  Since $x^2\to0$ as $x\to0$, by the Squeeze theorem $|f(x)|\to0$, i.e. $\lim_{x\to0}f(x)=0=f(0)$. Hence $f$ is \textbf{continuous} at $x=0$.
  \item We compute the derivative by definition:
  \[ f'(0)=\lim_{h\to0}\frac{f(h)-f(0)}{h}=\lim_{h\to0}\frac{h^2\sin\frac3h}{h}=\lim_{h\to0}h\sin\frac3h. \]
  Again $0\le\left|h\sin\frac3h\right|\le|h|\to0$, and by the Squeeze theorem the limit exists and equals $0$ (a function tending to zero times a bounded function). Hence $f$ is \textbf{differentiable} at $x=0$ and $f'(0)=0$.
\end{enumerate}
\end{solution}

\subsection*{Question 3}

\begin{question}
Analyze the following function (local extremum points, intervals of convexity/concavity, asymptotes):
\[ f(x)=\frac{2-x-3x^2}{x^2-1} \]
Sketch its graph.
\end{question}

\begin{hint}
Factor the numerator and the denominator -- there is a common factor.
\end{hint}

\begin{hint}
After cancelling, one can write $f(x)=-3-\frac{1}{x-1}$ (for $x\neq\pm1$).
\end{hint}

\begin{solution}
\textbf{Domain:} $x^2-1\neq0$, i.e. $x\neq\pm1$.

\textbf{Simplification:} $2-x-3x^2=-(3x^2+x-2)=-(3x-2)(x+1)$ and $x^2-1=(x-1)(x+1)$. Hence for every $x\neq\pm1$:
\[ f(x)=\frac{-(3x-2)(x+1)}{(x-1)(x+1)}=\frac{2-3x}{x-1}=-3-\frac{1}{x-1}. \]
(Check: $-3(x-1)-1=2-3x$.)

\textbf{The point $x=-1$:} $\lim_{x\to-1}f(x)=\frac{2+3}{-2}=-\frac52$ -- a finite limit, hence this is a removable discontinuity (a ``hole'' in the graph at the point $(-1,-\frac52)$), and there is no vertical asymptote there.

\textbf{Intersections with the axes:} $f(0)=-2$; $f(x)=0\iff x=\frac23$. The points $(0,-2)$, $(\frac23,0)$.

\textbf{First derivative:} $f'(x)=\frac{1}{(x-1)^2}>0$ for every $x$ in the domain. Hence $f$ is increasing on each of the intervals $(-\infty,-1)$, $(-1,1)$, $(1,\infty)$, and there are \textbf{no local extremum points} (the derivative does not vanish anywhere, and by Fermat's theorem there is no interior extremum).

\textbf{Second derivative:} $f''(x)=-\frac{2}{(x-1)^3}$.
\begin{itemize}
  \item For $x<1$ ($x\neq-1$): $(x-1)^3<0$, hence $f''>0$ -- $f$ is \textbf{convex} ($\cup$) on $(-\infty,-1)$ and on $(-1,1)$.
  \item For $x>1$: $f''<0$ -- $f$ is \textbf{concave} ($\cap$) on $(1,\infty)$.
\end{itemize}
There are no inflection points: the change of convexity happens only at $x=1$, which is not in the domain.

\textbf{Asymptotes:}
\begin{itemize}
  \item Vertical: $x=1$. $\lim_{x\to1^-}f(x)=-3-\frac{1}{0^-}=+\infty$, $\lim_{x\to1^+}f(x)=-\infty$.
  \item Horizontal: $\lim_{x\to\pm\infty}f(x)=-3$ (the degrees of the numerator and the denominator are equal, the ratio of the leading coefficients is $\frac{-3}{1}$). Hence $y=-3$ is a horizontal asymptote in both directions; there is no oblique asymptote.
\end{itemize}

\textbf{The graph:} the hyperbola $y=-3-\frac1{x-1}$ (a shift of $-\frac1x$) with a hole at the point $(-1,-\frac52)$. To the left of $x=1$ the graph rises from the asymptote $y=-3$ (above it) through the hole $(-1,-2.5)$, $(0,-2)$, $(\frac23,0)$ and tends to $+\infty$ as $x\to1^-$, and it is convex. To the right of $x=1$ the graph rises from $-\infty$ and approaches $y=-3$ from below, and it is concave.
\end{solution}

\subsection*{Question 4a}

\begin{question}
Write the Maclaurin polynomial of order 3 of $f(x)=\ln(1+x)$.
\end{question}

\begin{solution}
We compute the derivatives at $0$:
\[ f(x)=\ln(1+x),\quad f'(x)=\frac1{1+x},\quad f''(x)=-\frac1{(1+x)^2},\quad f'''(x)=\frac{2}{(1+x)^3}, \]
hence $f(0)=0,\ f'(0)=1,\ f''(0)=-1,\ f'''(0)=2$. The Maclaurin polynomial:
\[ P_3(x)=f(0)+f'(0)x+\frac{f''(0)}{2!}x^2+\frac{f'''(0)}{3!}x^3=\boxed{x-\frac{x^2}{2}+\frac{x^3}{3}}. \]
\end{solution}

\subsection*{Question 4b}

\begin{question}
Find a linear approximation of $\sqrt{15}$
\end{question}

\begin{hint}
Use the tangent line to the graph of $g(x)=\sqrt x$ at a point close to $15$ where the root is known.
\end{hint}

\begin{solution}
We use the function $g(x)=\sqrt x$ around the point $a=16$ (close to $15$, where the root is known). $g(16)=4$, $g'(x)=\frac{1}{2\sqrt x}$, $g'(16)=\frac18$. The linear approximation (the equation of the tangent line):
\[ g(x)\approx g(16)+g'(16)(x-16)=4+\frac{x-16}{8}. \]
For $x=15$: $\sqrt{15}\approx4-\frac18=\boxed{3.875}$.
(For comparison, $\sqrt{15}\approx3.873$. Since $g''<0$ -- $g$ is concave -- the tangent line lies above the graph, hence this is an approximation from above; by the Lagrange remainder the error is $\frac{|g''(c)|}{2}\cdot1^2=\frac{1}{8c^{3/2}}<\frac{1}{8\cdot 15^{3/2}}<0.003$.)
\end{solution}

\subsection*{Question 5a}

\begin{question}
Compute
\[ \int\frac{x^3+2x^2-3x+5}{x^2+x-6}\,dx \]
\end{question}

\begin{hint}
Polynomial long division, then partial fractions ($x^2+x-6=(x+3)(x-2)$).
\end{hint}

\begin{solution}
The degree of the numerator is greater than the degree of the denominator, so we divide the polynomials:
\[ x^3+2x^2-3x+5=(x+1)(x^2+x-6)+(2x+11), \]
(check: $(x+1)(x^2+x-6)=x^3+2x^2-5x-6$, and the remainder is $x^3+2x^2-3x+5-(x^3+2x^2-5x-6)=2x+11$.) Hence
\[ \frac{x^3+2x^2-3x+5}{x^2+x-6}=x+1+\frac{2x+11}{(x+3)(x-2)}. \]
Partial fraction decomposition: $\frac{2x+11}{(x+3)(x-2)}=\frac{A}{x+3}+\frac{B}{x-2}$, i.e. $2x+11=A(x-2)+B(x+3)$. Substituting $x=2$: $15=5B\Rightarrow B=3$; substituting $x=-3$: $5=-5A\Rightarrow A=-1$. Hence
\[ \int\frac{x^3+2x^2-3x+5}{x^2+x-6}\,dx=\int\left(x+1-\frac1{x+3}+\frac3{x-2}\right)dx=\boxed{\frac{x^2}{2}+x-\ln|x+3|+3\ln|x-2|+C}. \]
\end{solution}

\subsection*{Question 5b}

\begin{question}
Compute
\[ \int_1^e(\ln(x))^2\,dx \]
\end{question}

\begin{hint}
Integrate by parts twice, with $u=(\ln x)^2,\ v'=1$.
\end{hint}

\begin{solution}
Integrate by parts with $u=(\ln x)^2$, $v'=1$, i.e. $u'=\frac{2\ln x}{x}$, $v=x$:
\[ \int_1^e(\ln x)^2dx=\Big[x(\ln x)^2\Big]_1^e-\int_1^e 2\ln x\,dx=e-2\int_1^e\ln x\,dx. \]
Again by parts ($u=\ln x,\ v'=1$): $\int_1^e\ln x\,dx=\big[x\ln x\big]_1^e-\int_1^e1\,dx=e-(e-1)=1$. Hence
\[ \int_1^e(\ln x)^2dx=\boxed{e-2}\approx0.718. \]
\end{solution}

\section*{2023/24 Semester B Moed B}

\subsection*{Question 1a}

\begin{question}
Compute
\[ \lim_{x\to0}\sqrt[x]{1-2x} \]
\end{question}

\begin{hint}
Write $\sqrt[x]{1-2x}=(1-2x)^{1/x}=e^{\frac{\ln(1-2x)}{x}}$.
\end{hint}

\begin{solution}
For every $x\neq0$ close to $0$ we have $1-2x>0$, and therefore
\[ \sqrt[x]{1-2x}=(1-2x)^{1/x}=e^{\frac{\ln(1-2x)}{x}}. \]
We compute the limit of the exponent. This is a limit of the form $\frac00$, and by L'Hôpital's rule:
\[ \lim_{x\to0}\frac{\ln(1-2x)}{x}=\lim_{x\to0}\frac{\frac{-2}{1-2x}}{1}=-2. \]
By continuity of the exponential function we get
\[ \lim_{x\to0}\sqrt[x]{1-2x}=e^{-2}=\boxed{\frac1{e^2}}. \]
(This is also a form of the fundamental limit $\lim_{t\to0}(1+t)^{1/t}=e$ with $t=-2x$.)
\end{solution}

\subsection*{Question 1b}

\begin{question}
Compute
\[ \lim_{x\to0}\frac{\tan(x)-\sin(x)}{\sin^3(x)} \]
\end{question}

\begin{hint}
Factor out $\tan x$ in the numerator: $\tan x-\sin x=\tan x\,(1-\cos x)$.
\end{hint}

\begin{solution}
In a punctured neighborhood of $0$ we have $\sin x\neq0$, $\cos x\neq0$, and
\[ \tan x-\sin x=\frac{\sin x}{\cos x}-\sin x=\frac{\sin x\,(1-\cos x)}{\cos x}. \]
Hence, using $1-\cos^2x=\sin^2x$:
\[ \frac{\tan x-\sin x}{\sin^3x}=\frac{1-\cos x}{\cos x\,\sin^2x}=\frac{1-\cos x}{\cos x\,(1-\cos x)(1+\cos x)}=\frac{1}{\cos x\,(1+\cos x)}. \]
(The cancellation is allowed since $\cos x\neq1$ for $0<|x|<2\pi$.) The resulting expression is continuous at $0$, and therefore
\[ \lim_{x\to0}\frac{\tan x-\sin x}{\sin^3x}=\frac{1}{1\cdot2}=\boxed{\frac12}. \]
\end{solution}

\subsection*{Question 2}

\begin{question}
For every $a,b\in\R$ define the following function:
\[ f(x)=\begin{cases} \sin(ax) & x<0\\ b-2 & x=0\\ x^2\ln(x) & 0>x\end{cases} \]
\begin{enumerate}
  \item For which $a,b\in\R$ is the function $f(x)$ continuous at the point $x=0$?
  \item Do there exist $a,b\in\R$ for which the function $f(x)$ is differentiable at the point $x=0$?
\end{enumerate}
\end{question}

\begin{hint}
Compute the one-sided limits at $0$. For the limit $\lim_{x\to0^+}x^2\ln x$ (and also $x\ln x$) use L'Hôpital's rule after writing it as a quotient.
\end{hint}

\begin{hint}
Differentiability implies continuity. Compute the one-sided derivatives by the definition.
\end{hint}

\begin{solution}
\textbf{Remark:} the condition in the third row is written as ``$0>x$'', which is of course a typo: the first row already covers $x<0$, and for $x<0$ the expression $\ln x$ is undefined. The intended condition is $0<x$, and we solve accordingly.
\begin{enumerate}
  \item \textbf{Left limit:} $\lim_{x\to0^-}\sin(ax)=\sin0=0$ (continuity of the sine), for every $a$.

  \textbf{Right limit:} $\lim_{x\to0^+}x^2\ln x$ is of the form $0\cdot(-\infty)$. We write it as a quotient and use L'Hôpital's rule:
  \[ \lim_{x\to0^+}\frac{\ln x}{x^{-2}}=\lim_{x\to0^+}\frac{1/x}{-2x^{-3}}=\lim_{x\to0^+}\left(-\frac{x^2}{2}\right)=0. \]
  Hence the limit of $f$ at $0$ exists and equals $0$ for every $a$, and $f$ is continuous at $0$ if and only if $f(0)=b-2=0$.
  \textbf{Answer:} $f$ is continuous at $0$ for $b=2$ and any $a\in\R$.
  \item If $f$ is differentiable at $0$, then in particular it is continuous there, so we must have $b=2$, i.e. $f(0)=0$. We compute the one-sided derivatives by the definition:
  \[ f'_-(0)=\lim_{h\to0^-}\frac{\sin(ah)-0}{h}=a\lim_{h\to0^-}\frac{\sin(ah)}{ah}=a \]
  (for $a=0$ the quotient is identically $0$, and then too the limit is $0=a$), and
  \[ f'_+(0)=\lim_{h\to0^+}\frac{h^2\ln h-0}{h}=\lim_{h\to0^+}h\ln h=\lim_{h\to0^+}\frac{\ln h}{1/h}\overset{\text{L'Hôpital}}{=}\lim_{h\to0^+}\frac{1/h}{-1/h^2}=\lim_{h\to0^+}(-h)=0. \]
  $f$ is differentiable at $0$ if and only if the two one-sided derivatives are equal, i.e. $a=0$.
  \textbf{Answer:} yes: $f$ is differentiable at $0$ exactly for $a=0,\ b=2$, and then $f'(0)=0$.
\end{enumerate}
\end{solution}

\subsection*{Question 3}

\begin{question}
Analyze the following function (local extrema, intervals of increase/decrease, asymptotes):
\[ f(x)=\frac{e^x}{x^2-3} \]
Sketch its graph.
\end{question}

\begin{hint}
In the derivative, the numerator is $e^x(x^2-2x-3)=e^x(x-3)(x+1)$. Note that the extremum points must lie in the domain, and that $-\sqrt3<-1<\sqrt3<3$.
\end{hint}

\begin{solution}
\textbf{Domain:} $x^2\neq3$, i.e. $x\neq\pm\sqrt3$.

\textbf{Sign and intersections with the axes:} $e^x>0$, so there is no intersection with the $x$-axis, and $f>0$ for $|x|>\sqrt3$, $f<0$ for $|x|<\sqrt3$. Intersection with the $y$-axis: $f(0)=-\frac13$.

\textbf{Derivative:} by the quotient rule
\[ f'(x)=\frac{e^x(x^2-3)-e^x\cdot2x}{(x^2-3)^2}=\frac{e^x(x^2-2x-3)}{(x^2-3)^2}=\frac{e^x(x-3)(x+1)}{(x^2-3)^2}. \]
The denominator and $e^x$ are positive, so the sign of $f'$ is the sign of $(x-3)(x+1)$. We have $-\sqrt3\approx-1.73<-1<\sqrt3\approx1.73<3$:
\begin{itemize}
  \item $(-\infty,-\sqrt3)$ and $(-\sqrt3,-1)$: $f'>0$, $f$ is \textbf{increasing}.
  \item $(-1,\sqrt3)$ and $(\sqrt3,3)$: $f'<0$, $f$ is \textbf{decreasing}.
  \item $(3,\infty)$: $f'>0$, $f$ is \textbf{increasing}.
\end{itemize}

\textbf{Extremum points:} $f'$ changes sign from $+$ to $-$ at $x=-1$, so this is a \textbf{local maximum} point: $f(-1)=\frac{e^{-1}}{-2}=-\frac1{2e}\approx-0.18$. At $x=3$ the sign changes from $-$ to $+$, so it is a \textbf{local minimum}: $f(3)=\frac{e^3}{6}\approx3.35$.

\textbf{Asymptotes:}
\begin{itemize}
  \item Vertical: $x=\pm\sqrt3$ (the numerator $e^{\pm\sqrt3}\neq0$ and the denominator vanishes):
  $\lim_{x\to-\sqrt3^-}f=+\infty$, $\lim_{x\to-\sqrt3^+}f=-\infty$, $\lim_{x\to\sqrt3^-}f=-\infty$, $\lim_{x\to\sqrt3^+}f=+\infty$ (according to the sign of $x^2-3$ on each side).
  \item At $-\infty$: $e^x\to0$ and $x^2-3\to\infty$, so $f\to0$: horizontal asymptote $y=0$ (the graph lies above it).
  \item At $+\infty$: by L'Hôpital's rule (twice) $\lim_{x\to\infty}\frac{e^x}{x^2-3}=\lim\frac{e^x}{2}=\infty$, and also $\frac{f(x)}{x}=\frac{e^x}{x^3-3x}\to\infty$; hence there is no horizontal or oblique asymptote at $+\infty$.
\end{itemize}

\textbf{The graph:} on the left, the graph starts above the $x$-axis close to $y=0$ and rises to $+\infty$ as $x\to-\sqrt3^-$. Between $-\sqrt3$ and $\sqrt3$ the graph is negative: it rises from $-\infty$ up to the local maximum $(-1,-\frac1{2e})$, then descends through $(0,-\frac13)$ to $-\infty$ as $x\to\sqrt3^-$. To the right of $\sqrt3$ the graph descends from $+\infty$ to the local minimum $(3,\frac{e^3}{6})$ and then rises to $+\infty$.
\end{solution}

\subsection*{Question 4a}

\begin{question}
Write the Maclaurin polynomial of order 4 of $f(x)=\sin(x)\cos(x)$
\end{question}

\begin{hint}
$\sin x\cos x=\frac12\sin(2x)$.
\end{hint}

\begin{solution}
By the double-angle identity $f(x)=\sin x\cos x=\frac12\sin(2x)$. The Maclaurin polynomial of order 4 of $\sin t$ is $t-\frac{t^3}{6}$ (the coefficient of $t^4$ is $0$). Substitute $t=2x$:
\[ P_4(x)=\frac12\left(2x-\frac{(2x)^3}{6}\right)=\boxed{x-\frac{2}{3}x^3}. \]
(Directly: $f'(x)=\cos2x$, $f''=-2\sin2x$, $f'''=-4\cos2x$, $f^{(4)}=8\sin2x$; at the point $0$: $f=0,\ f'=1,\ f''=0,\ f'''=-4,\ f^{(4)}=0$, and therefore $P_4=x-\frac{4}{6}x^3$.)
\end{solution}

\subsection*{Question 4b}

\begin{question}
Find a linear approximation of $5^{\frac{2}{3}}$
\end{question}

\begin{hint}
Choose $g(x)=x^{2/3}$ and a point close to $5$ at which the value is known.
\end{hint}

\begin{solution}
We use $g(x)=x^{2/3}$ around $a=8$ (where $8^{2/3}=4$). $g'(x)=\frac23x^{-1/3}$, and therefore $g'(8)=\frac23\cdot\frac12=\frac13$. The linear approximation:
\[ g(x)\approx g(8)+g'(8)(x-8)=4+\frac{x-8}{3}. \]
For $x=5$: $5^{2/3}\approx4+\frac{-3}{3}=\boxed{3}$.
(The true value is $\approx2.924$; the error is relatively large because $5$ is far from $8$. Another choice, e.g. $a=1$, gives $1+\frac23\cdot4\approx3.67$, which is less accurate.)
\end{solution}

\subsection*{Question 5a}

\begin{question}
Compute
\[ \int\frac{2x^3+7x^2-5x-19}{x^2+x-6}\,dx \]
\end{question}

\begin{hint}
Polynomial long division, then partial fractions, $x^2+x-6=(x+3)(x-2)$.
\end{hint}

\begin{solution}
Polynomial long division:
\[ 2x^3+7x^2-5x-19=(2x+5)(x^2+x-6)+(2x+11), \]
(check: $(2x+5)(x^2+x-6)=2x^3+7x^2-7x-30$, and the remainder is $(-5x-19)-(-7x-30)=2x+11$.) Partial fractions: $\frac{2x+11}{(x+3)(x-2)}=\frac{A}{x+3}+\frac{B}{x-2}$, i.e. $2x+11=A(x-2)+B(x+3)$. For $x=2$: $15=5B$, $B=3$. For $x=-3$: $5=-5A$, $A=-1$. Hence
\[ \int\frac{2x^3+7x^2-5x-19}{x^2+x-6}dx=\int\left(2x+5-\frac1{x+3}+\frac3{x-2}\right)dx=\boxed{x^2+5x-\ln|x+3|+3\ln|x-2|+C}. \]
\end{solution}

\subsection*{Question 5b}

\begin{question}
Compute
\[ \int_0^{\pi}e^x\sin(x)\,dx \]
\end{question}

\begin{hint}
Integrate by parts twice; the original integral will appear again on the right-hand side, and one solves an equation for it.
\end{hint}

\begin{solution}
Denote $I=\int_0^\pi e^x\sin x\,dx$. By parts with $u=\sin x$, $v'=e^x$:
\[ I=\big[e^x\sin x\big]_0^\pi-\int_0^\pi e^x\cos x\,dx=0-\int_0^\pi e^x\cos x\,dx. \]
Again by parts with $u=\cos x$, $v'=e^x$:
\[ \int_0^\pi e^x\cos x\,dx=\big[e^x\cos x\big]_0^\pi+\int_0^\pi e^x\sin x\,dx=(-e^\pi-1)+I. \]
Hence $I=e^\pi+1-I$, i.e.
\[ I=\boxed{\frac{e^\pi+1}{2}}\approx12.07. \]
\end{solution}

\section*{2023/24 Semester B Moed C}

\subsection*{Question 1a}

\begin{question}
Compute
\[ \lim_{x\longrightarrow\infty}\left(\frac{x+3}{x-4}\right)^x \]
\end{question}

\begin{hint}
Write $\frac{x+3}{x-4}=1+\frac{7}{x-4}$ and use the limit $\lim_{t\to\infty}\left(1+\frac1t\right)^t=e$.
\end{hint}

\begin{solution}
For every $x>4$: $\frac{x+3}{x-4}=1+\frac{7}{x-4}$. Set $t=\frac{x-4}{7}\to\infty$; then $x=7t+4$ and
\[ \left(\frac{x+3}{x-4}\right)^x=\left(1+\frac1t\right)^{7t+4}=\left[\left(1+\frac1t\right)^{t}\right]^7\cdot\left(1+\frac1t\right)^4\xrightarrow[t\to\infty]{}e^7\cdot1. \]
(Alternatively: $x\ln\frac{x+3}{x-4}=\frac{\ln(x+3)-\ln(x-4)}{1/x}$, and by L'Hôpital's rule the limit is $\lim\frac{\frac1{x+3}-\frac1{x-4}}{-1/x^2}=\lim\frac{7x^2}{(x+3)(x-4)}=7$.)
The answer: $\boxed{e^7}$.
\end{solution}

\subsection*{Question 1b}

\begin{question}
Compute
\[ \lim_{x\longrightarrow-8}\frac{\sqrt{1+x}-3}{2+\sqrt[3]{x}} \]
\end{question}

\begin{hint}
First check the domain of $\sqrt{1+x}$ near $x=-8$.
\end{hint}

\begin{solution}
\textbf{As the question is printed}, the expression $\sqrt{1+x}$ is defined (over the reals) only for $x\ge-1$, so the function is not defined in any punctured neighborhood of $x=-8$, and the limit is \textbf{undefined}. (Substituting $x=-8$ also gives $\sqrt{-7}$.)

Most likely there is a typo, and the intended expression is $\frac{\sqrt{1-x}-3}{2+\sqrt[3]{x}}$: then at the point $x=-8$ the numerator is $\sqrt9-3=0$ and the denominator is $2+(-2)=0$, a limit of the form $\frac00$. We solve this limit. Multiply by ``conjugates'': in the numerator $\sqrt{1-x}-3=\frac{(1-x)-9}{\sqrt{1-x}+3}=\frac{-(x+8)}{\sqrt{1-x}+3}$, and in the denominator, by $a^3+b^3=(a+b)(a^2-ab+b^2)$ with $a=2,\ b=\sqrt[3]x$:
\[ 2+\sqrt[3]{x}=\frac{8+x}{4-2\sqrt[3]{x}+\sqrt[3]{x^2}}. \]
Hence for $x\neq-8$:
\[ \frac{\sqrt{1-x}-3}{2+\sqrt[3]{x}}=-\frac{4-2\sqrt[3]{x}+\sqrt[3]{x^2}}{\sqrt{1-x}+3}\xrightarrow[x\to-8]{}-\frac{4+4+4}{3+3}=-2, \]
by continuity of the last expression at $x=-8$. (The same result follows from L'Hôpital's rule: $\frac{-\frac{1}{2\sqrt{1-x}}}{\frac{1}{3\sqrt[3]{x^2}}}\to\frac{-1/6}{1/12}=-2$.)
The answer: as printed, the limit is undefined; in the corrected version with $\sqrt{1-x}$, the limit is $\boxed{-2}$.
\end{solution}

\subsection*{Question 2}

\begin{question}
For every $a,b\in\R$ define the following function:
\[ f(x)=\begin{cases} ax^2+bx+1 & x<0\\ b-2 & x=0\\ bx^2+ax+1 & 0<x\end{cases} \]
\begin{enumerate}
  \item For which $a,b\in\R$ is the function $f(x)$ continuous at the point $x=0$?
  \item Do there exist $a,b\in\R$ for which the function $f(x)$ is differentiable at the point $x=0$?
\end{enumerate}
\end{question}

\begin{hint}
Differentiability implies continuity, so the condition from part (a) must hold. Compute the one-sided derivatives by the definition.
\end{hint}

\begin{solution}
\begin{enumerate}
  \item Polynomials are continuous, so
  \[ \lim_{x\to0^-}f(x)=a\cdot0+b\cdot0+1=1,\qquad \lim_{x\to0^+}f(x)=1. \]
  The limit at $0$ exists and equals $1$ for all $a,b$, and $f$ is continuous at $0$ if and only if $f(0)=b-2=1$.
  \textbf{Answer:} $b=3$ and any $a\in\R$.
  \item If $f$ is differentiable at $0$ then it is continuous there, so $b=3$ and $f(0)=1$. We compute the one-sided derivatives by the definition:
  \[ f'_-(0)=\lim_{h\to0^-}\frac{ah^2+bh+1-1}{h}=\lim_{h\to0^-}(ah+b)=b, \qquad f'_+(0)=\lim_{h\to0^+}\frac{bh^2+ah+1-1}{h}=\lim_{h\to0^+}(bh+a)=a. \]
  $f$ is differentiable at $0$ if and only if $f'_-(0)=f'_+(0)$, i.e. $a=b$. Together with $b=3$:
  \textbf{Answer:} yes, exactly for $a=b=3$, and then $f'(0)=3$.
\end{enumerate}
\end{solution}

\subsection*{Question 3}

\begin{question}
Analyze the following function (local extrema, intervals of increase/decrease, asymptotes):
\[ f(x)=e^{-(x^3-3x)} \]
Sketch its graph.
\end{question}

\begin{hint}
$f=e^{g}$ with $g(x)=3x-x^3$; since $e^t$ is increasing, $f$ increases/decreases exactly where $g$ increases/decreases.
\end{hint}

\begin{solution}
\textbf{Domain:} all of $\R$; $f$ is continuous and differentiable everywhere, and $f(x)>0$ for every $x$ (no intersection with the $x$-axis). Intersection with the $y$-axis: $f(0)=e^0=1$.

\textbf{Derivative:} by the chain rule
\[ f'(x)=e^{-(x^3-3x)}\cdot(-(3x^2-3))=3(1-x^2)\,e^{-(x^3-3x)}=3(1-x)(1+x)\,e^{3x-x^3}. \]
The exponential is positive, so the sign of $f'$ is the sign of $1-x^2$:
\begin{itemize}
  \item $(-\infty,-1)$: $f'<0$, $f$ is \textbf{decreasing}.
  \item $(-1,1)$: $f'>0$, $f$ is \textbf{increasing}.
  \item $(1,\infty)$: $f'<0$, $f$ is \textbf{decreasing}.
\end{itemize}

\textbf{Extremum points:} at $x=-1$ the derivative changes from $-$ to $+$: \textbf{local minimum} $f(-1)=e^{-(-1+3)}=e^{-2}\approx0.135$. At $x=1$ the derivative changes from $+$ to $-$: \textbf{local maximum} $f(1)=e^{-(1-3)}=e^{2}\approx7.39$.

\textbf{Asymptotes:}
\begin{itemize}
  \item Vertical: none, $f$ is continuous on the whole line.
  \item $x\to+\infty$: $3x-x^3=-x^3\left(1-\frac{3}{x^2}\right)\to-\infty$, so $f(x)\to0$: horizontal asymptote $y=0$ (the graph lies above it).
  \item $x\to-\infty$: $3x-x^3\to+\infty$, so $f(x)\to+\infty$; also $\frac{f(x)}{x}\to-\infty$ (numerator $\to+\infty$, denominator $\to-\infty$), and therefore there is no horizontal or oblique asymptote at $-\infty$.
\end{itemize}

\textbf{The graph:} it comes from $+\infty$ on the left, decreases to the local minimum $(-1,e^{-2})$, increases through $(0,1)$ to the local maximum $(1,e^2)$, and then decreases and tends to $0$ (above the $x$-axis) as $x\to\infty$. The whole graph lies above the $x$-axis.
\end{solution}

\subsection*{Question 4a}

\begin{question}
Write the Maclaurin polynomial of order 3 of $f(x)=\sin(\cos(x))$
\end{question}

\begin{hint}
$f$ is even, so the coefficients of $x$ and of $x^3$ vanish; it remains to compute $f(0)$ and $f''(0)$.
\end{hint}

\begin{solution}
We compute derivatives by the chain rule:
\[ f'(x)=\cos(\cos x)\cdot(-\sin x)=-\sin x\cos(\cos x), \]
\[ f''(x)=-\cos x\cos(\cos x)-\sin x\cdot\sin(\cos x)\cdot\sin x=-\cos x\cos(\cos x)-\sin^2x\,\sin(\cos x). \]
At the point $0$: $f(0)=\sin1$, $f'(0)=0$, $f''(0)=-\cos1$. Since $f(-x)=\sin(\cos(-x))=f(x)$, the function is even, so all its derivatives of odd order vanish at $0$ (the derivative of an even function is odd), and in particular $f'''(0)=0$. Hence
\[ P_3(x)=f(0)+f'(0)x+\frac{f''(0)}{2}x^2+\frac{f'''(0)}{6}x^3=\boxed{\sin1-\frac{\cos1}{2}x^2}. \]
(Approximately: $0.8415-0.2702x^2$.)
\end{solution}

\subsection*{Question 4b}

\begin{question}
Find a linear approximation of $\sqrt[3]{7^2}$ using the function $f(x)=\sqrt[3]{x^2}$
\end{question}

\begin{hint}
Use the tangent line to the graph of $f$ at a point close to $7$ where the value is known.
\end{hint}

\begin{solution}
$f(x)=\sqrt[3]{x^2}=x^{2/3}$ around $a=8$, where $f(8)=4$. $f'(x)=\frac23x^{-1/3}$, $f'(8)=\frac23\cdot\frac12=\frac13$. The linear approximation:
\[ f(x)\approx4+\frac13(x-8), \qquad \sqrt[3]{7^2}=f(7)\approx4-\frac13=\boxed{\frac{11}{3}\approx3.667}. \]
(The true value is $\sqrt[3]{49}\approx3.659$.)
\end{solution}

\subsection*{Question 5a}

\begin{question}
Compute
\[ \int\frac{x^3}{x^2+x-2}\,dx \]
\end{question}

\begin{hint}
Polynomial long division and partial fractions, $x^2+x-2=(x+2)(x-1)$.
\end{hint}

\begin{solution}
Polynomial long division:
\[ x^3=(x-1)(x^2+x-2)+(3x-2), \]
(check: $(x-1)(x^2+x-2)=x^3-3x+2$, and the remainder is $x^3-(x^3-3x+2)=3x-2$.) Partial fractions: $\frac{3x-2}{(x+2)(x-1)}=\frac{A}{x+2}+\frac{B}{x-1}$, i.e. $3x-2=A(x-1)+B(x+2)$. For $x=1$: $1=3B$, $B=\frac13$. For $x=-2$: $-8=-3A$, $A=\frac83$. Hence
\[ \int\frac{x^3}{x^2+x-2}dx=\int\left(x-1+\frac{8/3}{x+2}+\frac{1/3}{x-1}\right)dx=\boxed{\frac{x^2}{2}-x+\frac83\ln|x+2|+\frac13\ln|x-1|+C}. \]
\end{solution}

\subsection*{Question 5b}

\begin{question}
Compute
\[ \int_0^{\ln2}xe^{-x}\,dx \]
\end{question}

\begin{hint}
Integrate by parts with $u=x$.
\end{hint}

\begin{solution}
Integrate by parts with $u=x$, $v'=e^{-x}$, i.e. $u'=1$, $v=-e^{-x}$:
\[ \int_0^{\ln2}xe^{-x}dx=\Big[-xe^{-x}\Big]_0^{\ln2}+\int_0^{\ln2}e^{-x}dx=-\frac{\ln2}{2}+\Big[-e^{-x}\Big]_0^{\ln2}=-\frac{\ln2}{2}+\left(1-\frac12\right), \]
since $e^{-\ln2}=\frac12$. Hence
\[ \int_0^{\ln2}xe^{-x}dx=\boxed{\frac{1-\ln2}{2}}\approx0.153. \]
\end{solution}

\section*{2024/25 Semester A Sample quiz}

\subsection*{Part A, Question 1}

\begin{question}
Compute the following limit:
\[
\lim_{x\to\infty}\left(\frac{x^2+x+2}{x^2+x+1}\right)^{2x^2+4}
\]
\end{question}

\begin{hint}
This is a limit of the form $1^\infty$. Write the base in the form $1+\frac{1}{x^2+x+1}$.
\end{hint}

\begin{solution}
The base tends to 1 and the exponent to $\infty$: a limit of the form $1^\infty$. Write the base as
\[
\frac{x^2+x+2}{x^2+x+1}=\frac{x^2+x+1+1}{x^2+x+1}=1+\frac{1}{x^2+x+1},
\]
and multiply and divide the exponent by $x^2+x+1$:
\[
\left(1+\frac{1}{x^2+x+1}\right)^{2x^2+4}
=\left[\left(1+\frac{1}{x^2+x+1}\right)^{x^2+x+1}\right]^{\frac{2x^2+4}{x^2+x+1}} .
\]
\textbf{The inner base:} substitute $y=x^2+x+1$; as $x\to\infty$ also $y\to\infty$, and by the fundamental limit
\[
\lim_{y\to\infty}\left(1+\frac1y\right)^y=e
\]
(and the theorem on the limit of a composition) the inner base tends to $e$.

\textbf{The exponent:} dividing by $x^2$ gives
\[
\lim_{x\to\infty}\frac{2x^2+4}{x^2+x+1}=\lim_{x\to\infty}\frac{2+\frac4{x^2}}{1+\frac1x+\frac1{x^2}}=2 .
\]
\textbf{Summary:} if $u(x)\to e>0$ and $v(x)\to2$ then $u^v=e^{v\ln u}\to e^{2\ln e}=e^2$, by continuity of $\ln$ and $\exp$ (limit of a quotient/product and continuity). Hence
\[
\lim_{x\to\infty}\left(\frac{x^2+x+2}{x^2+x+1}\right)^{2x^2+4}=\boxed{e^2}.
\]
\end{solution}

\subsection*{Part A, Question 2}

\begin{question}
Let $a,b\in\R$. Define a function by
\[
f(x)=\begin{cases}
\frac{\sin(8x)}{x} & x<0\\
ax+b & 0\le x\le 4\\
\frac{x^2-3x-4}{2-\sqrt{x}} & x>4
\end{cases}
\]
For which values of $a,b$ is the function continuous on $\R$ ? Justify.
\end{question}

\begin{hint}
At every point $x\neq0,4$ the function is continuous regardless of $a,b$. Check only the junction points $x=0$ and $x=4$.
\end{hint}

\begin{hint}
At $x=4$: $x^2-3x-4=(x-4)(x+1)$; multiply and divide by $2+\sqrt x$.
\end{hint}

\begin{solution}
\textbf{Points that are not junction points.} At every point $x<0$ the function is a quotient of continuous functions with a nonzero denominator; at every point $0<x<4$ it is a polynomial; and at every point $x>4$ it is a quotient of continuous functions with denominator $2-\sqrt x\neq0$ (since $\sqrt x>2$). Since at each of these points $f$ coincides with the formula on a whole neighborhood, $f$ is continuous there for all $a,b$. It remains to check $x=0$ and $x=4$.

\textbf{The point $x=0$.} $f(0)=b$. From the left, by $\lim_{y\to0}\frac{\sin y}{y}=1$ with $y=8x$:
\[
\lim_{x\to0^-}\frac{\sin(8x)}{x}=8\lim_{x\to0^-}\frac{\sin(8x)}{8x}=8\lim_{y\to0^-}\frac{\sin y}{y}=8\cdot1=8 .
\]
From the right: $\lim_{x\to0^+}(ax+b)=b=f(0)$ (continuity of a polynomial). Hence $f$ is continuous at 0 if and only if $\boxed{b=8}$.

\textbf{The point $x=4$.} $f(4)=4a+b=4a+8$, and from the left $\lim_{x\to4^-}(ax+8)=4a+8=f(4)$. From the right, a limit of the form $\frac00$. Multiply by the conjugate:
\[
\frac{x^2-3x-4}{2-\sqrt x}\cdot\frac{2+\sqrt x}{2+\sqrt x}=\frac{(x-4)(x+1)(2+\sqrt x)}{4-x}=-(x+1)(2+\sqrt x)\qquad(x\neq4),
\]
and therefore, by continuity of the last expression at 4,
\[
\lim_{x\to4^+}f(x)=-(4+1)(2+\sqrt4)=-5\cdot4=-20 .
\]
$f$ is continuous at 4 if and only if $4a+8=-20$, i.e. $\boxed{a=-7}$.

\textbf{Answer:} $f$ is continuous on $\R$ if and only if $a=-7,\ b=8$.

(Remark: in the attached handwritten solution the ``$-$'' sign was omitted in the line $-\lim(x+1)(2+\sqrt x)=5\cdot4=20$, but the next line correctly reads $4a+8=-20$ and the answer $a=-7$ is correct.)
\end{solution}

\subsection*{Part A, Question 3}

\begin{question}
Prove that the following equation has a real root and find an interval of length at most 1 that contains the root:
\[
x^2+\ln(x)=0
\]
\end{question}

\begin{hint}
Define $f(x)=x^2+\ln x$ on $(0,\infty)$ and look for points where $\ln$ is easy to compute, such as $x=1$ and $x=\frac1e$.
\end{hint}

\begin{solution}
Define $f(x)=x^2+\ln x$ on $(0,\infty)$. $f$ is continuous there as the sum of a polynomial and the logarithm function, which is continuous on its domain.

We compute:
\[
f\!\left(\tfrac1e\right)=\frac1{e^2}+\ln\frac1e=\frac{1}{e^2}-1<0,\qquad f(1)=1^2+\ln1=1>0 ,
\]
since $e^2>1$.

$f$ is continuous on the closed interval $\left[\frac1e,1\right]\subset(0,\infty)$ and changes sign at its endpoints. By the Intermediate Value Theorem (Bolzano) there exists $c\in\left(\frac1e,1\right)$ with $f(c)=0$, i.e. a real root of the equation.

\textbf{The required interval:} $\left[\frac1e,1\right]$, whose length is $1-\frac1e<1$.

(Remark: the root is unique, since $f$ is strictly increasing as the sum of $x^2$ and $\ln x$, both strictly increasing on $(0,\infty)$. Numerically $c\approx0.653$. One can also choose an interval that does not require $e$, for example $\left[\frac12,1\right]$: $f(\frac12)=\frac14-\ln2<0$.)
\end{solution}

\subsection*{Part B, Question 4}

\begin{question}
Let $f(x)$ be an even function. Which of the following statements is true?
\begin{enumerate}
  \item The function $g(x)=x^2f(x)-\frac{1}{|x|+1}$ is an even function.
  \item The function $g(x)=x^3f(x)-x|x|$ is an even function.
  \item The function $g(x)=x^3f(x)+1$ is an even function.
  \item The function $g(x)=\frac{f(x)}{x^2+x}$ is an odd function.
\end{enumerate}
\end{question}

\begin{hint}
A product of an even function and an even function is even; a product of an odd function and an even function is odd. To disprove, one example suffices, e.g. $f\equiv1$.
\end{hint}

\begin{solution}
\textbf{The correct answer: (a)}.

\textbf{(a) is true:} the domain of $g$ is the domain of $f$, which is symmetric (since $f$ is even). For every $x$ in the domain, since $f(-x)=f(x)$, $(-x)^2=x^2$ and $|-x|=|x|$:
\[
g(-x)=(-x)^2f(-x)-\frac{1}{|-x|+1}=x^2f(x)-\frac{1}{|x|+1}=g(x).
\]

\textbf{Why the others are wrong} (the statements must hold for every even $f$, so a single counterexample suffices; take $f(x)\equiv1$, which is even):
\begin{itemize}
  \item (b): $g(x)=x^3-x|x|$: $g(2)=8-4=4$ whereas $g(-2)=-8+4=-4\neq g(2)$. (In general this $g$ is odd, as the difference of two odd functions.)
  \item (c): $g(x)=x^3+1$: $g(1)=2$ whereas $g(-1)=0$.
  \item (d): $g(x)=\frac{1}{x^2+x}=\frac{1}{x(x+1)}$ is defined at $x=1$ but not at $x=-1$, so the domain is not symmetric and $g$ is not odd. (Even without the domain argument: $g(2)=\frac16$ and $g(-2)=\frac{1}{4-2}=\frac12\neq-\frac16=-g(2)$.)
\end{itemize}
\end{solution}

\subsection*{Part B, Question 5}

\begin{question}
Let
\[
f(x)=\begin{cases}
3x-1 & x<0\\
x^2 & x>0
\end{cases}
\]
Which of the following statements is true?
\begin{enumerate}
  \item The function is not monotone.
  \item The function is strictly decreasing.
  \item The function is strictly increasing.
  \item The function is monotone but not strictly.
\end{enumerate}
\end{question}

\begin{hint}
Check each piece separately, then compare some value to the left of 0 with some value to the right of 0.
\end{hint}

\begin{solution}
\textbf{The correct answer: (c)}.

(Note: as the function is written, it is not defined at $x=0$; the attached solution writes $x\le0$ in the first piece. The answer is the same in both cases.)

Let $x_1<x_2$ be in the domain. We distinguish three cases:
\begin{itemize}
  \item $x_1<x_2<0$: $f(x_2)-f(x_1)=3(x_2-x_1)>0$.
  \item $0<x_1<x_2$: $f(x_2)-f(x_1)=x_2^2-x_1^2=(x_2-x_1)(x_2+x_1)>0$, since both factors are positive.
  \item $x_1<0<x_2$: $f(x_1)=3x_1-1<-1<0<x_2^2=f(x_2)$.
\end{itemize}
(If one defines $f(0)=3\cdot0-1=-1$, the cases $x_1<0=x_2$ and $x_1=0<x_2$ also satisfy $f(x_1)<f(x_2)$ in the same way.)
In all cases $f(x_1)<f(x_2)$, and therefore $f$ is strictly increasing.

\textbf{Why the others are wrong:} (a): we showed that $f$ is monotone. (b): for example $f(1)=1>f(-1)=-4$, so it is not decreasing. (d): ``monotone but not strictly'' means that there exist $x_1<x_2$ with $f(x_1)=f(x_2)$, and this does not happen since $f$ is strictly increasing.
\end{solution}

\subsection*{Part B, Question 6}

\begin{question}
Let $f(x)$ be the function defined in the previous question. Which of the following statements is true?
\begin{enumerate}
  \item The function is one-to-one and onto $\R$.
  \item The function is one-to-one but not onto $\R$.
  \item The function is onto $\R$ but not one-to-one.
  \item The function is neither one-to-one nor onto $\R$.
\end{enumerate}
\end{question}

\begin{solution}
\textbf{The correct answer: (b)}.

(The function from the previous question: $f(x)=3x-1$ for $x<0$ and $f(x)=x^2$ for $x>0$.)

\textbf{One-to-one:} in the previous question we showed that $f$ is strictly increasing, and a strictly increasing function is one-to-one: if $x_1\neq x_2$, say $x_1<x_2$, then $f(x_1)<f(x_2)$ and in particular $f(x_1)\neq f(x_2)$.

\textbf{Not onto $\R$:} for $x<0$ we have $f(x)=3x-1<-1$, and for $x>0$ we have $f(x)=x^2>0$. Hence no value in the interval $[-1,0]$ is attained; for example there is no $x$ with $f(x)=-\frac12$. (The image is $(-\infty,-1)\cup(0,\infty)$; if one defines $f(0)=-1$ as in the attached solution, the image is $(-\infty,-1]\cup(0,\infty)$, and $-\frac12$ is still not attained.)

\textbf{Why the others are wrong:} (a) and (c) claim that $f$ is onto $\R$, which we disproved. (c) and (d) claim that $f$ is not one-to-one, which we disproved.
\end{solution}

\subsection*{Part B, Question 7}

\begin{question}
Check the continuity of the following function at the point 0:
\[
f(x)=\begin{cases}
\sin(x)\cos(\frac{1}{x}) & x<0\\
\frac{1}{1+x^2} & x\ge 0
\end{cases}
\]
Which of the following statements is true?
\begin{enumerate}
  \item The function is continuous at 0.
  \item The function has a removable discontinuity at 0.
  \item The function has a jump discontinuity at 0.
  \item The function has an essential discontinuity at 0.
\end{enumerate}
\end{question}

\begin{hint}
From the left: a function tending to 0 times a bounded function.
\end{hint}

\begin{solution}
\textbf{The correct answer: (c)}.

\textbf{From the left:} $\lim_{x\to0^-}\sin x=0$ (continuity of $\sin$) and $\left|\cos\frac1x\right|\le1$ for every $x\neq0$. By the theorem ``a function tending to 0 times a bounded function tends to 0'' (or by the squeeze $-|\sin x|\le\sin x\cos\frac1x\le|\sin x|$):
\[
\lim_{x\to0^-}f(x)=\lim_{x\to0^-}\sin(x)\cos\left(\tfrac1x\right)=0 .
\]
\textbf{From the right:} $\frac{1}{1+x^2}$ is continuous, so $\lim_{x\to0^+}f(x)=\frac{1}{1+0}=1=f(0)$.

Both one-sided limits exist and are finite but different ($0\neq1$), so there is a jump discontinuity at 0.

\textbf{Why the others are wrong:} (a) and (b): both require the two-sided limit at 0 to exist, and it does not exist. (d): although $\cos\frac1x$ alone has no limit at 0, the product with $\sin x$ does tend to 0, so both one-sided limits exist and are finite; an essential discontinuity requires one of them not to exist as a finite limit.
\end{solution}

\section*{2024/25 Semester A Quiz}

\subsection*{Part A, Question 1}

\begin{question}
Compute the following limit:
\[
\lim_{x\to\infty}\left(\sqrt{x^2+5x}-\sqrt{x^2-5x}\right)
\]
\end{question}

\begin{hint}
This is a limit of the form $\infty-\infty$. Multiply and divide by the conjugate $\sqrt{x^2+5x}+\sqrt{x^2-5x}$.
\end{hint}

\begin{solution}
For $x>5$ both roots are defined and positive. This is a limit of the form $\infty-\infty$, so we multiply and divide by the (positive) conjugate expression:
\[
\sqrt{x^2+5x}-\sqrt{x^2-5x}
=\frac{(x^2+5x)-(x^2-5x)}{\sqrt{x^2+5x}+\sqrt{x^2-5x}}
=\frac{10x}{\sqrt{x^2+5x}+\sqrt{x^2-5x}} .
\]
Divide the numerator and denominator by $x$. Since $x>0$ we have $x=\sqrt{x^2}$, and therefore
\[
=\frac{10}{\sqrt{1+\frac5x}+\sqrt{1-\frac5x}} .
\]
As $x\to\infty$: $\frac5x\to 0$, and by continuity of the square root function at 1 we get $\sqrt{1\pm\frac5x}\to 1$. By the arithmetic of limits (the denominator tends to $2\neq 0$):
\[
\lim_{x\to\infty}\left(\sqrt{x^2+5x}-\sqrt{x^2-5x}\right)=\frac{10}{1+1}=\boxed{5}.
\]
\end{solution}

\subsection*{Part A, Question 2}

\begin{question}
Let $a,b\in\R$. Define a function by
\[
f(x)=\begin{cases}
\frac{x\sin(ax)}{1-\cos x} & x<0\\
b & x=0\\
\frac{e^x-1}{xe^x} & x>0
\end{cases}
\]
For which values of $a,b$ is the function continuous at $0$ ? Justify.
\end{question}

\begin{hint}
$f$ is continuous at 0 if and only if both one-sided limits at 0 exist and are equal to $f(0)=b$.
\end{hint}

\begin{hint}
Use the fundamental limits $\frac{\sin t}{t}\to1$, $\frac{1-\cos x}{x^2}\to\frac12$, $\frac{e^x-1}{x}\to1$ (L'Hôpital's rule is not allowed in the quiz).
\end{hint}

\begin{solution}
The function is continuous at 0 if and only if
\[
\lim_{x\to0^-}f(x)=\lim_{x\to0^+}f(x)=f(0)=b .
\]

\textbf{The left limit.} We use the fundamental limit
\[
\lim_{x\to0}\frac{1-\cos x}{x^2}=\frac12 ,
\]
which follows from the identity $1-\cos x=2\sin^2\frac x2$: $\frac{1-\cos x}{x^2}=\frac12\left(\frac{\sin(x/2)}{x/2}\right)^2\to\frac12$.
If $a\neq0$ we write, for $x<0$ close to 0:
\[
\frac{x\sin(ax)}{1-\cos x}=a\cdot\frac{\sin(ax)}{ax}\cdot\frac{x^2}{1-\cos x}\xrightarrow[x\to0^-]{}a\cdot1\cdot2=2a ,
\]
by $\lim_{t\to0}\frac{\sin t}{t}=1$ (substitution $t=ax\to0$) and the arithmetic of limits.
If $a=0$ then $f(x)=0$ for every $x<0$ and the left limit is $0=2a$. In any case
\[
\lim_{x\to0^-}f(x)=2a .
\]

\textbf{The right limit.}
\[
\frac{e^x-1}{xe^x}=\frac{e^x-1}{x}\cdot\frac{1}{e^x}\xrightarrow[x\to0^+]{}1\cdot\frac{1}{e^0}=1 ,
\]
by the fundamental limit $\lim_{x\to0}\frac{e^x-1}{x}=\ln e=1$ (from the formula sheet) and continuity of $e^x$ at 0.

\textbf{Conclusion.} $f$ is continuous at 0 if and only if $2a=1=b$, i.e.
\[
\boxed{a=\tfrac12,\quad b=1.}
\]
\end{solution}

\subsection*{Part A, Question 3}

\begin{question}
Compute the inverse function of the function $f(x)=\ln(1+\sqrt{x})$ and write down the domain of the inverse function.
\end{question}

\begin{hint}
The domain of the inverse function is the image of $f$. Find the image, then solve the equation $y=\ln(1+\sqrt x)$ for $x$.
\end{hint}

\begin{solution}
\textbf{The domain and image of $f$.} $\sqrt x$ is defined for $x\ge0$, and then $1+\sqrt x\ge1>0$, so $\ln(1+\sqrt x)$ is defined. Hence $D_f=[0,\infty)$.

$f$ is strictly increasing: $\sqrt x$ is strictly increasing on $[0,\infty)$ and $\ln$ is strictly increasing, and a composition of strictly increasing functions is strictly increasing. In particular $f$ is one-to-one and therefore invertible on its image.

The image: $f(0)=\ln 1=0$, so $f(x)\ge0$ for every $x\ge0$. For every $y\ge0$ we will find below an $x\ge0$ with $f(x)=y$, and therefore $\operatorname{Im}f=[0,\infty)$.

\textbf{Computing the inverse.} Let $y\ge0$. Solve $y=\ln(1+\sqrt x)$:
\[
e^y=1+\sqrt x\iff \sqrt x=e^y-1 .
\]
Since $y\ge0$ we have $e^y-1\ge0$, so there is a unique solution $x=(e^y-1)^2\ge0$. (For $y<0$ we would get $\sqrt x<0$, so there is no solution, as expected since $y$ is not in the image.)

Hence
\[
\boxed{f^{-1}(x)=\left(e^x-1\right)^2,\qquad D_{f^{-1}}=[0,\infty).}
\]
\textbf{Check:} for $x\ge0$: $f(f^{-1}(x))=\ln\left(1+\sqrt{(e^x-1)^2}\right)=\ln\left(1+|e^x-1|\right)=\ln(e^x)=x$, since $e^x-1\ge0$. Note that restricting the domain to $[0,\infty)$ is necessary: the formula $(e^x-1)^2$ is defined for every $x$, but for $x<0$ it is not the inverse of $f$.
\end{solution}

\subsection*{Part B, Question 4}

\begin{question}
Let
\[
f(x)=e^x,\qquad g(x)=\cos(2x)
\]
Which of the following statements is true?
\begin{enumerate}
  \item The function $f\circ g$ is periodic.
  \item The function $f\circ g$ is even.
  \item The function $f\circ g$ is bounded on $\R$.
  \item All the answers are correct
\end{enumerate}
\end{question}

\begin{solution}
\textbf{The correct answer: (d)}.

$(f\circ g)(x)=f(g(x))=e^{\cos 2x}$, defined for every $x\in\R$. We check each statement:
\begin{itemize}
  \item \textbf{(a) is true:} $\cos$ is periodic with period $2\pi$, so $\cos(2(x+\pi))=\cos(2x+2\pi)=\cos 2x$, and hence $(f\circ g)(x+\pi)=e^{\cos 2x}=(f\circ g)(x)$ for every $x$. That is, $f\circ g$ is periodic (with period $\pi$).
  \item \textbf{(b) is true:} $\cos$ is even, so $(f\circ g)(-x)=e^{\cos(-2x)}=e^{\cos 2x}=(f\circ g)(x)$ for every $x$.
  \item \textbf{(c) is true:} $-1\le\cos 2x\le 1$ and $e^t$ is strictly increasing, so $e^{-1}\le e^{\cos 2x}\le e$ for every $x\in\R$.
\end{itemize}
Since all three statements are true, the answer is (d). (Each of (a), (b), (c) alone is not the complete answer, since the others are true as well.)
\end{solution}

\subsection*{Part B, Question 5}

\begin{question}
Let $f(x)=3^{-3x+3}$. Which of the following statements is true?
\begin{enumerate}
  \item The function is not monotone.
  \item The function is strictly decreasing.
  \item The function is strictly increasing.
  \item The function is monotone but not strictly.
\end{enumerate}
\end{question}

\begin{solution}
\textbf{The correct answer: (b)}.

$f$ is the composition $f=h\circ u$ of $u(x)=-3x+3$ and $h(t)=3^t$. The function $u$ is strictly decreasing (a linear function with slope $-3<0$) and the function $h$ is strictly increasing (an exponential function with base $3>1$).

Hence if $x_1<x_2$ then $u(x_1)>u(x_2)$, and so $3^{u(x_1)}>3^{u(x_2)}$, i.e. $f(x_1)>f(x_2)$: $f$ is strictly decreasing on $\R$.
(Alternatively: $f(x)=27\cdot\left(\frac{1}{27}\right)^{x}$, an exponential function with base $\frac1{27}<1$ times a positive constant.)

\textbf{Why the others are wrong:} (a): we showed that it is monotone. (c): for example $f(0)=27>f(1)=1$, so it is not increasing. (d): ``monotone but not strictly'' means that there are $x_1<x_2$ with $f(x_1)=f(x_2)$, but $f$ is strictly decreasing and therefore one-to-one.
\end{solution}

\subsection*{Part B, Question 6}

\begin{question}
Let $f(x)=x^3+x^2$. Which of the following statements is true?
\begin{enumerate}
  \item The function is one-to-one and onto $\R$.
  \item The function is one-to-one but not onto $\R$.
  \item The function is onto $\R$ but not one-to-one.
  \item The function is neither one-to-one nor onto $\R$.
\end{enumerate}
\end{question}

\begin{hint}
Factor: $x^3+x^2=x^2(x+1)$. How many roots are there?
\end{hint}

\begin{solution}
\textbf{The correct answer: (c)}.

\textbf{Not one-to-one:} $f(x)=x^2(x+1)$, so $f(0)=0=f(-1)$ with $0\neq-1$.

\textbf{Onto $\R$:} $f$ is a polynomial and therefore continuous on $\R$. Also
\[
\lim_{x\to\infty}x^3\left(1+\tfrac1x\right)=\infty,\qquad \lim_{x\to-\infty}x^3\left(1+\tfrac1x\right)=-\infty .
\]
Let $y\in\R$. By the definition of infinite limits there exist $x_1<x_2$ with $f(x_1)<y<f(x_2)$. By the Intermediate Value Theorem ($f$ is continuous on $[x_1,x_2]$) there exists $c\in(x_1,x_2)$ with $f(c)=y$. Hence $f$ is onto $\R$.

\textbf{Why the others are wrong:} (a) and (b) claim that $f$ is one-to-one, which we disproved. (d) claims that $f$ is not onto, but we proved that it is onto.
\end{solution}

\subsection*{Part B, Question 7}

\begin{question}
Check the continuity of the following function at the point 2:
\[
f(x)=\begin{cases}
\frac{1}{1+4^{\frac{1}{x-2}}} & x\neq 2\\
0 & x=2
\end{cases}
\]
Which of the following statements is true?
\begin{enumerate}
  \item The function is continuous at 2.
  \item The function has a removable discontinuity at 2.
  \item The function has a jump discontinuity at 2.
  \item The function has an essential discontinuity at 2.
\end{enumerate}
\end{question}

\begin{hint}
Compute the one-sided limits separately: where does $\frac{1}{x-2}$ tend as $x\to2^+$ and as $x\to2^-$?
\end{hint}

\begin{solution}
\textbf{The correct answer: (c)}.

We compute the one-sided limits at 2.

\textbf{From the right:} as $x\to2^+$, $x-2\to0^+$ and therefore $\frac1{x-2}\to+\infty$. Since $4>1$, $\lim_{t\to+\infty}4^t=\infty$, so $1+4^{\frac{1}{x-2}}\to\infty$ and
\[
\lim_{x\to2^+}f(x)=0 .
\]
\textbf{From the left:} as $x\to2^-$, $\frac1{x-2}\to-\infty$ and $\lim_{t\to-\infty}4^t=0$, therefore
\[
\lim_{x\to2^-}f(x)=\frac{1}{1+0}=1 .
\]
Both one-sided limits exist (and are finite) but are different, so the limit $\lim_{x\to2}f(x)$ does not exist, and this is exactly a jump discontinuity (of the first kind, not removable).

\textbf{Why the others are wrong:} (a): the limit at 2 does not exist, so $f$ is not continuous at 2 (the fact that $f(2)=0$ equals the right limit gives only continuity from the right). (b): a removable discontinuity requires the two-sided limit to exist. (d): an essential discontinuity (of the second kind) means that at least one of the one-sided limits does not exist as a finite limit, and here both exist.
\end{solution}

\section*{2024/25 Semester A Sample exam}

\subsection*{Question 1a}

\begin{question}
Define a function by
\[ f(x) = \begin{cases} x^3 \sin(\frac{5}{x}) & x \ne 0 \\ 0 & x = 0 \end{cases} \]
Is the function continuous and differentiable at the point $x = 0$? Justify your answer.
\end{question}

\begin{hint}
Use the fact that $\sin$ is bounded: ``a null function times a bounded function'' tends to zero. For differentiability, compute the limit of the difference quotient by definition.
\end{hint}

\begin{solution}
\textbf{Continuity.} For every $x \ne 0$ we have $\left|\sin\frac{5}{x}\right| \le 1$, and hence
\[ 0 \le |f(x)| = |x|^3\left|\sin\tfrac{5}{x}\right| \le |x|^3 \xrightarrow[x\to0]{} 0. \]
By the Squeeze theorem, $\lim_{x\to0} f(x) = 0 = f(0)$, so $f$ is continuous at $x = 0$.

\textbf{Differentiability.} By the definition of the derivative:
\[ \frac{f(x) - f(0)}{x - 0} = \frac{x^3\sin\frac5x}{x} = x^2\sin\frac5x, \qquad 0 \le \left| x^2 \sin\tfrac5x \right| \le x^2 \xrightarrow[x\to0]{} 0. \]
Again by the Squeeze theorem (null times bounded), the limit exists and equals $0$. Hence $f$ is differentiable at $x = 0$ and $f'(0) = 0$.

\textbf{Answer:} $f$ is continuous and differentiable at $x = 0$, and $f'(0) = 0$. (Differentiability implies continuity, so continuity also follows from differentiability.)
\end{solution}

\subsection*{Question 1b}

\begin{question}
Compute the limit
\[ \lim_{x\to\frac{\pi}{4}} (\tan x)^{\tan(2x)} \]
\end{question}

\begin{hint}
This is a $1^\infty$ form. Write $(\tan x)^{\tan 2x} = e^{\tan(2x)\ln(\tan x)}$ and compute the limit of the exponent (for example, with the substitution $t = \tan x$ and L'Hôpital's rule).
\end{hint}

\begin{solution}
In a punctured neighborhood of $\frac{\pi}{4}$ we have $\tan x > 0$, so the expression is defined and
\[ (\tan x)^{\tan 2x} = e^{\tan(2x)\cdot\ln(\tan x)}. \]
As $x \to \frac{\pi}{4}$: $\tan x \to 1$ and $|\tan 2x| \to \infty$ (a $1^\infty$ form). We compute the limit of the exponent. By the double-angle formula, $\tan 2x = \frac{2\tan x}{1 - \tan^2 x}$; substitute $t = \tan x$, so that $t \to 1$ as $x \to \frac{\pi}{4}$ ($\tan$ is continuous), and $t \ne 1$ in a punctured neighborhood:
\[ \lim_{x\to\frac{\pi}{4}} \tan(2x)\ln(\tan x) = \lim_{t\to1} \frac{2t\ln t}{1 - t^2}. \]
This is a $\frac00$ form; by L'Hôpital's rule:
\[ \lim_{t\to1} \frac{2t\ln t}{1 - t^2} = \lim_{t\to1} \frac{2\ln t + 2}{-2t} = \frac{0 + 2}{-2} = -1. \]
(Alternatively: $\frac{2t\ln t}{1 - t^2} = -\frac{2t}{1 + t}\cdot\frac{\ln t}{t - 1} \to -1\cdot 1$.)
By continuity of the exponential function:
\[ \boxed{\lim_{x\to\frac{\pi}{4}} (\tan x)^{\tan(2x)} = e^{-1} = \frac1e}. \]
Note that the limit is two-sided: although $\tan 2x \to +\infty$ from the left and $\tan 2x \to -\infty$ from the right, the product $\tan(2x)\ln(\tan x)$ tends to $-1$ from both sides.
\end{solution}

\subsection*{Question 2a}

\begin{question}
Find all asymptotes of the function
\[ f(x) = \frac{2x^2 - 3x + \ln(x)}{x} \]
Justify your answer.
\end{question}

\begin{hint}
Determine the domain, and check for a vertical asymptote at $x = 0^+$ and an oblique asymptote $y = ax + b$ as $x \to +\infty$.
\end{hint}

\begin{solution}
We write
\[ f(x) = \frac{2x^2 - 3x + \ln x}{x} = 2x - 3 + \frac{\ln x}{x}. \]
\textbf{Domain:} $x > 0$ (because of $\ln x$; in any case $x \ne 0$). $f$ is continuous on $(0, \infty)$.

\textbf{Vertical asymptote.} The only suspicious point is the endpoint of the domain, $x = 0$. As $x \to 0^+$: $\ln x \to -\infty$ and $\frac1x \to +\infty$, hence $\frac{\ln x}{x} = \ln x\cdot\frac1x \to -\infty$, while $2x - 3 \to -3$. Therefore
\[ \lim_{x\to0^+} f(x) = -\infty, \]
and $x = 0$ is a vertical asymptote (from the right). At every other point $f$ is continuous, so there is no vertical asymptote there.

\textbf{Oblique asymptote.} The domain is unbounded only at $+\infty$. We compute
\[ a = \lim_{x\to\infty}\frac{f(x)}{x} = \lim_{x\to\infty}\left( 2 - \frac3x + \frac{\ln x}{x^2} \right) = 2, \]
\[ b = \lim_{x\to\infty}\bigl(f(x) - 2x\bigr) = \lim_{x\to\infty}\left( -3 + \frac{\ln x}{x} \right) = -3, \]
where we used $\lim_{x\to\infty}\frac{\ln x}{x} = 0$ (by L'Hôpital: $\lim \frac{1/x}{1} = 0$), and hence also $\frac{\ln x}{x^2} \to 0$. Therefore $y = 2x - 3$ is an oblique asymptote at $+\infty$, and there is no horizontal asymptote.

\textbf{Answer:} vertical asymptote $x = 0$, and oblique asymptote $y = 2x - 3$ as $x \to +\infty$.
\end{solution}

\subsection*{Question 2b}

\begin{question}
Find the intervals of convexity and the inflection points of the function
\[ f(x) = x^2 \cdot e^{-x} \]
Justify your answer.
\end{question}

\begin{hint}
$f''(x) = e^{-x}(x^2 - 4x + 2)$; check the sign of the quadratic factor.
\end{hint}

\begin{solution}
$f(x) = x^2 e^{-x}$ is defined and infinitely differentiable on all of $\R$.
\[ f'(x) = 2xe^{-x} - x^2e^{-x} = e^{-x}(2x - x^2), \]
\[ f''(x) = -e^{-x}(2x - x^2) + e^{-x}(2 - 2x) = e^{-x}(x^2 - 4x + 2). \]
Since $e^{-x} > 0$, the sign of $f''$ is the sign of $x^2 - 4x + 2$, whose roots are
\[ x_{1,2} = \frac{4 \pm \sqrt{16 - 8}}{2} = 2 \pm \sqrt2. \]
This is an upward-opening parabola, hence:
\begin{itemize}
  \item $f'' > 0$ on $(-\infty, 2 - \sqrt2)$ and on $(2 + \sqrt2, \infty)$: there $f$ is \textbf{convex};
  \item $f'' < 0$ on $(2 - \sqrt2, 2 + \sqrt2)$: there $f$ is \textbf{concave}.
\end{itemize}
At the points $x = 2 \pm \sqrt2$ the function is differentiable and $f''$ changes sign, so these are inflection points. The values:
\[ f(2 \pm \sqrt2) = (2 \pm \sqrt2)^2 e^{-(2 \pm \sqrt2)} = (6 \pm 4\sqrt2)\,e^{-2 \mp \sqrt2}. \]

\textbf{Answer:} $f$ is convex on $(-\infty, 2 - \sqrt2)$ and on $(2 + \sqrt2, \infty)$, concave on $(2 - \sqrt2, 2 + \sqrt2)$, and the inflection points are $x = 2 - \sqrt2$ and $x = 2 + \sqrt2$.
\end{solution}

\subsection*{Question 3}

\begin{question}
Using a linear approximation (a first-degree Taylor polynomial), estimate $\sqrt[4]{18}$ and give a bound on the error.
\end{question}

\begin{hint}
Expand $f(x) = \sqrt[4]{x}$ around $a = 16$, since $\sqrt[4]{16} = 2$, and bound the Lagrange remainder $R_1$.
\end{hint}

\begin{solution}
Define $f(x) = \sqrt[4]{x} = x^{1/4}$ and expand around $a = 16$, the point closest to $18$ at which the fourth root is known. For $x > 0$:
\[ f'(x) = \frac14 x^{-3/4}, \qquad f''(x) = -\frac{3}{16} x^{-7/4}, \]
hence $f(16) = 2$ and $f'(16) = \frac{1}{4\cdot 16^{3/4}} = \frac{1}{4\cdot 8} = \frac{1}{32}$.

\textbf{Linear approximation:}
\[ P_1(x) = f(16) + f'(16)(x - 16) = 2 + \frac{x - 16}{32}, \qquad
   \sqrt[4]{18} \approx P_1(18) = 2 + \frac{2}{32} = 2 + \frac{1}{16} = \frac{33}{16} = 2.0625. \]

\textbf{Error bound.} By Taylor's theorem with the Lagrange remainder, there exists $c \in (16, 18)$ such that
\[ R_1(18) = \frac{f''(c)}{2!}(18 - 16)^2 = \frac{-\frac{3}{16}c^{-7/4}}{2}\cdot 4 = -\frac{3}{8\,c^{7/4}}. \]
Since $c > 16$, we have $c^{7/4} > 16^{7/4} = 2^7 = 128$, and hence
\[ |R_1(18)| = \frac{3}{8c^{7/4}} < \frac{3}{8\cdot128} = \frac{3}{1024} \approx 0.0029. \]

\textbf{Answer:} $\sqrt[4]{18} \approx \frac{33}{16} = 2.0625$, with error less than $\frac{3}{1024}$. Since $R_1 < 0$, the approximation is larger than the true value: $\frac{33}{16} - \frac{3}{1024} < \sqrt[4]{18} < \frac{33}{16}$.

(Check: $\sqrt[4]{18} \approx 2.05977$, and the actual error is about $0.00273 < \frac{3}{1024}$.)
\end{solution}

\subsection*{Question 4a}

\begin{question}
Compute the integral
\[ \int \frac{6x^2 + 13x + 8}{x^3 + 2x^2 + 2x}\,dx \]
\end{question}

\begin{hint}
$x^3 + 2x^2 + 2x = x(x^2 + 2x + 2)$ and the quadratic factor is irreducible. Decompose into $\frac{A}{x} + \frac{Bx + C}{x^2 + 2x + 2}$.
\end{hint}

\begin{solution}
\textbf{Factoring the denominator:} $x^3 + 2x^2 + 2x = x(x^2 + 2x + 2)$, and the discriminant of $x^2 + 2x + 2$ is $4 - 8 < 0$, so the quadratic factor is irreducible ($x^2 + 2x + 2 = (x + 1)^2 + 1 > 0$). The degree of the numerator is less than the degree of the denominator.

\textbf{Partial fractions:}
\[ \frac{6x^2 + 13x + 8}{x(x^2 + 2x + 2)} = \frac{A}{x} + \frac{Bx + C}{x^2 + 2x + 2}, \]
i.e. $6x^2 + 13x + 8 = A(x^2 + 2x + 2) + (Bx + C)x$.
\begin{itemize}
  \item $x = 0$: $8 = 2A$, hence $A = 4$.
  \item Coefficients of $x^2$: $6 = A + B$, hence $B = 2$.
  \item Coefficients of $x$: $13 = 2A + C$, hence $C = 5$.
\end{itemize}
Therefore
\[ \int \frac{6x^2 + 13x + 8}{x^3 + 2x^2 + 2x}\,dx = \int \left( \frac4x + \frac{2x + 2}{x^2 + 2x + 2} + \frac{3}{(x + 1)^2 + 1} \right) dx, \]
where we wrote $2x + 5 = (2x + 2) + 3$.
\begin{itemize}
  \item $\int \frac4x\,dx = 4\ln|x| + C$.
  \item With the substitution $t = x^2 + 2x + 2$, $dt = (2x + 2)\,dx$: $\int \frac{2x + 2}{x^2 + 2x + 2}\,dx = \ln(x^2 + 2x + 2) + C$.
  \item With the substitution $u = x + 1$: $\int \frac{3\,dx}{(x + 1)^2 + 1} = 3\arctan(x + 1) + C$.
\end{itemize}
\textbf{Answer:}
\[ \boxed{\int \frac{6x^2 + 13x + 8}{x^3 + 2x^2 + 2x}\,dx = 4\ln|x| + \ln(x^2 + 2x + 2) + 3\arctan(x + 1) + C.} \]
\end{solution}

\subsection*{Question 4b}

\begin{question}
Find the general solution of the differential equation
\[ (x^2 + 4)y' + 2xy^2 = 0 \]
\end{question}

\begin{hint}
The equation is separable: $\frac{dy}{y^2} = -\frac{2x}{x^2 + 4}\,dx$. Do not forget the solution $y \equiv 0$.
\end{hint}

\begin{solution}
Write the equation in the form ($x^2 + 4 > 0$ for every $x$):
\[ \frac{dy}{dx} = -\frac{2x}{x^2 + 4}\, y^2, \]
which is a separable equation.

\textbf{The constant solution:} $y \equiv 0$ satisfies the equation, so it is a solution.

\textbf{For $y \ne 0$:}
\[ \int \frac{dy}{y^2} = -\int \frac{2x}{x^2 + 4}\,dx. \]
The left-hand side is $-\frac1y$. On the right-hand side, with the substitution $t = x^2 + 4$, $dt = 2x\,dx$: $\int \frac{2x}{x^2 + 4}\,dx = \ln(x^2 + 4) + C$. Hence
\[ -\frac{1}{y} = -\ln(x^2 + 4) + C \quad\Longrightarrow\quad y = \frac{1}{\ln(x^2 + 4) - C}. \]

\textbf{Answer:} the general solution is
\[ \boxed{y = \frac{1}{\ln(x^2 + 4) + C}}, \quad C \in \R \]
(we renamed the constant; on every interval on which the denominator does not vanish), together with the solution $y \equiv 0$.
\end{solution}

\subsection*{Question 5a}

\begin{question}
Compute the volume of the solid of revolution of the function $f(x) = \sqrt{\cos x} \cdot \sin^2 x$ on the interval $[0, \frac{\pi}{2}]$.
\end{question}

\begin{hint}
For revolution about the $x$-axis: $V = \pi\int_0^{\pi/2} f^2(x)\,dx = \pi\int_0^{\pi/2}\cos x\,\sin^4 x\,dx$; substitute $t = \sin x$.
\end{hint}

\begin{solution}
We interpret the question as revolving the graph of $f$ about the $x$-axis. $f$ is defined and continuous on $[0, \frac{\pi}{2}]$ (where $\cos x \ge 0$), and the volume of the solid of revolution is
\[ V = \pi\int_0^{\pi/2} f^2(x)\,dx = \pi\int_0^{\pi/2} \cos x\,\sin^4 x\,dx. \]
Substitute $t = \sin x$, $dt = \cos x\,dx$; when $x = 0$ we have $t = 0$, and when $x = \frac{\pi}{2}$ we have $t = 1$:
\[ V = \pi\int_0^1 t^4\,dt = \pi\left[\frac{t^5}{5}\right]_0^1 = \frac{\pi}{5}. \]
\textbf{Answer:} $\boxed{V = \frac{\pi}{5}}$.
\end{solution}

\subsection*{Question 5b}

\begin{question}
Compute the improper integral $\int_0^\infty xe^{-x}\,dx$, or show that it does not converge.
\end{question}

\begin{hint}
Find an antiderivative using integration by parts, then compute the limit as the upper limit of integration tends to infinity.
\end{hint}

\begin{solution}
The integrand is continuous on $[0, \infty)$, and by definition
\[ \int_0^\infty xe^{-x}\,dx = \lim_{A\to\infty}\int_0^A xe^{-x}\,dx. \]
\textbf{Antiderivative} by integration by parts, $u = x$, $v' = e^{-x}$ ($u' = 1$, $v = -e^{-x}$):
\[ \int xe^{-x}\,dx = -xe^{-x} + \int e^{-x}\,dx = -xe^{-x} - e^{-x} + C = -(x + 1)e^{-x} + C. \]
Hence, by the Newton--Leibniz formula,
\[ \int_0^A xe^{-x}\,dx = \Bigl[-(x + 1)e^{-x}\Bigr]_0^A = 1 - \frac{A + 1}{e^A}. \]
As $A \to \infty$, $\frac{A + 1}{e^A} \to 0$ (a $\frac{\infty}{\infty}$ form; by L'Hôpital, $\lim \frac{1}{e^A} = 0$).

\textbf{Answer:} the integral converges and $\boxed{\int_0^\infty xe^{-x}\,dx = 1}$.
\end{solution}

\section*{2024/25 Semester A Moed A}

\subsection*{Question 1a}

\begin{question}
Define a function by
\[ f(x) = \begin{cases} ax^2 + b & x \le 0 \\ \dfrac{\ln(e^x - 1)}{2\ln x} & x > 0 \end{cases} \]
For which values of the parameters $a, b$ is the function continuous at the point $x=0$? Justify.
\end{question}

\begin{hint}
Write $e^x - 1 = x \cdot \frac{e^x - 1}{x}$ and use the known limit $\lim_{x\to 0}\frac{e^x-1}{x} = 1$.
\end{hint}

\begin{solution}
For all $a,b$ we have $f(0) = a\cdot 0 + b = b$, and to the left of $0$ the function is the polynomial $ax^2+b$, so
$\lim_{x\to 0^-} f(x) = b = f(0)$. Hence $f$ is continuous at $0$ if and only if $\lim_{x\to 0^+} f(x) = b$.

We compute the right limit. For $x>0$ we have $e^x - 1 > 0$, and therefore
\[ \ln(e^x - 1) = \ln\left( x \cdot \frac{e^x-1}{x} \right) = \ln x + \ln\frac{e^x - 1}{x}. \]
Hence
\[ f(x) = \frac{\ln x + \ln\frac{e^x-1}{x}}{2\ln x} = \frac12 + \frac{\ln\frac{e^x-1}{x}}{2\ln x}. \]
Since $\lim_{x\to 0}\frac{e^x - 1}{x} = 1$ (this is the derivative of $e^x$ at $0$) and $\ln$ is continuous at $1$, the numerator of the second fraction tends to $\ln 1 = 0$,
while $2\ln x \to -\infty$. Hence the second fraction tends to $0$ and we get
\[ \lim_{x\to 0^+} f(x) = \frac12. \]
(One can also use L'Hôpital's rule in the case $\frac{\infty}{\infty}$:
$\lim_{x\to0^+}\frac{\ln(e^x-1)}{2\ln x} = \lim_{x\to0^+}\frac{e^x/(e^x-1)}{2/x} = \lim_{x\to 0^+}\frac{e^x}{2}\cdot\frac{x}{e^x-1} = \frac12$.)

\textbf{Answer:} $f$ is continuous at $x=0$ if and only if $b = \frac12$, with $a$ arbitrary ($a \in \R$).
\end{solution}

\subsection*{Question 1b}

\begin{question}
How many real solutions does the following equation have?
\[ x + \sqrt{x} = 17 \]
Justify your answer carefully.
\end{question}

\begin{hint}
Define $g(x) = x + \sqrt{x} - 17$ on $[0,\infty)$: existence of a solution follows from the Intermediate Value Theorem, and uniqueness from monotonicity.
\end{hint}

\begin{solution}
Define $g(x) = x + \sqrt{x} - 17$. The equation is defined only for $x \ge 0$, and $g$ is continuous on $[0,\infty)$.

\textbf{Existence:} $g(0) = -17 < 0$ and $g(17) = \sqrt{17} > 0$. By the Intermediate Value Theorem there exists $c \in (0,17)$ with $g(c) = 0$.

\textbf{Uniqueness:} for $x > 0$ we have $g'(x) = 1 + \frac{1}{2\sqrt{x}} > 0$, and therefore (a corollary of Lagrange's theorem) $g$ is strictly increasing on $[0,\infty)$
(in fact this is also clear directly: it is a sum of two strictly increasing functions). A strictly increasing function attains each value at most once, so there is at most one solution.

\textbf{Answer:} the equation has \textbf{exactly one} real solution. (In fact, substituting $t = \sqrt{x}\ge 0$ gives $t^2 + t - 17 = 0$, i.e. $t = \frac{-1+\sqrt{69}}{2}$, and $x = \left(\frac{-1+\sqrt{69}}{2}\right)^2 \approx 13.35$.)
\end{solution}

\subsection*{Question 2a}

\begin{question}
Find all the asymptotes of the function
\[ f(x) = \frac{x^3 - 3x^2 + 3x + \sqrt{x}}{x^2 + x - 2} \]
Justify.
\end{question}

\begin{hint}
Pay attention to the domain: because of $\sqrt{x}$, the function is defined only for $x \ge 0$. Which zeros of the denominator lie in the domain?
\end{hint}

\begin{solution}
\textbf{Domain.} $\sqrt{x}$ is defined for $x\ge0$, and the denominator factors: $x^2 + x - 2 = (x-1)(x+2)$. Hence the domain is $[0,1)\cup(1,\infty)$
(the point $x=-2$ is not in the domain at all).

\textbf{Vertical asymptotes.} $f$ is continuous at every point of its domain (a quotient of continuous functions), so a vertical asymptote can occur only at $x=1$.
There the numerator tends to $1 - 3 + 3 + 1 = 2 \neq 0$ and the denominator tends to $0$, where $x^2+x-2 < 0$ for $0\le x<1$ and $>0$ for $x>1$. Hence
\[ \lim_{x\to1^-} f(x) = -\infty, \qquad \lim_{x\to1^+} f(x) = +\infty, \]
and $x = 1$ is a vertical asymptote. At $x=0$ the function is defined ($f(0)=0$) and continuous from the right, so there is no asymptote there.

\textbf{Oblique/horizontal asymptotes.} Only $x\to+\infty$ needs to be checked (the function is not defined for $x<0$). We look for $y = mx + n$:
\[ m = \lim_{x\to\infty}\frac{f(x)}{x} = \lim_{x\to\infty}\frac{x^3 - 3x^2 + 3x + \sqrt{x}}{x^3 + x^2 - 2x}
= \lim_{x\to\infty}\frac{1 - \frac3x + \frac3{x^2} + x^{-5/2}}{1 + \frac1x - \frac2{x^2}} = 1. \]
By polynomial long division $x^3 - 3x^2 + 3x = (x^2+x-2)(x-4) + (9x - 8)$, and therefore
\[ f(x) - x = -4 + \frac{9x - 8 + \sqrt{x}}{x^2 + x - 2} \xrightarrow[x\to\infty]{} -4 + 0 = -4, \]
since in the numerator of the fraction the highest power is $x$ and in the denominator it is $x^2$. Hence $n=-4$.

\textbf{Answer:} vertical asymptote $x = 1$, and oblique asymptote $y = x - 4$ as $x \to +\infty$. There are no other asymptotes.
\end{solution}

\subsection*{Question 2b}

\begin{question}
Find the absolute minimum and maximum of the function
\[ f(x) = e^x \cdot (x-1)^2 \]
or explain why they do not exist. Justify.
\end{question}

\begin{hint}
Note that $f(x) \ge 0$ for every $x$. What happens as $x \to \infty$?
\end{hint}

\begin{solution}
$f$ is defined and differentiable on all of $\R$.
\[ f'(x) = e^x(x-1)^2 + 2e^x(x-1) = e^x(x-1)(x+1). \]
Critical points: $x = \pm 1$. $f' > 0$ on $(-\infty,-1)$, $f'<0$ on $(-1,1)$, $f'>0$ on $(1,\infty)$.
Hence $x=-1$ is a local maximum ($f(-1) = 4/e$) and $x=1$ is a local minimum ($f(1)=0$).

\textbf{Absolute minimum:} for every $x$ we have $e^x>0$ and $(x-1)^2 \ge 0$, so $f(x) \ge 0 = f(1)$. Hence the absolute minimum is $0$, attained at $x=1$.

\textbf{Absolute maximum:} $\lim_{x\to\infty} e^x(x-1)^2 = \infty$ (a product of two factors tending to $\infty$), so $f$ is not bounded above and has no absolute maximum.
(The local maximum $4/e$ at $x=-1$ is not absolute; for example $f(3) = 4e^3 > 4/e$.)

\textbf{Answer:} absolute minimum $0$ at the point $x=1$; an absolute maximum does not exist since $f(x)\to\infty$ as $x\to\infty$.
\end{solution}

\subsection*{Question 3}

\begin{question}
Compute an approximation of $\sqrt[3]{2}$ with an error smaller than $\frac{1}{10}$. Justify.
\end{question}

\begin{hint}
Use the Taylor polynomial of $f(x) = \sqrt[3]{x}$ around $x_0 = 1$, and estimate the error using the Lagrange remainder.
\end{hint}

\begin{hint}
Check: is a linear approximation (order 1) enough to guarantee an error smaller than $\frac1{10}$? If not, move to order 2.
\end{hint}

\begin{solution}
Define $f(x) = x^{1/3}$ and expand around $x_0 = 1$ (where the values are known). The derivatives:
\[ f'(x) = \tfrac13 x^{-2/3},\quad f''(x) = -\tfrac29 x^{-5/3},\quad f'''(x) = \tfrac{10}{27}x^{-8/3}, \]
and therefore $f(1) = 1$, $f'(1) = \frac13$, $f''(1) = -\frac29$.

\textbf{Attempt with order 1.} $P_1(2) = 1 + \frac13 = \frac43$, and by the Lagrange remainder there exists $c\in(1,2)$ with
$R_1 = \frac{f''(c)}{2!}(2-1)^2 = -\frac19 c^{-5/3}$. From this we get only $|R_1| < \frac19$, which does \emph{not} guarantee an error smaller than $\frac1{10}$. Hence we move to order 2.

\textbf{Order 2.} The Taylor polynomial of order 2 around $1$:
\[ P_2(x) = 1 + \tfrac13(x-1) - \tfrac19(x-1)^2, \qquad P_2(2) = 1 + \tfrac13 - \tfrac19 = \tfrac{11}{9}. \]
By Taylor's theorem with the Lagrange remainder, there exists $c \in (1,2)$ such that
\[ \sqrt[3]{2} - \tfrac{11}{9} = R_2(2) = \frac{f'''(c)}{3!}(2-1)^3 = \frac{10}{27\cdot 6}\,c^{-8/3} = \frac{5}{81}\,c^{-8/3}. \]
Since $c > 1$ we have $c^{-8/3} < 1$, and therefore
\[ 0 < R_2(2) < \frac{5}{81} < \frac{1}{10} \qquad (\text{since } 50 < 81). \]

\textbf{Answer:} $\sqrt[3]{2} \approx \frac{11}{9} \approx 1.222$, with an error smaller than $\frac{5}{81}<\frac1{10}$ (and in particular $\frac{11}{9} < \sqrt[3]{2} < \frac{11}{9} + \frac{5}{81}$).
As a check: $\sqrt[3]{2} \approx 1.2599$, and the actual error is about $0.038$.

\textbf{Alternative method (linear approximation at a closer point).} $\left(\frac54\right)^3 = \frac{125}{64}$, which is close to $2$. Linear approximation around $x_0 = \frac{125}{64}$:
$f(2) \approx \frac54 + \frac13\cdot\frac{16}{25}\cdot\frac{3}{64} = \frac54 + \frac1{100} = 1.26$, and the error is bounded by $\frac19 c^{-5/3}\left(\frac3{64}\right)^2 < \frac{1}{9}\cdot\frac{9}{4096}<\frac1{10}$ (since $c>1$).
\end{solution}

\subsection*{Question 4a}

\begin{question}
Compute the integral
\[ \int \frac{2x^2 + 8}{x^3 - 4x^2 + 8x}\,dx \]
\end{question}

\begin{hint}
Factor the denominator: $x^3 - 4x^2 + 8x = x(x^2 - 4x + 8)$, where the quadratic factor has no real roots, and decompose into partial fractions.
\end{hint}

\begin{solution}
\textbf{Factoring the denominator.} $x^3 - 4x^2 + 8x = x(x^2 - 4x + 8)$, and $x^2 - 4x + 8 = (x-2)^2 + 4$ has no real roots (the discriminant is $16 - 32 < 0$).

\textbf{Partial fractions.} We look for
\[ \frac{2x^2 + 8}{x(x^2 - 4x + 8)} = \frac{A}{x} + \frac{Bx + C}{x^2 - 4x + 8}, \]
i.e. $2x^2 + 8 = A(x^2 - 4x + 8) + x(Bx + C)$. Substituting $x = 0$ gives $8 = 8A$, i.e. $A = 1$. Then
$2x^2 + 8 = x^2 - 4x + 8 + Bx^2 + Cx$, and comparing coefficients gives $B = 1$, $C = 4$.

\textbf{Integration.}
\[ \int\frac{dx}{x} = \ln|x|. \]
For the second term write $x + 4 = \frac12(2x - 4) + 6$, where $(x^2 - 4x + 8)' = 2x - 4$:
\[ \int\frac{x+4}{x^2-4x+8}\,dx = \frac12\int\frac{2x-4}{x^2-4x+8}\,dx + 6\int\frac{dx}{(x-2)^2 + 4}
= \frac12\ln(x^2-4x+8) + 6\cdot\frac12\arctan\frac{x-2}{2}, \]
where in the last integral we used $\int\frac{dt}{t^2 + a^2} = \frac1a\arctan\frac ta$ with $t = x-2$, $a = 2$.

\textbf{Answer:}
\[ \int \frac{2x^2 + 8}{x^3 - 4x^2 + 8x}\,dx = \ln|x| + \frac12\ln(x^2 - 4x + 8) + 3\arctan\frac{x-2}{2} + C. \]
(Check: differentiating the result gives back the integrand.)
\end{solution}

\subsection*{Question 4b}

\begin{question}
Find the general solution of the differential equation
\[ y' = e^{x+y} \]
\end{question}

\begin{hint}
$e^{x+y} = e^x e^y$, so this is a separable equation.
\end{hint}

\begin{solution}
$y' = e^x e^y$. Since $e^y > 0$ always, there are no constant solutions (for $y\equiv c$ we would get $0 = e^{x+c}$, a contradiction), and we may divide by $e^y$:
\[ e^{-y}\,y' = e^x \quad\Longrightarrow\quad \int e^{-y}\,dy = \int e^x\,dx \quad\Longrightarrow\quad -e^{-y} = e^x + c. \]
Set $C = -c$: $e^{-y} = C - e^x$. The left-hand side is positive, so the solution is defined only for $x$ satisfying $e^x < C$ (in particular $C > 0$ is required), and there
\[ y = -\ln\left(C - e^x\right), \qquad C > 0,\ x < \ln C. \]
\textbf{Check:} $y' = \frac{e^x}{C - e^x}$, and on the other hand $e^{x+y} = e^x\cdot\frac{1}{C - e^x}$. Equal.

\textbf{Answer:} the general solution is $y(x) = -\ln(C - e^x)$ (equivalently: $e^{-y} + e^x = C$), with a constant $C>0$, on the interval $x < \ln C$.
\end{solution}

\subsection*{Question 5a}

\begin{question}
Compute the area bounded between the function
\[ f(x) = \sin^2 x \cdot \cos^3 x \]
and the $x$-axis on the interval $[0, \frac{\pi}{2}]$.
\end{question}

\begin{hint}
Write $\cos^3 x = (1 - \sin^2 x)\cos x$ and substitute $t = \sin x$.
\end{hint}

\begin{solution}
On the interval $[0,\frac\pi2]$ we have $\sin^2 x \ge 0$ and $\cos x\ge 0$, so $f(x) \ge 0$ and the area is simply $\int_0^{\pi/2} f(x)\,dx$.
Write $\sin^2 x\cos^3 x = \sin^2 x(1 - \sin^2 x)\cos x$ and substitute $t = \sin x$, $dt = \cos x\,dx$; as $x$ goes from $0$ to $\frac\pi2$, $t$ goes from $0$ to $1$:
\[ S = \int_0^{\pi/2}\sin^2 x\cos^3 x\,dx = \int_0^1 (t^2 - t^4)\,dt = \left[\frac{t^3}{3} - \frac{t^5}{5}\right]_0^1 = \frac13 - \frac15 = \frac{2}{15}. \]
\textbf{Answer:} $S = \frac{2}{15}$.
\end{solution}

\subsection*{Question 5b}

\begin{question}
Compute the following improper integral, or show that it does not converge:
\[ \int_1^\infty \frac{1}{(1+x)\cdot\sqrt{x}}\,dx \]
\end{question}

\begin{hint}
Substitute $t = \sqrt{x}$.
\end{hint}

\begin{solution}
By definition, $\int_1^\infty g = \lim_{R\to\infty}\int_1^R g$. We compute the integral over $[1,R]$ with the substitution $t = \sqrt{x}$, i.e. $x = t^2$, $dx = 2t\,dt$; as $x$ goes from $1$ to $R$, $t$ goes from $1$ to $\sqrt R$:
\[ \int_1^R \frac{dx}{(1+x)\sqrt{x}} = \int_1^{\sqrt R}\frac{2t\,dt}{(1+t^2)\,t} = 2\int_1^{\sqrt R}\frac{dt}{1+t^2} = 2\left(\arctan\sqrt R - \arctan 1\right). \]
As $R \to \infty$ also $\sqrt R \to \infty$, and $\arctan\sqrt R \to \frac\pi2$. Hence
\[ \int_1^\infty \frac{dx}{(1+x)\sqrt{x}} = 2\left(\frac\pi2 - \frac\pi4\right) = \frac{\pi}{2}. \]
(Convergence can also be seen in advance by the comparison test: $0 < \frac{1}{(1+x)\sqrt x} < \frac{1}{x^{3/2}}$, and $\int_1^\infty x^{-3/2}dx$ converges.)

\textbf{Answer:} the integral converges and its value is $\frac{\pi}{2}$.
\end{solution}

\section*{2024/25 Semester A Moed B}

\subsection*{Question 1a}

\begin{question}
For a real number $a$ define a function by
\[ f(x) = \begin{cases} \dfrac{\sqrt{x+4} - 2}{x^2 + 2x} & x < 0 \\ a\cdot|x - 1| & x \ge 0 \end{cases} \]
(11 pts) For which value of $a$ is the function continuous at the point $x = 0$? Justify.
\end{question}

\begin{hint}
Multiply the numerator and denominator by the conjugate $\sqrt{x+4} + 2$.
\end{hint}

\begin{solution}
$f(0) = a\cdot|0 - 1| = a$, and to the right of $0$ (in a right neighborhood) $f(x) = a|x-1|$ is continuous, so $\lim_{x\to0^+}f(x) = a = f(0)$.
Hence $f$ is continuous at $0$ if and only if $\lim_{x\to 0^-} f(x) = a$.

For $-2 < x < 0$ we multiply by the conjugate:
\[ \frac{\sqrt{x+4} - 2}{x^2 + 2x} = \frac{(x + 4) - 4}{x(x+2)\left(\sqrt{x+4} + 2\right)} = \frac{1}{(x+2)\left(\sqrt{x+4}+2\right)}. \]
The last expression is continuous at $x=0$, and therefore
\[ \lim_{x\to 0^-} f(x) = \frac{1}{2\cdot(2 + 2)} = \frac18. \]
\textbf{Answer:} $f$ is continuous at $0$ if and only if $a = \frac18$.
\end{solution}

\subsection*{Question 1b}

\begin{question}
For a real number $a$ define a function by
\[ f(x) = \begin{cases} \dfrac{\sqrt{x+4} - 2}{x^2 + 2x} & x < 0 \\ a\cdot|x - 1| & x \ge 0 \end{cases} \]
(10 pts) For which value of $a$ is the function differentiable at the point $x = 1$? Justify.
\end{question}

\begin{hint}
Compute the one-sided derivatives of $a|x-1|$ at $x=1$.
\end{hint}

\begin{solution}
In a neighborhood of $x = 1$ (for example for $x>0$) $f(x) = a|x - 1|$, and $f(1) = 0$. We compute the one-sided derivatives:
\[ f'_+(1) = \lim_{h\to0^+}\frac{a|h| - 0}{h} = a, \qquad f'_-(1) = \lim_{h\to0^-}\frac{a|h|}{h} = \lim_{h\to 0^-}\frac{-ah}{h} = -a. \]
$f$ is differentiable at $1$ if and only if both one-sided derivatives exist and are equal: $a = -a$, i.e. $a = 0$.
In this case $f \equiv 0$ on $[0,\infty)$ and $f'(1) = 0$.

\textbf{Answer:} $f$ is differentiable at $x=1$ if and only if $a = 0$.
\end{solution}

\subsection*{Question 2}

\begin{question}
Define a function
\[ f(x) = \frac{x^2}{e^x} \]
\begin{enumerate}
  \item (11 pts) Compute the local extremum points and the intervals of increase and decrease of the function. Justify.
  \item (10 pts) Compute the inflection points and the intervals of convexity of the function. Justify.
\end{enumerate}
\end{question}

\begin{solution}
$f(x) = x^2e^{-x}$ is defined and infinitely differentiable on all of $\R$.
\begin{enumerate}
  \item
  \[ f'(x) = 2xe^{-x} - x^2e^{-x} = x(2 - x)e^{-x}. \]
  Since $e^{-x} > 0$, the sign of $f'$ is the sign of $x(2-x)$:
  \begin{itemize}
    \item $x < 0$: $f' < 0$, $f$ is decreasing;
    \item $0 < x < 2$: $f' > 0$, $f$ is increasing;
    \item $x > 2$: $f' < 0$, $f$ is decreasing.
  \end{itemize}
  \textbf{Answer:} $f$ is decreasing on $(-\infty,0]$, increasing on $[0,2]$ and decreasing on $[2,\infty)$.
  By the first derivative test (sign change of $f'$): $x = 0$ is a local minimum point, $f(0) = 0$; $x = 2$ is a local maximum point, $f(2) = \frac{4}{e^2}$.

  \item
  \[ f''(x) = (2 - 2x)e^{-x} - (2x - x^2)e^{-x} = (x^2 - 4x + 2)e^{-x}. \]
  $f''(x) = 0 \iff x^2 - 4x + 2 = 0 \iff x = 2 \pm \sqrt{2}$. The sign of $f''$ is the sign of the parabola $x^2 - 4x + 2$ (which opens upward):
  \begin{itemize}
    \item $x < 2 - \sqrt2$ or $x > 2+\sqrt2$: $f'' > 0$, $f$ is convex (upward, $\cup$);
    \item $2-\sqrt2 < x < 2 + \sqrt 2$: $f'' < 0$, $f$ is concave (downward, $\cap$).
  \end{itemize}
  At both points $x = 2\pm\sqrt2$, $f''$ changes sign, so they are inflection points.

  \textbf{Answer:} $f$ is convex on $(-\infty, 2 - \sqrt2]$ and on $[2+\sqrt2, \infty)$, and concave on $[2-\sqrt2, 2+\sqrt2]$.
  The inflection points: $x = 2 \pm \sqrt2$, with values $f(2\pm\sqrt2) = (6 \pm 4\sqrt2)\,e^{-(2\pm\sqrt2)}$
  (i.e. approximately $(0.586,\ 0.191)$ and $(3.414,\ 0.384)$).
\end{enumerate}
\end{solution}

\subsection*{Question 3}

\begin{question}
Use a linear approximation to compute $\ln(1.2)$, and estimate the error. Justify.
\end{question}

\begin{hint}
Use the Taylor polynomial of order 1 of $\ln x$ around $x_0 = 1$, and estimate the error using the Lagrange remainder.
\end{hint}

\begin{solution}
Define $f(x) = \ln x$ and choose $x_0 = 1$ (a point close to $1.2$ where the values are known). $f(1) = 0$, $f'(x) = \frac1x$, $f'(1) = 1$, $f''(x) = -\frac{1}{x^2}$.

\textbf{Linear approximation:} $P_1(x) = f(1) + f'(1)(x - 1) = x - 1$, and therefore
\[ \ln(1.2) \approx P_1(1.2) = 0.2. \]

\textbf{Error estimate:} by Taylor's theorem with the Lagrange remainder, there exists $c\in(1, 1.2)$ such that
\[ \ln(1.2) - 0.2 = R_1(1.2) = \frac{f''(c)}{2!}(1.2 - 1)^2 = -\frac{1}{2c^2}\cdot 0.04 = -\frac{0.02}{c^2}. \]
Since $1 < c < 1.2$ we have $1 < c^2 < 1.44$, and therefore
\[ \frac{0.02}{1.44} < |R_1| < 0.02, \qquad\text{i.e.}\qquad 0.0138 < |R_1| < 0.02, \]
and the error is negative, i.e. the approximation $0.2$ is larger than the true value.

\textbf{Answer:} $\ln(1.2) \approx 0.2$, with an error smaller than $0.02$ (and in particular $0.18 < \ln(1.2) < 0.2$).
As a check: $\ln(1.2) \approx 0.1823$, and the actual error is about $0.0177$.
\end{solution}

\subsection*{Question 4a}

\begin{question}
Compute the integral
\[ \int_0^{\frac{\sqrt2}{2}} \frac{x^2}{\sqrt{1 - x^2}}\,dx \]
Reminder: $\sin(\frac\pi4) = \cos(\frac\pi4) = \frac{\sqrt2}{2}$.
\end{question}

\begin{hint}
Substitute $x = \sin t$, and then use the identity $\sin^2 t = \frac{1 - \cos 2t}{2}$.
\end{hint}

\begin{solution}
Substitute $x = \sin t$ with $t\in[0,\frac\pi4]$ (a differentiable increasing function, $\sin 0 = 0$, $\sin\frac\pi4 = \frac{\sqrt2}2$), $dx = \cos t\,dt$.
On this interval $\cos t > 0$, and therefore $\sqrt{1 - \sin^2 t} = \cos t$:
\[ \int_0^{\frac{\sqrt2}{2}}\frac{x^2}{\sqrt{1-x^2}}\,dx = \int_0^{\pi/4}\frac{\sin^2 t}{\cos t}\cos t\,dt = \int_0^{\pi/4}\sin^2 t\,dt
= \int_0^{\pi/4}\frac{1 - \cos 2t}{2}\,dt = \left[\frac t2 - \frac{\sin 2t}{4}\right]_0^{\pi/4}. \]
Hence
\[ \int_0^{\frac{\sqrt2}{2}}\frac{x^2}{\sqrt{1-x^2}}\,dx = \frac\pi8 - \frac{\sin\frac\pi2}{4} = \frac{\pi}{8} - \frac14. \]
(The integrand is continuous on the closed interval, since $1 - x^2 \ge \frac12$ there, so this is an ordinary integral.)

\textbf{Answer:} $\frac\pi8 - \frac14 \approx 0.1427$.
\end{solution}

\subsection*{Question 4b}

\begin{question}
Find the general solution of the differential equation
\[ (x^4 + 4)yy' = x^3 \]
\end{question}

\begin{hint}
This is a separable equation: $y\,dy = \frac{x^3}{x^4 + 4}\,dx$.
\end{hint}

\begin{solution}
$x^4 + 4 > 0$ for every $x$, so we can divide: $y y' = \frac{x^3}{x^4 + 4}$. This is a separable equation. Integrate both sides with respect to $x$:
\[ \int y\,dy = \int\frac{x^3}{x^4 + 4}\,dx \quad\Longrightarrow\quad \frac{y^2}{2} = \frac14\ln(x^4 + 4) + c, \]
where on the right-hand side we substituted $u = x^4 + 4$, $du = 4x^3dx$. Multiplying by $2$ and setting $C = 2c$:
\[ y^2 = \frac12\ln(x^4 + 4) + C. \]
(Note that $y\equiv0$ is not a solution, since then we would get $0 = x^3$, which does not hold for every $x$. Also, no division by $y$ was needed: we only used the fact that $yy' = \left(\frac{y^2}{2}\right)'$.)

\textbf{Answer:} the general solution is given implicitly by $y^2 = \frac12\ln(x^4+4) + C$, or explicitly
\[ y = \pm\sqrt{\tfrac12\ln(x^4 + 4) + C}, \]
on the interval where the expression under the root is positive. \textbf{Check:} differentiating $y^2 = \frac12\ln(x^4+4)+C$ gives $2yy' = \frac{2x^3}{x^4+4}$, i.e. $(x^4+4)yy' = x^3$.
\end{solution}

\subsection*{Question 5a}

\begin{question}
Compute the volume of the solid of revolution about the $x$-axis on the interval $[1,4]$ of the function
\[ f(x) = \frac{\sqrt{2x+1}}{x^2 + x} \]
\end{question}

\begin{hint}
The volume of the solid of revolution is $V = \pi\int_a^b f^2(x)\,dx$. Note that $(x^2 + x)' = 2x + 1$.
\end{hint}

\begin{solution}
The volume of the solid obtained by revolving the graph of $f$ about the $x$-axis over $[a,b]$ is $V = \pi\int_a^b f(x)^2\,dx$. Here
$f(x)^2 = \frac{2x+1}{(x^2 + x)^2}$, which is continuous on $[1,4]$. Substitute $u = x^2 + x$, $du = (2x + 1)\,dx$; as $x$ goes from $1$ to $4$, $u$ goes from $2$ to $20$:
\[ V = \pi\int_1^4\frac{2x+1}{(x^2+x)^2}\,dx = \pi\int_2^{20}\frac{du}{u^2} = \pi\left[-\frac1u\right]_2^{20} = \pi\left(\frac12 - \frac1{20}\right) = \frac{9\pi}{20}. \]
\textbf{Answer:} $V = \frac{9\pi}{20}$.
\end{solution}

\subsection*{Question 5b}

\begin{question}
Compute the following improper integral, or show that it does not converge:
\[ \int_1^\infty \frac{\ln x}{x}\,dx \]
\end{question}

\begin{hint}
Find an antiderivative using the substitution $u = \ln x$.
\end{hint}

\begin{solution}
By definition, $\int_1^\infty\frac{\ln x}{x}dx = \lim_{R\to\infty}\int_1^R\frac{\ln x}{x}dx$. With the substitution $u = \ln x$, $du = \frac{dx}{x}$:
\[ \int_1^R\frac{\ln x}{x}\,dx = \int_0^{\ln R}u\,du = \frac{(\ln R)^2}{2} \xrightarrow[R\to\infty]{} \infty. \]
(Alternatively, by the comparison test: for $x \ge e$ we have $\frac{\ln x}{x} \ge \frac1x \ge 0$, and $\int_e^\infty\frac{dx}x$ diverges.)

\textbf{Answer:} the integral does not converge (it diverges to $\infty$).
\end{solution}

\section*{2024/25 Semester A Special Moed}

\subsection*{Question 1a}

\begin{question}
Compute the limit
\[ \lim_{x\to\infty}\quad x\cdot\left(\sqrt{x^2 + 4} - \sqrt{x^2 - 4}\right) \]
\end{question}

\begin{hint}
Multiply and divide by the conjugate $\sqrt{x^2+4} + \sqrt{x^2-4}$.
\end{hint}

\begin{solution}
For $x \ge 2$, multiply and divide by the conjugate (which is positive):
\[ x\left(\sqrt{x^2+4} - \sqrt{x^2-4}\right) = x\cdot\frac{(x^2 + 4) - (x^2 - 4)}{\sqrt{x^2+4} + \sqrt{x^2-4}} = \frac{8x}{\sqrt{x^2+4} + \sqrt{x^2-4}}. \]
Divide the numerator and the denominator by $x$ (for $x>0$, $\sqrt{x^2 + 4} = x\sqrt{1 + 4/x^2}$):
\[ = \frac{8}{\sqrt{1 + \frac4{x^2}} + \sqrt{1 - \frac4{x^2}}} \xrightarrow[x\to\infty]{} \frac{8}{1 + 1} = 4, \]
by the arithmetic of limits and continuity of the square root.

\textbf{Answer:} the limit equals $4$.
\end{solution}

\subsection*{Question 1b}

\begin{question}
Define a function by
\[ f(x) = \begin{cases} x^3\cdot\cos(\frac1x) & x \neq 0 \\ 0 & x = 0 \end{cases} \]
Is the function differentiable at $x = 0$? \\
If so, compute its derivative at this point. If not, justify.
\end{question}

\begin{hint}
Use the definition of the derivative and the fact that $|\cos(\frac1x)| \le 1$.
\end{hint}

\begin{solution}
By the definition of the derivative, for $h \neq 0$:
\[ \frac{f(h) - f(0)}{h} = \frac{h^3\cos(\frac1h)}{h} = h^2\cos\left(\tfrac1h\right). \]
Since $|\cos(\frac1h)| \le 1$, we have $0 \le \left|h^2\cos(\frac1h)\right| \le h^2 \to 0$, and hence by the Squeeze theorem (a bounded sequence/function times one tending to zero)
\[ \lim_{h\to 0}\frac{f(h) - f(0)}{h} = 0. \]
\textbf{Answer:} yes, $f$ is differentiable at $x=0$ and $f'(0) = 0$.
\end{solution}

\subsection*{Question 2}

\begin{question}
Define a function
\[ f(x) = \frac{x}{x^4 + 3} \]
\begin{enumerate}
  \item (11 pts) Compute the local extremum points and the intervals of increase and decrease of the function on its domain. Justify your answer.
  \item (10 pts) Compute all asymptotes of the function. Justify your answer.
\end{enumerate}
\end{question}

\begin{solution}
Domain: $x^4 + 3 \ge 3 > 0$ for every $x$, so $f$ is defined, continuous and differentiable on all of $\R$.
\begin{enumerate}
  \item By the quotient rule:
  \[ f'(x) = \frac{(x^4 + 3) - x\cdot 4x^3}{(x^4+3)^2} = \frac{3 - 3x^4}{(x^4 + 3)^2} = \frac{3(1 - x^2)(1 + x^2)}{(x^4+3)^2}. \]
  The denominator and $1 + x^2$ are positive, so the sign of $f'$ is the sign of $1 - x^2$:
  \begin{itemize}
    \item $|x| < 1$: $f' > 0$, $f$ is increasing on $[-1, 1]$;
    \item $|x| > 1$: $f' < 0$, $f$ is decreasing on $(-\infty, -1]$ and on $[1, \infty)$.
  \end{itemize}
  By the first derivative test: $x=-1$ is a local minimum point, $f(-1) = -\frac14$; $x = 1$ is a local maximum point, $f(1) = \frac14$.
  (Since $f\to0$ at infinity, these are also the absolute extrema.)

  \item \textbf{Vertical:} $f$ is continuous on all of $\R$, so at every point $x_0$ we have $\lim_{x\to x_0}f(x) = f(x_0)$, which is finite, and there are no vertical asymptotes.

  \textbf{Horizontal/oblique:}
  \[ \lim_{x\to\pm\infty}\frac{x}{x^4 + 3} = \lim_{x\to\pm\infty}\frac{1/x^3}{1 + 3/x^4} = 0. \]
  Hence $y = 0$ is a horizontal asymptote both at $+\infty$ and at $-\infty$ (so there is no additional oblique asymptote: $m = \lim \frac{f(x)}{x} = 0$, $n = \lim f(x) = 0$).

  \textbf{Answer:} the only asymptote is $y = 0$ (horizontal, in both directions); there are no vertical asymptotes.
\end{enumerate}
\end{solution}

\subsection*{Question 3}

\begin{question}
Using a linear approximation, compute $\sqrt[3]{10}$ and estimate the error. Justify your answer.
\end{question}

\begin{hint}
Choose a point close to $10$ at which the cube root is known: $x_0 = 8$. Estimate the error using the Lagrange remainder.
\end{hint}

\begin{solution}
Define $f(x) = x^{1/3}$ and choose $x_0 = 8$ (since $\sqrt[3]{8} = 2$). The derivatives:
\[ f'(x) = \tfrac13x^{-2/3}, \qquad f''(x) = -\tfrac29x^{-5/3}, \]
hence $f(8) = 2$, $f'(8) = \frac13\cdot\frac14 = \frac1{12}$.

\textbf{Linear approximation:} $P_1(x) = 2 + \frac{1}{12}(x - 8)$, and hence
\[ \sqrt[3]{10} \approx P_1(10) = 2 + \frac{2}{12} = \frac{13}{6} \approx 2.1667. \]

\textbf{Error estimate:} by Taylor's theorem with the Lagrange remainder, there exists $c \in (8, 10)$ such that
\[ \sqrt[3]{10} - \frac{13}{6} = R_1(10) = \frac{f''(c)}{2!}(10 - 8)^2 = -\frac19c^{-5/3}\cdot 4 = -\frac{4}{9}c^{-5/3}. \]
Since $c > 8$, we have $c^{-5/3} < 8^{-5/3} = \frac1{32}$, and hence
\[ |R_1(10)| < \frac49\cdot\frac{1}{32} = \frac{1}{72} \approx 0.0139, \]
and the error is negative, i.e. the approximation is larger than the true value.

\textbf{Answer:} $\sqrt[3]{10} \approx \frac{13}{6} \approx 2.1667$, with error less than $\frac{1}{72}$; in particular $\frac{13}{6} - \frac1{72} < \sqrt[3]{10} < \frac{13}{6}$.
Check: $\sqrt[3]{10} \approx 2.1544$, and the actual error is about $0.0122$.
\end{solution}

\subsection*{Question 4a}

\begin{question}
Compute the integral
\[ \int\frac{7x^2 + 6}{x^3 + 2x^2 + 2x}\,dx \]
\end{question}

\begin{hint}
$x^3 + 2x^2 + 2x = x\left((x+1)^2 + 1\right)$. Decompose into partial fractions.
\end{hint}

\begin{solution}
\textbf{Factoring the denominator.} $x^3 + 2x^2 + 2x = x(x^2 + 2x + 2)$, and $x^2 + 2x + 2 = (x+1)^2 + 1$ has no real roots.

\textbf{Partial fractions.}
\[ \frac{7x^2 + 6}{x(x^2 + 2x + 2)} = \frac{A}{x} + \frac{Bx + C}{x^2 + 2x + 2}, \qquad 7x^2 + 6 = A(x^2 + 2x + 2) + x(Bx + C). \]
Substituting $x = 0$: $6 = 2A$, i.e. $A = 3$. Then $7x^2 + 6 = 3x^2 + 6x + 6 + Bx^2 + Cx$, and comparing coefficients gives $B = 4$, $C = -6$.

\textbf{Integration.} Write $4x - 6 = 2(2x + 2) - 10$, where $(x^2 + 2x + 2)' = 2x + 2$:
\[ \int\frac{4x - 6}{x^2 + 2x + 2}\,dx = 2\int\frac{2x+2}{x^2+2x+2}\,dx - 10\int\frac{dx}{(x+1)^2 + 1} = 2\ln(x^2 + 2x + 2) - 10\arctan(x + 1). \]

\textbf{Answer:}
\[ \int\frac{7x^2 + 6}{x^3 + 2x^2 + 2x}\,dx = 3\ln|x| + 2\ln(x^2 + 2x + 2) - 10\arctan(x+1) + C. \]
(Check: differentiating the result gives back the integrand.)
\end{solution}

\subsection*{Question 4b}

\begin{question}
Find the general solution of the differential equation
\[ y' - xy^2 = 4x \]
\end{question}

\begin{hint}
Rearrange: $y' = x(y^2 + 4)$ --- this is a separable equation.
\end{hint}

\begin{solution}
Rearrange: $y' = xy^2 + 4x = x(y^2 + 4)$. This is a separable equation. Since $y^2 + 4 > 0$ always, there are no constant solutions (for $y\equiv c$: $0 = x(c^2+4)$ does not hold for every $x$), and we may divide:
\[ \frac{y'}{y^2 + 4} = x \quad\Longrightarrow\quad \int\frac{dy}{y^2 + 4} = \int x\,dx \quad\Longrightarrow\quad \frac12\arctan\frac y2 = \frac{x^2}{2} + c. \]
Multiply by $2$ and set $C = 2c$: $\arctan\frac y2 = x^2 + C$, and hence
\[ y = 2\tan\left(x^2 + C\right), \]
on the region where $x^2 + C \in (-\frac\pi2, \frac\pi2)$ (since $\arctan$ takes values only in this interval).

\textbf{Check:} $y' = 2\cdot 2x\left(1 + \tan^2(x^2 + C)\right) = 4x + x\cdot 4\tan^2(x^2+C) = 4x + xy^2$. It holds.

\textbf{Answer:} $y(x) = 2\tan(x^2 + C)$, $C \in \R$.
\end{solution}

\subsection*{Question 5a}

\begin{question}
Compute the finite area enclosed between the graph of the function $f(x) = xe^x$ and the line $y = e\cdot x$.
\end{question}

\begin{hint}
First find the intersection points: $xe^x = ex \iff x(e^x - e) = 0$. Compute $\int xe^x\,dx$ by integration by parts.
\end{hint}

\begin{solution}
\textbf{Intersection points:} $xe^x = ex \iff x(e^x - e) = 0 \iff x = 0$ or $e^x = e$, i.e. $x = 0$ or $x = 1$.

\textbf{Which graph is on top:} $ex - xe^x = x(e - e^x)$. For $0 < x < 1$: $x > 0$ and $e^x < e$, hence $ex > xe^x$ --- the line lies above the graph.
(For $x < 0$ the two functions do not intersect again, and the region between them is unbounded with infinite area, since $ex \to -\infty$ while $xe^x \to 0$; therefore ``the finite area'' is the one between $0$ and $1$.)

\textbf{Computation:} by integration by parts ($u = x$, $dv = e^xdx$): $\int xe^x\,dx = xe^x - e^x + C$. Hence
\[ S = \int_0^1\left(ex - xe^x\right)dx = \left[\frac{ex^2}{2}\right]_0^1 - \left[xe^x - e^x\right]_0^1 = \frac e2 - \left((e - e) - (0 - 1)\right) = \frac e2 - 1. \]
\textbf{Answer:} $S = \frac e2 - 1 \approx 0.359$.
\end{solution}

\subsection*{Question 5b}

\begin{question}
Compute the following improper integral, or show that it does not converge:
\[ \int_e^\infty\frac{1}{x\ln^2 x}\,dx \]
\end{question}

\begin{hint}
Substitute $u = \ln x$.
\end{hint}

\begin{solution}
By definition, $\int_e^\infty = \lim_{R\to\infty}\int_e^R$. With the substitution $u = \ln x$, $du = \frac{dx}{x}$ ($x = e \mapsto u = 1$, $x = R\mapsto u = \ln R$):
\[ \int_e^R\frac{dx}{x\ln^2x} = \int_1^{\ln R}\frac{du}{u^2} = \left[-\frac1u\right]_1^{\ln R} = 1 - \frac{1}{\ln R} \xrightarrow[R\to\infty]{} 1, \]
since $\ln R \to \infty$.

\textbf{Answer:} the integral converges and its value is $1$.
\end{solution}

\section*{2025/26 Semester A Quiz}

\subsection*{Part A, Question 1}

\begin{question}
Define a function by
\[
f(x)=\begin{cases}
x^3 & x<1\\
2x+1 & x\ge 1
\end{cases}
\]
\begin{enumerate}
  \item (5 pts) Sketch the graph of the function.
  \item (15 pts) Is this function invertible as a function from $\R$ to $\R$?
  If so, compute its inverse function $f^{-1}:\R\to\R$.
  If not, justify.
\end{enumerate}
\end{question}

\begin{hint}
A function that is invertible from $\R$ to $\R$ must be both one-to-one and onto $\R$. What is the image of each of the two pieces?
\end{hint}

\begin{solution}
\begin{enumerate}
  \item \textbf{The graph:} for $x<1$ this is the graph of the cubic parabola $y=x^3$: strictly increasing, passing through the points $(-1,-1)$, $(0,0)$ (an inflection point with a horizontal tangent), and approaching the point $(1,1)$ --- this point does \emph{not} belong to the graph (open circle). For $x\ge1$ it is a ray with slope 2 starting at the point $(1,3)$ (filled circle) and passing through $(2,5)$. At the point $x=1$ there is a ``jump'' from height 1 to height 3.
  \item \textbf{The function is not invertible as a function from $\R$ to $\R$}, since it is not onto $\R$.

  We find the image. For $x<1$: $x^3$ is strictly increasing, hence $x^3<1^3=1$; i.e. the values of the first piece are less than 1 (and in fact cover $(-\infty,1)$). For $x\ge1$: $2x+1\ge3$; i.e. the values of the second piece are greater than or equal to 3 (and cover $[3,\infty)$). Therefore
  \[
  \operatorname{Im}f=(-\infty,1)\cup[3,\infty).
  \]
  In particular, there is no $x\in\R$ with $f(x)=2$: if $x<1$ then $f(x)<1$, and if $x\ge1$ then $f(x)\ge3$. Hence $f$ is not onto $\R$, and therefore it has no inverse $f^{-1}:\R\to\R$.

  (Remark: $f$ is one-to-one --- it is strictly increasing, since each piece is strictly increasing and every value of the first piece is less than every value of the second piece. Hence $f$ is invertible as a function from $\R$ onto its image $(-\infty,1)\cup[3,\infty)$, with $f^{-1}(y)=\sqrt[3]{y}$ for $y<1$ and $f^{-1}(y)=\frac{y-1}{2}$ for $y\ge3$; but not as a function from $\R$ to $\R$.)
\end{enumerate}
\end{solution}

\subsection*{Part A, Question 2}

\begin{question}
Let $a,b\in\R$. Define a function by
\[
f(x)=\begin{cases}
\frac{x^2-4x+3}{x-1} & x<1\\
a & x=1\\
2b\cdot\frac{\ln x}{x-1} & x>1
\end{cases}
\]
For which values of $a,b$ is the function continuous at $1$ ? Justify your answer.
\end{question}

\begin{hint}
For the left limit, factor the numerator. For the right limit, substitute $y=x-1$ and use the basic limit $\frac{\ln(1+y)}{y}\to1$.
\end{hint}

\begin{solution}
The function is continuous at 1 if and only if the one-sided limits at 1 exist and are equal to $f(1)=a$.

\textbf{The left limit.} $x^2-4x+3=(x-1)(x-3)$, and for $x\neq1$ we may cancel:
\[
\lim_{x\to1^-}f(x)=\lim_{x\to1^-}\frac{(x-1)(x-3)}{x-1}=\lim_{x\to1^-}(x-3)=1-3=-2 ,
\]
where in the last equality we used the continuity of the polynomial $x-3$.

\textbf{The right limit.} By the arithmetic of limits (the constant $2b$ comes out of the limit), and with the substitution $y=x-1$ ($y\to0^+$ as $x\to1^+$):
\[
\lim_{x\to1^+}f(x)=2b\cdot\lim_{x\to1^+}\frac{\ln x}{x-1}=2b\cdot\lim_{y\to0^+}\frac{\ln(1+y)}{y}=2b\cdot1=2b ,
\]
by the basic limit $\lim_{y\to0}\frac{\ln(1+y)}{y}=1$ (from the formula sheet: $\lim_{y\to0}\frac{\log_a(1+y)}{y}=\frac{1}{\ln a}$ with $a=e$).

\textbf{Conclusion.} $f$ is continuous at 1 if and only if $-2=a=2b$, i.e.
\[
\boxed{a=-2,\quad b=-1.}
\]
\end{solution}

\subsection*{Part A, Question 3}

\begin{question}
Prove that the equation
\[
x^5+e^x=0
\]
has a real solution, and find an interval of length at most 1 that contains the solution. Justify your answer.
\end{question}

\begin{hint}
Define $f(x)=x^5+e^x$ and check its sign at integer points close to 0.
\end{hint}

\begin{solution}
Define $f(x)=x^5+e^x$. $f$ is continuous on $\R$ as the sum of a polynomial and an exponential function, both of which are continuous.

We compute:
\[
f(0)=0^5+e^0=1>0,\qquad f(-1)=(-1)^5+e^{-1}=-1+\frac1e<0 ,
\]
since $e>1$ and hence $\frac1e<1$.

$f$ is continuous on the closed interval $[-1,0]$ and has opposite signs at its endpoints. By the Intermediate Value Theorem (Bolzano's theorem) there exists $c\in(-1,0)$ such that $f(c)=0$, i.e. $c^5+e^c=0$, and $c$ is a real solution of the equation.

\textbf{The required interval:} $[-1,0]$, whose length is 1.

(Remark: the solution is unique --- $f$ is strictly increasing as the sum of the two strictly increasing functions $x^5$ and $e^x$ --- so ``the solution'' is well defined. Numerically, $c\approx-0.845$.)
\end{solution}

\subsection*{Part B, Question 4}

\begin{question}
Let
\[
f(x)=\sqrt[3]{x},\qquad g(x)=\cos(x)-\sin(x)
\]
Which of the following statements is true?
\begin{enumerate}
  \item The function $f\circ g$ is even.
  \item The function $f\circ g$ is odd.
  \item The function $f\circ g$ is monotone.
  \item The function $f\circ g$ is periodic.
\end{enumerate}
\end{question}

\begin{solution}
\textbf{The correct answer: (d)}.

$(f\circ g)(x)=f(g(x))=\sqrt[3]{\cos x-\sin x}$, defined for every $x\in\R$ (the cube root is defined for every real number).

\textbf{(d) is true:} $\sin$ and $\cos$ are periodic with period $2\pi$, hence for every $x$:
\[
(f\circ g)(x+2\pi)=\sqrt[3]{\cos(x+2\pi)-\sin(x+2\pi)}=\sqrt[3]{\cos x-\sin x}=(f\circ g)(x).
\]

\textbf{Why the others are false:} denote $h=f\circ g$.
\begin{itemize}
  \item (a) --- $h(-x)=\sqrt[3]{\cos x+\sin x}$. For example, $h\!\left(\frac\pi4\right)=\sqrt[3]{0}=0$ but $h\!\left(-\frac\pi4\right)=\sqrt[3]{\sqrt2}\neq0$, so $h$ is not even.
  \item (b) --- an odd function defined at 0 satisfies $h(0)=-h(0)$, i.e. $h(0)=0$; but $h(0)=\sqrt[3]{1}=1\neq0$.
  \item (c) --- a non-constant periodic function is not monotone: for example, $h(0)=1$, $h\!\left(\frac\pi4\right)=0$, $h(2\pi)=1$, i.e. $h(0)>h(\frac\pi4)<h(2\pi)$ --- the function decreases and then increases.
\end{itemize}
\end{solution}

\subsection*{Part B, Question 5}

\begin{question}
Compute the limit
\[
\lim_{x\to\infty}\left(\frac{x^2+2x}{3x^2+7x}\right)^{x^2}
\]
\begin{enumerate}
  \item The limit exists in the wide sense and equals $\infty$.
  \item The limit exists and equals $e$.
  \item The limit exists and equals $0$.
  \item The limit exists and equals $1$.
\end{enumerate}
\end{question}

\begin{hint}
What does the base tend to? This is \emph{not} a limit of the form $1^\infty$.
\end{hint}

\begin{solution}
\textbf{The correct answer: (c)}.

The base tends to $\frac13$ (divide by $x^2$), not to 1, so this is not a limit of the form $1^\infty$. We use the Squeeze theorem. For $x>0$:
\[
0\le\frac{x^2+2x}{3x^2+7x}\le\frac{x^2+2x}{3x^2+6x}=\frac{x(x+2)}{3x(x+2)}=\frac13 ,
\]
since the numerator is positive and decreasing the positive denominator ($7x>6x$) increases the fraction. Since the function $t\mapsto t^{x^2}$ is increasing on $[0,\infty)$ (for $x^2>0$):
\[
0\le\left(\frac{x^2+2x}{3x^2+7x}\right)^{x^2}\le\left(\frac13\right)^{x^2}.
\]
As $x\to\infty$ also $x^2\to\infty$, and since $0<\frac13<1$ we have $\left(\frac13\right)^{x^2}\to0$. By the Squeeze theorem
\[
\lim_{x\to\infty}\left(\frac{x^2+2x}{3x^2+7x}\right)^{x^2}=\boxed{0}.
\]

\textbf{Why the others are false:} (a) --- the expression is bounded by $\left(\frac13\right)^{x^2}\le1$, so it does not tend to $\infty$. (b) and (d) --- these values come from wrongly treating the expression as if it were of the form $1^\infty$; here the base tends to $\frac13<1$ and the limit is 0.
\end{solution}

\subsection*{Part B, Question 6}

\begin{question}
Let $f(x)=x+e^x$. Which of the following statements is true?
\begin{enumerate}
  \item The function is one-to-one and onto $\R$.
  \item The function is one-to-one but not onto $\R$.
  \item The function is onto $\R$ but not one-to-one.
  \item The function is neither one-to-one nor onto $\R$.
\end{enumerate}
\end{question}

\begin{solution}
\textbf{The correct answer: (a)}.

\textbf{One-to-one:} let $x_1<x_2$. Since $e^x$ is strictly increasing, $e^{x_1}<e^{x_2}$. Adding the two inequalities we get $x_1+e^{x_1}<x_2+e^{x_2}$, i.e. $f$ is strictly increasing, and in particular one-to-one.

\textbf{Onto $\R$:} $f$ is continuous on $\R$ as a sum of continuous functions, and
\[
\lim_{x\to-\infty}(x+e^x)=-\infty+0=-\infty,\qquad \lim_{x\to\infty}(x+e^x)=\infty+\infty=\infty .
\]
Let $y\in\R$. By these limits there exist $x_1<x_2$ with $f(x_1)<y<f(x_2)$, and by the Intermediate Value Theorem on $[x_1,x_2]$ there exists $c$ with $f(c)=y$. Hence $f$ is onto $\R$.

\textbf{Why the others are false:} (b) and (d) claim that $f$ is not onto, (c) and (d) claim that it is not one-to-one --- both claims were refuted above.
\end{solution}

\subsection*{Part B, Question 7}

\begin{question}
Check the continuity of the following function at the point 3:
\[
f(x)=\begin{cases}
\frac{\sqrt{x+7}-\sqrt{10}}{\sqrt{x+4}-\sqrt{7}} & x>3\\
\sqrt{x^2+1} & x\le 3
\end{cases}
\]
Which of the following statements is true?
\begin{enumerate}
  \item The function is continuous at 3.
  \item The function has a removable discontinuity at 3.
  \item The function has a jump discontinuity at 3.
  \item The function has an essential discontinuity at 3.
\end{enumerate}
\end{question}

\begin{hint}
For the right limit, multiply the numerator and the denominator by the conjugates of both expressions: $\sqrt{x+7}+\sqrt{10}$ and $\sqrt{x+4}+\sqrt7$.
\end{hint}

\begin{solution}
\textbf{The correct answer: (c)}.

\textbf{From the left:} $\sqrt{x^2+1}$ is continuous (a composition of a polynomial and a square root), hence
\[
\lim_{x\to3^-}f(x)=\sqrt{3^2+1}=\sqrt{10}=f(3).
\]
\textbf{From the right:} a limit of the form $\frac00$. Multiply by the conjugates:
\[
\frac{\sqrt{x+7}-\sqrt{10}}{\sqrt{x+4}-\sqrt7}\cdot\frac{\sqrt{x+4}+\sqrt7}{\sqrt{x+4}+\sqrt7}\cdot\frac{\sqrt{x+7}+\sqrt{10}}{\sqrt{x+7}+\sqrt{10}}
=\frac{\bigl((x+7)-10\bigr)\left(\sqrt{x+4}+\sqrt7\right)}{\bigl((x+4)-7\bigr)\left(\sqrt{x+7}+\sqrt{10}\right)}
=\frac{\sqrt{x+4}+\sqrt7}{\sqrt{x+7}+\sqrt{10}},
\]
where we cancelled $x-3\neq0$. The last expression is continuous at 3, hence
\[
\lim_{x\to3^+}f(x)=\frac{\sqrt7+\sqrt7}{\sqrt{10}+\sqrt{10}}=\frac{2\sqrt7}{2\sqrt{10}}=\frac{\sqrt7}{\sqrt{10}} .
\]
Both one-sided limits exist and are finite, but $\frac{\sqrt7}{\sqrt{10}}\neq\sqrt{10}$. Hence the limit at 3 does not exist, and there is a jump discontinuity at 3.

\textbf{Why the others are false:} (a) --- the limit at 3 does not exist, so the function is not continuous there (it is only left-continuous). (b) --- a removable discontinuity requires the two-sided limit to exist. (d) --- an essential discontinuity requires at least one of the one-sided limits not to exist as a finite limit, and here both exist.
\end{solution}

\section*{2025/26 Semester A Moed A}

\subsection*{Question 1a}

\begin{question}
How many real solutions does the following equation have?
\[ x^4 + 2x^2 = 1 \]
Justify your answer carefully.
\end{question}

\begin{hint}
Substitute $t = x^2$ to get a quadratic equation in $t$; remember that $t = x^2$ must be non-negative.
\end{hint}

\begin{solution}
\textbf{Method 1 (substitution).} Move everything to one side: $x^4 + 2x^2 - 1 = 0$. Substitute $t = x^2 \ge 0$ to get
\[ t^2 + 2t - 1 = 0 \quad\Longrightarrow\quad t_{1,2} = \frac{-2 \pm \sqrt{4 + 4}}{2} = -1 \pm \sqrt{2}. \]
The solution $t = -1 - \sqrt{2} < 0$ is rejected, since $x^2 \ge 0$ for every real $x$, so it gives no real solution.
The solution $t = -1 + \sqrt{2} > 0$ (since $\sqrt 2 > 1$) gives
\[ x^2 = \sqrt{2} - 1 \quad\Longrightarrow\quad x_{1,2} = \pm\sqrt{\sqrt{2} - 1}, \]
two distinct real numbers.

\textbf{Method 2 (function analysis).} Define $f(x) = x^4 + 2x^2 - 1$, a function that is continuous and differentiable on all of $\R$. We have
\[ f'(x) = 4x^3 + 4x = 4x(x^2 + 1), \]
and since $x^2 + 1 > 0$, the sign of $f'$ is the sign of $x$: $f' < 0$ on $(-\infty, 0)$ and $f' > 0$ on $(0, \infty)$.
Hence $f$ is strictly decreasing on $(-\infty, 0]$ and strictly increasing on $[0, \infty)$, so it has at most one zero in each of these intervals.
Now $f(0) = -1 < 0$ and $f(1) = f(-1) = 2 > 0$. By the Intermediate Value Theorem, $f$ has a zero in $(-1, 0)$ and a zero in $(0, 1)$, and by monotonicity these are its only zeros.

\textbf{Answer:} the equation has exactly \textbf{two} real solutions, $x = \pm\sqrt{\sqrt 2 - 1}$.
\end{solution}

\subsection*{Question 1b}

\begin{question}
Compute the following limit
\[ \lim_{x\to 0} \frac{x - \sin x}{\cos x - 1} \]
\end{question}

\begin{hint}
Identify the type of the expression ($\frac{0}{0}$) and use L'Hôpital's rule.
\end{hint}

\begin{solution}
As $x \to 0$ the numerator $x - \sin x \to 0$ and the denominator $\cos x - 1 \to 0$, i.e. an expression of the form $\frac{0}{0}$.
The numerator and the denominator are differentiable in a neighborhood of $0$, and the derivative of the denominator satisfies $-\sin x \ne 0$ in a punctured neighborhood of $0$. By L'Hôpital's rule (provided the new limit exists):
\[ \lim_{x\to 0} \frac{x - \sin x}{\cos x - 1} = \lim_{x\to 0} \frac{1 - \cos x}{-\sin x}. \]
This is also a $\frac{0}{0}$ expression. We compute it using $1 - \cos x = 2\sin^2\frac{x}{2}$:
\[ \frac{1 - \cos x}{-\sin x} = \frac{2\sin^2\frac{x}{2}}{-\sin x}
   = -\frac{2\left(\dfrac{\sin\frac{x}{2}}{\frac{x}{2}}\right)^2 \cdot \dfrac{x^2}{4}}{\dfrac{\sin x}{x} \cdot x}
   = -\frac{\left(\dfrac{\sin\frac{x}{2}}{\frac{x}{2}}\right)^2}{\dfrac{\sin x}{x}} \cdot \frac{x}{2}
   \xrightarrow[x\to 0]{} -\frac{1^2}{1}\cdot 0 = 0, \]
where we used the basic limit $\lim_{u\to 0}\frac{\sin u}{u} = 1$ and the arithmetic of limits.
(Alternatively, applying L'Hôpital once more: $\lim_{x\to0}\frac{\sin x}{-\cos x} = \frac{0}{-1} = 0$.)
Since the limit after L'Hôpital exists, the original limit also exists and is equal to it:
\[ \boxed{\lim_{x\to 0} \frac{x - \sin x}{\cos x - 1} = 0}. \]
\end{solution}

\subsection*{Question 2a}

\begin{question}
Compute all asymptotes of the function
\[ f(x) = \frac{x^2 - 2\ln(x)}{x - 1} \]
\end{question}

\begin{hint}
Start from the domain: the ``suspicious'' points for a vertical asymptote are the endpoints of the domain and the points where the denominator vanishes. An oblique asymptote $y = ax + b$ needs to be checked only in a direction in which the domain is unbounded.
\end{hint}

\begin{solution}
\textbf{Domain:} $\ln x$ is defined for $x > 0$, and the denominator vanishes at $x = 1$. Hence the domain is $(0,1)\cup(1,\infty)$, and $f$ is continuous on it.

\textbf{Vertical asymptotes.} The suspicious points are $x = 0$ (an endpoint of the domain) and $x = 1$.
\begin{itemize}
  \item $x \to 0^+$: the numerator $x^2 - 2\ln x \to 0 - (-\infty) = +\infty$ and the denominator $x - 1 \to -1$, hence
  \[ \lim_{x\to 0^+} f(x) = -\infty. \]
  Therefore $x = 0$ is a vertical asymptote (from the right).
  \item $x \to 1$: the numerator $\to 1 - 2\ln 1 = 1 > 0$ and the denominator $\to 0$, where $x - 1 < 0$ on the left and $x - 1 > 0$ on the right. Hence
  \[ \lim_{x\to 1^-} f(x) = -\infty, \qquad \lim_{x\to 1^+} f(x) = +\infty, \]
  and $x = 1$ is a vertical asymptote (from both sides).
\end{itemize}

\textbf{Horizontal/oblique asymptotes.} The domain is unbounded only in the direction $+\infty$, so we check only $x \to +\infty$. We look for $y = ax + b$:
\[ a = \lim_{x\to\infty} \frac{f(x)}{x} = \lim_{x\to\infty} \frac{x^2 - 2\ln x}{x^2 - x}
     = \lim_{x\to\infty} \frac{x^2\left(1 - \frac{2\ln x}{x^2}\right)}{x^2\left(1 - \frac{1}{x}\right)} = \frac{1 - 0}{1 - 0} = 1, \]
since $\frac{\ln x}{x^2} \to 0$ (for example by L'Hôpital: $\lim \frac{1/x}{2x} = 0$).
\[ b = \lim_{x\to\infty} \bigl(f(x) - x\bigr) = \lim_{x\to\infty} \frac{x^2 - 2\ln x - x^2 + x}{x - 1}
     = \lim_{x\to\infty} \frac{x\left(1 - \frac{2\ln x}{x}\right)}{x\left(1 - \frac{1}{x}\right)} = 1, \]
since $\frac{\ln x}{x} \to 0$. Since both limits are finite, the line $y = x + 1$ is an oblique asymptote at $+\infty$ (and there is no horizontal asymptote).

\textbf{Answer:} vertical asymptotes $x = 0$ and $x = 1$; oblique asymptote $y = x + 1$ as $x \to +\infty$.
\end{solution}

\subsection*{Question 2b}

\begin{question}
Find the absolute minimum and maximum of the function
\[ f(x) = \frac{\sqrt{x}}{e^x} \]
on its domain, or explain why they do not exist.
\end{question}

\begin{hint}
Find the critical points, determine the intervals of increase and decrease, and check the behavior of $f$ at the ends of the domain ($x = 0$ and $x \to \infty$).
\end{hint}

\begin{solution}
\textbf{Domain:} $x \ge 0$ (because of $\sqrt{x}$). $f$ is continuous on $[0,\infty)$, $f(x) \ge 0$ for every $x$, and $f(0) = 0$.

\textbf{Derivative} (for $x > 0$):
\[ f'(x) = \frac{\frac{1}{2\sqrt{x}}\, e^x - \sqrt{x}\, e^x}{e^{2x}} = \frac{e^x(1 - 2x)}{2\sqrt{x}\, e^{2x}} = \frac{1 - 2x}{2\sqrt{x}\, e^x}. \]
The denominator is positive, hence $f'(x) = 0 \iff x = \frac12$, $f' > 0$ on $(0, \frac12)$ and $f' < 0$ on $(\frac12, \infty)$.
That is, $f$ is strictly increasing on $[0, \frac12]$ and strictly decreasing on $[\frac12, \infty)$ (including the endpoint, since $f$ is continuous there).

\textbf{Absolute maximum:} by monotonicity, $f(x) \le f(\frac12)$ for every $x \ge 0$, hence
\[ \max f = f\left(\tfrac12\right) = \frac{\sqrt{1/2}}{e^{1/2}} = \frac{1}{\sqrt{2e}}, \]
attained at $x = \frac12$.

\textbf{Absolute minimum:} $f(x) \ge 0 = f(0)$ for every $x$ in the domain (non-negative numerator and positive denominator), hence
\[ \min f = f(0) = 0, \]
attained at $x = 0$. (Note: although $\lim_{x\to\infty} f(x) = 0$, the value $0$ is actually attained at $x = 0$, so the minimum exists.)

\textbf{Answer:} absolute minimum $0$ at $x = 0$; absolute maximum $\frac{1}{\sqrt{2e}}$ at $x = \frac12$.
\end{solution}

\subsection*{Question 3}

\begin{question}
Using a second-order Taylor approximation, compute $\ln\left(\frac{6}{5}\right) = \ln(1.2)$, and estimate the error.

Note that the approximation and the error bound must be rational numbers.
\end{question}

\begin{hint}
Expand $f(x) = \ln x$ around $x_0 = 1$ (where the values of $f$ and its derivatives are rational) and use the second-order Lagrange remainder.
\end{hint}

\begin{solution}
Define $f(x) = \ln x$ and expand around $x_0 = 1$, a point close to $1.2$ at which all the values are known. $f$ is infinitely differentiable on $(0,\infty)$:
\[ f(x) = \ln x,\quad f'(x) = \frac{1}{x},\quad f''(x) = -\frac{1}{x^2},\quad f'''(x) = \frac{2}{x^3}, \]
hence $f(1) = 0$, $f'(1) = 1$, $f''(1) = -1$.

\textbf{The second-order Taylor polynomial around $1$:}
\[ P_2(x) = f(1) + f'(1)(x - 1) + \frac{f''(1)}{2!}(x - 1)^2 = (x - 1) - \frac{(x - 1)^2}{2}. \]
Hence
\[ \ln(1.2) \approx P_2(1.2) = 0.2 - \frac{0.04}{2} = 0.18 = \frac{9}{50}. \]

\textbf{Error estimate.} By Taylor's theorem with the Lagrange remainder, there exists a point $c$ between $1$ and $1.2$ such that
\[ R_2(1.2) = \ln(1.2) - P_2(1.2) = \frac{f'''(c)}{3!}(1.2 - 1)^3 = \frac{2}{6c^3}\cdot (0.2)^3 = \frac{0.008}{3c^3}. \]
Since $1 < c < 1.2$, we have $c^3 > 1$, hence $\frac{1}{c^3} < 1$ and
\[ 0 < R_2(1.2) < \frac{0.008}{3} = \frac{8}{3000} = \frac{1}{375}. \]

\textbf{Answer:} $\ln(1.2) \approx 0.18 = \frac{9}{50}$, and the error is less than $\frac{1}{375}$ (in particular, since the remainder is positive, $0.18 < \ln(1.2) < 0.18 + \frac{1}{375}$).

(Check: $\ln 1.2 \approx 0.18232$, and the actual error is about $0.0023 < \frac{1}{375} \approx 0.00267$.)
\end{solution}

\subsection*{Question 4a}

\begin{question}
Compute the integral
\[ \int \frac{-3x - 7}{x^3 + x^2 - 2}\,dx \]
Hint: note that the denominator vanishes at $x = 1$.
\end{question}

\begin{hint}
Divide the denominator by $(x - 1)$, check that the remaining quadratic factor is irreducible, and decompose into partial fractions of the form $\frac{A}{x-1} + \frac{Bx + C}{x^2 + 2x + 2}$.
\end{hint}

\begin{solution}
\textbf{Factoring the denominator.} $1^3 + 1^2 - 2 = 0$, hence $(x - 1)$ divides the denominator. Polynomial division gives
\[ x^3 + x^2 - 2 = (x - 1)(x^2 + 2x + 2). \]
The discriminant of $x^2 + 2x + 2$ is $4 - 8 < 0$, so the quadratic factor is irreducible over $\R$ (and in particular $x^2 + 2x + 2 = (x+1)^2 + 1 > 0$).

\textbf{Partial fractions.} The degree of the numerator is less than the degree of the denominator, so we look for
\[ \frac{-3x - 7}{(x - 1)(x^2 + 2x + 2)} = \frac{A}{x - 1} + \frac{Bx + C}{x^2 + 2x + 2}, \]
i.e. $-3x - 7 = A(x^2 + 2x + 2) + (Bx + C)(x - 1)$ for every $x$.
\begin{itemize}
  \item $x = 1$: $-10 = 5A$, hence $A = -2$.
  \item $x = 0$: $-7 = 2A - C$, hence $C = 2A + 7 = 3$.
  \item Comparing the coefficients of $x^2$: $0 = A + B$, hence $B = 2$.
\end{itemize}
Therefore
\[ \int \frac{-3x - 7}{x^3 + x^2 - 2}\,dx = \int \left( \frac{-2}{x - 1} + \frac{2x + 3}{x^2 + 2x + 2} \right) dx
   = \int \left( \frac{-2}{x - 1} + \frac{2x + 2}{x^2 + 2x + 2} + \frac{1}{(x + 1)^2 + 1} \right) dx. \]
\begin{itemize}
  \item $\int \frac{-2}{x - 1}\,dx = -2\ln|x - 1| + C$.
  \item In the second term the numerator is the derivative of the denominator: with the substitution $t = x^2 + 2x + 2$, $dt = (2x + 2)\,dx$,
  \[ \int \frac{2x + 2}{x^2 + 2x + 2}\,dx = \int \frac{dt}{t} = \ln|t| + C = \ln(x^2 + 2x + 2) + C. \]
  \item By completing the square and substituting $u = x + 1$: $\int \frac{dx}{(x + 1)^2 + 1} = \arctan(x + 1) + C$.
\end{itemize}
\textbf{Answer:}
\[ \boxed{\int \frac{-3x - 7}{x^3 + x^2 - 2}\,dx = -2\ln|x - 1| + \ln(x^2 + 2x + 2) + \arctan(x + 1) + C.} \]
(Check: differentiating the result gives back the integrand.)
\end{solution}

\subsection*{Question 4b}

\begin{question}
Find the general solution of the differential equation
\[ y' = y^2 x \sin x \]
\end{question}

\begin{hint}
This is a separable equation: move $y^2$ to the side of $dy$, and compute $\int x\sin x\,dx$ by integration by parts. Do not forget the solution $y \equiv 0$.
\end{hint}

\begin{solution}
This is a separable equation: $\frac{dy}{dx} = y^2 \cdot x\sin x$.

\textbf{The constant solution:} $y \equiv 0$ satisfies $y' = 0 = 0^2 \cdot x\sin x$, so it is a solution.

\textbf{For $y \ne 0$:} separate the variables and integrate:
\[ \int \frac{dy}{y^2} = \int x\sin x\,dx. \]
On the left-hand side, $\int y^{-2}\,dy = -\frac{1}{y}$. On the right-hand side, integrate by parts with $u = x$, $dv = \sin x\,dx$ (i.e. $du = dx$, $v = -\cos x$):
\[ \int x\sin x\,dx = -x\cos x + \int \cos x\,dx = -x\cos x + \sin x + C. \]
Hence
\[ -\frac{1}{y} = -x\cos x + \sin x + C \quad\Longrightarrow\quad \boxed{y = \frac{1}{x\cos x - \sin x - C}}, \]
where $C$ is an arbitrary real constant (on every interval on which the denominator does not vanish).

\textbf{Answer:} the general solution is $y = \dfrac{1}{x\cos x - \sin x + C}$ ($C \in \R$, we renamed the constant), together with the solution $y \equiv 0$.
\end{solution}

\subsection*{Question 5a}

\begin{question}
Compute the volume of the solid of revolution about the $x$-axis of the function
\[ f(x) = \sin x + \cos x \]
on the interval $[0, \frac{\pi}{2}]$.
\end{question}

\begin{hint}
Expand $(\sin x + \cos x)^2$ and use the identities $\sin^2 x + \cos^2 x = 1$ and $2\sin x\cos x = \sin 2x$.
\end{hint}

\begin{solution}
The volume of the solid obtained by revolving the graph of $f$ about the $x$-axis on the interval $[a,b]$ is $V = \pi\int_a^b f^2(x)\,dx$. Here
\[ f^2(x) = \sin^2 x + 2\sin x\cos x + \cos^2 x = 1 + \sin 2x, \]
and hence, by the Fundamental Theorem of Calculus (Newton--Leibniz),
\[ V = \pi\int_0^{\pi/2} (1 + \sin 2x)\,dx = \pi\left[ x - \frac{\cos 2x}{2} \right]_0^{\pi/2}
     = \pi\left[ \left(\frac{\pi}{2} - \frac{\cos\pi}{2}\right) - \left(0 - \frac{\cos 0}{2}\right) \right]
     = \pi\left[ \frac{\pi}{2} + \frac12 + \frac12 \right]. \]
\textbf{Answer:} $\boxed{V = \pi\left(\frac{\pi}{2} + 1\right) = \frac{\pi^2}{2} + \pi}$.
\end{solution}

\subsection*{Question 5b}

\begin{question}
Compute the following improper integral, or show that it diverges:
\[ \int_e^\infty \frac{dx}{x \cdot \ln^3 x} \]
\end{question}

\begin{hint}
Substitute $t = \ln x$.
\end{hint}

\begin{solution}
The integrand $\frac{1}{x\ln^3 x}$ is continuous and positive on $[e, \infty)$, so the integral is improper only because of the infinite upper limit. By definition
\[ \int_e^\infty \frac{dx}{x\ln^3 x} = \lim_{A\to\infty} \int_e^A \frac{dx}{x\ln^3 x}. \]
We find an antiderivative with the substitution $t = \ln x$, $dt = \frac{dx}{x}$:
\[ \int \frac{dx}{x\ln^3 x} = \int \frac{dt}{t^3} = -\frac{1}{2t^2} + C = -\frac{1}{2\ln^2 x} + C. \]
Hence, for $A > e$:
\[ \int_e^A \frac{dx}{x\ln^3 x} = \left[ -\frac{1}{2\ln^2 x} \right]_e^A = -\frac{1}{2\ln^2 A} + \frac{1}{2\ln^2 e} = \frac12 - \frac{1}{2\ln^2 A}. \]
As $A \to \infty$ we have $\ln A \to \infty$, and hence $\frac{1}{2\ln^2 A} \to 0$.

\textbf{Answer:} the integral converges and $\boxed{\int_e^\infty \frac{dx}{x\ln^3 x} = \frac12}$.
\end{solution}

\section*{2025/26 Semester A Moed B}

\subsection*{Question 1a}

\begin{question}
Compute the image of the function
\[ f(x) = x^3 - 12x + 8 \]
on the interval $[-3, 5]$. That is, compute the set
\[ \{ f(x) : -3 \le x \le 5 \} \]
Justify your answer carefully.
\end{question}

\begin{hint}
A continuous function on a closed interval attains a minimum and a maximum (Weierstrass) and every value between them (Intermediate Value Theorem). It remains to find the minimum and the maximum.
\end{hint}

\begin{solution}
$f$ is a polynomial, hence continuous on the closed interval $[-3, 5]$. By Weierstrass's second theorem, $f$ attains on the interval a minimum $c$ and a maximum $d$, and by the Intermediate Value Theorem it attains every value between $c$ and $d$. Hence the image is exactly the closed interval $[c, d]$, and it remains to compute $c$ and $d$.

The absolute extremum points of a differentiable function on a closed interval are at interior critical points or at the endpoints.
\[ f'(x) = 3x^2 - 12 = 3(x^2 - 4) = 0 \iff x = \pm 2, \]
and both points lie inside the interval. We compute the values of $f$ at the candidate points:
\[ f(-3) = -27 + 36 + 8 = 17,\quad f(-2) = -8 + 24 + 8 = 24,\quad f(2) = 8 - 24 + 8 = -8,\quad f(5) = 125 - 60 + 8 = 73. \]
Hence the minimum is $c = -8$ (at $x = 2$) and the maximum is $d = 73$ (at $x = 5$).

\textbf{Answer:} $\{ f(x) : -3 \le x \le 5 \} = [-8, 73]$.
\end{solution}

\subsection*{Question 1b}

\begin{question}
Compute the following limit
\[ \lim_{x\to\infty} \left( \frac{x^3 + 1}{x^3 - 2} \right)^{x^3 + 3} \]
\end{question}

\begin{hint}
This is an expression of the form $1^\infty$. Write $\frac{x^3 + 1}{x^3 - 2} = 1 + \frac{3}{x^3 - 2}$ and use the limit $\left(1 + \frac{1}{u}\right)^u \to e$.
\end{hint}

\begin{solution}
As $x \to \infty$ the base tends to $1$ and the exponent to $\infty$ (a $1^\infty$ expression). Write
\[ \frac{x^3 + 1}{x^3 - 2} = \frac{(x^3 - 2) + 3}{x^3 - 2} = 1 + \frac{3}{x^3 - 2}, \qquad x^3 + 3 = (x^3 - 2) + 5, \]
and hence (for large $x$, so that $x^3 - 2 > 0$)
\[ \left( \frac{x^3 + 1}{x^3 - 2} \right)^{x^3 + 3} = \left( 1 + \frac{3}{x^3 - 2} \right)^{x^3 - 2} \cdot \left( 1 + \frac{3}{x^3 - 2} \right)^{5}. \]
Set $u = \frac{x^3 - 2}{3} \to \infty$. Then the first factor equals $\left[\left(1 + \frac{1}{u}\right)^{u}\right]^3 \to e^3$, by the basic limit $\lim_{u\to\infty}\left(1 + \frac1u\right)^u = e$ and continuity of $t \mapsto t^3$.
The second factor tends to $1^5 = 1$, since $\frac{3}{x^3 - 2} \to 0$.
By the arithmetic of limits:
\[ \boxed{\lim_{x\to\infty} \left( \frac{x^3 + 1}{x^3 - 2} \right)^{x^3 + 3} = e^3 \cdot 1 = e^3}. \]
(Alternative method: the $\ln$ of the expression is $(x^3 + 3)\ln\left(1 + \frac{3}{x^3 - 2}\right)$, and since $\ln(1 + t) \sim t$ as $t \to 0$, it tends to $\lim \frac{3(x^3 + 3)}{x^3 - 2} = 3$.)
\end{solution}

\subsection*{Question 2a}

\begin{question}
Compute the intervals of convexity and the inflection points of the function
\[ f(x) = \sqrt[3]{6x + 12} \]
\end{question}

\begin{hint}
Write $f(x) = (6x + 12)^{1/3}$, differentiate twice and check the sign of $f''$. Pay attention to the point where $f''$ is not defined.
\end{hint}

\begin{solution}
\textbf{Domain:} the cube root is defined for every real number, so $f$ is defined and continuous on all of $\R$. Write $f(x) = (6x + 12)^{1/3}$. For $x \ne -2$ (i.e. $6x + 12 \ne 0$), by the chain rule:
\[ f'(x) = \frac13 (6x + 12)^{-2/3} \cdot 6 = 2(6x + 12)^{-2/3}, \]
\[ f''(x) = 2\cdot\left(-\frac23\right)(6x + 12)^{-5/3} \cdot 6 = -8(6x + 12)^{-5/3} = \frac{-8}{(6x + 12)^{5/3}}. \]
At $x = -2$ the function is not differentiable: $\frac{f(x) - f(-2)}{x + 2} = \frac{(6(x+2))^{1/3}}{x + 2} = \frac{6^{1/3}}{(x + 2)^{2/3}} \to +\infty$ (a vertical tangent).

\textbf{Sign of $f''$:} the power $(6x + 12)^{5/3}$ has the same sign as $6x + 12$ (an odd power of a cube root). Hence
\begin{itemize}
  \item $x < -2 \iff 6x + 12 < 0 \iff f''(x) > 0$, so $f$ is convex on $(-\infty, -2]$;
  \item $x > -2 \iff 6x + 12 > 0 \iff f''(x) < 0$, so $f$ is concave on $[-2, \infty)$.
\end{itemize}

\textbf{Inflection point:} at $x = -2$ the function is continuous, the graph has a (vertical) tangent, and the convexity changes on the two sides of the point. Hence $(-2, f(-2)) = (-2, 0)$ is an inflection point, even though $f$ is not differentiable there. There are no other inflection points, since $f'' \ne 0$ at every other point.

\textbf{Answer:} $f$ is convex on $(-\infty, -2)$, concave on $(-2, \infty)$, and the inflection point is $(-2, 0)$.
\end{solution}

\subsection*{Question 2b}

\begin{question}
Prove that for every $x \ge 0$
\[ 0 \le \frac{x^2}{e^x} \le 1 \]
\end{question}

\begin{hint}
Investigate $g(x) = \frac{x^2}{e^x}$ on $[0, \infty)$ and find its maximum.
\end{hint}

\begin{solution}
Define $g(x) = \frac{x^2}{e^x} = x^2 e^{-x}$ on $[0, \infty)$; $g$ is continuous and differentiable there.

\textbf{The left inequality:} $x^2 \ge 0$ and $e^x > 0$, hence $g(x) \ge 0$ for every $x$.

\textbf{The right inequality:} we compute
\[ g'(x) = \frac{2x e^x - x^2 e^x}{e^{2x}} = \frac{x(2 - x)}{e^x}. \]
On $(0, 2)$ we have $g' > 0$, and on $(2, \infty)$ we have $g' < 0$. Hence $g$ is increasing on $[0, 2]$ and decreasing on $[2, \infty)$, and $x = 2$ is an absolute maximum point of $g$ on $[0, \infty)$. Therefore, for every $x \ge 0$:
\[ g(x) \le g(2) = \frac{4}{e^2} = \left( \frac{2}{e} \right)^2 < 1, \]
since $0 < 2 < e$ and hence $0 < \frac{2}{e} < 1$.

Thus, for every $x \ge 0$: $0 \le \frac{x^2}{e^x} \le \frac{4}{e^2} < 1$, and in particular $0 \le \frac{x^2}{e^x} \le 1$. $\blacksquare$
\end{solution}

\subsection*{Question 3}

\begin{question}
Using a Taylor approximation, compute $\sqrt[3]{11}$ with an error less than $\frac{1}{10}$. Justify your answer.

Note that the approximation and the error bound must be rational numbers.
\end{question}

\begin{hint}
Expand $f(x) = \sqrt[3]{x}$ around $a = 8$ (since $\sqrt[3]{8} = 2$). Start with a first-order polynomial and check whether the Lagrange remainder is already less than $\frac{1}{10}$.
\end{hint}

\begin{solution}
Define $f(x) = \sqrt[3]{x} = x^{1/3}$ and expand around $a = 8$, the point closest to $11$ at which the cube root is known ($\sqrt[3]{8} = 2$). We look for the minimal $n$ such that $|R_n(11)| < \frac{1}{10}$.

\textbf{Derivatives} (for $x > 0$):
\[ f'(x) = \frac13 x^{-2/3}, \qquad f''(x) = -\frac29 x^{-5/3}, \]
hence $f(8) = 2$ and $f'(8) = \frac{1}{3\cdot 8^{2/3}} = \frac{1}{3\cdot 4} = \frac{1}{12}$.

\textbf{Lagrange remainder for $n = 1$:} by Taylor's theorem, there exists $c$ between $8$ and $11$ such that
\[ R_1(11) = \frac{f''(c)}{2!}(11 - 8)^2 = \frac{-\frac29 c^{-5/3}}{2}\cdot 9 = -\frac{1}{c^{5/3}}. \]
Since $8 < c < 11$, we have $c^{5/3} > 8^{5/3} = 2^5 = 32$, and hence
\[ |R_1(11)| = \frac{1}{c^{5/3}} < \frac{1}{32} < \frac{1}{10}. \]
Thus the first-order Taylor polynomial suffices ($n = 1$).

\textbf{The approximation:}
\[ P_1(x) = f(8) + f'(8)(x - 8) = 2 + \frac{1}{12}(x - 8), \]
\[ \sqrt[3]{11} = f(11) \approx P_1(11) = 2 + \frac{3}{12} = 2\frac14 = \frac94. \]

\textbf{Answer:} $\sqrt[3]{11} \approx \frac{9}{4} = 2.25$, with error less than $\frac{1}{32} < \frac{1}{10}$. (Since $R_1(11) < 0$, the approximation is larger than the true value: $\frac94 - \frac{1}{32} < \sqrt[3]{11} < \frac94$.)

(Check: $\sqrt[3]{11} \approx 2.22398$, and the actual error is about $0.026 < \frac{1}{32} \approx 0.031$.)
\end{solution}

\subsection*{Question 4a}

\begin{question}
Compute the integral
\[ \int \frac{6x^3 - 12x + 24}{x^4 - 4x^3 + 8x^2}\,dx \]
\end{question}

\begin{hint}
Factor the denominator: $x^4 - 4x^3 + 8x^2 = x^2(x^2 - 4x + 8)$, and the quadratic factor is irreducible. The partial fraction decomposition has the form $\frac{A}{x} + \frac{B}{x^2} + \frac{Cx + D}{x^2 - 4x + 8}$.
\end{hint}

\begin{solution}
This is a rational function, and we follow the usual steps.
\begin{itemize}
  \item \textbf{Degrees:} the degree of the numerator ($3$) is less than the degree of the denominator ($4$), so no polynomial division is needed.
  \item \textbf{Factoring the denominator:} $x^4 - 4x^3 + 8x^2 = x^2(x^2 - 4x + 8)$, and the discriminant of $x^2 - 4x + 8$ is $16 - 32 < 0$, so the quadratic factor is irreducible; in fact $x^2 - 4x + 8 = (x - 2)^2 + 4$.
  \item \textbf{Partial fractions:}
  \[ \frac{6x^3 - 12x + 24}{x^2(x^2 - 4x + 8)} = \frac{A}{x} + \frac{B}{x^2} + \frac{Cx + D}{x^2 - 4x + 8}, \]
  i.e. $6x^3 - 12x + 24 = Ax(x^2 - 4x + 8) + B(x^2 - 4x + 8) + (Cx + D)x^2$. Comparing coefficients:
  \begin{align*}
    x^0 &: \ 8B = 24 \ \Rightarrow\ B = 3, \\
    x^1 &: \ 8A - 4B = -12 \ \Rightarrow\ 8A = 0 \ \Rightarrow\ A = 0, \\
    x^2 &: \ -4A + B + D = 0 \ \Rightarrow\ D = -3, \\
    x^3 &: \ A + C = 6 \ \Rightarrow\ C = 6.
  \end{align*}
  Hence
  \[ \frac{6x^3 - 12x + 24}{x^4 - 4x^3 + 8x^2} = \frac{3}{x^2} + \frac{6x - 3}{x^2 - 4x + 8}. \]
\end{itemize}
\textbf{Integration.} Write $6x - 3 = 3(2x - 4) + 9$, so that in the first part the numerator is the derivative of the denominator, and in the second part we complete the square:
\[ \int \left( \frac{3}{x^2} + 3\cdot\frac{2x - 4}{x^2 - 4x + 8} + \frac{9}{(x - 2)^2 + 4} \right) dx. \]
\begin{itemize}
  \item $\int \frac{3}{x^2}\,dx = -\frac{3}{x} + C$.
  \item With the substitution $t = x^2 - 4x + 8$, $dt = (2x - 4)\,dx$: $\int \frac{2x - 4}{x^2 - 4x + 8}\,dx = \ln|t| = \ln(x^2 - 4x + 8) + C$ (the denominator is always positive).
  \item With the substitution $x - 2 = 2u$, $dx = 2\,du$: $\int \frac{9\,dx}{(x - 2)^2 + 4} = \int \frac{18\,du}{4u^2 + 4} = \frac92\arctan u + C = \frac92 \arctan\frac{x - 2}{2} + C$.
\end{itemize}
\textbf{Answer:}
\[ \boxed{\int \frac{6x^3 - 12x + 24}{x^4 - 4x^3 + 8x^2}\,dx = -\frac{3}{x} + 3\ln(x^2 - 4x + 8) + \frac92\arctan\left(\frac{x - 2}{2}\right) + C.} \]
\end{solution}

\subsection*{Question 4b}

\begin{question}
Find the general solution of the differential equation
\[ \sqrt{1 - x^2} \cdot y' = y^2 \arcsin x \]
\end{question}

\begin{hint}
The equation is separable. After separating, on the $x$ side substitute $t = \arcsin x$.
\end{hint}

\begin{solution}
The equation is defined for $-1 < x < 1$, and there $\sqrt{1 - x^2} > 0$. Write it in the form
\[ \frac{dy}{dx} = y^2 \cdot \frac{\arcsin x}{\sqrt{1 - x^2}}, \]
which is a separable equation.

\textbf{The constant solution:} $y \equiv 0$ satisfies the equation ($0 = 0$), so it is a solution.

\textbf{For $y \ne 0$:} separate the variables:
\[ \int \frac{dy}{y^2} = \int \frac{\arcsin x}{\sqrt{1 - x^2}}\,dx. \]
The left-hand side gives $-\frac{1}{y}$. On the right-hand side substitute $t = \arcsin x$, $dt = \frac{dx}{\sqrt{1 - x^2}}$:
\[ \int \frac{\arcsin x}{\sqrt{1 - x^2}}\,dx = \int t\,dt = \frac{t^2}{2} + C = \frac{\arcsin^2 x}{2} + C. \]
Hence
\[ -\frac{1}{y} = \frac{\arcsin^2 x}{2} + C \quad\Longrightarrow\quad y = \frac{-1}{\frac{\arcsin^2 x}{2} + C} = \frac{-2}{\arcsin^2 x + 2C}. \]

\textbf{Answer:} the general solution is
\[ \boxed{y = \frac{-2}{\arcsin^2 x + C}}, \quad C \in \R \]
(on every interval on which the denominator does not vanish), together with the solution $y \equiv 0$.
\end{solution}

\subsection*{Question 5a}

\begin{question}
Compute the area enclosed between the $x$-axis and the graph of the function
\[ f(x) = xe^x \]
on the interval $[-1, 1]$.
\end{question}

\begin{hint}
Note that $f$ changes sign at $x = 0$: the area is $\int_{-1}^1 |f(x)|\,dx$. Find the antiderivative using integration by parts.
\end{hint}

\begin{solution}
Since $e^x > 0$, the sign of $f(x) = xe^x$ is the sign of $x$: $f \le 0$ on $[-1, 0]$ and $f \ge 0$ on $[0, 1]$. Hence the area is
\[ S = \int_{-1}^1 |f(x)|\,dx = -\int_{-1}^0 f(x)\,dx + \int_0^1 f(x)\,dx. \]
\textbf{Antiderivative} by integration by parts, with $u = x$, $v' = e^x$ (i.e. $u' = 1$, $v = e^x$):
\[ F(x) = \int xe^x\,dx = xe^x - \int e^x\,dx = xe^x - e^x + C = (x - 1)e^x + C. \]
By the Newton--Leibniz formula:
\[ S = -\bigl(F(0) - F(-1)\bigr) + \bigl(F(1) - F(0)\bigr) = F(1) - 2F(0) + F(-1). \]
Here $F(1) = 0$, $F(0) = -1$, $F(-1) = -2e^{-1}$, and hence
\[ S = 0 - 2\cdot(-1) - \frac{2}{e} = 2 - \frac{2}{e}. \]
\textbf{Answer:} $\boxed{S = 2\left(1 - \frac1e\right)}$.
\end{solution}

\subsection*{Question 5b}

\begin{question}
Determine whether the following improper integral converges or diverges:
\[ \int_1^\infty \frac{\cos x}{x^2 - \ln x} \]
Remark: you may assume that the denominator is positive on the domain of integration.
\end{question}

\begin{hint}
The integrand does not have a constant sign. Check absolute convergence: bound $|f(x)| \le \frac{1}{x^2 - \ln x}$ and compare (using the limit comparison test) with $\frac{1}{x^2}$.
\end{hint}

\begin{solution}
Denote $f(x) = \frac{\cos x}{x^2 - \ln x}$. The denominator is positive on $[1, \infty)$ (as we may assume; in fact $\ln x \le x - 1 < x \le x^2$ for $x \ge 1$), so $f$ is continuous there and the integral is improper only because of the infinite limit. $f$ changes sign infinitely many times, so we check \textbf{absolute convergence}.

\textbf{Bound:} for every $x \ge 1$,
\[ 0 \le |f(x)| = \frac{|\cos x|}{x^2 - \ln x} \le \frac{1}{x^2 - \ln x} =: h(x). \]

\textbf{Convergence of $\int_1^\infty h$:} compare with $g(x) = \frac{1}{x^2}$. Both functions are positive, and
\[ \lim_{x\to\infty} \frac{g(x)}{h(x)} = \lim_{x\to\infty} \frac{x^2 - \ln x}{x^2} = \lim_{x\to\infty} \left( 1 - \frac{\ln x}{x^2} \right) = 1, \]
since $\frac{\ln x}{x^2} \to 0$ (L'Hôpital: $\frac{1/x}{2x} \to 0$). The limit is finite and positive, so by the limit comparison test $\int_1^\infty h$ and $\int_1^\infty g$ either both converge or both diverge. Since $\int_1^\infty \frac{dx}{x^\alpha}$ converges for $\alpha > 1$, and in particular $\int_1^\infty \frac{dx}{x^2} = 1$, $\int_1^\infty h$ also converges.

\textbf{Conclusion:} by the comparison test (since $0 \le |f| \le h$), $\int_1^\infty |f|$ converges. An absolutely convergent improper integral converges, and hence
\[ \int_1^\infty \frac{\cos x}{x^2 - \ln x}\,dx \ \text{converges (even absolutely).} \]
\end{solution}

\end{document}
